\documentclass[UTF8,fontset=none,11pt,oneside,openany,a4paper]{ctexbook}

\setCJKmainfont{Songti SC}
\setCJKsansfont{Heiti SC}
\setCJKmonofont{Heiti SC}
\setCJKfamilyfont{zhsong}{Songti SC}
\setCJKfamilyfont{zhhei}{Heiti SC}
\usepackage{geometry}
\geometry{left=2.6cm,right=2.6cm,top=2.6cm,bottom=2.6cm}

\usepackage{amsmath,amssymb,amsthm,bm}
\usepackage{mathtools}
\usepackage{tikz}
\usetikzlibrary{arrows.meta,calc,decorations.pathmorphing,decorations.markings,patterns,angles,quotes,3d}
\usepackage{pgfplots}
\pgfplotsset{compat=1.17}
\usepackage[most]{tcolorbox}
\usepackage{booktabs,enumitem,array,multirow}
\usepackage{caption,needspace}
\usepackage{hyperref}
\hypersetup{colorlinks=true,linkcolor=blue!50!black,urlcolor=blue!50!black}

\setcounter{secnumdepth}{2}
\setcounter{tocdepth}{2}

\ctexset{
  chapter/format={\centering\LARGE\bfseries},
  section/format={\Large\bfseries},
  subsection/format={\large\bfseries}
}

\newtheorem{definition}{定义}[chapter]
\newtheorem{theorem}{定理}[chapter]
\newtheorem{example}{例}[chapter]
\newtheorem{remark}{注}[chapter]

\newtcolorbox{guidebox}{
  colback=gray!4,colframe=gray!50!black,title={本章导读},
  fonttitle=\bfseries,boxrule=.6pt,arc=2pt,breakable
}
\newlist{bookexercises}{enumerate}{1}
\setlist[bookexercises]{label=\textbf{\thechapter.\arabic*},ref=\thechapter.\arabic*,leftmargin=3.2em}
\newlist{bookreading}{enumerate}{1}
\setlist[bookreading]{label={[\arabic*]},leftmargin=2em}
\newcommand{\exlevel}[1]{\textsf{〔#1〕}\,}

\newtcolorbox{keybox}[1][]{
  colback=blue!4,colframe=blue!55!black,
  fonttitle=\bfseries,title={核心结论},
  boxrule=0.6pt,arc=2pt,left=6pt,right=6pt,top=4pt,bottom=4pt,breakable,#1
}
\newtcolorbox{physicsbox}[1][]{
  colback=green!4,colframe=green!45!black,
  fonttitle=\bfseries,title={物理图像与理解（Physical Insight）},
  boxrule=0.6pt,arc=2pt,left=6pt,right=6pt,top=4pt,bottom=4pt,breakable,#1
}
\newtcolorbox{mathbox}[1][]{
  colback=orange!4,colframe=orange!60!black,
  fonttitle=\bfseries,title={数学注记},
  boxrule=0.6pt,arc=2pt,left=6pt,right=6pt,top=4pt,bottom=4pt,breakable,#1
}
\newtcolorbox{modernbox}[1][]{
  colback=violet!4,colframe=violet!55!black,
  fonttitle=\bfseries,title={进一步探索（Further Exploration）},
  boxrule=0.6pt,arc=2pt,left=6pt,right=6pt,top=4pt,bottom=4pt,breakable,#1
}

\newcommand{\dd}{\mathrm{d}}
\newcommand{\br}{\boldsymbol{r}}
\newcommand{\bv}{\boldsymbol{v}}
\newcommand{\ba}{\boldsymbol{a}}
\newcommand{\bF}{\boldsymbol{F}}
\newcommand{\bE}{\boldsymbol{E}}
\newcommand{\bB}{\boldsymbol{B}}
\newcommand{\bR}{\boldsymbol{R}}
\newcommand{\bJ}{\boldsymbol{J}}
\newcommand{\bM}{\boldsymbol{M}}
\newcommand{\bp}{\boldsymbol{p}}
\newcommand{\bP}{\boldsymbol{P}}
\newcommand{\bvec}[1]{\boldsymbol{#1}}
\newcommand{\bu}{\boldsymbol{u}}
\newcommand{\bg}{\boldsymbol{g}}
\newcommand{\bV}{\boldsymbol{V}}
\newcommand{\bT}{\boldsymbol{T}}
\newcommand{\bq}{\boldsymbol{q}}
\newcommand{\bA}{\boldsymbol{A}}
\newcommand{\bN}{\boldsymbol{N}}
\newcommand{\bI}{\boldsymbol{I}}
\newcommand{\bom}{\boldsymbol{\omega}}
\newcommand{\bOm}{\boldsymbol{\Omega}}
\newcommand{\bth}{\boldsymbol{\theta}}
\newcommand{\be}{\boldsymbol{e}}
\newcommand{\er}{\hat{\be}_r}
\newcommand{\etheta}{\hat{\be}_\theta}
\newcommand{\ephi}{\hat{\be}_\varphi}
\newcommand{\et}{\hat{\be}_t}
\newcommand{\en}{\hat{\be}_n}
\newcommand{\eb}{\hat{\be}_b}
\newcommand{\ii}{\hat{\boldsymbol{\imath}}}
\newcommand{\jj}{\hat{\boldsymbol{\jmath}}}
\newcommand{\kk}{\hat{\boldsymbol{k}}}
\newcommand{\dv}[2]{\frac{\dd #1}{\dd #2}}
\newcommand{\dvv}[3][]{\frac{\dd^{#1} #2}{\dd #3^{#1}}}
\newcommand{\pdv}[2]{\frac{\partial #1}{\partial #2}}
\newcommand{\ddv}[2]{\frac{\dd^2 #1}{\dd #2^2}}

\captionsetup{font=small,labelfont=bf}
\renewcommand{\arraystretch}{1.25}

\begin{document}

\frontmatter
\begin{titlepage}
\thispdfpagelabel{Cover}
\centering
\vspace*{3cm}
{\Huge\bfseries 理论力学\par}
\vspace{0.8cm}
{\large 第1章\quad 牛顿力学\par}
\vspace{0.15cm}
{\large 第2章\quad 变分法\par}
\vspace{0.15cm}
{\large 第3章\quad 拉格朗日力学\par}
\vspace{0.15cm}
{\large 第4章\quad 小振动\par}
\vspace{0.15cm}
{\large 第5章\quad 刚体力学\par}
\vspace{0.15cm}
{\large 第6章\quad 哈密顿力学\par}
\vspace{0.15cm}
{\large 第7章\quad 狭义相对论\par}
\vspace{0.15cm}
\vspace{1.2cm}
{\large 全书统稿版\par}
\vfill
{\large 2026\par}
\end{titlepage}

\tableofcontents
\chapter*{读者指南与记号约定}
\phantomsection
\addcontentsline{toc}{chapter}{读者指南与记号约定}

初读时，每节先弄清正在研究的物体、观察方式和要算的量。遇到新符号，回到它第一次出现的等式，找出它由哪些已知量定义。遇到较长的推导，可以先用一句话说出每一步在做什么，例如投影、积分、用初值定常数，再自己写出中间的等式。

例题可以先盖住解答，试着画图并写第一条方程。章节末尾先做基础题，再做综合题。拓展和计算部分可以在掌握核心模型之后回读。下表可用来查找跨章重复使用的符号。

\medskip
\begin{center}
\begin{tabular}{@{}p{3.3cm}p{10.7cm}@{}}
\toprule
记号 & 含义与使用条件\\
\midrule
$\br,\bv,\ba$ & 位置、速度和加速度；比较分量前先注明参考系和基矢。\\
$T,V,E_{\rm mech}$ & 经典动能、位置势能、机械能；在相应模型中 $E_{\rm mech}=T+V$。\\
$L,S,E_L,H$ & 拉格朗日量、作用量、能量函数、哈密顿量。能否将 $E_L$ 或 $H$ 认作机械能，须核对模型与变量。\\
$\bp,\boldsymbol{\pi}$ & 电磁耦合中的广义动量（第6章称正则动量）和机械动量；自由质点两者相同。\\
$e,e_r$ & 电磁模型的带符号电荷、碰撞模型的恢复系数；轨道例中的 $e$ 也可局部表示偏心率。\\
$\theta,\varphi$ & 球坐标中为极角、方位角；第1章的平面极角 $\theta$ 在第3、6章的平面中心力模型中也记作 $\varphi$。\\
$m,M,E,K$ & 第7章分别为粒子不变质量、系统不变质量、粒子总能量与动能；$K=E-mc^2$。本书取时空度规号差 $(+,-,-,-)$。\\
\bottomrule
\end{tabular}
\end{center}

\textbf{怎样核对结果。}
先读题设与假设，再看推导和量纲；若使用守恒量，先说明系统边界。核心结论框归纳公式及条件，物理图像框解释机制，数学注记说明判据，进一步探索框给出延伸主题。各章习题按“章号.题号”编号，标出基础、综合或拓展；解答使用同一题号。计算题可用随书脚本检查残差、守恒或步长误差。

\mainmatter

\chapter{牛顿力学}\label{ch:particle}

\begin{guidebox}
想象你在分析一辆正在过弯的车。先选参考系，说明车在哪里、速度朝哪，再问什么力使速度改变。若题目只问到达某处时有多快，能量关系往往能省去求整条轨迹的步骤。研究多个物体时，先圈出研究对象，再分清内力和外力。

本章采用低速、宏观的经典模型，各惯性系共用同一时间。除火箭等专门说明的开放系统外，质点质量取常量。读到转动参考系和中心力时，可以回到三个问题检查自己的方程。观察者跟谁一起运动？哪些条件限制了运动？什么量能够守恒？
\end{guidebox}

\section{运动的描述与分量表达式}

先问怎样记录运动。若把物体的每个时刻都拍下来，照片上的位置点能画出轨迹，但只有给每张照片标上时刻，才能进一步算速度和加速度。

同一运动可用不同坐标记录。走廊里沿两面墙量距离很方便，绕固定中心运动时，量半径和转角通常更省事。接下来用几种坐标描述同一类位置、速度和加速度。

\subsection{参考系、坐标系与质点}

车厢里放着一只静止的行李箱。坐在车里的人说它没动，站台上的人却看到它随车经过。先说明跟着谁观察，再选择怎样标记位置。若我们只关心箱子走了多远，还可以把它当作一个有质量的点。

运动是相对的。描述一个物体的位置及其变化，必须先指定一个作为基准的物体或物体群，称为\textbf{参考系}（也称参照系）。在参考系上建立适当的数学框架，用以定量地标定空间位置，称为\textbf{坐标系}。参考系是物理的选择，坐标系是数学的选择；同一参考系中可以选用直角坐标、极坐标、球坐标等不同坐标系。

\begin{definition}[质点]
当物体的形状和大小对所研究的问题影响可以忽略时，可以把物体抽象为具有全部质量的几何点，称为\textbf{质点}。质点是一个理想模型，而不是“很小的物体”。
\end{definition}

\begin{physicsbox}
质点模型是否适用，取决于“尺度比”而不是绝对大小。研究地球公转时，地球可视为质点；研究地球自转或潮汐时，地球的大小和形状不可忽略，质点模型失效。判断标准是：物体的线度远小于问题中的特征长度，且转动、形变等内部自由度不参与动力学。
\end{physicsbox}

在直角坐标系中，质点 $P$ 的位置由三个坐标 $(x,y,z)$ 给出。从坐标原点 $O$ 指向 $P$ 的有向线段称为\textbf{位置矢量}，简称\textbf{位矢}：
\begin{equation}
\br=x\ii+y\jj+z\kk .
\label{eq:position}
\end{equation}
位矢的大小为
\begin{equation}
r=|\br|=\sqrt{x^2+y^2+z^2}.
\end{equation}

\begin{figure}[htbp]
\centering
\begin{tikzpicture}[>=Stealth,x={(1cm,0cm)},y={(-0.55cm,-0.4cm)},z={(0cm,1cm)}]
\coordinate (O) at (0,0,0);\coordinate (Q) at (3,1.5,0);\coordinate (P) at (3,1.5,2.5);
\draw[->] (O)--(4,0,0) node[right]{$x$};
\draw[->] (O)--(0,3,0) node[below]{$y$};
\draw[->] (O)--(0,0,3.4) node[above]{$z$};
\draw[dashed] (3,0,0)--(Q)--(0,1.5,0);
\draw[dashed] (Q)--(P)--(0,0,2.5);
\draw[->,very thick,blue!70!black] (O)--(P) node[midway,left]{$\br$};
\fill (P) circle (2pt) node[right]{$P(x,y,z)$};
\fill (Q) circle (1.5pt) node[below right]{$Q(x,y,0)$};
\node[below] at (3,0,0) {$x$};\node[left] at (0,1.5,0) {$y$};\node[left] at (0,0,2.5) {$z$};
\end{tikzpicture}
\caption{位矢与三个坐标分量。图为同一三维坐标系的平行投影，$Q$ 为 $P$ 在 $xy$ 平面的垂足。}
\label{fig:cartesian}
\end{figure}

\subsection{运动方程、位移、速度与加速度}

知道一辆车经过哪些地方，还不能判断它在哪里快、在哪里慢。给位置加上时间，就可以用相邻时刻的位置变化定义速度，再用速度的变化定义加速度。

当质点运动时，其坐标随时间变化。位矢作为时间的函数
\begin{equation}
\br(t)=x(t)\ii+y(t)\jj+z(t)\kk
\label{eq:motion}
\end{equation}
称为\textbf{运动方程}。消去时间 $t$，得到坐标之间的约束关系
\begin{equation}
F(x,y,z)=0,
\label{eq:orbit}
\end{equation}
通常只能描述一个曲面；一般三维轨迹需两条独立关系 $F_1(x,y,z)=F_2(x,y,z)=0$，或保留参数形式。对平面运动，一条平面内的关系才足以称为\textbf{轨道方程}。运动方程给出“何时到何处”，轨道方程只给出“经过哪些点”。


设质点在 $t$ 时刻位于 $A$，在 $t+\Delta t$ 时刻位于 $B$，则
\begin{equation}
\Delta\br=\br_B-\br_A
\end{equation}
称为 $\Delta t$ 时间内的\textbf{位移}。位移是矢量，描述位置的净变化；它与路程 $\Delta s$ 不同，后者是轨迹弧长的增量。

平均速度为 $\Delta\br/\Delta t$。令 $\Delta t\to0$，得到\textbf{瞬时速度}
\begin{equation}
\bv=\lim_{\Delta t\to0}\frac{\Delta\br}{\Delta t}
=\dv{\br}{t}=\dot{\br}.
\label{eq:velocity}
\end{equation}
在直角坐标中
\begin{equation}
\bv=\dot x\ii+\dot y\jj+\dot z\kk .
\end{equation}
速度的大小称为\textbf{速率}：
\begin{equation}
v=|\bv|=\dv{s}{t}=\sqrt{\dot x^2+\dot y^2+\dot z^2}.
\end{equation}
加速度定义为
\begin{equation}
\ba=\dv{\bv}{t}=\ddv{\br}{t}
=\ddot x\ii+\ddot y\jj+\ddot z\kk .
\label{eq:acceleration}
\end{equation}

\begin{figure}[htbp]
\centering
\begin{tikzpicture}[>=Stealth]
\draw[->] (-.6,0)--(4.1,0) node[right]{$x$};\draw[->] (-.3,-.3)--(-.3,2.5) node[above]{$y$};
\draw[thick,blue!65!black,domain=-.3:3.6,samples=80] plot (\x,{1+.5*\x-.12*\x*\x});
\coordinate (A) at (0,1);\coordinate (B) at (2,1.52);
\fill (A) circle (2pt) node[left]{$A$};\fill (B) circle (2pt) node[above]{$B$};
\draw[->,thick,red!70!black] (A)--(B) node[midway,below]{$\Delta\br$};
\draw[->,thick,green!40!black] (A)--(1,1.5) node[above]{$\bv_A$};
\end{tikzpicture}
\caption{解析曲线上两点的位移沿弦方向，而瞬时速度沿切线。图中弦与切线由同一曲线计算，只有在两事件趋近时方向才一致。}
\label{fig:displacement}
\end{figure}

\begin{physicsbox}
速度是“位置对时间的变化率”，加速度是“速度对时间的变化率”。二者都是矢量，但方向一般不同：匀速圆周运动的加速度指向圆心，而速度沿切线；一般曲线运动中，加速度总可以分解为切向分量（改变速率）与法向分量（改变速度方向）。
\end{physicsbox}


同一运动在不同坐标系中有不同的分量形式。选择坐标系的原则是：让约束和对称性尽可能简单。

\subsection{直角坐标与平面极坐标}

沿墙量长宽时，坐标轴一直朝同一方向。跟着转盘上的一点看，径向和切向却一直在转。对极坐标位置求导时，要把长度的变化和方向的变化都计入。

在直角坐标系中，基矢量 $\ii,\jj,\kk$ 不随时间变化，因此
\begin{equation}
\bv=\dot x\ii+\dot y\jj+\dot z\kk,\qquad
\ba=\ddot x\ii+\ddot y\jj+\ddot z\kk .
\end{equation}
这是最简单的情形：对分量分别求导即可。

在平面极坐标 $(r,\theta)$ 中，取径向单位矢量 $\er$ 沿位矢方向，横向单位矢量 $\etheta$ 垂直于 $\er$ 且指向 $\theta$ 增大的方向：
\begin{equation}
\er=\cos\theta\,\ii+\sin\theta\,\jj,\qquad
\etheta=-\sin\theta\,\ii+\cos\theta\,\jj .
\end{equation}
位矢的直角坐标参数表示为
\begin{equation}
\br=r\cos\theta\,\ii+r\sin\theta\,\jj .
\label{eq:position-polar-cartesian}
\end{equation}
固定 $r$，对 $\theta$ 求偏导：
\begin{equation}
\left.\pdv{\br}{\theta}\right|_r
=r(-\sin\theta\,\ii+\cos\theta\,\jj)=r\etheta .
\label{eq:polar-tangent}
\end{equation}
因此
\[
\etheta=\frac{1}{r}\left.\pdv{\br}{\theta}\right|_r,
\]
即 $\etheta$ 是圆周 $r=\text{常量}$ 的切向方向。类似地，固定 $\theta$ 对 $r$ 求偏导，得到
\[
\left.\pdv{\br}{r}\right|_\theta
=\cos\theta\,\ii+\sin\theta\,\jj=\er,
\]
即径向方向。可见极坐标的两个基矢量分别是位矢对 $r$、$\theta$ 的偏导数；其中 $\etheta$ 需除以 $r$ 才能归一化为单位矢量。

当 $\theta$ 变化时，基矢量也随之转动：
\begin{equation}
\dv{\er}{\theta}=\etheta,\qquad
\dv{\etheta}{\theta}=-\er .
\label{eq:polar-unit}
\end{equation}
因此
\begin{equation}
\dv{\er}{t}=\dot\theta\,\etheta,\qquad
\dv{\etheta}{t}=-\dot\theta\,\er .
\end{equation}
位矢为 $\br=r\er$，求导得
\begin{equation}
\bv=\dot r\,\er+r\dot\theta\,\etheta .
\label{eq:polar-velocity}
\end{equation}
再求导，并利用基矢量导数，得
\begin{equation}
\ba=(\ddot r-r\dot\theta^2)\er+
\left(r\ddot\theta+2\dot r\dot\theta\right)\etheta .
\label{eq:polar-acceleration}
\end{equation}
其中 $a_r=\ddot r-r\dot\theta^2$ 称为\textbf{径向加速度}，$a_\theta=r\ddot\theta+2\dot r\dot\theta$ 称为\textbf{横向加速度}。横向项也可写成
\begin{equation}
a_\theta=\frac{1}{r}\dv{}{t}\left(r^2\dot\theta\right),
\label{eq:polar-transverse}
\end{equation}
这一形式在角动量守恒问题中特别有用。

\begin{figure}[htbp]
\centering
\begin{tikzpicture}[>=Stealth,scale=1.1]
  \draw[->] (-0.5,0)--(4.2,0) node[below]{$x$};
  \draw[->] (0,-0.5)--(0,3.4) node[left]{$y$};
  \coordinate (P) at (2.7,1.9);
  \draw[->,thick] (0,0)--(P) node[midway,above left]{$\br$};
  \draw[->,very thick,blue!70!black] (P)--++(0.95,0.67) node[right]{$\er$};
  \draw[->,very thick,red!75!black] (P)--++(-0.67,0.95) node[above]{$\etheta$};
  \draw[dashed] (0,0)--(2.7,0) node[midway,below]{$r\cos\theta$};
  \draw[dashed] (2.7,0)--(P) node[midway,right]{$r\sin\theta$};
  \draw (1.0,0) arc (0:35:1.0);
  \node at (1.2,0.25) {$\theta$};
  \fill (P) circle (1.6pt) node[right]{$P$};
\end{tikzpicture}
\caption{平面极坐标。$\er$ 沿位矢方向，$\etheta$ 垂直于 $\er$ 指向 $\theta$ 增大的方向。}
\label{fig:polar}
\end{figure}

\begin{physicsbox}
式 \eqref{eq:polar-acceleration} 中，$-\,r\dot\theta^2\er$ 是向心加速度项，$2\dot r\dot\theta\,\etheta$ 是科里奥利型项。它们不是“新的力”，而是基矢量随质点转动所产生的几何效应。只要把 $\er,\etheta$ 看成随位置变化的标架，这些项就会自然出现。
\end{physicsbox}

\begin{example}[同一椭圆，恒角速度与开普勒运动有何不同？]
取平面椭圆 $r=p/(1+e\cos\theta)$，$p=a(1-e^2)$，$0<e<1$。先把它当作给定路径，若外部驱动强制 $\dot\theta=\omega$ 恒定，则
\[
\dot r=\frac{pe\sin\theta}{(1+e\cos\theta)^2}\omega,
\qquad v_\theta=r\omega.
\]
此时面积速度 $r^2\omega/2$ 随位置改变，并非开普勒运动，通常需要额外切向力。

若改为平方反比引力下的真实轨道，面积常量 $h=r^2\dot\theta$ 给出 $\dot\theta=h/r^2$。用同一个 $\dd r/\dd\theta$，得到
\[
v_r=\frac{he}{p}\sin\theta,\qquad
v_\theta=\frac{h}{r}=\frac{h}{p}(1+e\cos\theta),
\]
\[
v=\frac{h}{p}\sqrt{1+2e\cos\theta+e^2}.
\]
近心点与远心点处 $v_r=0$，且 $v_p/v_a=r_a/r_p=(1+e)/(1-e)>1$：真实引力轨道在近心点更快。$e\to0$ 时两种描述都退化为圆周运动。给定一条几何轨迹并不等于已经给定它的时间演化。
\end{example}

\subsection{自然坐标系与内禀方程}

车过弯时，即使速度表读数不变，速度的方向也在变。沿路面的切线方向看速率怎样变，沿指向弯内的法线方向看车怎样转弯。

当轨道已知时，用弧长 $s$ 作为参数最方便。定义切向单位矢量
\begin{equation}
\et=\dv{\br}{s},
\end{equation}
它沿轨道切线指向 $s$ 增大的方向。在光滑正则曲线上，曲率定义为
\begin{equation}
\kappa=\left|\dv{\et}{s}\right|=\frac{1}{\rho},
\end{equation}
其中 $\rho$ 为曲率半径。平面曲线还可用切线转角表示 $\kappa=|\dd\varphi/\dd s|$；空间曲线一般不能用单个切线角代替。以下要求速率非零、曲率非零；直线段可取 $\rho\to\infty$，但主法向不唯一。主法向单位矢量 $\en$ 指向曲率中心，满足
\begin{equation}
\dv{\et}{s}=\kappa\en=\frac{1}{\rho}\en .
\end{equation}
于是
\begin{equation}
\bv=\dv{\br}{s}\dv{s}{t}=v\et,
\end{equation}
\begin{equation}
\ba=\dv{\bv}{t}=\dot v\,\et+\frac{v^2}{\rho}\en .
\label{eq:natural-acceleration}
\end{equation}
这就是自然坐标下的\textbf{内禀方程}：切向加速度 $\dot v$ 改变速率，法向加速度 $v^2/\rho$ 改变速度方向。

\begin{figure}[htbp]
\centering
\begin{tikzpicture}[>=Stealth]
\draw[very thick,blue!65!black,domain=-1.8:1.8,samples=100] plot (\x,{\x*\x/4});
\draw[dashed,gray] (0,2) circle (2);
\fill (0,0) circle (2pt) node[below]{$P$};
\fill (0,2) circle (2pt) node[above right]{$C$};
\draw[->,thick,red!70!black] (0,0)--(1.5,0) node[right]{$\et$};
\draw[->,thick,blue!70!black] (0,0)--(0,1.2) node[left]{$\en$};
\draw[<->] (-0.5,0)--(-0.5,2) node[midway,left]{$\rho=2$};
\node at (2.5,0.3) {$\eb\odot$};
\end{tikzpicture}
\caption{以无量纲曲线 $y=x^2/4$ 的顶点为例：曲率半径为2，密切圆与曲线在 $P$ 处具有相同切向和曲率。副法向垂直纸面向外，不能画成平面内的第三个独立方向。}
\label{fig:natural}
\end{figure}

对于平面曲线 $y=y(x)$，曲率半径的常用公式为
\begin{equation}
\rho=\frac{\left[1+(y')^2\right]^{3/2}}{|y''|}.
\label{eq:curvature-plane}
\end{equation}
对于一般空间曲线 $\br(t)$，
\begin{equation}
\kappa=\frac{|\br'\times\br''|}{|\br'|^3}.
\label{eq:curvature-space}
\end{equation}

\begin{example}[加速度与速度夹角恒定的圆周运动]
质点沿半径为 $R$ 的圆周运动，其加速度 $\ba$ 与速度 $\bv$ 的夹角 $\alpha$ 保持不变，初速率为 $v_0$。求速率 $v(t)$，并讨论 $\alpha=90^\circ$ 与 $\alpha<90^\circ$ 的物理含义。

\begin{center}
\begin{tikzpicture}[>=Stealth,scale=0.82]
  \draw[thick] (0,0) circle (1.35);
  \coordinate (P) at (0.955,0.955);
  \fill (P) circle (2.2pt) node[above right]{$P$};
  \draw[->,very thick,blue!70!black] (P)--++(-0.66,0.66) node[above]{$\bv$};
  \draw[->,very thick,red!75!black] (P)--++(-0.78,-0.21) node[below]{$\ba$};
  \draw[dashed] (P)--(0,0);
  \draw[->] (P)++(-0.62,0.62) arc (135:195:0.58);
  \node at ($(P)+(-0.88,0.20)$) {$\alpha$};
\end{tikzpicture}
\end{center}

解：在自然坐标系中，加速度分解为切向与法向分量：
\[
a_t=\dot v,\qquad a_n=\frac{v^2}{R}.
\]
由“$\ba$ 与 $\bv$ 的夹角 $\alpha$ 恒定”得
\[
\tan\alpha=\frac{a_n}{a_t}
=\frac{v^2}{R\dot v}.
\]
于是
\[
\frac{\dd v}{v^2}=\frac{\dd t}{R\tan\alpha}.
\]
积分并代入 $v(0)=v_0$：
\[
-\frac1v+\frac1{v_0}=\frac{t}{R\tan\alpha},
\]
即
\[
v(t)=\frac{v_0}{1-\dfrac{v_0t}{R\tan\alpha}} .
\]
当 $\alpha=90^\circ$ 时，$\tan\alpha\to\infty$，$v(t)\equiv v_0$，退化为匀速圆周运动。当 $0<\alpha<90^\circ$ 时，加速度有沿速度方向的分量，速率增大；形式上在 $t_*=R\tan\alpha/v_0$ 处发散，说明不存在延拓到 $t_*$ 的有限速率运动；实际驱动力与相对论效应会限制这一理想模型。当 $90^\circ<\alpha<180^\circ$ 时，切向分量与速度相反，质点减速。
\end{example}

\subsection{柱坐标与球坐标}

绕竖直轴转动的物体，可以用离轴多远、转过多大角、升到多高来定位。若问题围绕一个中心点，则改用离中心的距离和两个方向角。

先从柱坐标 $(\rho,\varphi,z)$ 开始。它在水平面使用极坐标，沿轴仍用直角坐标：$x=\rho\cos\varphi,y=\rho\sin\varphi$。令 $\hat{\be}_\rho=(\cos\varphi,\sin\varphi,0)$，$\hat{\be}_\varphi=(-\sin\varphi,\cos\varphi,0)$，则
\[
\br=\rho\hat{\be}_\rho+z\kk,\qquad
\dot{\hat{\be}}_\rho=\dot\varphi\hat{\be}_\varphi,\quad
\dot{\hat{\be}}_\varphi=-\dot\varphi\hat{\be}_\rho.
\]
逐项微分得
\begin{align}
\bv&=\dot\rho\hat{\be}_\rho+\rho\dot\varphi\hat{\be}_\varphi+\dot z\kk,\\
\ba&=(\ddot\rho-\rho\dot\varphi^2)\hat{\be}_\rho
+(\rho\ddot\varphi+2\dot\rho\dot\varphi)\hat{\be}_\varphi+\ddot z\kk.
\end{align}
若 $z$ 恒定，它恢复前面的平面极坐标；柱坐标特别适合轴对称问题。注意 $\rho$ 是到轴的距离，而下面的 $r$ 是到原点的距离。

在球坐标 $(r,\theta,\varphi)$ 中，
\begin{equation}
x=r\sin\theta\cos\varphi,\qquad
y=r\sin\theta\sin\varphi,\qquad
z=r\cos\theta .
\end{equation}
三个单位矢量为
\begin{equation}
\er=\sin\theta\cos\varphi\,\ii+\sin\theta\sin\varphi\,\jj+\cos\theta\,\kk,
\end{equation}
\begin{equation}
\etheta=\cos\theta\cos\varphi\,\ii+\cos\theta\sin\varphi\,\jj-\sin\theta\,\kk,
\end{equation}
\begin{equation}
\ephi=-\sin\varphi\,\ii+\cos\varphi\,\jj .
\end{equation}
它们随时间的导数为
\begin{equation}
\dot{\er}=\dot\theta\,\etheta+\dot\varphi\sin\theta\,\ephi,
\end{equation}
\begin{equation}
\dot{\etheta}=-\dot\theta\,\er+\dot\varphi\cos\theta\,\ephi,
\end{equation}
\begin{equation}
\dot{\ephi}=-\dot\varphi\left(\sin\theta\,\er+\cos\theta\,\etheta\right).
\end{equation}
于是速度为
\begin{equation}
\bv=\dot r\,\er+r\dot\theta\,\etheta+r\dot\varphi\sin\theta\,\ephi .
\label{eq:spherical-velocity}
\end{equation}
为看清交叉项的来源，将 $\dot{\bv}$ 按三个乘积求导：
\[
\ba=\ddot r\er+\dot r\dot{\er}
+(\dot r\dot\theta+r\ddot\theta)\etheta+r\dot\theta\dot{\etheta}
+\frac{\dd}{\dd t}(r\dot\varphi\sin\theta)\ephi
+r\dot\varphi\sin\theta\dot{\ephi}.
\]
代入基矢导数，分别收集 $\er,\etheta,\ephi$ 的系数，得到加速度分量
\begin{align}
a_r&=\ddot r-r\dot\theta^2-r\dot\varphi^2\sin^2\theta,\\
a_\theta&=r\ddot\theta+2\dot r\dot\theta-r\dot\varphi^2\sin\theta\cos\theta,\\
a_\varphi&=\frac{1}{r\sin\theta}\dv{}{t}\left(r^2\dot\varphi\sin^2\theta\right).
\end{align}
其中 $a_\varphi$ 的写法体现了绕 $z$ 轴的角动量守恒结构。

\begin{figure}[htbp]
\centering
\begin{tikzpicture}[>=Stealth]
\begin{scope}
\draw[->] (0,0)--(3,0) node[right]{$x$};\draw[->] (0,0)--(0,2.6) node[above]{$y$};
\coordinate (Q) at (2.2,1.6);
\draw[->,thick,blue!70!black] (0,0)--(Q) node[midway,above]{$\rho$};
\draw (0.85,0) arc (0:36:0.85);\node at (1.05,0.3) {$\varphi$};
\draw[dashed] (Q)--(2.2,0);\draw[dashed] (Q)--(0,1.6);
\node at (1.2,-0.7) {$xy$ 平面投影};
\end{scope}
\begin{scope}[shift={(5.5,0)}]
\draw[->] (0,0)--(3.1,0) node[right]{$\rho$};\draw[->] (0,0)--(0,2.9) node[above]{$z$};
\draw[->,thick,blue!70!black] (0,0)--(2.2,2) node[midway,above left]{$r$};
\draw[dashed] (2.2,0)--(2.2,2)--(0,2);
\draw (0,1) arc (90:42:1);\node at (0.45,1.2) {$\theta$};
\node[below] at (1.1,0) {$r\sin\theta$};\node[left] at (0,2) {$r\cos\theta$};
\node at (1.3,-0.7) {过 $z$ 轴的子午截面};
\end{scope}
\end{tikzpicture}
\caption{球坐标的两个截面：方位角 $\varphi$ 在 $xy$ 平面内测量，极角 $\theta$ 从 $+z$ 向位矢测量。两个角的图示定义与正文一致。}
\label{fig:spherical}
\end{figure}

\begin{mathbox}
单位矢量的导数可以统一写成
\begin{equation}
\dv{\be_i}{t}=\bom\times\be_i,
\end{equation}
其中 $\bom$ 是标架转动的角速度。极坐标中 $\bom=\dot\theta\kk$，球坐标中 $\bom=\dot\varphi\kk+\dot\theta\,\ephi$（注意 $\ephi$ 本身也在转）。用这个观点可以快速检查式 \eqref{eq:polar-acceleration} 和球坐标加速度的各项。
\end{mathbox}


最后，运动学问题通常分为两类：
\begin{enumerate}[label=(\arabic*),leftmargin=2.2em]
  \item 已知运动方程 $\br(t)$，求速度、加速度：对时间求导；
  \item 已知加速度 $\ba(t)$ 或速度 $\bv(t)$，求运动方程：积分并利用初始条件。
\end{enumerate}
第二类问题的数学形式为
\begin{equation}
\bv(t)=\bv_0+\int_{t_0}^{t}\ba(t')\,\dd t',\qquad
\br(t)=\br_0+\int_{t_0}^{t}\bv(t')\,\dd t'.
\end{equation}
初始条件 $(\br_0,\bv_0)$ 的选取决定了具体的运动轨道。

\begin{example}[月球车车轮上的点：摆线]
月球车车轮半径为 $R$，在月面作无滑动滚动，轮缘上一点 $P$ 的轨迹称为\textbf{摆线}。设 $t=0$ 时 $P$ 与月面接触，圆心以速度 $v_0=R\omega$ 匀速前进。求 $P$ 的运动方程、速度和加速度，并解释轨迹尖点处的运动学。

\begin{center}
\begin{tikzpicture}[>=Stealth,scale=0.82]
  \draw[->] (-0.4,0)--(6.6,0) node[below]{$x$};
  \draw[->] (0,-0.25)--(0,2.35) node[left]{$y$};
  \draw[very thick,blue!70!black,domain=0:6.283,samples=220]
    plot ({0.8*(\x - sin(\x r))},{0.8*(1-cos(\x r))});
  \coordinate (C) at (2.513,0.8);
  \draw[thick] (C) circle (0.8);
  \draw[dashed] (C)--++(0,0.8) node[midway,left]{$R$};
  \fill (C) circle (1.6pt) node[below left]{$C$};
  \coordinate (P) at (2.513,1.6);
  \fill (P) circle (2.2pt) node[above right]{$P$};
  \draw[->,thick,red!75!black] (C)++(0,0.18)--++(0.72,0) node[below]{$v_0$};
  \draw[->,thick,green!45!black] (P)--++(0.0,-0.72) node[right]{$\ba$};
\end{tikzpicture}
\end{center}

解：取滚动方向为 $x$ 轴，月面为 $y=0$。设车轮转角的正大小为 $\theta=\omega t$。无滑动条件给出圆心坐标 $(R\theta,R)$；点 $P$ 相对圆心的坐标为 $(-R\sin\theta,-R\cos\theta)$。因此
\[
x=R(\theta-\sin\theta),\qquad y=R(1-\cos\theta).
\]
对时间求导，注意 $\dot\theta=\omega$：
\[
\dot x=R\omega(1-\cos\theta),\qquad
\dot y=R\omega\sin\theta,
\]
速率为
\[
v=\sqrt{\dot x^2+\dot y^2}
=R\omega\sqrt{2(1-\cos\theta)}
=2R\omega\left|\sin\frac{\theta}{2}\right|.
\]
再求导得
\[
\ddot x=R\omega^2\sin\theta,\qquad
\ddot y=R\omega^2\cos\theta,
\]
故加速度大小恒为 $R\omega^2$，方向始终指向圆心。

在 $\theta=0$ 处，$P$ 与月面接触，$v=0$，但 $\ba=(0,R\omega^2)\neq0$。轨迹在此出现\textbf{尖点}（cusp）：速度为零只是瞬时静止，并不表示加速度为零。在 $\theta=\pi$ 处，$P$ 位于最高点，$v=2R\omega$，加速度向下。这个例子说明：运动学描述必须同时看位置、速度和加速度，不能由“某一时刻速度为零”推断“没有加速度”。月面重力只影响动力学，不改变车轮滚动的几何轨迹。
\end{example}

\begin{example}[空间站机械臂：正运动学与雅可比矩阵]
考虑基座固定的平面二连杆教学模型：两段刚性连杆长度分别为 $l_1,l_2$，关节角为 $\theta_1,\theta_2$，末端执行器为 $P$。求末端位置、速度和加速度，并讨论奇异位形。本例用于训练复合运动求导，不代表实际空间机械臂的完整结构；若基座自由漂浮，还须计入基座反冲。

\begin{center}
\begin{tikzpicture}[>=Stealth,scale=1.0]
  \draw[->] (-0.5,0)--(5.2,0) node[below]{$x$};
  \draw[->] (0,-0.5)--(0,4.2) node[left]{$y$};
  \coordinate (O) at (0,0);
  \coordinate (A) at (2.718,1.269);
  \coordinate (P) at (3.470,3.336);
  \draw[very thick,blue!70!black] (O)--(A) node[midway,below left]{$l_1$};
  \draw[very thick,red!75!black] (A)--(P) node[midway,above left]{$l_2$};
  \fill (O) circle (2.2pt) node[below left]{$O$};
  \fill (A) circle (2.6pt) node[below left]{关节 1};
  \fill (P) circle (2.8pt) node[right]{末端执行器 $P$};
  \draw[dashed] (A)--++(25:1.5);
  \draw (1.1,0) arc (0:25:1.1);
  \node at (1.28,0.32) {$\theta_1$};
  \draw (A) ++(25:0.85) arc (25:70:0.85);
  \coordinate (L2) at ($(A)+(0.95,0.78)$);
  \node at (L2) {$\theta_2$};
\end{tikzpicture}
\end{center}

解：设基座 $O$ 为原点，第一段连杆与 $x$ 轴夹角为 $\theta_1$，第二段连杆相对第一段的夹角为 $\theta_2$。末端位置为
\begin{equation}
x=l_1\cos\theta_1+l_2\cos(\theta_1+\theta_2),\qquad
y=l_1\sin\theta_1+l_2\sin(\theta_1+\theta_2).
\label{eq:robot-fk}
\end{equation}
这就是\textbf{正运动学}：给定关节角 $(\theta_1,\theta_2)$，求末端位置 $(x,y)$。

对时间求导，得速度映射
\begin{equation}
\begin{pmatrix}\dot x\\ \dot y\end{pmatrix}
=
\underbrace{
\begin{pmatrix}
-l_1\sin\theta_1-l_2\sin(\theta_1+\theta_2) & -l_2\sin(\theta_1+\theta_2)\\
l_1\cos\theta_1+l_2\cos(\theta_1+\theta_2) & l_2\cos(\theta_1+\theta_2)
\end{pmatrix}
}_{J(\theta_1,\theta_2)}
\begin{pmatrix}\dot\theta_1\\ \dot\theta_2\end{pmatrix}.
\label{eq:robot-jacobian}
\end{equation}
矩阵 $J$ 称为\textbf{雅可比矩阵}，它把关节速度映射为末端速度。再求导得
\begin{equation}
\begin{pmatrix}\ddot x\\ \ddot y\end{pmatrix}
=J\begin{pmatrix}\ddot\theta_1\\ \ddot\theta_2\end{pmatrix}
+\dot J\begin{pmatrix}\dot\theta_1\\ \dot\theta_2\end{pmatrix},
\label{eq:robot-acceleration}
\end{equation}
其中 $\dot J$ 由 $J$ 对时间逐项求导得到。加速度不仅包含关节加速度经 $J$ 的贡献，还包含速度二次项（离心与科里奥利型项），这正是机器人动力学中非线性项的来源。

雅可比行列式为
\begin{equation}
\det J=l_1l_2\sin\theta_2 .
\label{eq:robot-singularity}
\end{equation}
当 $\theta_2=0$ 或 $\theta_2=\pi$ 时，$\det J=0$，机械臂处于\textbf{奇异位形}：末端在某个方向上瞬时失去运动能力。$\theta_2=0$ 对应手臂完全伸直，$\theta_2=\pi$ 对应完全折叠。奇异位形在轨迹规划、力控制和强化学习中都需要特别处理。

反过来，给定末端位置 $(x,y)$，求关节角 $(\theta_1,\theta_2)$ 称为\textbf{逆运动学}。由余弦定理
\[
\cos\theta_2=\frac{x^2+y^2-l_1^2-l_2^2}{2l_1l_2},
\]
一般有两组解（“肘上”和“肘下”），对应机器人的不同构型。
\end{example}

\begin{modernbox}
关节坐标与末端坐标间的雅可比是局部映射。由本例继续学习机器人动力学，需要多元微积分以及第三、五章的约束与刚体知识。注意一个关节角的梯度给出一个输出的偏导；计算完整多输出雅可比一般需要组织多次求导，不能把它简单等同于一次标量梯度计算。进一步阅读可先完成本章运动学题，再学习第三章的广义坐标。
\end{modernbox}

\section{质点运动的基本规律与微分方程}

前面从已知位置求出了加速度。现在把这条路倒过来走，从受力找加速度，再用出发时的位置和速度求之后的运动。

牛顿定律把“运动”与“相互作用”联系起来，是质点动力学的公理基础。有了受力分析，牛顿第二定律就给出关于位置或速度的微分方程；求解这个微分方程并代入初始条件，就得到运动规律。本节把牛顿三定律、惯性系概念和各种典型的运动微分方程放在同一条逻辑线上，重点是学会“建方程—定初始条件—解方程—作物理讨论”的完整流程。

\subsection{牛顿三定律与惯性系}

比较两种观察方式。放手后的滑块在近似静止的地面实验室中可保持匀速，但在突然制动的车厢中，看起来会向前加速。写牛顿方程之前，先判断观察用的参考系。

牛顿力学以三个基本定律为基础。第一定律指出惯性参考系的存在；第二定律给出力与加速度的定量关系；第三定律说明物体间相互作用的对称性。

\begin{enumerate}[label=(\arabic*),leftmargin=2.2em]
  \item \textbf{惯性定律}：任何物体在不受外力作用时，保持静止或匀速直线运动状态。它定义了惯性参考系，也说明“力不是维持运动的原因，而是改变运动状态的原因”。
  \item \textbf{运动定律}：质点的动量对时间的变化率等于作用在质点上的合力，
  \begin{equation}
  \bF=\dv{\bp}{t}=\dv{}{t}(m\bv)=m\ba .
  \label{eq:newton2}
  \end{equation}
  这里已经使用了恒定质量条件。对于火箭这样的开放系统，不能将当前火箭的 $\dd(m\bv)/\dd t$ 直接等同于外力，必须计入跨边界的动量通量。
  \item \textbf{作用与反作用定律}：对接触力、万有引力和静电力等常见情形，两物体间的作用力与反作用力大小相等、方向相反、沿同一直线，分别作用在两个物体上：
  \begin{equation}
  \bF_{12}=-\bF_{21}.
  \end{equation}
  对运动电荷之间的磁力，这一强形式一般不成立，下文的例子会说明其适用边界。
\end{enumerate}

\paragraph{牛顿第三定律的适用边界：沿 $x,y$ 轴运动的两个电荷.}
考虑两个点电荷：$q_1$ 位于原点，沿 $x$ 轴以速度 $v_1$ 运动；$q_2$ 位于 $(0,d,0)$，沿 $y$ 轴以速度 $v_2$ 运动。取这一瞬时，在 \(v\ll c\) 的非相对论近似下计算它们之间的磁力。

\begin{center}
\begin{tikzpicture}[>=Stealth,scale=1.0]
  \draw[->] (-0.5,0)--(3.6,0) node[below]{$x$};
  \draw[->] (0,-0.5)--(0,3.4) node[left]{$y$};
  \fill (0,0) circle (2pt) node[below left]{$q_1$};
  \draw[->,very thick,red!75!black] (0,0)--(1.15,0) node[above]{$\bv_1$};
  \fill (0,2.2) circle (2pt) node[left]{$q_2$};
  \draw[->,very thick,blue!70!black] (0,2.2)--(0,3.15) node[right]{$\bv_2$};
  \draw[dashed] (0,0)--(0,2.2) node[midway,right]{$d$};
  \node at (0.85,2.55) {$\bB_1\odot$};
  \draw[->,very thick,green!45!black] (0,2.2)--(1.2,2.2) node[below]{$\bF_{21}^{(m)}$};
\end{tikzpicture}
\end{center}

$q_1$ 在 $q_2$ 处产生的磁场为
\[
\bB_1=\frac{\mu_0}{4\pi}\frac{q_1\bv_1\times(d\hat{\jj})}{d^3}
=\frac{\mu_0q_1v_1}{4\pi d^2}\hat{\kk},
\]
方向垂直纸面向外。于是 $q_2$ 受到的磁力为
\[
\bF_{21}^{(m)}=q_2v_2\hat{\jj}\times\bB_1
=\frac{\mu_0q_1q_2v_1v_2}{4\pi d^2}\hat{\ii},
\]
沿 $x$ 轴正方向。反过来，$q_2$ 在原点产生的磁场为
\[
\bB_2=\frac{\mu_0}{4\pi}\frac{q_2\bv_2\times(-d\hat{\jj})}{d^3}=0,
\]
因为 $\bv_2\parallel \hat{\jj}$。所以 $q_1$ 在这一瞬时不受磁力：
\[
\bF_{12}^{(m)}=0.
\]
上述只比较低速磁场贡献；不能在保留 $v^2/c^2$ 阶磁力的同时，把同阶电场修正和迟滞效应全部忽略，借此求精确的总受力。磁力部分本身已显示
\[
\bF_{21}^{(m)}\neq -\bF_{12}^{(m)}.
\]
因此，牛顿第三定律的强形式对运动电荷之间的磁力不成立。

粒子的机械动量可以与场交换。电磁场的动量密度为 $\epsilon_0\bE\times\bB$；孤立系统的总动量须计入场，若边界有动量通量，则应写出相应的收支。在本例的低速两电荷模型中，磁相互作用相对静电尺度带有 $v_1v_2/c^2$ 因子。这个估计不适用于独立给定的外磁场。第\ref{ch:relativity}章再讨论完整系统的能动量。

\begin{physicsbox}
牛顿第二定律是矢量方程，等价于三个分量方程。解题时先选参考系、再画受力图、最后写分量方程。切忌把“合力”与“某个力”混为一谈，也切忌在非惯性系中直接套用 $\bF=m\ba$ 而不加惯性力。
\end{physicsbox}

牛顿定律的成立依赖于参考系的选择。若在一个参考系中，不受力的质点保持静止或匀速直线运动，则该参考系称为\textbf{惯性参考系}。相对于某一惯性系作匀速直线平动的参考系仍是惯性系；作加速运动或转动的参考系则是\textbf{非惯性系}。

伽利略相对性原理断言：一切惯性系对力学规律等价。数学上，这表现为牛顿方程在伽利略变换下形式不变。因此，我们不能通过力学实验判断自己是“静止”还是“匀速运动”；只有相对于其他物体的运动才有意义。

\begin{modernbox}
“惯性系”并不是一个绝对的概念。在广义相对论中，引力与惯性力在局部不可区分，惯性系只在局部时空邻域内定义。经典力学把惯性系当作全局存在的背景，这是宏观低速条件下的理想化。
\end{modernbox}


\subsection{自由质点与约束质点}

珠子套在铁丝上时，铁丝会推着它改变方向。若先只求珠子沿铁丝怎样走，可把方程投影到切向。若还要知道铁丝推了多大力，再用法向方程补回来。

若质点不受约束，只受已知外力 $\bF(\br,\dot{\br},t)$，则运动由二阶矢量微分方程
\begin{equation}
m\ddot{\br}=\bF(\br,\dot{\br},t)
\label{eq:free-particle}
\end{equation}
描述。若 $\bF$ 对状态变量局部满足利普希茨条件，例如在所考察区域连续可微，给定初始位置 $\br_0$ 和初始速度 $\bv_0$ 可局部唯一确定运动；奇异力心处不能直接引用这一结论。直角坐标下为三个耦合的标量方程：
\begin{equation}
m\ddot x=F_x(x,y,z,\dot x,\dot y,\dot z,t),\quad
m\ddot y=F_y(\cdots),\quad
m\ddot z=F_z(\cdots).
\end{equation}

下面考虑受约束的情形。若质点被限制在曲线或曲面上运动，约束会对质点施加\textbf{约束力} $\bR$。此时
\begin{equation}
m\ddot{\br}=\bF+\bR .
\label{eq:constrained}
\end{equation}
约束力通常未知，需要与约束方程联立求解。若约束是光滑的，约束力垂直于约束面（或曲线）；若存在摩擦，还需补充摩擦定律。

在自然坐标系中，式 \eqref{eq:constrained} 分解为
\begin{equation}
m\dot v=F_t+R_t,\qquad
m\frac{v^2}{\rho}=F_n+R_n,\qquad
0=F_b+R_b .
\label{eq:natural-eom}
\end{equation}
此处自然标架公式针对固定曲线。光滑固定曲线上 $R_t=0$；若沿 $s$ 正向滑动，库仑动摩擦模型给出 $R_t=-\mu_k R_N$，其中 $R_N=\sqrt{R_n^2+R_b^2}$。静止接触时用 $|R_t|\le\mu_s R_N$，方向由平衡需求决定。移动曲面上的光滑反力不一定与实际速度垂直，因而也可能做功。

\begin{physicsbox}
约束力不是“外加的已知力”，而是约束条件对运动的响应。求解约束力时，先用运动方程求出法向加速度，再由式 \eqref{eq:natural-eom} 反推 $R_n$。这也是判断约束是否“离开”的依据：对于只能推、不能拉的单侧接触，若沿允许支持方向求得负反力，接触解就失效，应改为脱离运动；双侧导轨可提供两侧反力，不能仅凭某个分量为负就判定脱离。
\end{physicsbox}

\subsection{线性微分方程与受迫振动}

推秋千之后放手，和继续按节拍推，是两种不同的题。前者只由初值决定，后者还受外部推力控制。线性方程允许把这两种运动叠加起来。

在平衡点附近，许多系统可线性化为
\begin{equation}
\ddot x+2\gamma\dot x+\omega_0^2x=0,
\label{eq:damped}
\end{equation}
其中 $\gamma$ 为阻尼系数，$\omega_0$ 为固有频率。试探解 $x=e^{\lambda t}$ 给出特征方程
\begin{equation}
\lambda^2+2\gamma\lambda+\omega_0^2=0,
\end{equation}
其根为 $\lambda=-\gamma\pm\sqrt{\gamma^2-\omega_0^2}$。于是有三类解：
\begin{align}
\gamma<\omega_0 &: \quad x=e^{-\gamma t}\left(A\cos\omega t+B\sin\omega t\right),\quad \omega=\sqrt{\omega_0^2-\gamma^2},\\
\gamma=\omega_0 &: \quad x=e^{-\gamma t}\left(A+Bt\right),\\
\gamma>\omega_0 &: \quad x=Ae^{-(\gamma-\sqrt{\gamma^2-\omega_0^2})t}
+B e^{-(\gamma+\sqrt{\gamma^2-\omega_0^2})t}.
\end{align}
分别称为欠阻尼、临界阻尼和过阻尼。

若外加周期驱动力 $F_0\cos\omega t$，则
\begin{equation}
\ddot x+2\gamma\dot x+\omega_0^2x=\frac{F_0}{m}\cos\omega t .
\end{equation}
稳态解为 $x=A\cos(\omega t-\delta)$，其中
\begin{equation}
A=\frac{F_0/m}{\sqrt{(\omega_0^2-\omega^2)^2+4\gamma^2\omega^2}},\qquad
\tan\delta=\frac{2\gamma\omega}{\omega_0^2-\omega^2}.
\label{eq:resonance}
\end{equation}
当 $0<\gamma<\omega_0/\sqrt2$ 时，位移振幅的峰值频率是 $\omega_{\rm peak}=\sqrt{\omega_0^2-2\gamma^2}$，只有弱阻尼下才近似为 $\omega_0$。相位须取连续分支 $\delta=\operatorname{atan2}(2\gamma\omega,\omega_0^2-\omega^2)\in[0,\pi]$。稳态表达式不适用于无阻尼且恰好共振的情况；那时特解含随时间增长的项。详细的振动分析留到第四章。

\begin{figure}[htbp]
\centering
\begin{tikzpicture}
\begin{axis}[
  width=10.5cm,height=5.6cm,
  axis lines=middle,
  xlabel={$\omega/\omega_0$},ylabel={$A/(F_0/m\omega_0^2)$},
  xmin=0,xmax=3,ymin=0,ymax=11,
  xtick={0,0.5,1,1.5,2,2.5,3},
  ytick={0,2,4,6,8,10},
  legend style={at={(0.97,0.97)},anchor=north east,font=\small},
  samples=300,domain=0.05:3
]
\addplot[very thick,blue!70!black] {1/sqrt((1-x^2)^2+4*0.05^2*x^2)};
\addplot[very thick,red!75!black] {1/sqrt((1-x^2)^2+4*0.15^2*x^2)};
\addplot[very thick,green!45!black] {1/sqrt((1-x^2)^2+4*0.3^2*x^2)};
\legend{$\gamma=0.05\omega_0$,$\gamma=0.15\omega_0$,$\gamma=0.30\omega_0$}
\end{axis}
\end{tikzpicture}
\caption{受迫振动的幅频曲线。阻尼越小，共振峰越高、越窄；当 $\omega=\omega_0$ 时相位差为 $\pi/2$。}
\label{fig:resonance}
\end{figure}

\subsection{典型示例}

下面的模型各有不同的受力。先选坐标和正方向，再判断哪些力控制运动。若要从方程求出随时间变化的轨迹，还要写下出发时的位置和速度。

\begin{example}[空间站微重力实验中的受迫振动与三维谐振子]
(1) 在空间站微重力实验平台上，带电粒子（电荷 $-e$）在振荡电场 $E_x=E_0\cos(\omega t+\theta)$ 中运动，$m\ddot x=-eE_x$。积分两次得
\[
x=x_0-\frac{eE_0}{m\omega^2}\cos\theta
+\left(v_0+\frac{eE_0}{m\omega}\sin\theta\right)t
+\frac{eE_0}{m\omega^2}\cos(\omega t+\theta),
\]
即匀速漂移与同频受迫振荡的叠加。

(2) 作为给定外力的教学模型，设三维恢复力可近似为三个独立谐振子（本例不声称仅凭静电场就能形成三维稳定阱）：
\[
m\ddot x_i=-k_i x_i,\qquad x_i=A_i\cos(\omega_i t+\phi_i),\qquad \omega_i=\sqrt{k_i/m}.
\]
对所有被激发的方向，若频率存在共同周期，轨道闭合；若三个频率有理独立，相位轨道在三维相位环面上稠密。仅有某一频率比无理并不足以推出三维环面上的稠密性；位置空间轨迹则是该相位轨道的投影。
\end{example}

\begin{example}[火星探测器进入段的阻力下落模型]
质量为 $m$ 的火星探测器在降落伞展开后的简化模型中，从高度 $h$ 处近似无初速下落，阻力与速度成正比，$\bR=-mk\bv$。取竖直向上为 $x$ 轴正方向，初始位置为 $x(0)=h$，则重力为 $-mg$，阻力为 $-mk\dot x$，运动方程为
\[
m\ddot x=-mg-mk\dot x,
\]
即
\[
\ddot x=-g-k\dot x.
\]
分离变量并积分，利用 $\dot x(0)=0$，
\[
\ln\frac{g+k\dot x}{g}=-kt,
\]
故
\[
\dot x=-\frac{g}{k}\left(1-e^{-kt}\right).
\]
再积分，利用 $x(0)=h$，得
\[
x=h+\frac{g}{k^2}\left(1-e^{-kt}\right)-\frac{g}{k}t.
\]
在形式上的长时间延拓中，$\dot x\to -g/k$，探测器趋于匀速下落，称为\textbf{收尾速度}。若取竖直向下为 $x$ 轴，则速度与坐标的符号整体反转。本模型仅在触地之前使用。线性阻力适用于相应的低雷诺数或经验证的线性范围；实际降落伞及进入段常有速度平方阻力，不能据本例定量预测真实降落过程。

取 $k>0$，将方程写为 $\dd(\dot x+g/k)/\dd t=-k(\dot x+g/k)$。推广到初速 $\dot x(0)=v_0$ 时可得 $\dot x+g/k=(v_0+g/k)e^{-kt}$，再积分一次得到位移。
\end{example}

\begin{example}[月球摆的大振幅周期]
在月球表面，长为 $l$ 的单摆，摆角 $\theta$ 满足
\[
\ddot\theta+\frac{g}{l}\sin\theta=0.
\]
两边乘以 $\dot\theta$ 并积分，得能量积分
\[
\frac12\dot\theta^2-\frac{g}{l}\cos\theta
=\frac12\dot\theta_0^2-\frac{g}{l}\cos\theta_0.
\]
若从最大摆角 $\theta_0$ 处由静止释放，则
\[
\dot\theta^2=\frac{2g}{l}\left(\cos\theta-\cos\theta_0\right).
\]
于是周期为
\[
T=4\sqrt{\frac{l}{g}}\int_0^{\theta_0}
\frac{\dd\theta}{\sqrt{2(\cos\theta-\cos\theta_0)}}.
\]
作 $\sin(\theta/2)=\sin(\theta_0/2)\sin\phi$ 代换，可将它化成第一类完全椭圆积分；章末给出代换细节。小角近似 $\sin\theta\approx\theta$ 给出简谐振动，$T_0=2\pi\sqrt{l/g}$；随着振幅增大，周期变长。按月球平均参数的近似 $g_M\approx1.62\,\mathrm{m/s^2}$，同样长度的摆周期约为地面的 $\sqrt{g/g_M}\approx2.46$ 倍。

\begin{figure}[htbp]
\centering
\begin{tikzpicture}[>=Stealth]
\coordinate (O) at (0,3);\coordinate (B) at (1.2,.921539);
\draw[thick] (-1,3)--(1,3);\draw[thick] (O)--(B);
\draw[dashed] (O)--(0,.5);\fill (B) circle (3pt) node[right]{$m$};
\draw (0,2.1) arc (-90:-60:.9);\node at (.3,1.9) {$\theta$};
\draw[->,thick,red!70!black] (B)--++(0,-.9) node[below]{$m\bg$};
\draw[->,thick,blue!70!black] (B)--++(-.5,.866025) node[left]{$\bT$};
\draw[decorate,decoration={brace,amplitude=4pt}] (.2,3.11547)--(1.4,1.037009) node[midway,right=5pt]{$l$};
\end{tikzpicture}
\caption{固定摆长的单摆：拉力沿绳并垂直瞬时速度，切向重力为 $-mg\sin\theta$。长度标记沿摆线段绘制。}
\label{fig:pendulum}
\end{figure}
\end{example}



\begin{example}[霍尔推进器中的电子]
在霍尔推进器中，电子在互相垂直的均匀电场 $\bE$ 与磁场 $\bB$ 中运动。取 $\bE=E\ii$、$\bB=B\kk$，初速度 $\bv(0)=v_0\jj$，初始位置为原点。求电子的运动规律，并说明漂移速度的物理意义。

\begin{center}
\begin{tikzpicture}[>=Stealth,scale=0.95]
  \draw[->] (0,0)--(4.2,0) node[below]{$x$};
  \draw[->] (0,0)--(0,2.8) node[left]{$y$};
  \draw[->,very thick,red!75!black] (0.6,2.2)--(1.7,2.2) node[above]{$\bE$};
  \node at (0.9,1.25) {$\bB\odot$};
  \draw[->,very thick,green!45!black] (2.8,1.2)--(2.8,2.2) node[right]{$\bv_0$};
  \draw[->,very thick,blue!70!black] (3.5,1.2)--(3.5,0.2) node[right]{$\bv_d=\dfrac{\bE\times\bB}{B^2}$};
\end{tikzpicture}
\end{center}

解：电子电荷为 $-e$，运动方程为
\[
m\dot{\bv}=-e(\bE+\bv\times\bB).
\]
写成分量：
\[
m\dot v_x=-eE-eBv_y,\qquad
m\dot v_y=eBv_x,\qquad
m\dot v_z=0.
\]
令回旋频率 $\omega_c=eB/m$，并定义复速度 $w=v_x+\mathrm{i}v_y$，则
\[
\dot w=-\frac{eE}{m}+\mathrm{i}\omega_c w .
\]
其特解给出漂移速度 $v_x=0$、$v_y=-E/B$，即
\[
\bv_d=\frac{\bE\times\bB}{B^2}.
\]
以初始条件 $w(0)=\mathrm{i}v_0$ 积分，得
\[
v_x(t)=-\left(v_0+\frac{E}{B}\right)\sin\omega_ct,\qquad
v_y(t)=\left(v_0+\frac{E}{B}\right)\cos\omega_ct-\frac{E}{B}.
\]
再积分一次：
\[
x(t)=\frac{v_0+E/B}{\omega_c}(\cos\omega_ct-1),\qquad
y(t)=\frac{v_0+E/B}{\omega_c}\sin\omega_ct-\frac{E}{B}t .
\]
因此电子的运动是“回旋”与“匀速漂移”的叠加：在垂直于 $\bE$ 与 $\bB$ 的方向上，它以速度 $E/B$ 漂移；在以 $\bv_d$ 平动的低速参考系中，它作匀速圆周运动。本模型说明交叉场漂移的单粒子机制。真实霍尔推进器还涉及电离、碰撞及电子输运，离子束电荷中和主要由外置阴极提供电子，不能从这条漂移轨迹直接推出。进一步资料见本章末 NASA 技术介绍。
\end{example}

\section{功、能与势能}

若只想知道滑到底有多快，可以先问各个力在途中做了多少功。这个问题有时比逐时刻求加速度更容易回答。

功和能提供了描述运动的另一套语言：力沿路径的累积效应可以用一个标量——功——来度量，保守力又可以进一步用势能表示。对许多问题，能量守恒比直接求解矢量微分方程更简洁，也更容易揭示运动的全局特征。本节从功的定义出发，引入势能、机械能守恒和势能曲线，为后续的守恒律与中心力问题作准备。

\subsection{功与功率}

拉行李箱时，沿拉力方向移动的那一部分位移决定做功。用同样的力做同样多的功，花的时间越短，平均功率越大。

力 $\bF$ 在微小位移 $\dd\br$ 上做的功为
\begin{equation}
\dd W=\bF\cdot\dd\br=F\cos\alpha\,\dd s,
\end{equation}
其中 $\alpha$ 是力与位移的夹角。这里的 $\dd W$ 表示过程中的元功，一般不是某个仅依赖状态的函数的全微分。沿路径 $L$ 从 $A$ 到 $B$ 的总功为
\begin{equation}
W_{A\to B}=\int_L \bF\cdot\dd\br .
\label{eq:work}
\end{equation}
功率为
\begin{equation}
P=\dv{W}{t}=\bF\cdot\bv .
\end{equation}

\begin{figure}[htbp]
\centering
\begin{tikzpicture}[>=Stealth]
\draw[thick,blue!65!black,domain=-1:2,samples=80] plot (\x,{0.25*\x*\x});
\fill (0,0) circle (2pt) node[below left]{$P$};
\draw[->,very thick,green!40!black] (0,0)--(1.7,0) node[below]{$\dd\br$};
\draw[->,very thick,red!70!black] (0,0)--(1.3,1.3) node[above]{$\bF$};
\draw (0.65,0) arc (0:45:0.65);\node at (0.9,0.3) {$\alpha$};
\end{tikzpicture}
\caption{功只取决于力在瞬时位移上的投影。绿色箭头沿 $P$ 处切线，红色力矢量与它夹角为 $\alpha$。}
\label{fig:work}
\end{figure}

\subsection{保守力与势能}

把书从桌上抬到同一高度，若只算重力的功，走哪条路线结果都相同。这种只由起点和终点决定的功，可以用势能差来记账。

若力沿任意闭合路径做的功为零，
\begin{equation}
\oint_L \bF\cdot\dd\br=0,
\label{eq:conservative-loop}
\end{equation}
则称该力为\textbf{保守力}。等价地，保守力沿 $A\to B$ 的功只与始末位置有关，与路径无关。于是可以定义\textbf{势能} $V(\br)$：
\begin{equation}
V(\br)-V(\br_0)=-\int_{\br_0}^{\br}\bF\cdot\dd\br',
\label{eq:potential}
\end{equation}
或微分形式
\begin{equation}
\bF=-\nabla V .
\label{eq:force-potential}
\end{equation}
在直角坐标中，
\begin{equation}
F_x=-\pdv{V}{x},\qquad F_y=-\pdv{V}{y},\qquad F_z=-\pdv{V}{z}.
\end{equation}
对定义在连通区域内的光滑位置力场，闭路功为零等价于全局单值势能存在。若势能存在，则旋度必为零；反向推论还须区域单连通等拓扑条件。在这样的区域内才可用判据
\begin{equation}
\nabla\times\bF=0
\quad\Longleftrightarrow\quad
\pdv{F_z}{y}=\pdv{F_y}{z},\quad
\pdv{F_x}{z}=\pdv{F_z}{x},\quad
\pdv{F_y}{x}=\pdv{F_x}{y}.
\label{eq:curl-zero}
\end{equation}

\begin{physicsbox}
重力、静电力、弹簧弹力都是保守力；摩擦力、空气阻力一般不是保守力。当势能不显含时间且其余力不做功时，保守力功可由势能变化描述，因此机械能守恒；非保守力的功则把机械能转化为内能等其他形式。
\end{physicsbox}

\subsection{机械能守恒与势能曲线}

把总能量画成水平线，与势能曲线比较。水平线高于势能的区间才可能到达，交点处动能为零；是否真会在交点折返，还要看该处的受力。

质点动能为
\begin{equation}
T=\frac12m\bv^2=\frac12m\dot{\br}^2.
\end{equation}
由牛顿第二定律，
\begin{equation}
\dd T=\bF\cdot\dd\br=\dd W .
\end{equation}
若只有保守力做功，则 $\dd W=-\dd V$，于是
\begin{equation}
\dd(T+V)=0\quad\Longrightarrow\quad T+V=E=\text{常量}.
\label{eq:energy-conservation}
\end{equation}
这就是\textbf{机械能守恒定律}。
更一般地，若 $\bF=-\nabla V(\br,t)+\bF_{\rm nc}$，则
\[
\dot T=-\nabla V\cdot\bv+\bF_{\rm nc}\cdot\bv,\qquad
\dot V=\nabla V\cdot\bv+\pdv{V}{t},
\]
\begin{equation}\label{eq:ch1-energy-general}
\boxed{\frac{\dd}{\dd t}(T+V)=\pdv{V}{t}+\bF_{\rm nc}\cdot\bv.}
\end{equation}
这明确了守恒的条件：时间变化的势或做功的驱动约束都可向质点传递能量。磁力可以不保守却不做功，不能把所有非保守力都等同于耗散力。
\begin{example}[无旋不一定全局保守]
在去掉 $z$ 轴的空间中，取 $\bF=K(-y,x,0)/(x^2+y^2)$，其中 $K$ 的量纲为能量。直接求偏导得到 $\nabla\times\bF=0$。但在半径 $a$ 的圆上，$\br=(a\cos\phi,a\sin\phi,0)$，有 $\bF\cdot\dd\br=K\dd\phi$，故闭路功为 $2\pi K\ne0$。局部可写 $V=-K\phi$，绕轴一周却不单值。这说明使用旋度判据前必须先说明力场定义域。
\end{example}

一维运动的势能曲线 $V(x)$ 提供了运动的定性图像（图中用固定能量与长度单位无量纲化）：由 $E=T+V(x)\ge V(x)$ 可知，运动只能发生在 $V(x)\le E$ 的区间；$V(x)=E$ 的普通单根对应\textbf{转折点}（速度为零并反向）；重根处还可能是平衡或需无限时间才能到达的渐近位置，应单独判断；$V'(x)=0$ 的点为平衡点，$V''(x)>0$ 为稳定平衡，$V''(x)<0$ 为不稳定平衡。若二阶导数为零，须检查更高阶，不能由二阶检验直接定论。

\begin{figure}[htbp]
\centering
\begin{tikzpicture}
\begin{axis}[width=0.82\textwidth,height=6cm,axis lines=left,
 xlabel={$X$},ylabel={$U$},xmin=-1.6,xmax=1.6,ymin=-.85,ymax=1.2,
 xtick={-1.376382,-.324920,.324920,1.376382},xticklabels={$X_1$,$X_2$,$X_3$,$X_4$}]
\addplot[very thick,blue!70!black,domain=-1.55:1.55,samples=200] {x^4-2*x^2+.4};
\addplot[dashed,red!70!black,domain=-1.55:1.55] {.2};
\addplot[only marks,mark=*,blue!70!black] coordinates {(-1,-.6) (0,.4) (1,-.6)};
\pgfplotsinvokeforeach{-1.376382,-.324920,.324920,1.376382}{\draw[dashed,gray] (axis cs:#1,-.85)--(axis cs:#1,.2);}
\draw[<->,thick,green!40!black] (axis cs:-1.376382,-.73)--(axis cs:-.324920,-.73);
\draw[<->,thick,green!40!black] (axis cs:.324920,-.73)--(axis cs:1.376382,-.73);
\node[red!70!black] at (axis cs:1.15,.35) {$\mathcal E=0.2$};
\end{axis}
\end{tikzpicture}
\caption{无量纲双阱势 $U=X^4-2X^2+0.4$。在能量 $\mathcal E=0.2$ 下，两段绿色区间允许运动，中央势垒不能跨过；四个端点由 $X^2=1\pm\sqrt{0.8}$ 精确确定。}
\label{fig:potential-curve}
\end{figure}

\begin{example}[保守力场的势能]
以下坐标和力按固定长度、力的单位给出数值，势能按相应的力乘长度单位表示。判断 $\bF=(x+2y+z+5)\ii+(2x+y+z)\jj+(x+y+z-6)\kk$ 是否保守，并求势能。

解：三个交叉偏导分别满足 $\partial_yF_x=\partial_xF_y=2$、$\partial_zF_x=\partial_xF_z=1$ 和 $\partial_zF_y=\partial_yF_z=1$。因此在整个 $\mathbb R^3$ 上旋度为零，存在全局势能。先积分 $V_x=-F_x$，得
\[
V=-x^2/2-2xy-xz-5x+g(y,z).
\]
再由 $V_y=-F_y$ 得 $g_y=-y-z$，所以 $g=-y^2/2-yz+h(z)$。最后由 $V_z=-F_z$ 得 $h'=-z+6$。合并得
\[
V=-\left[\frac{x^2}{2}+2xy+xz+5x+\frac{y^2}{2}+yz+\frac{z^2}{2}-6z\right]+C.
\]
取 $V(0,0,0)=0$，则 $C=0$。
\end{example}

\begin{example}[月球车沿抛物线月坑下滑]
月球车（视为质点）质量为 $m$，沿光滑抛物线形月坑 $x^2=4ay$ 从 $x=2a$ 处由静止下滑。求到达坑底时的速度和月面支持力。

解：由 $y=x^2/(4a)$ 得 $y'=x/(2a)$，$y''=1/(2a)$。在 $x=2a$ 处，$y=a$，$y'=1$，切线与水平方向成 $45^\circ$。由机械能守恒，
\[
\frac12mv^2=mg(a-y),
\]
故到达坑底 $y=0$ 时
\[
v=\sqrt{2ga}.
\]
坑底处曲率半径为
\[
\rho=\frac{[1+(y')^2]^{3/2}}{|y''|}\bigg|_{x=0}
=\frac{1}{1/(2a)}=2a.
\]
在坑底，法向（指向曲率中心，即竖直向上）方程为
\[
m\frac{v^2}{\rho}=R-mg.
\]
代入 $v^2=2ga$、$\rho=2a$，得
\[
R=mg+\frac{m(2ga)}{2a}=2mg.
\]
\end{example}
\begin{example}[月球车翻越球形月丘]
月球车（视为质点）质量为 $m$，从半径为 $R$ 的光滑球形月丘顶部附近、极小角位移处由静止释放，并取该位移趋于零的下滑极限。求脱离月丘表面时的位置和速率。

\begin{center}
\begin{tikzpicture}[>=Stealth,scale=1.0]
  \draw[thick] (0,0) circle (1.6);
  \draw[dashed] (0,0)--(0,1.6);
  \coordinate (P) at (1.226,1.028);
  \fill (P) circle (2.2pt) node[above left]{$m$};
  \draw[thick] (0,0)--(P) node[midway,below]{$R$};
  \draw (0,0.7) arc (90:40:0.7);
  \node at (0.22,1.20) {$\theta$};
  \draw[->,thick,red!75!black] (P)--++(0,-0.72) node[below]{$m\bg$};
  \draw[->,thick,blue!70!black] (P)--++(0.52,0.44) node[above right]{$\bN$};
  \draw[->,thick,green!45!black] (P)--++(0.48,-0.40) node[below right]{$\bv$};
\end{tikzpicture}
\end{center}

解：设月球车与竖直方向的夹角为 $\theta$。由机械能守恒，
\[
mgR(1-\cos\theta)=\frac12mv^2 .
\]
沿径向（指向球心）应用牛顿第二定律：
\[
mg\cos\theta-N=m\frac{v^2}{R}.
\]
代入 $v^2=2gR(1-\cos\theta)$，得
\[
N=mg(3\cos\theta-2).
\]
脱离月丘表面的条件是 $N=0$，故
\[
\cos\theta_0=\frac23,\qquad
\theta_0\approx48.2^\circ .
\]
此时速率为
\[
v_0=\sqrt{2gR(1-\cos\theta_0)}
=\sqrt{\frac{2gR}{3}} .
\]
若恰在最高点静止且毫无扰动，它将保持不稳定平衡；上面的下滑结果是微扰极限。在顶部极限 $N=mg$；随着下滑，支持力单调减小，在“赤道”以上就已降为零。此后月球车脱离月丘表面，在月面引力下作弹道运动。
\end{example}


\begin{example}[矩阵力场的势能条件]
设力场由常矩阵 $A=(a_{ij})$ 给出：
\[
\bF=A\br,\qquad F_i=\sum_{j=1}^{3}a_{ij}x_j .
\]
问 $A$ 满足什么条件时 $\bF$ 保守？并求势能。

解：由 $\nabla\times\bF=0$ 得
\[
(\nabla\times\bF)_x=\partial_yF_z-\partial_zF_y=a_{32}-a_{23}=0,
\]
以及其他两个分量循环对称的条件。因此 $\bF$ 保守的充要条件是
\[
a_{ij}=a_{ji},
\]
即 $A$ 为对称矩阵。此时
\[
V(\br)=-\frac12\br^T A\br
=-\frac12\sum_{i,j=1}^{3}a_{ij}x_ix_j
\]
（相差一个常数），且 $-\nabla V=A\br=\bF$。若 $A$ 不对称，只有其对称部分对势能有贡献，反对称部分对应一个不能由势能导出的涡旋力。

把 $A$ 分为对称部分 $S$ 和反对称部分 $W$，有 $\br^{\mathsf T}W\br=0$，因为此标量取转置后反号。逐分量微分给 $\nabla(\br^{\mathsf T}A\br)=(A+A^{\mathsf T})\br$，所以该二次型只能产生 $S\br$。
\end{example}


\section{质点动力学基本定理与守恒定律}

动量、角动量和动能回答不同问题。分别检查外力、绕选定点的外力矩和各力做功，便能知道哪一种量的变化最好算。

\subsection{动量定理与动量守恒}

碰撞力在短时间内可以很大，而且难以逐时刻测定。把力随时间加起来，得到的冲量直接告诉我们动量改变了多少。

由牛顿第二定律 $\bF=\dd\bp/\dd t$，在时间区间 $[t_1,t_2]$ 上积分得
\begin{equation}
\bI=\int_{t_1}^{t_2}\bF\,\dd t=\bp_2-\bp_1,
\label{eq:momentum-theorem}
\end{equation}
其中 $\bI$ 称为\textbf{冲量}。若合外力为零，则
\begin{equation}
\bp=m\bv=\text{常量}.
\label{eq:momentum-conservation}
\end{equation}
这就是\textbf{动量守恒定律}。

\subsection{角动量定理与角动量守恒}

推门时，力的大小相同，推在门把手和推在铰链边的效果不同。选定矩心后，力的大小、方向和力臂共同决定力矩。

质点相对于原点 $O$ 的角动量为
\begin{equation}
\bJ=\br\times\bp=m\,\br\times\bv .
\end{equation}
力矩为
\begin{equation}
\bM=\br\times\bF .
\end{equation}
由
\begin{equation}
\dv{\bJ}{t}=\dv{}{t}(\br\times m\bv)
=\br\times m\ddot{\br}=\br\times\bF=\bM,
\label{eq:angular-momentum-theorem}
\end{equation}
得\textbf{角动量定理}：
\begin{equation}
\dv{\bJ}{t}=\bM .
\end{equation}
若 $\bM=0$，则 $\bJ=\text{常量}$。这里 $O$ 是惯性系中的固定点。若改用移动点 $A(t)$，且角动量仍用绝对动量定义为 $\bJ_A=(\br-\br_A)\times m\bv$，则乘积法则给出
\[
\dot{\bJ}_A=(\bv-\bv_A)\times m\bv+(\br-\br_A)\times\bF
=\bM_A-\bv_A\times\bp.
\]
因此不能不加条件地把固定点角动量定理套到移动参考点。第五章会说明为何质心是一个特别方便的参考点。特别地，若力始终通过某一固定点（中心力），则对该点的角动量守恒。

\begin{physicsbox}
角动量守恒的核心是“力矩为零”，而不是“力为零”。中心力 $\bF=F(r)\er$ 满足 $\br\times\bF=0$，所以角动量守恒；但这并不意味着质点做匀速运动，也不意味着能量一定守恒。
\end{physicsbox}

\subsection{动能定理与守恒律的对称性}

若关心物体快了还是慢了，沿实际位移算功通常最直接。至于动量方向和轨道形状，还要用其他方程或守恒量确定。

由
\begin{equation}
\dd T=m\bv\cdot\dd\bv=\bF\cdot\dd\br=\dd W,
\end{equation}
得\textbf{动能定理}：
\begin{equation}
W_{A\to B}=T_B-T_A=\frac12mv_B^2-\frac12mv_A^2 .
\label{eq:work-energy}
\end{equation}
若只有保守力做功，则 $W=-\Delta V$，从而 $T+V=E$。

进一步看，三个守恒律分别对应时空对称性：空间平移不变 $\Rightarrow$ 动量守恒；空间转动不变 $\Rightarrow$ 角动量守恒；时间平移不变 $\Rightarrow$ 能量守恒。这一对应关系在分析力学中由诺特定理精确表述，是“对称性决定相互作用”的思想源头。

\begin{example}[中心力作用下的平面运动]
质点在中心力 $\bF=F(r)\er$ 作用下运动。由于 $\bM=\br\times\bF=0$，角动量
\[
\bJ=\br\times m\bv
\]
守恒。若 $\bJ\ne0$，因 $\br\cdot\bJ=0$，位矢始终垂直同一个常矢量，轨道位于过力心的固定平面内。若 $\bJ=0$，运动为沿力心直线的径向运动，需单独处理；下式的角度参数化假设 $J\ne0$。

在极坐标中，$\bJ=mr^2\dot\theta\,\kk$，故
\[
mr^2\dot\theta=J=\text{常量}.
\]
扫过的面积速度为
\[
\dv{A}{t}=\frac12r^2\dot\theta=\frac{J}{2m}=\text{常量},
\]
这正是开普勒第二定律的力学基础。
\end{example}

\section{质点系与质心运动}

观察一辆两人推着走的车，可先把人和车合起来研究，再看各部分相对整体怎样运动。系统边界一旦选定，才能分清哪一些力是内力。

本节从每个质点的牛顿第二定律出发，逐步导出质点系的动量、角动量和动能定理。设系统由 $N$ 个质点组成，第 $i$ 个质点的质量为 $m_i$、位矢为 $\br_i$。对每个质点写出牛顿第二定律：
\begin{equation}
m_i\ddot{\br}_i=\bF_i^{(e)}+\sum_{j\ne i}\bF_{ij},
\qquad i=1,\dots,N,
\label{eq:newton-system}
\end{equation}
其中 $\bF_i^{(e)}$ 是外力，$\bF_{ij}$ 是第 $j$ 个质点对第 $i$ 个质点的内力。由牛顿第三定律
\begin{equation}
\bF_{ij}=-\bF_{ji}.
\label{eq:third-law-system}
\end{equation}
以下结论由式 \eqref{eq:newton-system} 和 \eqref{eq:third-law-system} 推出；在讨论内力矩时，还需要假设内力满足强形式。

\subsection{内力、外力与质心}

质心是按质量加权的位置平均。两个人坐跷跷板时，质量较大的一侧对质心位置影响更大，所以平衡时它通常要坐得更靠近支点。

先证明内力之和为零。把内力按 $i<j$ 成对求和：
\begin{equation}
\sum_{i=1}^{N}\sum_{j\ne i}\bF_{ij}
=\sum_{i<j}\left(\bF_{ij}+\bF_{ji}\right)=0.
\label{eq:internal-force-zero}
\end{equation}
若进一步假设内力沿两质点的连线（强形式），则
\begin{equation}
\bF_{ij}\parallel \br_i-\br_j .
\end{equation}
于是每一对内力对固定点 $O$ 的力矩之和为
\begin{equation}
\br_i\times\bF_{ij}+\br_j\times\bF_{ji}
=(\br_i-\br_j)\times\bF_{ij}=0.
\label{eq:internal-torque-zero}
\end{equation}
因此所有内力矩之和也为零。以下各质点质量及系统成员固定；开放系统另行分析。

注意：若内力是运动电荷之间的磁力，强形式一般不成立，必须把电磁场的角动量计入；本节的质点系角动量定理默认内力为中心力。

定义系统的\textbf{质心}
\begin{equation}
\bR=\frac{1}{M}\sum_{i=1}^{N}m_i\br_i,\qquad
M=\sum_{i=1}^{N}m_i .
\label{eq:cm-ch1}
\end{equation}
对连续体，取 $m_i\to\dd m$ 的极限，求和化为积分：
\begin{equation}
\bR=\frac{1}{M}\int\br\,\dd m,\qquad
M=\int\dd m .
\end{equation}
质心不一定在物体上，但它把系统的整体运动浓缩为一个代表点。

\begin{example}[球形贮箱球冠端盖的质心]
航天器中的一个薄壁球形贮箱，被距球心 $b$ 的平面隔板截出一个球冠形端盖（$0<b<a$）。设球面面密度为 $\sigma$，求端盖质心到球心的距离。

\begin{center}
\begin{tikzpicture}[>=Stealth,scale=1.0]
  \fill[blue!8] (0,0) circle (1.5);
  \begin{scope}
    \clip (0,0) circle (1.5);
    \fill[blue!22] (-1.6,0.6) rectangle (1.6,1.6);
  \end{scope}
  \draw[thick,blue!50!black] (0,0) circle (1.5);
  \draw[thick,blue!60!black] (-1.5,0.6)--(1.5,0.6);
  \draw[dashed] (0,0)--(1.5,0) node[midway,below]{$a$};
  \draw[->,thick] (0,0)--(0,2.1) node[above]{$z$};
  \fill (0,1.05) circle (2pt) node[right]{$z_c$};
  \node at (1.0,0.35) {$b$};
  \node at (-1.0,1.15) {球冠端盖};
\end{tikzpicture}
\end{center}

解：设球面面密度为 $\sigma$，取球心为原点、对称轴为 $z$ 轴。球面面积元为
\[
\dd S=2\pi a^2\sin\theta\,\dd\theta,\qquad
z=a\cos\theta,
\]
其中 $\theta$ 从 $0$ 到 $\theta_0$，且 $\cos\theta_0=b/a$。球冠的面积为
\[
S=2\pi a^2\int_0^{\theta_0}\sin\theta\,\dd\theta
=2\pi a(a-b),
\]
质量为 $M=\sigma S$。质心的 $z$ 坐标为
\[
z_c=\frac{1}{S}\int z\,\dd S
=\frac{1}{2\pi a(a-b)}
\int_0^{\theta_0}a\cos\theta\,2\pi a^2\sin\theta\,\dd\theta.
\]
计算积分：
\[
\int_0^{\theta_0}\cos\theta\sin\theta\,\dd\theta
=\frac12\sin^2\theta_0
=\frac12\left(1-\frac{b^2}{a^2}\right).
\]
代入得
\[
z_c=\frac{a^2}{a-b}\cdot\frac12\left(1-\frac{b^2}{a^2}\right)
=\frac{a+b}{2}.
\]
本例忽略端盖厚度且面密度均匀。当 $b=0$ 时，球冠退化为半球面，$z_c=a/2$，这是熟悉的半球壳结果。质心位置直接影响贮箱和航天器的质量特性与姿态控制。
\end{example}

对式 \eqref{eq:cm-ch1} 求导两次，并利用式 \eqref{eq:newton-system}：
\begin{equation}
M\ddot{\bR}
=\sum_i m_i\ddot{\br}_i
=\sum_i\bF_i^{(e)}+\sum_i\sum_{j\ne i}\bF_{ij}.
\end{equation}
由式 \eqref{eq:internal-force-zero}，内力和为零，于是
\begin{equation}
M\ddot{\bR}=\sum_i\bF_i^{(e)}=\bF^{(e)}.
\label{eq:cm-theorem-ch1}
\end{equation}
这就是\textbf{质心运动定理}：质心像一个质量为 $M$、受全部外力矢量和作用的质点。烟花爆炸、航天器分离时，只要外力变化不大，系统质心的轨迹就与爆炸前相同；内部碎片如何飞散不影响质心运动。

\subsection{质点系动量定理与动量守恒}

系统内部互推时，某一部分获得的动量变化会由另一部分抵消。整个系统的动量变化由外界给出的合冲量决定。

由质心定义，系统的总动量为
\begin{equation}
\bP=\sum_i m_i\dot{\br}_i
=\frac{\dd}{\dd t}\sum_i m_i\br_i
=M\dot{\bR}.
\label{eq:system-momentum-def}
\end{equation}
对时间求导，并利用质心运动定理 \eqref{eq:cm-theorem-ch1}：
\begin{equation}
\dv{\bP}{t}=M\ddot{\bR}=\bF^{(e)}.
\label{eq:sys-mom-ch1}
\end{equation}
写成冲量形式：
\begin{equation}
\bP(t_2)-\bP(t_1)=\int_{t_1}^{t_2}\bF^{(e)}\,\dd t .
\end{equation}
若外力和为零，则
\begin{equation}
\bP=M\bV_c=\text{常量}.
\end{equation}
这就是质点系的\textbf{动量守恒定律}。内力可以很大，却不会改变系统的总动量；动量守恒是对整个系统而言的，单个质点的动量一般并不守恒。

\begin{example}[火箭方程]
取竖直向上为正，在所考察高度范围内把 $g$ 近似视为常量。设在 \(\dd t\) 内火箭质量变化为 \(\dd m<0\)，喷出气体相对火箭的速度为 \(-v_e\)（\(v_e>0\)）。喷出前系统动量为 \(mv\)；喷出后火箭动量为 \((m+\dd m)(v+\dd v)\)，气体动量为 \((-\dd m)(v-v_e)\)。外力冲量为 \(-mg\,\dd t\)。由动量定理，
\[
(m+\dd m)(v+\dd v)+(-\dd m)(v-v_e)-mv=-mg\,\dd t .
\]
展开并保留一阶小量：
\[
m\,\dd v+v_e\,\dd m=-mg\,\dd t,
\]
即
\begin{equation}
m\dot v=-v_e\dot m-mg.
\end{equation}
忽略重力时，分离变量并积分：
\begin{equation}
\Delta v=v_e\ln\frac{m_0}{m_f},
\label{eq:rocket-ch1}
\end{equation}
其中 \(m_0,m_f\) 分别为初始质量和燃料耗尽后的质量。若 \(v_e=3.0\,\mathrm{km/s}\)、\(m_0/m_f=10\)，则
\[
\Delta v=3.0\ln10\approx6.9\,\mathrm{km/s}.
\]
质量比越大、排气速度越高，速度增量越大；这就是多级火箭和电推进技术追求高比冲的原因。
\end{example}

\subsection{质点系角动量定理与角动量守恒}

一群物体的总角动量既可以来自整个质心绕外点运动，也可以来自各部分绕质心运动。机器人在自由圆环上爬行时，这两部分都要算。

系统对固定点 $O$ 的总角动量为
\begin{equation}
\bJ=\sum_i\br_i\times m_i\dot{\br}_i .
\end{equation}
对时间求导，并利用式 \eqref{eq:newton-system}：
\begin{equation}
\dv{\bJ}{t}
=\sum_i\br_i\times m_i\ddot{\br}_i
=\sum_i\br_i\times\bF_i^{(e)}
+\sum_i\br_i\times\sum_{j\ne i}\bF_{ij}.
\end{equation}
由式 \eqref{eq:internal-torque-zero}，内力矩成对抵消，于是
\begin{equation}
\dv{\bJ}{t}=\bM^{(e)},\qquad
\bM^{(e)}=\sum_i\br_i\times\bF_i^{(e)}.
\label{eq:sys-ang-ch1}
\end{equation}
若外力对 $O$ 的总力矩为零，则 $\bJ$ 守恒。

为了分离整体转动与内部转动，令
\begin{equation}
\br_i=\bR+\br_i',\qquad
\bv_i=\bV_c+\bv_i'.
\end{equation}
由质心定义，
\begin{equation}
\sum_i m_i\br_i'=0.
\end{equation}
对时间求导得
\begin{equation}
\sum_i m_i\bv_i'=0.
\end{equation}
把 $\br_i,\bv_i$ 代入 $\bJ$：
\begin{equation}
\bJ=\sum_i(\bR+\br_i')\times m_i(\bV_c+\bv_i').
\end{equation}
展开后，含 $\sum m_i\br_i'$ 和 $\sum m_i\bv_i'$ 的交叉项均为零，故
\begin{equation}
\bJ=\bR\times\bP+\bJ',\qquad
\bJ'=\sum_i\br_i'\times m_i\bv_i',
\end{equation}
其中 $\bP=M\bV_c$。第一项是质心的轨道角动量，第二项是系统相对质心的角动量。

\begin{example}[受冲量的双质点杆]
两个质量均为 $m$ 的质点 $A,B$ 由长为 $a$ 的轻杆相连，静止放在光滑水平桌面上。在 $B$ 点施加一个垂直于杆的冲量 $\bI$。求此后质心的速度和杆的角速度。

\begin{center}
\begin{tikzpicture}[>=Stealth,scale=1.0]
  \draw[thick] (0,0)--(4,0);
  \fill (0,0) circle (2.2pt) node[below]{$A$};
  \fill (4,0) circle (2.2pt) node[below]{$B$};
  \fill (2,0) circle (2.2pt) node[below]{$C$};
  \draw[->,very thick,red!75!black] (4,0)--++(0,1.0) node[right]{$\bI$};
  \draw[->,thick,blue!70!black] (2,0)--++(0,0.65) node[right]{$\bV_c$};
  \draw[->] (2.7,0.05) arc (0:120:0.7);
  \node at (3.15,0.55) {$\omega$};
\end{tikzpicture}
\end{center}

解：系统所受外力冲量为 $\bI$，由动量定理
\[
2m\bV_c=\bI\quad\Longrightarrow\quad
\bV_c=\frac{\bI}{2m}.
\]
质心 $C$ 位于杆的中点。对质心的转动惯量为
\[
I_C=2m\left(\frac{a}{2}\right)^2=\frac{ma^2}{2}.
\]
冲量对质心的角冲量为 $(a/2)I$，故
\[
I_C\omega=\frac{a}{2}I
\quad\Longrightarrow\quad
\omega=\frac{I}{ma}.
\]
此后质心以 $\bV_c$ 匀速平动，杆绕质心以 $\omega$ 匀速转动；两质点相对质心的位置矢量作匀速圆周运动。
\end{example}

\begin{example}[自由圆环上的机器人：先找对质心]
细圆环质量 $m$、半径 $R$，质量同为 $m$ 的点状机器人沿环爬行。系统初始静止，无外力及外力矩。令环的转角为 $\varphi$，机器人相对环的角为 $\theta$，其相对环心连线的空间方位为 $\alpha=\varphi+\theta$。机器人相对圆环爬一周时，环转过多少？

解：用 $C$ 表示环心，$G$ 表示总质心。令 $\be(\alpha)$ 为由环心指向机器人的单位矢量，则
\[
\br_b=\br_C+R\be,\quad
\bR_G=\frac{\br_C+\br_b}{2},\quad
\br_C-\bR_G=-\frac R2\be,\quad
\br_b-\bR_G=\frac R2\be.
\]
两者都距总质心 $R/2$，不能把机器人的距质心距离误当作 $R$。关于 $G$ 的角动量是环自转与两部分轨道角动量之和：
\[
J_G=mR^2\dot\varphi+2m(R/2)^2\dot\alpha
=mR^2\left(\dot\varphi+\tfrac12\dot\alpha\right)=0.
\]
因此 $\dd\varphi=-\dd\alpha/2$，又 $\dd\theta=\dd\alpha-\dd\varphi=3\dd\alpha/2$。当 $\Delta\theta=2\pi$ 时，
\[
\boxed{\Delta\alpha=\frac{4\pi}{3},\qquad
\Delta\varphi=-\frac{2\pi}{3}=-120^\circ.}
\]
总质心在空间静止，但环心在动。能量可由机器人内部驱动提供，不能另加“系统动能恒为零”的条件。计算只依赖总爬行角，与具体速度过程无关。
\begin{center}
\begin{tikzpicture}[>=Stealth]
\draw[thick,blue!65!black] (0,0) circle (1.6);
\fill (0,0) circle (2pt) node[below]{$C$};
\fill (0,0.8) circle (2pt) node[right]{$G$};
\fill (0,1.6) circle (3pt) node[above]{机器人};
\draw[<->] (-0.35,0)--(-0.35,0.8) node[midway,left]{$R/2$};
\draw[<->] (-0.35,0.8)--(-0.35,1.6) node[midway,left]{$R/2$};
\draw[->,thick] (1.6,0) arc (0:65:1.6);
\node[right] at (1.7,0.8) {$\alpha=\varphi+\theta$};
\end{tikzpicture}
\end{center}
\end{example}

\subsection{质点系动能定理与柯尼希定理}

向前滚的轮子一边带着质心移动，一边绕质心转。把每个质点的速度分成这两部分后，动能的混合项会因质心定义而消掉。

系统总动能为
\begin{equation}
T=\sum_i\frac12 m_i\dot{\br}_i^2 .
\end{equation}
把 $\dot{\br}_i=\bV_c+\bv_i'$ 代入：
\begin{equation}
T=\frac12\sum_i m_i(\bV_c+\bv_i')^2
=\frac12MV_c^2+\frac12\sum_i m_i v_i'^2
+\bV_c\cdot\sum_i m_i\bv_i'.
\end{equation}
最后一项因 $\sum_i m_i\bv_i'=0$ 而消失，于是
\begin{equation}
T=\frac12MV_c^2+T',\qquad
T'=\sum_i\frac12 m_i v_i'^2 .
\label{eq:koenig-ch1}
\end{equation}
这就是\textbf{柯尼希定理}：总动能等于质心整体运动的动能与相对质心的内部动能之和。

下面从牛顿第二定律导出动能定理。对每个质点，
\begin{equation}
m_i\ddot{\br}_i\cdot\dd\br_i=\bF_i\cdot\dd\br_i,
\end{equation}
对 $i$ 求和：
\begin{equation}
\dd T=\sum_i m_i\ddot{\br}_i\cdot\dd\br_i
=\sum_i\bF_i\cdot\dd\br_i
=\dd W^{(e)}+\dd W^{(i)}.
\label{eq:work-energy-system}
\end{equation}
内力功一般不为零。对一对内力，
\begin{equation}
\dd W_{ij}
=\bF_{ij}\cdot\dd\br_i+\bF_{ji}\cdot\dd\br_j
=\bF_{ij}\cdot\dd(\br_i-\br_j).
\end{equation}
对于定距离的理想中心约束，$\dd(\br_i-\br_j)$ 与 $\br_i-\br_j$ 垂直，而 $\bF_{ij}$ 沿 $\br_i-\br_j$，故内力功为零；对弹簧等系统，内力功则等于对应势能的减少量。

若作用可由不显含时间的总势能描述，可以引入总势能
\begin{equation}
V=\sum_i V_i^{(e)}+\sum_{i<j}V_{ij},
\end{equation}
其中 $V_i^{(e)}$ 是外势能，$V_{ij}$ 是第 $i,j$ 个质点间的相互作用势能。由保守力做功与势能的关系 $\dd W=-\dd V$，结合式 \eqref{eq:work-energy-system} 得
\begin{equation}
\dd(T+V)=0,
\end{equation}
即
\begin{equation}
E=T+V=\text{常量}.
\end{equation}
这就是质点系的\textbf{机械能守恒定律}。

\subsection{两体问题与约化质量}

研究两颗星时，可以先问两星的连线怎样变，再用质心的位置找回各自的位置。约化质量就是把这段相对运动写成一个等效质点所需的质量。

考虑两个质点 $m_1,m_2$，只受彼此的内力作用。由牛顿第二定律，
\begin{equation}
m_1\ddot{\br}_1=\bF_{12},\qquad
m_2\ddot{\br}_2=\bF_{21},\qquad
\bF_{21}=-\bF_{12}.
\end{equation}
定义相对位矢和质心
\begin{equation}
\br=\br_1-\br_2,\qquad
\bR=\frac{m_1\br_1+m_2\br_2}{m_1+m_2}.
\end{equation}
两式相减：
\begin{equation}
\ddot{\br}=\ddot{\br}_1-\ddot{\br}_2
=\frac{\bF_{12}}{m_1}-\frac{\bF_{21}}{m_2}
=\left(\frac1{m_1}+\frac1{m_2}\right)\bF_{12}.
\end{equation}
定义约化质量
\begin{equation}
\mu=\frac{m_1m_2}{m_1+m_2},
\end{equation}
则
\begin{equation}
\mu\ddot{\br}=\bF_{12}.
\label{eq:reduced-ch1}
\end{equation}
同时，
\begin{equation}
(m_1+m_2)\ddot{\bR}
=m_1\ddot{\br}_1+m_2\ddot{\br}_2
=\bF_{12}+\bF_{21}=0,
\end{equation}
所以质心作匀速直线运动。若两体内力进一步仅依赖相对距离且沿连线，两体问题因此化为一个质量为 \(\mu\) 的质点在中心力 \(\bF_{12}\) 中的单体问题；相对运动的具体轨道可直接引用后续中心力部分的结果，无需重新推导开普勒定律。

相对坐标与质心坐标还能反解原位置：
\[
\br_1=\bR+\frac{m_2}{m_1+m_2}\br,\qquad
\br_2=\bR-\frac{m_1}{m_1+m_2}\br.
\]
微分后代入总动能，交叉项因 $m_1m_2/(m_1+m_2)-m_2m_1/(m_1+m_2)=0$ 抵消，剩下
\begin{equation}\label{eq:ch1-twobody-energy}
T=\tfrac12(m_1+m_2)\dot{\bR}^{\,2}+\tfrac12\mu\dot{\br}^{\,2},\qquad
\bJ_G=\br\times\mu\dot{\br}.
\end{equation}
若势能为 $V(r)$，相对能量 $E_{\rm rel}=\mu\dot{\br}^{\,2}/2+V(r)$ 独立守恒；不能把一个粒子相对地面的速度直接当作相对速度。

\begin{example}[双星系统的动能]
两颗星质量分别为 $m_1,m_2$，相距 $r$，绕共同质心作匀速圆周运动。求系统的总动能。

\begin{center}
\begin{tikzpicture}[>=Stealth,scale=1.0]
  \fill (0,0) circle (3pt) node[below]{$m_1$};
  \fill (4,0) circle (2pt) node[below]{$m_2$};
  \fill (1.3,0) circle (2pt) node[below]{$C$};
  \draw[dashed] (0,0)--(4,0) node[midway,above]{$r$};
  \draw[->,thick,blue!70!black] (0,0)--++(0,0.8) node[left]{$\bv_1$};
  \draw[->,thick,red!75!black] (4,0)--++(0,-0.65) node[right]{$\bv_2$};
\end{tikzpicture}
\end{center}

解：质心静止，故由柯尼希定理 $T=T'$。相对运动满足式 \eqref{eq:reduced-ch1}。由动量守恒 $m_1v_1=m_2v_2$，且 $v=v_1+v_2$，可解出
\[
v_1=\frac{m_2}{m_1+m_2}v,\qquad
v_2=\frac{m_1}{m_1+m_2}v .
\]
相对运动是半径为 $r$ 的圆周运动，向心加速度为 $v^2/r$，引力大小为 $Gm_1m_2/r^2$，故
\[
\frac{m_1m_2}{m_1+m_2}\frac{v^2}{r}
=\frac{Gm_1m_2}{r^2},
\]
即 $v^2=G(m_1+m_2)/r$。代入动能：
\[
T=\frac12m_1v_1^2+\frac12m_2v_2^2
=\frac12\frac{m_1m_2}{m_1+m_2}v^2
=\frac{Gm_1m_2}{2r}.
\]
这个结果也可以看作“约化质量 $\mu$ 的动能”。
\end{example}

\begin{example}[空间站释放立方星：质心系中的反冲]
空间站质量为 $M$，一颗立方星质量为 $m$，二者由压缩弹簧连接，初始相对静止。弹簧释放能量 $E_s$ 后把立方星推出。忽略轨道曲率、潮汐力和空气阻力，求二者相对速度、各自速度，并说明质心运动。

\begin{center}
\begin{tikzpicture}[>=Stealth,scale=1.0]
  \draw[thick,fill=gray!15] (0,0.2) rectangle (2.2,1.0);
  \node at (1.1,1.2) {空间站 $M$};
  \fill (1.1,0.5) circle (2pt) node[below]{$C$};
  \draw[decorate,decoration={coil,amplitude=3pt,segment length=4pt}]
    (2.3,0.35)--(3.3,0.35);
  \node[above] at (2.8,0.55) {弹簧};
  \fill[fill=gray!45] (3.5,0.15) rectangle (3.9,0.75);
  \node at (3.7,-0.15) {$m$};
  \draw[->,thick,blue!70!black] (1.1,0.5)--++(-0.8,0) node[left]{$V$};
  \draw[->,thick,red!75!black] (3.9,0.5)--++(0.8,0) node[right]{$v$};
\end{tikzpicture}
\end{center}

解：在质心系中，系统总动量为零。设空间站和立方星的速度分别为 $V$ 和 $v$，方向相反，则
\[
MV+mv=0.
\]
设二者分离的相对速度为 $u=v-V$，解得
\[
v=\frac{M}{M+m}u,\qquad
V=-\frac{m}{M+m}u.
\]
弹簧释放的能量全部转化为质心系中的动能：
\[
E_s=\frac12MV^2+\frac12mv^2
=\frac12\frac{Mm}{M+m}u^2
=\frac12\mu u^2,
\]
其中 $\mu=Mm/(M+m)$ 为约化质量。因此
\[
u=\sqrt{\frac{2E_s}{\mu}},\qquad
v=\frac{M}{M+m}u,\qquad
V=-\frac{m}{M+m}u.
\]
质心不受额外外力，继续沿原轨道运动；释放过程只改变系统内部相对运动，不改变质心轨道。以下是教学设定，不是空间站实测数据。若取 $M=4.2\times10^5\,\mathrm{kg}$、$m=10\,\mathrm{kg}$、$E_s=10\,\mathrm{J}$，则
\[
u\approx1.4\,\mathrm{m/s},\qquad
|V|\approx3.4\times10^{-5}\,\mathrm{m/s}.
\]
空间站的反冲速度极小，而立方星获得约 $1.4\,\mathrm{m/s}$ 的分离速度。这里“质心不动”仅指所选质心系；在地心惯性系中质心仍运动。短时忽略潮汐的条件下，可把它近似为沿原轨道运动。
\end{example}

\begin{example}[地月系统的质心]
地球质量 \(M_E\approx5.97\times10^{24}\,\mathrm{kg}\)，月球质量 \(M_m\approx7.35\times10^{22}\,\mathrm{kg}\)，地月距离 \(d\approx3.84\times10^8\,\mathrm{m}\)。约化质量为
\[
\mu=\frac{M_EM_m}{M_E+M_m}\approx7.25\times10^{22}\,\mathrm{kg}.
\]
质心到地心的距离为
\[
d_E=\frac{M_m}{M_E+M_m}d\approx4.67\times10^6\,\mathrm{m},
\]
小于地球平均半径（约 $6.37\times10^6\,\mathrm{m}$），所以地月质心位于地球内部。严格说，地球和月球都绕这个共同质心运动；把月球看作绕静止地球只是近似。
\end{example}

\subsection{碰撞：质心系为什么能简化计算？}\label{sec:newton-collision}
两辆小车相撞，先减去它们共同的质心速度，就只剩相向接近或彼此分离的运动。恢复系数比较的正是接近和分离的相对速率。

碰撞时间很短，若外力冲量相对接触冲量可忽略，则总动量在碰撞前后近似守恒。不应把这一近似写成“碰撞时外力一定为零”。考虑一维正碰，初速度为 $u_1,u_2$，末速度为 $v_1,v_2$，取 $u_1-u_2>0$ 表示接近。质心速度
\[
V_G=\frac{m_1u_1+m_2u_2}{m_1+m_2}
=\frac{m_1v_1+m_2v_2}{m_1+m_2}
\]
在碰撞中不变。高速碰撞见第\ref{sec:sr-collision}节。用恢复系数描述分离相对速度：
\[
v_1-v_2=-e_r(u_1-u_2),\qquad 0\le e_r\le1.
\]
令 $w=u_1-u_2$，解动量及相对速度两式可得
\begin{align}
v_1&=V_G-\frac{m_2}{m_1+m_2}e_rw,\\
v_2&=V_G+\frac{m_1}{m_1+m_2}e_rw.
\end{align}
由式 \eqref{eq:ch1-twobody-energy}，质心平动能不变，损失的动能只来自相对运动：
\begin{equation}\label{eq:ch1-collision-loss}
T_{\rm before}-T_{\rm after}=\tfrac12\mu(1-e_r^2)w^2.
\end{equation}
完全弹性碰撞 $e_r=1$ 保持总动能；完全非弹性正碰 $e_r=0$ 时两物体同速。耗散的动能可变成形变和内能，总能量仍守恒。该单一恢复系数模型不包括斜碰中的摩擦和转动。
\begin{example}[一维碰撞的极限检查]
质量 $m$ 的粒子以速度 $u>0$ 撞上静止靶 $M$，取 $e_r=1$。上式给出
\[
v_1=\frac{m-M}{m+M}u,\qquad
v_2=\frac{2m}{m+M}u.
\]
若 $m=M$，二者交换速度；若 $M\gg m$，入射粒子近似以 $-u$ 反弹，而靶速度很小。若改取 $e_r=0$，共同速度为 $mu/(m+M)$，动能损失为 $mMu^2/[2(m+M)]$。质量相等且靶静止时，完全结合损失初始动能的一半，不能同时再要求动能守恒。

取 $m$ 为入射质量、$M$ 为静止靶。解 $mv_1+Mv_2=mu$、$v_2-v_1=e_ru$ 得 $v_1=(m-e_rM)u/(m+M)$、$v_2=(1+e_r)mu/(m+M)$。损失可直接用质心系动能 $\frac12[mM/(m+M)](v_2-v_1)^2$ 与碰前相对动能比较。
\end{example}
\begin{physicsbox}[title={物理图像与理解：把共同平移从碰撞中拿走}]
实验室看到两物体既共同前进又相互接近，容易混淆这两部分运动。在质心系中，二者动量大小相等、方向相反；一维弹性碰撞只将它们反向。解完内部碰撞再加回同一个 $V_G$，即可返回实验室系。质心系是惯性系的前提是该阶段合外力为零或可忽略。
\end{physicsbox}
\begin{figure}[htbp]
\centering
\begin{tikzpicture}[>=Stealth]
\node[anchor=west] at (-1.1,1.2) {碰前};
\node[anchor=west] at (-1.1,-0.2) {碰后};
\foreach \y in {1.2,-0.2}{
\fill[blue!70!black] (1,\y) circle (3pt);
\fill[red!70!black] (4,\y) circle (4pt);}
\draw[->,thick,blue!70!black] (1,1.2)--(2,1.2) node[midway,above]{$+p_*$};
\draw[->,thick,red!70!black] (4,1.2)--(3,1.2) node[midway,above]{$-p_*$};
\draw[->,thick,blue!70!black] (1,-0.2)--(0,-0.2) node[midway,below]{$-p_*$};
\draw[->,thick,red!70!black] (4,-0.2)--(5,-0.2) node[midway,below]{$+p_*$};
\end{tikzpicture}
\caption{质心系中的一维弹性碰撞。箭头表示动量，各物体速度还须分别除以自身质量；两物体的动量相消，但动能不为零。}
\label{fig:ch1-cm-collision}
\end{figure}

\subsection{位力定理}

物体绕行一周时，动能和势能可以一直变化，但它们的长时间平均可能有简单关系。这里的关键是区分有界运动的平均与某个时刻的瞬时值。

定义位力
\begin{equation}
G=\sum_i\bp_i\cdot\br_i .
\end{equation}
利用 \(\dot{\bp}_i=\bF_i\) 和 \(\bp_i=m_i\dot{\br}_i\)，对时间求导：
\begin{equation}
\dot G=\sum_i\dot{\bp}_i\cdot\br_i+\sum_i\bp_i\cdot\dot{\br}_i
=\sum_i\bF_i\cdot\br_i+2T .
\end{equation}
对时长 $\mathcal T$ 取平均，严格有
\[
\frac{G(t_0+\mathcal T)-G(t_0)}{\mathcal T}
=2\langle T\rangle_{\mathcal T}
+\left\langle\sum_i\bF_i\cdot\br_i\right\rangle_{\mathcal T}.
\]
只有当左边在 $\mathcal T\to\infty$ 时趋零，且右边平均存在，才得到\textbf{位力定理}
\begin{equation}
2\langle T\rangle+\left\langle\sum_i\bF_i\cdot\br_i\right\rangle=0.
\label{eq:virial-ch1}
\end{equation}
$G$ 有界是充分条件，位置有界本身并不自动保证动量有界，特别是存在碰撞奇点时。周期运动可直接取整数周期平均。

若全部力来自不显含时间的 $n$ 次齐次势 $V$，欧拉齐次函数定理给出
\[
\sum_i\br_i\cdot\nabla_iV=nV,\qquad 2\langle T\rangle=n\langle V\rangle.
\]
平方反比引力的势能次数为 $-1$，故 $2\langle T\rangle+\langle V\rangle=0$。对引力系统应使用质心内部动能，否则整体平动会污染估计。
\begin{modernbox}[title={进一步探索：位力质量估计的假设}]
若用模型 $V\simeq-CGM^2/R$，其中 $C$ 依赖密度分布，并定义三维内部均方速度 $\langle v^2\rangle$，则 $T=M\langle v^2\rangle/2$，估计为 $M\simeq R\langle v^2\rangle/(CG)$。若观测量是一维视向速度弥散 $\sigma_{\rm los}$，还需各向同性等假设才有 $\langle v^2\rangle=3\sigma_{\rm los}^2$。因此没有不注明速度定义与结构系数的通用质量公式。进一步学习需统计平均与引力势能积分，可先用均匀球完成章末习题。
\end{modernbox}

\section{非惯性参考系}

公交车突然起步时，乘客在车内看起来向后偏。若坚持在车里写方程，就要把车本身的加速带来的影响一起算进去。

前面讨论的参考系都默认是惯性系。实际观测中，我们常常身处加速或转动的参考系，例如加速的车厢、旋转的地球和太空电梯。只要在牛顿第二定律中引入相应的惯性力，非惯性系中的运动方程仍可写成“合力等于质量乘加速度”的形式。本节先建立相对运动的普遍关系，再讨论平动、转动参考系中的惯性力，并以地球自转和太空电梯为例说明其物理效应。

\subsection{相对运动与伽利略变换}

当两位观察者彼此匀速移动时，从一人的位置读数中减去另一坐标原点的运动，就能得到新读数。再微分一次得速度关系，两次得加速度关系。

设有两个参考系 $S$ 与 $S'$。$S'$ 的原点 $O'$ 在 $S$ 中的位矢为 $\bR(t)$。同一质点 $P$ 在 $S$ 中的位矢为 $\br$，在 $S'$ 中的位矢为 $\br'$，则
\begin{equation}
\br=\bR+\br'.
\label{eq:relative-position}
\end{equation}
对时间求导：
\begin{equation}
\bv=\dot{\bR}+\bv',\qquad
\ba=\ddot{\bR}+\ba'.
\label{eq:relative-velocity}
\end{equation}
若 $S'$ 相对 $S$ 作匀速直线平动，则 $\ddot{\bR}=0$，于是 $\ba=\ba'$。这说明加速度在相互作匀速直线平动的参考系中相同，为伽利略相对性原理的运动学基础。

\begin{figure}[htbp]
\centering
\begin{tikzpicture}[>=Stealth,scale=1.05]
  \draw[->] (0,0)--(5.2,0) node[below]{$x$};
  \draw[->] (0,0)--(0,3.5) node[left]{$y$};
  \coordinate (O) at (0,0);
  \coordinate (Op) at (1.8,0.75);
  \coordinate (P) at (4.0,2.7);
  \draw[->,very thick,blue!65!black] (O)--(Op) node[midway,below]{$\bR$};
  \draw[->] (Op)--++(2.6,0) node[below]{$x'$};
  \draw[->] (Op)--++(0,2.2) node[left]{$y'$};
  \draw[->,very thick,red!75!black] (O)--(P) node[pos=0.55,above left]{$\br$};
  \draw[->,very thick,green!45!black] (Op)--(P) node[midway,right]{$\br'$};
  \fill (P) circle (1.7pt) node[above right]{$P$};
  \node at (0.5,0.3) {$S$};
  \node at (2.2,1.4) {$S'$};
\end{tikzpicture}
\caption{两个参考系中的位矢关系 $\br=\bR+\br'$。若 $S'$ 相对 $S$ 平动，则三个矢量构成一个随时间变化的三角形。}
\label{fig:frames}
\end{figure}

最重要的特殊情形是伽利略变换。若 $S'$ 相对 $S$ 以恒定速度 $\bV$ 沿 $x$ 轴平动，且 $t=0$ 时两原点重合，则
\begin{equation}
x'=x-Vt,\qquad y'=y,\qquad z'=z,\qquad t'=t.
\label{eq:galilean}
\end{equation}
这是对同一事件换用运动参考系坐标的被动变换。第\ref{ch:relativity}章式 \eqref{eq:sr-lorentz} 在低速且同时性修正可忽略时恢复此式。
对时间求导得
\begin{equation}
v_x'=v_x-V,\qquad v_y'=v_y,\qquad v_z'=v_z,
\end{equation}
再求导得
\begin{equation}
a_x'=a_x,\qquad a_y'=a_y,\qquad a_z'=a_z .
\end{equation}
因此，在伽利略变换下加速度不变，牛顿第二定律的形式也不变。这就是\textbf{伽利略相对性原理}：一切惯性系对力学规律等价，不存在“绝对静止”的惯性系。

\begin{modernbox}
伽利略变换隐含了绝对时间 $t'=t$ 和绝对空间（长度不变）的假设。狭义相对论表明，当速度接近光速时，这些假设失效，应由洛伦兹变换替代。但在 $v\ll c$ 的经典力学范围内，伽利略变换是极好的近似。本书前六章除特别说明外均采用经典时空观。
\end{modernbox}

\begin{example}[理想径向收绳：角动量如何限制回收过程？]
在惯性平面内，一载荷以质量 $m$ 的质点建模，轻绳另一端固定于原点。绳始终绷紧，张力沿径向；卷扬机规定 $\dot r=-u$，$u>0$，初始 $r=r_0$，比角动量为 $h>0$。忽略其他力，求可行收绳阶段的轨道与张力。

解：力矩为零，故 $r^2\dot\theta=h$。而 $r=r_0-ut$，因此
\[
\theta(t)-\theta_0=\int_0^t\frac{h\,\dd t'}{(r_0-ut')^2}
=\frac hu\left(\frac1{r(t)}-\frac1{r_0}\right).
\]
由 $m(\ddot r-r\dot\theta^2)=-T$ 且 $\ddot r=0$ 得
\[
T=\frac{mh^2}{r^3},\qquad
\frac{\dd}{\dd t}\left[\frac m2(u^2+h^2/r^2)\right]
=\frac{mh^2u}{r^3}=Tu.
\]
收绳机构对载荷做正功，横向速率随 $r$ 减小而增大。模型形式上在 $r\to0$ 时要求无限张力与能量；实际必须在有限捕获半径停止、减速或转移动量。不存在只比较“横向速度与收绳速度”就能判断捕获的普遍准则；本题也不描述空间站轨道系中的完整相对动力学。
\end{example}



\subsection{加速平动参考系}

在向上加速的电梯里，人会感觉更重。把地面系中的加速度改写为相对于电梯的加速度，就能看出为什么电梯里出现一个额外的向下项。

设惯性系 $S$ 中质点受力 $\bF$，加速度为 $\ba$。另一参考系 $S'$ 相对 $S$ 以加速度 $\ba_0$ 平动，则质点在 $S'$ 中的加速度为
\[
\ba'=\ba-\ba_0.
\]
代入牛顿第二定律，
\begin{equation}
m\ba'=\bF-m\ba_0 .
\label{eq:accel-frame}
\end{equation}
若把 $-m\ba_0$ 称为\textbf{惯性力}，则形式上仍有“合力 = 质量 $\times$ 加速度”。惯性力不是物体间的相互作用，没有反作用力，只在非惯性系中出现。

\begin{physicsbox}
在加速的车厢中，乘客感到一个向后的“力”，这就是惯性力。它可以局域地等效于引力场：$-\,m\ba_0$ 与 $m\bg$（均匀引力）在动力学上不可区分。这一观察是爱因斯坦等效原理的起点。
\end{physicsbox}

\subsection{转动参考系}

坐标轴一旦转动，速度箭头即使相对这些轴分量不变，在地面看也可能转向。对转动基矢求导，正是各个附加项的来源。

设 $S'$ 相对惯性系 $S$ 以角速度 $\bom$ 转动，且两原点重合。取 $S'$ 的基矢量为 $\be_i'$，则基矢量随 $S'$ 转动：
\begin{equation}
\dv{\be_i'}{t}=\bom\times\be_i'.
\end{equation}
把任意矢量写成 $\bA=A_i'\be_i'$，在惯性系 $S$ 中求导：
\begin{equation}
\left(\dv{\bA}{t}\right)_S
=\dot A_i'\be_i'+A_i'\dot{\be}_i'
=\left(\dv{\bA}{t}\right)_{S'}+\bom\times\bA .
\label{eq:rotating-derivative}
\end{equation}
第一项是 $S'$ 中观察到的变化率，第二项来自基矢量的转动。

取 $\bA=\br$，得速度关系
\begin{equation}
\bv=\bv'+\bom\times\br .
\label{eq:rotating-velocity}
\end{equation}
再对 $\bv$ 应用式 \eqref{eq:rotating-derivative}：
\begin{align}
\ba&=\left(\dv{\bv}{t}\right)_S
=\left(\dv{\bv}{t}\right)_{S'}+\bom\times\bv \nonumber\\
&=\ba'+\dot{\bom}\times\br+\bom\times\bv'
+\bom\times(\bv'+\bom\times\br)\nonumber\\
&=\ba'+2\bom\times\bv'+\bom\times(\bom\times\br)
+\dot{\bom}\times\br .
\label{eq:rotating-acceleration}
\end{align}
代入牛顿第二定律，得到转动系中的运动方程
\begin{equation}
m\ba'=\bF
-2m\bom\times\bv'
-m\bom\times(\bom\times\br)
-m\dot{\bom}\times\br .
\label{eq:rotating-eom}
\end{equation}
后三项对应的惯性力依次为：
\begin{itemize}[leftmargin=2.2em]
  \item \textbf{科里奥利力} $\bF_C=-2m\bom\times\bv'$；
  \item \textbf{离心力} $\bF_{\rm cf}=-m\bom\times(\bom\times\br)$，当 $\bom\perp\br$ 时大小为 $m\omega^2\rho$，方向背离转轴；
  \item \textbf{欧拉力} $\bF_E=-m\dot{\bom}\times\br$，只在角速度变化时出现。
\end{itemize}

\begin{figure}[htbp]
\centering
\begin{tikzpicture}[>=Stealth]
\begin{scope}
\draw[thick,fill=gray!8] (0,0) rectangle (3,2);
\fill (1.5,1) circle (3pt) node[above]{$m$};
\draw[->,thick,blue!70!black] (1.5,1)--(.3,1) node[midway,below]{$-m\ba_0$};
\draw[->,thick] (1.8,2.3)--(3,2.3) node[right]{$\ba_0$};
\node at (1.5,-.6) {(a) 加速平动系};
\end{scope}
\begin{scope}[shift={(6,0)}]
\draw[->] (0,0)--(3.6,0) node[right]{$x'$};\draw[->] (0,-1.2)--(0,2.5) node[above]{$y'$};
\node[left] at (0,0) {$O$};\node at (.7,2.2) {$\bom\odot$};
\coordinate (P) at (2,0);\fill (P) circle (3pt);
\draw[->,thick,red!70!black] (P)--(2,1.2) node[above]{$\bv'$};
\draw[->,thick,green!40!black] (P)--(3.2,0) node[above]{$\bF_C$};
\draw[->,thick,blue!70!black] (2,-.45)--(3.2,-.45) node[below]{$\bF_{\rm cf}$};
\node at (1.6,-1.6) {(b) 在 $+x'$ 轴的粒子};
\end{scope}
\end{tikzpicture}
\caption{明确取 $\bom$ 垂直纸面向外。(b) 粒子位于 $+x'$ 轴且相对速度沿 $+y'$，于是 $-2m\bom\times\bv'$ 和离心力都沿 $+x'$；离心力箭头为便于辨认平移绘出。}
\label{fig:noninertial}
\end{figure}

\subsection{地球自转效应}

地面实验室随地球一起转。分析落体或摆时，先按东、北、竖直分解地球角速度，再判断哪些转动效应在题目的精度下需要保留。

在地球表面，取随地球转动的参考系。地球自转角速度记为 $\omega$。把 $\bom$ 分解为当地竖直分量和水平分量：
\begin{equation}
\bom=\omega\left(\cos\lambda\,\hat{\boldsymbol{\jmath}}_{\rm north}
+\sin\lambda\,\hat{\boldsymbol{k}}_{\rm up}\right),
\end{equation}
其中 $\lambda$ 为地理纬度。科里奥利力
\begin{equation}
\bF_C=-2m\bom\times\bv'
\end{equation}
对纯水平速度，科里奥利力的水平分量大小为
\begin{equation}
|F_{C,\rm horizontal}|=2m\omega v|\sin\lambda|,
\end{equation}
对纯水平运动，它使物体在北半球向右偏、南半球向左偏。其中 $\lambda$ 为地理纬度。

对傅科摆的小幅摆动，设摆锤在水平面内的坐标为 $(x,y)$，则科里奥利加速度的水平分量为 $(2\omega\sin\lambda\,\dot y,-2\omega\sin\lambda\,\dot x)$。于是
\begin{equation}
\ddot x-2\omega\sin\lambda\,\dot y+\omega_0^2x=0,\qquad
\ddot y+2\omega\sin\lambda\,\dot x+\omega_0^2y=0,
\end{equation}
其中 $\omega_0=\sqrt{g/l}$。这组方程的解表明，摆动平面以角速度 $\omega\sin\lambda$ 缓慢进动，进动周期为
\begin{equation}
T_{\rm F}=\frac{2\pi}{\omega|\sin\lambda|}.
\end{equation}
此处 $g$ 已吸收当地静态离心修正，忽略 $\omega^2$ 的其他影响；周期分子是恒星日，而非严格的24小时。赤道处该阶进动为零，周期趋于无穷；进动方向在两半球相反。

\begin{modernbox}
大尺度大气运动和长程弹道需要评估地球自转项。判断是否可以忽略科里奥利效应，可比较特征速度 $U$、长度 $L$ 与转速，构造罗斯贝数 $\mathrm{Ro}=U/(2\omega L)$。局部小尺度运动不能只凭“位于北半球”就断言自转项占主导。
\end{modernbox}

\begin{example}[绕竖直直径旋转的圆环与小珠]
在地面实验舱中，半径为 $R$ 的光滑圆形导轨以恒定角速度 $\omega$ 绕竖直直径转动，研究离心效应与重力竞争；小环质量为 $m$，套在导轨上。以最低点为 $\theta=0$，转动系中的有效势能为
\[
V_{\rm eff}(\theta)=-mgR\cos\theta-\frac12m\omega^2R^2\sin^2\theta.
\]
平衡条件 $V_{\rm eff}'=0$ 给出
\[
\sin\theta\left(g-\omega^2R\cos\theta\right)=0.
\]
故平衡位置为 $\theta=0,\pi$，以及（当 $\omega^2>g/R$ 时，左右各有一个对称位置）
\[
\cos\theta_0=\frac{g}{\omega^2R}.
\]
$\theta=0$ 在 $\omega^2<g/R$ 时稳定，在 $\omega^2>g/R$ 时失稳；$\theta=\pi$ 恒不稳定；$\theta_0$ 稳定。这是离心力导致的\emph{分岔}：转速增大到临界值时，最低点失稳，系统转向新的稳定平衡。本例假定地面均匀重力。失重旋转舱应取相应的微重力模型，不能直接照搬地面平衡角。

由 $V_{\rm eff}''=mR[g\cos\theta-\omega^2R\cos(2\theta)]$，最低点和最高点的曲率分别为 $mR(g-\omega^2R)$ 和 $-mR(g+\omega^2R)$，新平衡支的曲率为 $m\omega^2R^2\sin^2\theta_0$。恰在临界转速，原点的二次项消失，而 $V_{\rm eff}(\theta)-V_{\rm eff}(0)=\frac{mgR}{2}(1-\cos\theta)^2$，最低点仍是严格局部极小，恢复力的首项为三阶。
\end{example}

\begin{figure}[htbp]
\centering
\begin{tikzpicture}[>=Stealth,scale=1.0]
  \draw[thick] (0,3.3)--(0,-0.3);
  \draw[blue!45,thick] (0,1.5) circle (1.5);
  \draw[dashed] (0,1.5)--(0,3.0);
  \fill (1.06,0.44) circle (2.4pt) node[below right]{$m$};
  \draw[->,thick,red!75!black] (1.06,0.44)--++(0,-0.9) node[below]{$m\boldsymbol g$};
  \draw[->,thick,blue!70!black] (1.06,0.44)--++(1.0,0) node[above right]{$m\omega^2\rho$};
  \draw[->,thick,green!45!black] (1.06,0.44)--++(-0.65,0.65) node[midway,below left]{$\bN$};
  \draw (0,1.5)--(1.06,0.44);
  \draw (0,1.5) ++(-90:0.75) arc (-90:-40:0.75);
  \node at (0.30,0.95) {$\theta$};
  \draw[->] (1.75,2.9) arc (60:20:1.0);
  \node at (2.25,2.75) {$\omega$};
\end{tikzpicture}
\caption{绕竖直直径旋转的圆环与小珠。转动系中，重力、约束力与离心力共同决定平衡位置；当 $\omega^2>g/R$ 时，最低点失稳，出现偏离竖直方向的稳定平衡。}
\label{fig:rotating-hoop}
\end{figure}


\subsection{太空电梯：旋转系中的有效引力}

沿赤道上方的旋转缆绳移动时，引力随高度减弱，而随地球转动所需的向心加速度随半径增大。比较两者就能判断不同高度处缆绳怎样受拉。

太空电梯的设想是：从地球赤道向外延伸一根缆绳，顶端系一个配重，使整根缆绳随地球自转。缆绳上的每一点都在旋转系中静止，因此可以直接用离心力与引力的平衡来分析。这个例子把非惯性系、有效势能和材料强度联系起来，是现代空间工程中的经典力学问题。

\begin{example}[太空电梯的张力与视重]
设地球质量为 $M$、半径为 $R_E$、自转角速度为 $\Omega$。缆绳沿赤道径向向外，线密度为 $\mu$，顶端 $r_t$ 处连接配重。求缆绳中的张力分布、同步轨道半径，以及攀爬器在缆绳上的视重和科里奥利偏转。

\begin{center}
\begin{tikzpicture}[>=Stealth,scale=0.9]
  \shade[ball color=blue!20] (0,0) circle (0.75);
  \draw[thick] (0,0) circle (0.75);
  \node[font=\small] at (0,-1.0) {地球};
  \draw[very thick,black!70] (0,0.75)--(0,5.6);
  \fill (0,5.6) circle (3pt) node[right]{配重};
  \draw[dashed] (-1.2,3.7)--(1.2,3.7);
  \fill[red!75!black] (0,3.7) circle (2pt) node[left]{$r=r_g$};
  \node[font=\small,anchor=west] at (1.3,3.7) {同步轨道};
  \fill (0,2.1) circle (2.2pt) node[right]{$m$};
  \draw[->,thick,red!75!black] (0,2.1)--(0,1.35) node[midway,left]{$mg$};
  \draw[->,thick,blue!70!black] (0,2.1)--(0,2.85) node[midway,right]{$m\Omega^2r$};
  \draw[->] (-1.0,0.1) arc (180:120:1.0);
  \node at (-1.25,0.55) {$\Omega$};
  \node[font=\small,anchor=west] at (0.15,4.8) {缆绳};
\end{tikzpicture}
\end{center}

解：在随地球转动的参考系中，单位质量的等效势能为
\[
\Phi_{\rm eff}(r)=-\frac{GM}{r}-\frac12\Omega^2r^2,
\]
等效径向加速度（取向外为正）为
\[
a_{\rm eff}(r)=-\Phi_{\rm eff}'(r)=\Omega^2r-\frac{GM}{r^2}.
\]
令 $a_{\rm eff}=0$，得同步轨道半径
\[
r_g=\left(\frac{GM}{\Omega^2}\right)^{1/3}
\approx4.22\times10^7\ \mathrm{m},
\]
对应高度约 $3.58\times10^7\ \mathrm{m}$。在 $r<r_g$ 处引力占优，$a_{\rm eff}<0$，缆绳被向下拉；在 $r>r_g$ 处离心力占优，$a_{\rm eff}>0$，缆绳被向外拉。同步轨道处视重为零。

取 $r$ 到 $r+\dd r$ 的一段缆绳，其平衡条件为
\[
\dd T-\mu\frac{GM}{r^2}\dd r+\mu\Omega^2r\,\dd r=0,
\]
即
\[
\frac{\dd T}{\dd r}
=\mu\left(\frac{GM}{r^2}-\Omega^2r\right).
\]
顶端若挂有质量 $M_t$ 的配重，配重的径向平衡要求
\[
T_t=T(r_t)=M_t\left(\Omega^2r_t-\frac{GM}{r_t^2}\right).
\]
取 $r_t>r_g$ 才能在这个模型中提供正的顶端张力。由微分方程从 $r_t$ 积到 $r$，
\begin{align}
T(r)-T_t
&=\mu\int_{r_t}^{r}\left(\frac{GM}{s^2}-\Omega^2s\right)\dd s\\
&=\mu GM\left(\frac1{r_t}-\frac1r\right)
+\frac{\mu\Omega^2}{2}(r_t^2-r^2).
\end{align}
对最后一式求导可恢复原微分方程，这是检查符号的直接方法。顶端自由且无配重时才可取 $T_t=0$，并仍须检查全段 $T(r)\ge0$；负张力表示柔绳模型不可行。因为 $T'$ 在 $r_g$ 内侧为正、外侧为负，跨过 $r_g$ 的可行绳段在同步半径处达到张力最大值。本例只分析理想线密度恒定绳，未论证具体材料或工程可建性。

对固定在绳上的小质量 $m$，保持其静止所需的净径向支承力为 $-ma_{\rm eff}$，大小为 $m|a_{\rm eff}|$；超过同步半径后支承方向反转。因此“视重为零”并非沿整根绳成立。

若攀爬器以径向速度 $v$ 沿缆绳运动，它还受到科里奥利力
\[
\bF_C=-2m\bOm\times\bv',
\]
其大小为 $2m\Omega v$。在赤道处，向上爬升时科里奥利力指向西，向下返回时指向东。这个横向力必须由缆绳的约束或攀爬机构承担，是太空电梯设计中的另一个关键力学因素。
\end{example}

\section{中心力}\label{sec:central}

若力始终沿质点与固定中心的连线，它对该中心的力矩为零。我们可先用这个事实确定轨道平面，再讨论半径和方位角怎样变。

中心力是质点力学中最重要的一类力：行星绕日、电子绕核、散射实验都归结为中心力问题。由于力始终指向或背离一个固定点，角动量守恒，运动被限制在一个平面内；再结合能量守恒，就可以把二维问题化为一维的径向问题。本节从这些基本性质出发，导出比内方程、开普勒定律、有效势能、轨道稳定性、伯特兰定理、散射截面和潮汐力。

\subsection{中心力的基本性质与比内方程}

轨道方程可以把半径当作方位角的函数。用角动量把时间导数换成角度导数后，便能直接求轨道形状；若还要知道飞了多久，仍需用时间关系。

本书也将中心力称为有心力。以下取力心固定、径向力的大小只依赖到力心距离的模型。

若力始终指向或背离某一固定点，且大小只与到该点的距离有关，
\begin{equation}
\bF=F(r)\er,
\label{eq:central-force}
\end{equation}
则称为\textbf{中心力}。取力心为原点，则
\begin{equation}
\bM=\br\times\bF=F(r)\,\br\times\er=0,
\end{equation}
因此角动量守恒，运动限于一个平面内。取该平面为极坐标平面，则
\begin{equation}
J=mr^2\dot\theta=\text{常量},\qquad
h=\frac{J}{m}=r^2\dot\theta .
\label{eq:areal-constant}
\end{equation}
$h$ 称为\textbf{比角动量}。面积速度 ${\rm d}A/{\rm d}t=h/2$ 为常量。

径向运动方程为
\begin{equation}
m\left(\ddot r-r\dot\theta^2\right)=F(r).
\label{eq:radial-eom}
\end{equation}
利用 $\dot\theta=h/r^2$，可消去 $\dot\theta$：
\begin{equation}
m\ddot r-\frac{mh^2}{r^3}=F(r).
\label{eq:radial-eom2}
\end{equation}
机械能为
\begin{equation}
E=\frac12m\dot r^2+\frac{mh^2}{2r^2}+V(r)
=\frac12m\dot r^2+V_{\rm eff}(r),
\label{eq:central-energy}
\end{equation}
其中
\begin{equation}
V_{\rm eff}(r)=V(r)+\frac{mh^2}{2r^2}
=V(r)+\frac{J^2}{2mr^2}
\label{eq:effective-potential}
\end{equation}
$V_{\rm eff}$ 称为\textbf{有效势能}，其中第二项称为\textbf{离心势能}。若 $J\ne0$ 且 $r^2V(r)\to0$，则有效势在 $r\to0$ 时趋于正无穷。若 $V=-a/r^2+o(r^{-2})$，主导项系数为 $J^2/(2m)-a$；系数为负时没有这种势垒，恰为零时须检查次阶项。

\begin{figure}[htbp]
\centering
\begin{tikzpicture}
\begin{axis}[width=.85\textwidth,height=6.2cm,axis lines=left,
 xlabel={$R=r/r_0$},ylabel={$U=V_{\rm eff}/(k/r_0)$},xmin=0,xmax=5,ymin=-.65,ymax=.8,
 xtick={.6126,1,2.7208},xticklabels={$R_1$,$1$,$R_2$},ytick={-.5,-.3,0,.3}]
\addplot[thick,blue!70!black,domain=.35:5,samples=250] {-1/x+.5/x^2};
\addplot[dashed,red!70!black,domain=.35:5] {-.3};
\addplot[dashed,green!40!black,domain=.35:5] {.3};
\addplot[only marks,mark=*] coordinates {(1,-.5) (.612574,-.3) (2.720759,-.3) (.441518,.3)};
\node at (axis cs:3.8,-.43) {$\mathcal E=-0.3$};
\node[above left] at (axis cs:.441518,.3) {$R_3$};\node at (axis cs:3.8,.38) {$\mathcal E=0.3$};
\end{axis}
\end{tikzpicture}
\caption{平方反比吸引势的无量纲有效势 $U=-1/R+1/(2R^2)$，$r_0=J^2/(mk)$。负能量的两个交点给出径向振动范围，正能量只有一个近心转折点；$R=1$ 为圆轨道。标记均由所绘曲线与水平能量线的交点计算。}
\label{fig:effective-potential}
\end{figure}

为了求轨道，下面引入比内方程。令 $u=1/r$。利用
\begin{equation}
\dot r=\dv{r}{\theta}\dot\theta
=-\frac{1}{u^2}\dv{u}{\theta}\,hu^2
=-h\dv{u}{\theta},
\end{equation}
\begin{equation}
\ddot r=\dv{\dot r}{\theta}\dot\theta
=-h\dvv[2]{u}{\theta}\,hu^2
=-h^2u^2\dvv[2]{u}{\theta}.
\end{equation}
代入径向方程 \eqref{eq:radial-eom2}，得
\begin{equation}
-mh^2u^2\left(\dvv[2]{u}{\theta}+u\right)=F(1/u).
\end{equation}
于是
\begin{equation}
\boxed{\;\dvv[2]{u}{\theta}+u
=-\frac{F(1/u)}{mh^2u^2}
=-\frac{mF(1/u)}{J^2u^2}\;}
\label{eq:binet}
\end{equation}
这就是\textbf{比内方程}。它把“求轨道”的问题转化为关于 $u(\theta)$ 的二阶常微分方程。

\begin{mathbox}
比内方程的推导只用了角动量守恒和极坐标运动方程，不依赖于 $F(r)$ 的具体形式。因此在 $J\ne0$、可用单调角度参数的轨道段上，它适用于任意中心力；$J=0$ 的纯径向运动应直接用径向时间方程。若已知 $F(r)$，解出 $u(\theta)$ 即得轨道；反之，若已知轨道 $u(\theta)$，也可反推出中心力 $F(r)$。
\end{mathbox}

\subsection{平方反比引力与开普勒定律}

越靠近太阳，引力越强，行星一般也走得更快。把平方反比力代入轨道方程，可以看到椭圆、等面积和周期关系怎样来自同一模型。

考虑引力
\begin{equation}
F(r)=-\frac{k}{r^2}=-ku^2,\qquad k=GMm.
\end{equation}
代入比内方程：
\begin{equation}
\dvv[2]{u}{\theta}+u=\frac{km}{J^2}\equiv\text{常量}.
\end{equation}
其通解为
\begin{equation}
u=A\cos(\theta-\theta_0)+\frac{km}{J^2}.
\end{equation}
于是
\begin{equation}
r=\frac{p}{1+e\cos(\theta-\theta_0)},\qquad
p=\frac{J^2}{km},\qquad
e=\frac{AJ^2}{km}.
\label{eq:conic}
\end{equation}
以下取 $J\ne0$、$e\ge0$，这是以力心为焦点的圆锥曲线：
\begin{itemize}[leftmargin=2.2em]
  \item $e=0$：圆；
  \item $0<e<1$：椭圆；
  \item $e=1$：抛物线；
  \item $e>1$：双曲线。
\end{itemize}
在近心点 $\theta=\theta_0$，$\dot r=0$、$r_p=p/(1+e)$，故
\[
E=\frac{J^2}{2mr_p^2}-\frac{k}{r_p}
=\frac{k}{2p}\bigl[(1+e)^2-2(1+e)\bigr].
\]
用 $p=J^2/(km)$ 消去 $p$，得到总能量与偏心率的关系
\begin{equation}
E=\frac{mk^2}{2J^2}\left(e^2-1\right).
\label{eq:energy-eccentricity}
\end{equation}
因此 $E<0$ 对应椭圆（束缚态），$E=0$ 对应抛物线（临界逃逸），$E>0$ 对应双曲线（散射态）。

\begin{figure}[htbp]
\centering
\begin{tikzpicture}\begin{axis}[width=.31\textwidth,height=4.6cm,axis equal image,axis lines=middle,xmin=-3.8,xmax=1.6,ymin=-3,ymax=3,ticks=none,title={椭圆},xlabel={$x/p$},ylabel={$y/p$}]
\addplot[thick,blue!70!black,domain=-180:180,samples=200] ({cos(x)/(1+0.6*cos(x))},{sin(x)/(1+0.6*cos(x))});
\fill[orange!80!black] (axis cs:0,0) circle (2pt) node[below left]{$F$};\end{axis}\end{tikzpicture}
\begin{tikzpicture}\begin{axis}[width=.31\textwidth,height=4.6cm,axis equal image,axis lines=middle,xmin=-3.8,xmax=1.6,ymin=-3,ymax=3,ticks=none,title={抛物线},xlabel={$x/p$},ylabel={$y/p$}]
\addplot[thick,blue!70!black,domain=-130:130,samples=200] ({cos(x)/(1+1*cos(x))},{sin(x)/(1+1*cos(x))});
\fill[orange!80!black] (axis cs:0,0) circle (2pt) node[below left]{$F$};\end{axis}\end{tikzpicture}
\begin{tikzpicture}\begin{axis}[width=.31\textwidth,height=4.6cm,axis equal image,axis lines=middle,xmin=-3.8,xmax=1.6,ymin=-3,ymax=3,ticks=none,title={双曲线},xlabel={$x/p$},ylabel={$y/p$}]
\addplot[thick,blue!70!black,domain=-115:115,samples=200] ({cos(x)/(1+1.5*cos(x))},{sin(x)/(1+1.5*cos(x))});
\fill[orange!80!black] (axis cs:0,0) circle (2pt) node[below left]{$F$};\end{axis}\end{tikzpicture}

\caption{统一取焦点为原点、半通径为尺度，分别绘制 $e=0.6,1,1.5$ 的轨道。双曲线只画与该吸引力焦点对应的物理支，不能把另一支也当成同一次引力散射轨迹。}
\label{fig:conics}
\end{figure}

对椭圆，$p=a(1-e^2)$、短半轴 $b=a\sqrt{1-e^2}$，因此 $E=-k/(2a)$。面积速度为 $h/2$，一周面积为 $\pi ab$，于是
\[
T=\frac{2\pi ab}{h},\qquad
h^2=\frac{k}{m}p=GM a(1-e^2),\qquad
T^2=\frac{4\pi^2a^3}{GM}.
\]
这个推导先把力心视为固定。对真实两体相对运动，应把动力学质量换成约化质量、$k=Gm_1m_2$，从而 $k/\mu=G(m_1+m_2)$。此时 $a$ 是两体相对距离轨道的半长轴，而非任意一颗天体关于质心轨道的半长轴。

开普勒三定律可由上述结果导出：
\begin{enumerate}[label=(\arabic*),leftmargin=2.2em]
  \item 行星沿椭圆轨道运动，太阳位于椭圆的一个焦点上；
  \item 行星与太阳的连线在相等时间内扫过相等面积（面积速度守恒）；
  \item 行星公转周期 $T$ 的平方与椭圆半长轴 $a$ 的立方成正比：
  \begin{equation}
  T^2=\frac{4\pi^2a^3}{G(M+m)}.
  \label{eq:kepler3}
  \end{equation}
  当 $M\gg m$ 时，$T^2\approx4\pi^2a^3/(GM)$。
\end{enumerate}

\begin{figure}[htbp]
\centering
\begin{tikzpicture}[>=Stealth,scale=1.35]
\coordinate (F) at (1.2,0);
\draw[thick,blue!70!black] (0,0) ellipse (2 and 1.6);
\fill[blue!12] (F)--(2,0) plot[domain=0:65.712221235,samples=60] ({2*cos(\x)},{1.6*sin(\x)})--cycle;
\fill[red!12] (F)--(-2,0) plot[domain=180:201.678498205,samples=60] ({2*cos(\x)},{1.6*sin(\x)})--cycle;
\draw[thick] (F)--(2,0);\draw[thick] (F)--({2*cos(65.712221235)},{1.6*sin(65.712221235)});
\draw[thick] (F)--(-2,0);\draw[thick] (F)--({2*cos(201.678498205)},{1.6*sin(201.678498205)});
\fill[orange!80!black] (F) circle (2pt) node[below right]{$F$};
\node at (1.4,.4) {$A_1$};\node at (-1,-.25) {$A_2$};
\end{tikzpicture}
\caption{两块区域由椭圆弧和焦点连线围成。取 $a=2,b=1.6,e=0.6$，偏近点角区间按相等的 $\Delta(E-e\sin E)$ 计算，从而表示相等时间与相等扫掠面积，不能以任意直线三角形替代扇形。}
\label{fig:kepler-area}
\end{figure}

\subsection{有效势能、轨道分类与稳定性}

把同一角动量下的转动动能和原来的势能合在一起，半径便像一个一维坐标那样运动。有效势的最低点给出可能的稳定圆轨道。

由式 \eqref{eq:central-energy}，径向运动等价于一维问题：
\begin{equation}
\frac12m\dot r^2+V_{\rm eff}(r)=E.
\end{equation}
圆轨道满足
\begin{equation}
V_{\rm eff}'(r_0)=0
\quad\Longleftrightarrow\quad
F(r_0)=-\frac{J^2}{mr_0^3}.
\label{eq:circular-condition}
\end{equation}
这就是“中心力提供向心力”的条件。圆轨道的能量为
\begin{equation}
E_0=V_{\rm eff}(r_0).
\end{equation}
令 $r=r_0+\delta r$，在固定角动量下展开 $m\ddot r=-V_{\rm eff}'(r)$：
\[
m\ddot{\delta r}=-V_{\rm eff}''(r_0)\delta r+O(\delta r^2).
\]
所以线性径向振动频率为 $\omega_r^2=V_{\rm eff}''(r_0)/m$，稳定性要求
\begin{equation}
V_{\rm eff}''(r_0)>0.
\label{eq:stability}
\end{equation}
对于平方反比引力，
\begin{equation}
r_0=\frac{J^2}{mk},\qquad
E_0=-\frac{mk^2}{2J^2},\qquad
\omega_0=\frac{J}{mr_0^2}=\sqrt{\frac{k}{mr_0^3}}.
\end{equation}

\begin{example}[月球探测器的霍曼转移]
月球探测器在半径 $r_1$ 的环月圆轨道上运行，欲通过霍曼转移进入半径 $r_2>r_1$ 的共面圆轨道。求两次点火所需的速度增量，并估算从 $100\,\mathrm{km}$ 到 $1000\,\mathrm{km}$ 环月圆轨道所需的总 $\Delta v$。

\begin{center}
\begin{tikzpicture}[>=Stealth,scale=0.95]
  \fill[gray!25] (0,0) circle (0.55);
  \node at (0,-0.85) {月球};
  \draw[blue!55,thick] (0,0) circle (1.2);
  \draw[blue!35,thick] (0,0) circle (2.0);
  \draw[red!75!black,thick] (-0.4,0) ellipse (1.6 and 1.549193);
  \fill (1.2,0) circle (2pt);
  \fill (-2.0,0) circle (2pt);
  \draw[->,thick,green!45!black] (1.2,0)--++(0,0.55) node[right]{$\Delta v_1$};
  \draw[->,thick,green!45!black] (-2.0,0)--++(0,-0.45) node[left]{$\Delta v_2$};
  \node[blue!55] at (0.9,1.10) {$r_1$};
  \node[blue!35] at (1.45,1.80) {$r_2$};
\end{tikzpicture}
\end{center}

解：霍曼转移轨道是一个椭圆，近月点为 $r_1$、远月点为 $r_2$，半长轴为
\[
a=\frac{r_1+r_2}{2}.
\]
由椭圆能量 $E=-m\mu_M/(2a)$ 与任意点能量 $E=mv^2/2-m\mu_M/r$ 相等，消去 $E$ 得活力公式
\[
v^2=\mu\left(\frac{2}{r}-\frac{1}{a}\right),
\]
其中本例 $\mu=GM$ 为月球引力参数，与前文约化质量符号同形，但物理量、单位不同；以下数值用 $\mu_M$ 表示。第一次点火在 $r_1$ 处，把圆轨道速度
\[
v_1^c=\sqrt{\frac{\mu}{r_1}}
\]
加速到转移轨道近月点速度
\[
v_1^t=\sqrt{\mu\left(\frac{2}{r_1}-\frac{1}{a}\right)},
\]
故
\[
\Delta v_1=v_1^t-v_1^c .
\]
到达远月点 $r_2$ 后，转移轨道速度为
\[
v_2^t=\sqrt{\mu\left(\frac{2}{r_2}-\frac{1}{a}\right)},
\]
而目标圆轨道速度为
\[
v_2^c=\sqrt{\frac{\mu}{r_2}}.
\]
第二次点火使探测器加速到 $v_2^c$，于是
\[
\Delta v_2=v_2^c-v_2^t .
\]
总速度增量为 $\Delta v=\Delta v_1+\Delta v_2$。

转移飞行时间为椭圆半周期 $t_{\rm tr}=\pi\sqrt{a^3/\mu_M}$。本例计算约为 $4.91\times10^3\,\mathrm{s}$。

取月球半径 $R_M=1737.4\,\mathrm{km}$、$\mu_M=4.902800118\times10^{12}\,\mathrm{m^3/s^2}$。对 $r_1=R_M+100\,\mathrm{km}$、$r_2=R_M+1000\,\mathrm{km}$，计算得
\[
\Delta v_1\approx153\,\mathrm{m/s},\qquad
\Delta v_2\approx139\,\mathrm{m/s},
\]
故总 $\Delta v\approx0.29\,\mathrm{km/s}$。霍曼转移的两次点火都沿速度方向，在球对称两体模型、瞬时切向双脉冲、指定两圆轨道之间具有标准的低速度增量性质；不能推广为允许任意脉冲数及中间轨道时的全局最优。实际任务需另计摄动、有限推力和误差。采用的月球参数来源于 JPL，见章末来源；轨道高度是教学设定，不对应特定任务。
\end{example}

有效势能还决定圆轨道的稳定性。考虑一般幂次中心力
\begin{equation}
F(r)=-\frac{k}{r^n}=-ku^n .
\end{equation}
由比内方程，圆轨道满足
\begin{equation}
ku_0^{n-3}=mh^2 .
\end{equation}
令 $u=u_0+\xi$，线性化得
\begin{equation}
\dvv[2]{\xi}{\theta}+(3-n)\xi=0 .
\label{eq:stability-perturb}
\end{equation}
因此圆轨道稳定的条件是
\begin{equation}
\boxed{\;n<3\;}
\label{eq:stability-condition}
\end{equation}
当 $n<3$ 时，径向扰动表现为简谐振荡，轨道在圆附近振荡；当 $n>3$ 时出现指数增长；$n=3$ 时线性方程为 $\xi''=0$，是退化临界情况，扰动可线性增长，不能称为指数增长或据此断言稳定。平方反比引力 $n=2$ 是稳定情形，且扰动频率与角频率相同，轨道闭合（伯特兰定理的一个特例）。


\begin{example}[幂次势的圆轨道与径向频率]
质量为 $m$ 的质点在势能 $V(r)=\beta r^k$ 中运动，角动量为 $L$。求圆轨道半径与径向微振动频率，并验证拱点角公式。

解：有效势能为
\[
V_{\rm eff}(r)=\frac{L^2}{2mr^2}+\beta r^k .
\]
设 $L\ne0$、$\beta k>0$ 且 $k\ne-2$，圆轨道满足 $V_{\rm eff}'(r_0)=0$：
\[
-\frac{L^2}{mr_0^3}+\beta k r_0^{k-1}=0
\quad\Longrightarrow\quad
r_0=\left(\frac{L^2}{m\beta k}\right)^{1/(k+2)} .
\]
在 $r_0$ 附近作小振动，径向频率为
\[
\omega_r^2=\frac{V_{\rm eff}''(r_0)}{m}
=\frac{1}{m}\left[\frac{3L^2}{mr_0^4}+\beta k(k-1)r_0^{k-2}\right].
\]
利用圆轨道条件 $L^2/(mr_0^3)=\beta k r_0^{k-1}$，得
\[
\omega_r^2=\frac{\beta k(k+2)}{m}r_0^{k-2}
=(k+2)\omega_0^2,
\]
其中 $\omega_0=L/(mr_0^2)$ 是圆轨道的角频率。稳定性还要求 $k>-2$。在小扰动极限，从近心点到远心点的拱点角为
\[
\Psi=\pi\frac{\omega_0}{\omega_r}
=\frac{\pi}{\sqrt{k+2}} .
\]
当 $k=2$ 时 $\Psi=\pi/2$，当 $k=-1$ 时 $\Psi=\pi$；其他幂次也可能给出有理频率比，但近圆有理比只约束线性近似，不能保证有限振幅的所有轨道闭合。这个结果正是伯特兰定理证明的起点。
\end{example}

\subsection{进一步探索：伯特兰定理与轨道进动}
近圆轨道每次到达最近点时，方向未必一样。比较半径往复的周期与绕中心转动的周期，便可判断近心方向是否逐圈移动。

圆轨道稳定不等于所有轨道闭合。本小节超出首次学习所必需的核心内容，使用比内方程与小振幅展开，说明闭合要求有多严格；所需基础是泰勒展开和三角恒等式。
\begin{theorem}[伯特兰定理的通常正则情形]
对光滑吸引中心势，在可形成稳定圆轨道的连通半径范围内，若所有非碰撞束缚轨道都闭合，则允许的势只有 $V=-k/r$ 与 $V=kr^2/2$（$k>0$），相差常数不影响运动。这里要求一族稳定圆轨道与其邻近束缚轨道，不能用仅存在孤立闭轨道替代。
\end{theorem}

\paragraph{必要条件：近圆频率比必须保持不变。}
由前节径向小振动，近圆轨道的半径向周期转角为
\begin{equation}\label{eq:apsidal-angle}
\Psi_0=\pi\left[3+\frac{r_0V''(r_0)}{V'(r_0)}\right]^{-1/2}.
\end{equation}
闭合要求径向与角向周期可公度。若对连续变化的轨道参数始终闭合，连续变化的频率比又只能取有理数，故它必须是常数。令中括号为 $C>0$，则
\[
rV''=(C-3)V',\qquad V'=A r^{C-3}.
\]
在该连通范围内，势能为幂次 $V=\beta r^k+\mathrm{const}$，$k=C-2>-2$；$k=0$ 的极限形式为对数势 $A\ln r$。仅检验近圆运动还不能排除其他有理频率比。

\paragraph{再进一步：振幅变化不能改变频率比。}
对幂次势，令 $u=u_0(1+x)$，用圆轨道条件消去系数，比内方程变成
\[
x''+1+x=(1+x)^{-k-1},
\]
撇号表示对方位角 $\theta$ 求导。展开到三阶：
\[
x''+w_0^2x=a_2x^2+a_3x^3+O(x^4),\quad
w_0^2=k+2,
\]
\[
a_2=\frac{(k+1)(k+2)}2,\qquad
a_3=-\frac{(k+1)(k+2)(k+3)}6.
\]
设 $\zeta=w\theta$，用小振幅参数 $\epsilon$ 写为
\[
x=\epsilon\cos\zeta+\epsilon^2(c_0+c_2\cos2\zeta)+O(\epsilon^3),\quad
w=w_0(1+d\epsilon^2+\cdots).
\]
二阶常数项及倍频项给出 $c_0=a_2/(2w_0^2)$、$c_2=-a_2/(6w_0^2)$。三阶 $\cos\zeta$ 项若不消去，会产生非周期的世俗项；其系数条件为
\[
-2w_0^2d=2a_2(c_0+c_2/2)+3a_3/4,
\]
于是
\[
d=-\frac{5a_2^2}{12w_0^4}-\frac{3a_3}{8w_0^2}
=\frac{(k+1)(2-k)}{24}.
\]
若所有小而有限振幅轨道均闭合，频率比不能随连续振幅改变，必须有 $d=0$。故只剩 $k=-1$ 或 $k=2$；对数势对应的比内展开取 $k=0$，其 $d\ne0$，亦被排除。这给出了所述正则情形中的筛选条件。

\paragraph{充分性：两个候选确实闭合。}
平方反比吸引势的非碰撞束缚轨道已证明为焦点椭圆。对各向同性谐振子，$\ddot x+\omega^2x=\ddot y+\omega^2y=0$，通过旋转平面坐标并选择时间原点，可写成 $x=A\cos\omega t,y=B\sin\omega t$，轨道是以力心为中心的椭圆；径向退化轨道为线段。因此两个候选均满足闭合要求。

\begin{physicsbox}
开普勒椭圆的长轴方向固定，比单纯的角动量守恒更强。在 $V=-k/r$ 中，还存在沿近心方向的守恒量 $\bA=\bp\times\bJ-mk\er$。其守恒可在章末习题直接检验。隐藏守恒量说明为何这两类势具有特别简单的轨道；更一般的中心势通常允许拱点进动。
\end{physicsbox}
\begin{modernbox}[title={近圆拱点进动的适用范围}]
对幂次势，$\Psi_0=\pi/\sqrt{k+2}$ 只是在近圆极限的半径向周期转角。从一个近心点到下一个近心点的方位增量约为 $2\Psi_0$。对接近开普勒势的运动，常用进动约定为 $\Delta\varpi\simeq2\Psi_0-2\pi$。这不是一般偏心轨道的精确公式，也不能直接用来计算广义相对论进动。

有限振幅的精确拱点角应从能量积分计算：
\begin{equation}\label{eq:apsidal-integral}
\Psi(E,J)=\int_{r_-}^{r_+}
\frac{|J|\,\dd r}{r^2\sqrt{2m[E-V(r)]-J^2/r^2}}.
\end{equation}
这里从近心到远心只积分一次；再乘2才是一个完整径向周期的转角。进一步阅读可用本章参考的中心力讲义，第\ref{sec:kepler-action}节再用径向作用变量讨论周期。
\end{modernbox}

\subsection{散射截面与卢瑟福散射}

把一束粒子打向靶，瞄得越偏，经过靶附近时通常偏转得越少。用入射位置的面积去对应出射角的范围，就得到散射截面的含义。

考虑质量为 $m$ 的粒子以初速 $v_0$、瞄准距 $b$ 射向力心。角动量和能量为
\begin{equation}
J=mv_0b,\qquad E=\frac12mv_0^2 .
\end{equation}
粒子在中心力场中偏转，出射方向与入射方向的夹角 $\Theta$ 称为\textbf{散射角}。所有瞄准距在 $b\sim b+\dd b$ 内的粒子，被散射到 $\Theta\sim\Theta+\dd\Theta$ 的立体角内。在轴对称、每个散射角只对应一个瞄准距的分支上，入射环带面积为 $\dd\sigma=2\pi b|\dd b|$，出射角环的立体角为 $\dd\Omega=2\pi\sin\Theta|\dd\Theta|$。两者相除，定义\textbf{微分散射截面}
\begin{equation}
\dv{\sigma}{\Omega}=\frac{b}{\sin\Theta}\left|\dv{b}{\Theta}\right|,
\label{eq:cross-section}
\end{equation}
若多个不同瞄准距产生同一散射角，微分截面应把各分支贡献相加。总截面为
\begin{equation}
\sigma=\int \dv{\sigma}{\Omega}\,\dd\Omega
=2\pi\int_0^\pi \frac{b}{\sin\Theta}\left|\dv{b}{\Theta}\right|\sin\Theta\,\dd\Theta .
\end{equation}

\begin{figure}[htbp]
\centering
\begin{tikzpicture}[>=Stealth,scale=1.05]
\coordinate (C) at (0.666667,0);
\draw[dashed,gray] (C)--++(-60:3.5);
\draw[dashed,gray] (C)--++(60:3.5);
\draw[dashed,gray] (C)--++(120:1.6);
\draw[very thick,blue!65!black,domain=-48:48,samples=160]
 plot ({cos(\x)/(2*cos(\x)-1)},{sin(\x)/(2*cos(\x)-1)});
\fill[orange!80!black] (0,0) circle (2.5pt) node[left]{力心};
\draw[->,thick,red!70!black] (2.07,-2.43)--(1.62,-1.65) node[right]{入射};
\draw[->,thick,red!70!black] (1.62,1.65)--(2.07,2.43) node[right]{出射};
\draw[->,thick] (C)++(60:.9) arc (60:120:.9);
\node at (.67,1.2) {$\Theta$};
\draw[<->] (0,0)--(.5,-.288675) node[midway,below left]{$b$};
\end{tikzpicture}
\caption{解析排斥轨道 $r=p/(e\cos\theta-1)$ 的实例，取 $p=1,e=2$。虚线为真正的渐近线；散射角比较入射、出射的速度方向，本例为 $60^\circ$。$b$ 是力心到入射渐近线的垂距。}
\label{fig:scattering}
\end{figure}

对于库仑排斥力
\begin{equation}
F(r)=\frac{k}{r^2},\qquad k=\frac{q_1q_2}{4\pi\varepsilon_0},
\end{equation}
以下假设固定重靶、点电荷、非相对论且无屏蔽，取 $k>0$。以近心点方向为极轴，比内方程的排斥解可写为
\[
r=\frac{p}{e\cos\theta-1},\qquad
p=\frac{J^2}{mk},\quad
e^2=1+\frac{2EJ^2}{mk^2}.
\]
无穷远的渐近角满足 $\cos\theta_\infty=1/e$。入射速度指向近心区、出射速度离开，因此速度偏转角为 $\Theta=\pi-2\theta_\infty$，不是直接把两条无向直线的小夹角当作偏转角。于是
\[
\cot(\Theta/2)=\tan\theta_\infty=\sqrt{e^2-1}
=\frac{mv_0^2b}{k},
\]
散射角与瞄准距满足
\begin{equation}
b=\frac{k}{mv_0^2}\cot\frac{\Theta}{2}
=\frac{k}{2E}\cot\frac{\Theta}{2}.
\label{eq:rutherford-b}
\end{equation}
对 $b=(k/2E)\cot(\Theta/2)$ 求导，
\[
\left|\frac{\dd b}{\dd\Theta}\right|=\frac{k}{4E}\csc^2(\Theta/2).
\]
代入式 \eqref{eq:cross-section}，用 $\sin\Theta=2\sin(\Theta/2)\cos(\Theta/2)$ 化简，得\textbf{卢瑟福散射公式}
\begin{equation}
\boxed{\;
\dv{\sigma}{\Omega}
=\left(\frac{k}{4E}\right)^2
\frac{1}{\sin^4(\Theta/2)}
=\left(\frac{q_1q_2}{16\pi\varepsilon_0E}\right)^2
\frac{1}{\sin^4(\Theta/2)}
\;}
\label{eq:rutherford}
\end{equation}
总截面发散，反映库仑力是长程力；实验报告通常给定非零的最小散射角，屏蔽也可改变小角行为。对 $\Theta\ge\Theta_{\min}>0$，积分截面为 $\sigma=\pi[k/(2E)]^2\cot^2(\Theta_{\min}/2)$，与 $\pi b_{\max}^2$ 一致。不能只凭“探测器有限大小”就消除包含零角的长程发散。

\begin{example}[月球探测器与微陨石的刚球散射]
把月球探测器与微陨石的碰撞简化为半径为 $a$ 的刚性球：由 $b=a\sin\varphi$、$\Theta=\pi-2\varphi$ 得
\[
b=a\cos\frac{\Theta}{2},\qquad
\dv{\sigma}{\Omega}
=\frac{b}{\sin\Theta}\left|\dv{b}{\Theta}\right|
=\frac{a^2}{4}.
\]
总截面 $\sigma=\pi a^2$，恰为几何截面积，可用于估算微陨石撞击探测器的概率。

在轴对称入射束中，入射环带截面是 $2\pi b|\dd b|$，对应的散射立体角是 $2\pi\sin\Theta\,\dd\Theta$。所需导数为
\[
\left|\frac{\dd b}{\dd\Theta}\right|=\frac a2\sin\frac\Theta2 .
\]
用 $\sin\Theta=2\sin(\Theta/2)\cos(\Theta/2)$ 约去半角因子，比值即为 $a^2/4$。将常数对 $4\pi$ 立体角积分即可复核总截面。
\end{example}

\begin{modernbox}
卢瑟福散射是经典力学最著名的“判决性实验”之一。大角度散射对短距离相互作用敏感，因而可以用来检验原子内部电荷分布模型。今天，散射截面仍是粒子物理、核物理和冷原子物理中探测相互作用的基本工具；量子力学中的卢瑟福公式在库仑势下与经典结果形式上一致。
\end{modernbox}

\subsection{潮汐力}

月球对地球不同位置的引力略有差别。先减去地心也受到的那一份，再比较近端和远端，就能分清潮汐作用与整体加速。

潮汐力不是一种新的基本力，而是引力场在有限大小物体上的\emph{差分}效应。设外引力场为 $\bg(\br)$，物体上相距 $\Delta\br$ 的两点加速度之差为
\begin{equation}
\Delta\ba=\bg(\br+\Delta\br)-\bg(\br)
\approx(\Delta\br\cdot\nabla)\bg .
\label{eq:tidal-def}
\end{equation}
对于点质量 $M$ 的引力场 $\bg=-GM\br/r^3$，先求雅可比 $\partial g_i/\partial x_j=GM(3x_ix_j/r^5-\delta_{ij}/r^3)$，再乘以 $\Delta x_j$，得到
\begin{equation}
\Delta\ba=\frac{GM}{r^3}\left[3(\Delta\br\cdot\hat{\br})\hat{\br}-\Delta\br\right].
\label{eq:tidal-acceleration}
\end{equation}
沿径向（$\Delta\br\parallel\hat{\br}$）有
\[
\Delta a_r=+\frac{2GM}{r^3}\Delta r,
\]
物体被\emph{拉伸}；沿垂直方向（$\Delta\br\perp\hat{\br}$）有
\[
\Delta a_\perp=-\frac{GM}{r^3}\Delta r_\perp,
\]
物体被\emph{压缩}。这正是潮汐隆起的几何来源：地球面向月球和背向月球的两侧都被“拉出去”，给出理想平衡潮汐的拉伸方向；真实海潮还受海盆、海岸及耗散动力学影响。

\begin{figure}[htbp]
\centering
\begin{tikzpicture}[>=Stealth,scale=1.0]
  \draw[blue!20,thick] (0,0) circle (1.5);
  \draw[dashed] (-2.4,0)--(4.2,0);
  \draw[gray!60,thick] (5.4,0) circle (0.55);
  \node at (5.4,-1.0) {月球};
  \node at (0,-1.9) {地球};
  \fill (1.5,0) circle (2pt);
  \fill (-1.5,0) circle (2pt);
  \fill (0,1.5) circle (2pt);
  \fill (0,-1.5) circle (2pt);
  \draw[->,very thick,red!75!black] (1.5,0)--++(0.85,0);
  \draw[->,very thick,red!75!black] (-1.5,0)--++(-0.85,0);
  \draw[->,very thick,blue!70!black] (0,1.5)--++(0,-0.8);
  \draw[->,very thick,blue!70!black] (0,-1.5)--++(0,0.8);
  \node[red!75!black] at (2.45,0.35) {拉伸};
  \node[blue!70!black] at (0.55,0.95) {压缩};
\end{tikzpicture}
\caption{潮汐力的几何图像。沿地月连线方向，引力差把地球两侧向外拉（拉伸）；垂直方向则向内压（压缩），形成两个潮汐隆起。}
\label{fig:tidal}
\end{figure}

潮汐力的大小随距离迅速衰减：
\begin{equation}
|\Delta\ba|\sim\frac{GM}{r^3}\,\Delta r .
\end{equation}
因此潮汐效应在近距离天体系统中特别显著。

\begin{itemize}[leftmargin=2.2em]
  \item \textbf{潮汐锁定}：潮汐隆起受到引力转矩，使天体自转角速度逐渐趋向公转角速度。月球已被地球潮汐锁定，以近乎同一面朝向地球，观测中的天平动仍会改变可见区域；潮汐耗散也会交换自转与轨道角动量。
  \item \textbf{洛希极限}：当小物体的自身引力不足以抵抗主星的潮汐拉伸时，它会被撕裂。设小物体半径为 $R$、密度为 $\rho$，主星质量为 $M$，则潮汐撕裂的临界距离约为
  \begin{equation}
  d_{\rm Roche}\sim R\left(\frac{2M}{m}\right)^{1/3}
  =\left(\frac{3M}{2\pi\rho}\right)^{1/3}
  \propto \rho^{-1/3}.
  \end{equation}
  对理想同步流体卫星，常用模型写成 $d_{\rm Roche}\simeq2.44R_M(\rho_M/\rho)^{1/3}$，其中 $R_M,\rho_M$ 为主星半径、平均密度。这个系数对应另一套流体平衡模型，不能直接乘到上述点对点拉伸估算式上。土星环、彗星碎裂等现象都与洛希极限有关。
  \item \textbf{潮汐加热}：若轨道偏心率较大，潮汐形变随时间变化，机械能转化为内能。木卫二冰壳下的海洋、木卫一的火山活动，都与潮汐加热密切相关。
\end{itemize}

\begin{modernbox}
在广义相对论中，潮汐力对应时空曲率引起的\emph{测地线偏离}。等效原理只能在一个自由落体参考系的局部邻域内消去引力加速度，但无法消去潮汐力；因此潮汐力是“引力的不可消除部分”，也是探测引力波的基本原理：引力波通过时，空间中的自由检验质量会发生差分振荡。
\end{modernbox}


\subsection{拉格朗日点与限制性三体问题}

如果在随两个主天体一起转的坐标系中看，有些位置上的小物体可以相对两者保持不动。找出这些平衡点后，还要另问轻轻推离会发生什么。

拉格朗日点是限制性三体问题中最著名的结果。设两个大质量天体 $M_1,M_2$ 绕共同质心作圆轨道，间距为 $a$，角速度为 $\Omega$，满足
\begin{equation}
\Omega^2=\frac{G(M_1+M_2)}{a^3}.
\end{equation}
第三个物体质量 $m\ll M_1,M_2$，它对两个大天体的引力影响可以忽略，但仍在它们的引力场中运动。这就是\textbf{限制性三体问题}。

取随两个大天体一起转动的参考系，令 $M_1$ 位于 $x=0$，$M_2$ 位于 $x=a$。总质量 $M=M_1+M_2$，质心位于 $x=a\mu$，其中
\begin{equation}
\mu=\frac{M_2}{M_1+M_2}.
\end{equation}
第三体的坐标为 $(x,y)$，到两个大天体的距离分别为
\begin{equation}
r_1=\sqrt{x^2+y^2},\qquad
r_2=\sqrt{(x-a)^2+y^2}.
\end{equation}
在转动系中，第三体受到两个引力、离心力和科里奥利力。科里奥利力不做功，不贡献势能；单位质量的\textbf{有效势能}为
\begin{equation}
\Phi(x,y)=-\frac{GM_1}{r_1}-\frac{GM_2}{r_2}
-\frac12\Omega^2\left[(x-a\mu)^2+y^2\right].
\label{eq:lagrange-effective}
\end{equation}
第三体的运动方程为
\begin{equation}
\ddot x-2\Omega\dot y=-\pdv{\Phi}{x},\qquad
\ddot y+2\Omega\dot x=-\pdv{\Phi}{y}.
\label{eq:lagrange-eom}
\end{equation}
平衡点满足 $\dot x=\dot y=0$、$\ddot x=\ddot y=0$，因此
\begin{equation}
\pdv{\Phi}{x}=0,\qquad
\pdv{\Phi}{y}=0.
\label{eq:lagrange-equilibrium}
\end{equation}
这些平衡点就是\textbf{拉格朗日点}。

\paragraph{共线拉格朗日点 $L_1,L_2,L_3$.}
先令 $y=0$。在 $M_1,M_2$ 之间的 $L_1$ 点，$0<x<a$，有
\[
\pdv{\Phi}{x}
=\frac{GM_1}{x^2}-\frac{GM_2}{(a-x)^2}
-\Omega^2(x-a\mu)=0.
\]
令 $x=au$，并利用 $\Omega^2=GM/a^3$，得到无量纲方程
\begin{equation}
\frac{1-\mu}{u^2}-\frac{\mu}{(1-u)^2}-(u-\mu)=0,
\qquad 0<u<1.
\label{eq:L1-equation}
\end{equation}
这是关于 $u$ 的代数方程，一般没有简单的闭式解，但可用数值方法求解。对小的质量比 $\mu$，近似有
\begin{equation}
u_{L1}\approx1-\left(\frac{\mu}{3}\right)^{1/3},\qquad
u_{L2}\approx1+\left(\frac{\mu}{3}\right)^{1/3},\qquad
u_{L3}\approx-1+\frac{7}{12}\mu .
\end{equation}
因此 $L_1$ 位于两天体之间、靠近较小天体；$L_2$ 位于较小天体外侧；$L_3$ 位于较大天体另一侧。三个共线点都在两体连线上。

\paragraph{三角拉格朗日点 $L_4,L_5$.}
令
\begin{equation}
x=\frac a2,\qquad y=\pm\frac{\sqrt3}{2}a .
\end{equation}
此时 $r_1=r_2=a$，即 $L_4,L_5$ 与 $M_1,M_2$ 构成等边三角形。直接代入式 \eqref{eq:lagrange-equilibrium} 验证：
\[
\pdv{\Phi}{x}
=\frac{GM_1(a/2)}{a^3}
+\frac{GM_2(-a/2)}{a^3}
-\Omega^2\left(\frac a2-a\mu\right)
=0,
\]
\[
\pdv{\Phi}{y}
=\frac{GM_1(\sqrt3a/2)}{a^3}
+\frac{GM_2(\sqrt3a/2)}{a^3}
-\Omega^2\frac{\sqrt3a}{2}
=0.
\]
所以 $L_4,L_5$ 确实是平衡点，它们位于两体连线的两侧，与两天体构成等边三角形。

\begin{center}
\begin{tikzpicture}[>=Stealth,scale=.82]
\fill[blue!6] (-2,0)--(0,3.464102)--(2,0)--cycle;
\fill[blue!6] (-2,0)--(0,-3.464102)--(2,0)--cycle;
\draw[dashed] (-6,0)--(4.2,0);
\draw[dashed] (-2,0)--(0,3.464102)--(2,0)--(0,-3.464102)--cycle;
\fill[blue!65!black] (-2,0) circle (4pt) node[below]{$M_1$};
\fill[orange!85!black] (2,0) circle (3pt) node[below]{$M_2$};
\fill[green!40!black] (0,3.464102) circle (2.5pt) node[above]{$L_4$};
\fill[green!40!black] (0,-3.464102) circle (2.5pt) node[below]{$L_5$};
\fill[red!70!black] (0.8361404400928096,0) circle (2pt) node[above]{$L_1$};
\fill[red!70!black] (3.438799331609326,0) circle (2pt) node[above]{$L_2$};
\fill[red!70!black] (-5.76643563428424,0) circle (2pt) node[above]{$L_3$};
\node at (-1,-4) {$\mu=0.1$，共线位置由平衡方程数值求得};
\end{tikzpicture}
\end{center}

\paragraph{稳定性分析.}
平衡点是否存在与是否稳定是两个问题。令 $x=x_0+\xi$、$y=y_0+\eta$，在平衡点附近线性化式 \eqref{eq:lagrange-eom}：
\begin{equation}
\ddot\xi-2\Omega\dot\eta
=-\Phi_{xx}\xi-\Phi_{xy}\eta,
\end{equation}
\begin{equation}
\ddot\eta+2\Omega\dot\xi
=-\Phi_{xy}\xi-\Phi_{yy}\eta,
\end{equation}
其中二阶偏导数在平衡点取值。设指数模态 $\xi,\eta\propto e^{\lambda t}$，代入后得到关于 $\xi,\eta$ 的齐次线性方程组。要有非零解，系数行列式必须为零：
\begin{equation}
\begin{vmatrix}
\lambda^2+\Phi_{xx} & -2\Omega\lambda+\Phi_{xy}\\
2\Omega\lambda+\Phi_{xy} & \lambda^2+\Phi_{yy}
\end{vmatrix}=0 .
\end{equation}
展开行列式，就得到特征方程
\begin{equation}
\lambda^4+\left(\Phi_{xx}+\Phi_{yy}+4\Omega^2\right)\lambda^2
+\Phi_{xx}\Phi_{yy}-\Phi_{xy}^2=0.
\label{eq:lagrange-characteristic}
\end{equation}

对共线点 $L_1,L_2,L_3$，令 $A=GM_1/r_1^3+GM_2/r_2^3$。在 $y=0$ 直接微分可得 $\Phi_{xy}=0$、$\Phi_{xx}=-2A-\Omega^2$、$\Phi_{yy}=A-\Omega^2$。在两主星之间，$r_1,r_2<a$，故 $A>G(M_1+M_2)/a^3=\Omega^2$。对外侧平衡点，记 $x_C=a\mu$，将 $\Phi_x=0$ 写成
\[
(A-\Omega^2)(x-x_C)
=G M_2(a-x_C)(r_2^{-3}-r_1^{-3}).
\]
右边用到了 $M_1x_C=M_2(a-x_C)$。右侧外点有 $r_2<r_1$ 且 $x-x_C>0$，左侧外点两处符号同时反转，因而两处也有 $A>\Omega^2$。于是
\begin{equation}
\Phi_{xx}\Phi_{yy}-\Phi_{xy}^2<0,
\end{equation}
把特征方程视为关于 $u=\lambda^2$ 的二次式。常数项为负，两个根必为异号实数，正的 $u$ 给出增长指数 $\lambda=+\sqrt u$，因此平衡不稳定。这正是为什么航天器不能长期“停”在 $L_1,L_2,L_3$ 上，而需要主动轨道保持。

对三角点，设 $\boldsymbol d_i$ 是从第 $i$ 个主星指向第三体的位移。引力势 $-GM_i/r_i$ 的二阶导矩阵为 $GM_i(I/r_i^3-3\boldsymbol d_i\boldsymbol d_i^{\mathsf T}/r_i^5)$。因 $r_1=r_2=a$，两项的单位矩阵部分恰与离心势的 $-\Omega^2I$ 抵消。将 $\boldsymbol d_1/a=(1/2,\pm\sqrt3/2)$ 和 $\boldsymbol d_2/a=(-1/2,\pm\sqrt3/2)$ 代入余下的外积，得
\begin{equation}
\Phi_{xx}=-\frac34\Omega^2,\qquad
\Phi_{yy}=-\frac94\Omega^2,\qquad
\Phi_{xy}=\mp\frac{3\sqrt3}{4}(1-2\mu)\Omega^2.
\end{equation}
代入式 \eqref{eq:lagrange-characteristic}，得到
\begin{equation}
\lambda^4+\Omega^2\lambda^2
+\frac{27}{4}\mu(1-\mu)\Omega^4=0.
\label{eq:L4-characteristic}
\end{equation}
令 $u=\lambda^2$，则
\begin{equation}
u^2+\Omega^2u+\frac{27}{4}\mu(1-\mu)\Omega^4=0.
\end{equation}
两根之和为 $-\Omega^2<0$，乘积为 $\frac{27}{4}\mu(1-\mu)\Omega^4>0$，若判别式严格为正且 $0<\mu<1$，则两个根为不同的负实数。临界时 $u=-\Omega^2/2$ 为重根。对应 $\lambda=\pm i\Omega/\sqrt2$ 时，前面的二阶模态矩阵行列式为零且非零（非对角元含非零虚部），故每个重根只有一个独立振型。一阶系统因此有含 $t e^{\lambda t}$ 的广义模态，临界扰动并非全都保持有界。判别式为
\begin{equation}
\Omega^4\left[1-27\mu(1-\mu)\right],
\end{equation}
故线性有界稳定的条件为
\begin{equation}
1-27\mu(1-\mu)>0.
\end{equation}
如果没有科里奥利力，$L_4,L_5$ 是有效势能的极大点，运动不稳定；正是科里奥利力在质量比足够小时把运动“拧”成有界振荡。
因此在线性近似中稳定的条件为
\begin{equation}
\mu<\frac12\left(1-\sqrt{\frac{23}{27}}\right)
\approx0.03852
\quad\text{或}\quad
\mu>0.96148 .
\end{equation}
对于 $M_1>M_2$ 的情形，这等价于
\begin{equation}
\frac{M_1}{M_2}>24.96 .
\end{equation}
太阳--木星系统、地球--月球系统都满足这一条件，这说明圆型限制性模型中三角点具有线性稳定性，但不能仅凭质量比就断言真实系统中的尘埃可长期聚集；而 $L_1,L_2,L_3$ 是不稳定的。

\begin{modernbox}
拉格朗日点在航天中有重要应用：日地 $L_1$ 点的太阳观测器可以连续观测太阳；日地 $L_2$ 点是詹姆斯·韦布空间望远镜和许多深空探测器的理想位置；地月 $L_2$ 点可用于月球背面通信中继。木星轨道上的特洛伊小行星群则聚集在太阳--木星 $L_4,L_5$ 附近。由于共线点不稳定，航天器通常需要在 $L_1,L_2$ 附近做周期性轨道保持，或运行在环绕该点的晕轨道上。
\end{modernbox}

\section{计算与模拟：用守恒量检验轨道}\label{sec:orbit-compute}
画出一条像椭圆的轨迹只是第一步。再比较不同时刻的能量和角动量，并把步长缩小重算，才能区分图形看起来正确与积分确实收敛。

对平方反比引力，把长度单位取为初始圆轨道半径、时间单位取为其倒角频率，方程成为 $\ddot{\boldsymbol x}=-\boldsymbol x/|\boldsymbol x|^3$。取 $\boldsymbol x_0=(1,0)$、$\boldsymbol v_0=(0,0.8)$，它是一条束缚椭圆，不是圆轨道。可采用速度 Verlet 格式：
\begin{align}
\boldsymbol x_{n+1}&=\boldsymbol x_n+h\boldsymbol v_n+\tfrac12h^2\boldsymbol a(\boldsymbol x_n),\\
\boldsymbol v_{n+1}&=\boldsymbol v_n+\tfrac h2[\boldsymbol a(\boldsymbol x_n)+\boldsymbol a(\boldsymbol x_{n+1})].
\end{align}
每一步计算 $\mathcal E=v^2/2-1/r$、$j=xv_y-yv_x$，与初始 $\mathcal E_0=-0.68,j_0=0.8$ 比较。理论上 $a=-1/(2\mathcal E_0)=25/34$、$e=0.36$、$r_p=8/17$、$r_a=1$。作图像“像椭圆”远远不够，还应检查转折点、能量误差和步长减半后的误差变化。

式 \eqref{eq:ham-verlet} 从分裂流给出同一个算法。算法不能消除近心点分辨率不足，也不能跨过 $r=0$ 奇点。教材随附的 \texttt{review/verify\_ch01.py} 使用上述初值比较两个步长，并保存误差报告；这是理想教学模型，不是任务轨道预报。

\section*{关键概念与常见误区}
\begin{itemize}[leftmargin=1.8em]
\item 车在匀速过弯时，速度表读数不变，但车仍有指向弯内的加速度。上抛物体到最高点时，速度瞬间为零，重力加速度仍在。
\item 先确定观察者跟谁一起运动，再选坐标。用随位置转动的极坐标基矢描述运动时，几何项来自基矢求导。
\item 物体在系统内部彼此推拉，内力总和可以为零，内力仍能做功。两辆小车黏在一起就是动量守恒而动能减少的例子。
\item 能量关系可先求出速率或能到达的范围。若还要求轨道的支持力，再写沿法线的牛顿方程。
\item 算角动量前先标矩心，并判断该点是否运动。算火箭时，先说明系统是否包含已经喷出的气体。
\item 圆轨道稳定只说明附近的小扰动受到限制。若要判断轨道是否闭合，还要比较径向运动和绕行的周期。
\end{itemize}

\section*{本章小结}
\addcontentsline{toc}{section}{本章小结}

本章以受力方程为起点，用守恒量、参考系变换和轨道几何组织解题：
\begin{itemize}[leftmargin=2.2em]
  \item 一个模型：质点；
  \item 两个方程：运动学方程 $\br(t)$ 与动力学方程 $m\ddot{\br}=\bF$；
  \item 三个守恒律：动量、角动量、机械能；
  \item 一个核心问题：中心力场中的轨道与散射；
  \item 一个关键视角：非惯性系中的惯性力与太空电梯；
  \item 一个推广：质点系与质心运动、两体约化质量与位力定理；
  \item 三个现代延伸：伯特兰定理、潮汐力与拉格朗日点。
\end{itemize}

\begin{center}
\renewcommand{\arraystretch}{1.5}
\begin{tabular}{ll}
\toprule
物理量 & 表达式 \\
\midrule
速度 & $\bv=\dot{\br}=\dot r\er+r\dot\theta\etheta$ \\
加速度（极坐标） & $\ba=(\ddot r-r\dot\theta^2)\er+(r\ddot\theta+2\dot r\dot\theta)\etheta$ \\
自然坐标加速度 & $\ba=\dot v\et+\dfrac{v^2}{\rho}\en$ \\
转动系加速度 & $\ba=\ba'+2\bom\times\bv'+\bom\times(\bom\times\br)+\dot{\bom}\times\br$ \\
动量定理 & $\displaystyle\int_{t_1}^{t_2}\bF\,\dd t=\bp_2-\bp_1$ \\
角动量定理 & $\dfrac{\dd\bJ}{\dd t}=\br\times\bF$ \\
机械能守恒 & $T+V=E$ \\
比内方程 & $\dfrac{\dd^2u}{\dd\theta^2}+u=-\dfrac{mF(1/u)}{J^2u^2}$ \\
有效势能 & $V_{\rm eff}=V(r)+\dfrac{J^2}{2mr^2}$ \\
圆轨道条件 & $F(r_0)=-\dfrac{J^2}{mr_0^3}$ \\
卢瑟福截面 & $\dfrac{\dd\sigma}{\dd\Omega}=\left(\dfrac{q_1q_2}{16\pi\varepsilon_0E}\right)^2\dfrac{1}{\sin^4(\Theta/2)}$ \\
伯特兰定理 & $F=-\dfrac{k}{r^2}$ 或 $F=-kr$ \\
潮汐加速度 & $\Delta\ba=\dfrac{GM}{r^3}\left[3(\Delta\br\cdot\hat{\br})\hat{\br}-\Delta\br\right]$ \\
\bottomrule
\end{tabular}
\end{center}

\section*{分层习题}
\addcontentsline{toc}{section}{分层习题}
先在草图上标出参考系、正方向和力的作用点。若从守恒量入手，先说明研究的是单个质点还是整个系统。

以下除特别说明外使用理想模型，给定数值是练习参数。
\begin{bookexercises}
\item \exlevel{基础} \textbf{基础理解：旋转基矢。} 不引用极坐标加速度公式，从 $\br=r(\cos\theta,\sin\theta)$ 微分两次得到它。用 $r=R$、$\dot\theta=\omega$ 检查向心项的方向和量纲。
\item \exlevel{基础} \textbf{基础理解：力场定义域。} 对连续径向力 $F(r)\er$ 在 $r>0$ 写出全局径向势；再比较正文去轴无旋力场为何没有全局单值势。不能将后者的反例误解为所有去原点的中心势都不保守。
\item \exlevel{基础} \textbf{独立推导：大振幅摆。} 把周期积分用 $\sin(\theta/2)=k\sin\phi$、$k=\sin(\theta_0/2)$ 化为标准积分，推到 $T/T_0=1+\theta_0^2/16+O(\theta_0^4)$。讨论 $\theta_0\to\pi$。
\item \exlevel{基础} \textbf{基本应用：静摩擦。} 漏斗母线与水平夹角为 $0<\alpha<\pi/2$，在距轴 $r>0$ 处的物块随漏斗以 $\omega$ 匀速转动，静摩擦系数为 $\mu_s$。求静止接触条件，说明什么时候存在转速上界。
\item \exlevel{综合} \textbf{综合分析：有限深势阱。} 质点在 $V=-V_0e^{-r^2/a^2}$（$V_0,a>0$）中运动，角动量 $J\ne0$。求圆轨道的角动量上界及径向稳定性。特别区分严格稳定轨道与边界退化轨道，不能要求最大角动量处仍有严格稳定圆轨道。
\item \exlevel{综合} \textbf{势能图。} 以 $X=x/\ell$、$U=V/V_*$ 无量纲化，$U=X^4-2X^2+0.4$。求能量 $E/V_*=0.2$ 时的运动区间与全部平衡点；求两个势阱底的小振动频率。
\item \exlevel{综合} \textbf{地球自转。} 在赤道将物体以速度 $u$ 竖直上抛。$x$ 向东、$y$ 向北、$z$ 向上，忽略空气阻力并保留 $\Omega$ 的一次项，求返回初始高度时的东西偏移。说明为何不能只分析下落半程。
\item \exlevel{综合} \textbf{傅科摆。} 令 $f=\Omega\sin\lambda$，从正文两条水平运动方程出发，令 $w=x+iy$，再令 $w=e^{-ift}z$，求进动角速度的符号与周期。
\item \exlevel{综合} \textbf{潮汐。} 用点质量外场雅可比求地球近、远两端相对地心的潮汐加速度，并求其差。已知地球半径 $R_E$、地月距 $d\gg R_E$、月球质量 $M_m$。解释相对加速度与绝对引力大小的区别。
\item \exlevel{综合} \textbf{质心与开放系统。} (a) 求均匀半球壳质心。(b) 火箭有效排气速率 $v_e=3.5\,\mathrm{km/s}$，无外力下需 $\Delta v=8.0\,\mathrm{km/s}$，求质量比，并说明能量是否可只对火箭本体闭合守恒。
\item \exlevel{综合} \textbf{位力模型。} 均匀自引力球质量 $M$、半径 $R$，有内部各向同性速度弥散 $\sigma_{\rm los}$。从逐层组装球的势能积分推导 $V=-3GM^2/(5R)$，在位力平衡下求 $M(R,\sigma_{\rm los})$。是否可在未给可见质量和误差时直接声称发现暗物质？
\item \exlevel{综合} \textbf{碰撞。} 质量 $m$ 的粒子以 $u$ 撞上静止的同质量粒子，恢复系数 $e_r=1/2$。求两末速度及初始动能损失比例，并在质心系复算。
\item \exlevel{综合} \textbf{综合守恒。} 将机器人质量推广为 $m$、圆环质量推广为 $M$，均匀细环半径为 $R$。机器人相对环爬行一周，求环转角；再讨论 $M\gg m$ 极限。
\item \exlevel{拓展} \textbf{轨道设计。} 用正文给定月球参数重算两个点火增量和转移时间，再令 $r_2\to r_1$ 检查极限。说明两次点火时动量、能量的变化与中间弹道段有何区别。
\item \exlevel{拓展} \textbf{隐藏守恒量。} 对 $V=-k/r$，从 $\bA=\bp\times\bJ-mk\er$ 出发逐项求导，证明它守恒。该论证在哪些地方使用了力心固定与 $J$ 守恒？
\item \exlevel{拓展} \textbf{计算探索。} 按本节 Verlet 算法积分十个理论周期，比较 $h=0.01,0.005$ 的最大相对能量误差及角动量误差；记录每一步到力心的最小距离。解释为什么这些检查比只展示轨迹可靠。
\end{bookexercises}

\section*{习题解答}
先用同一题号找到解答。核对时，从第一条方程开始比较自己的做法，再检查最后的符号、单位和适用条件。
\begin{bookexercises}
\item 由 $\er=(\cos\theta,\sin\theta)$、$\etheta=(-\sin\theta,\cos\theta)$ 得 $\dot{\er}=\dot\theta\etheta$、$\dot{\etheta}=-\dot\theta\er$。先求 $\bv=\dot r\er+r\dot\theta\etheta$，再逐项微分，
\[
\ba=(\ddot r-r\dot\theta^2)\er+(r\ddot\theta+2\dot r\dot\theta)\etheta.
\]
匀速圆周时 $\dot r=\ddot r=\ddot\theta=0$，只剩指向圆心的 $-R\omega^2\er$，$R\omega^2$ 的量纲为长度除以时间平方。
\item $V(r)=-\int_{r_*}^{r}F(s)\dd s+C$，在可达的径向区间单值。去轴涡旋场则沿绕轴闭路做功 $2\pi K$，不能用单值径向积分消除。

由 $\nabla r=\er$ 可得 $-\nabla V=F(r)\er$，对任意闭路的积分为 $-\oint\dd V=0$。涡旋场取半径 $R$ 的正向圆周，切向力为 $K/R$，弧元为 $R\dd\varphi$，绕一周得到 $2\pi K$。
\item 恒等式 $\cos\theta-\cos\theta_0=2[k^2-\sin^2(\theta/2)]$ 给出 
\[T=4\sqrt{l/g}\int_0^{\pi/2}(1-k^2\sin^2\phi)^{-1/2}\dd\phi\]
展开积分，$T/T_0=1+k^2/4+\cdots=1+\theta_0^2/16+\cdots$。$\theta_0\to\pi$ 时积分发散，非小振动极限。

变量替换给出 $\dd\theta=2k\cos\phi\,\dd\phi/\sqrt{1-k^2\sin^2\phi}$，同时 $\sqrt{k^2-\sin^2(\theta/2)}=k\cos\phi$。展开后用 $\int_0^{\pi/2}\sin^2\phi\,\dd\phi=\pi/4$。当 $k=1$ 时被积函数为 $1/\cos\phi$，在上端点呈对数发散。
\item 取沿坡向外上升为切向正向。转动系中有效力为 $m\omega^2r$ 向外与 $mg$ 向下，故 \[
N=m(g\cos\alpha+\omega^2r\sin\alpha),\qquad
F_t=m(g\sin\alpha-\omega^2r\cos\alpha).
\]令 $q=\omega^2r/g$，条件是
\[
|\sin\alpha-q\cos\alpha|\le\mu_s(\cos\alpha+q\sin\alpha).
\]
下界为 
\[q\ge\max\{0,(\sin\alpha-\mu_s\cos\alpha)/(\cos\alpha+\mu_s\sin\alpha)\}\]
若 $\mu_s<\cot\alpha$，还需 
\[q\le(\sin\alpha+\mu_s\cos\alpha)/(\cos\alpha-\mu_s\sin\alpha)\]
否则无有限上界。$\mu_s\ge\tan\alpha$ 时允许 $q=0$。
\item 令 $x=r_0^2/a^2$，圆轨道条件为 $J^2=2mV_0a^2x^2e^{-x}$。右端在 $x=2$ 达最大，故 $|J|\le J_{\max}=\sqrt{8mV_0a^2}\,e^{-1}$。用圆轨道条件求二阶导数得到 $V_{\rm eff}''=4V_0e^{-x}(2-x)/a^2$：$x<2$ 稳定，$x>2$ 不稳定。$x=2$ 为退化临界，不能称为严格稳定。严格稳定圆轨道只取 $|J|<J_{\max}$ 的内支。

由 $V_{\rm eff}'=2V_0re^{-r^2/a^2}/a^2-J^2/(mr^3)=0$ 得到上式。$x^2e^{-x}$ 的导数为 $xe^{-x}(2-x)$。固定 $J=J_{\max}$ 再对原有效势求导，在退化半径有 $V_{\rm eff}'''=-8V_0r_0e^{-2}/a^4\ne0$，因此这里也不是局部极小。
\item 由 $E-V\ge0$ 得 $(X^2-1)^2\le0.8$。令 $b_\pm=\sqrt{1\pm\sqrt{0.8}}$，两个允许区间为 $[-b_+,-b_-]$ 和 $[b_-,b_+]$。平衡由 $U'=4X(X^2-1)=0$ 给出。$U''=12X^2-4$ 在 $X=0$ 为负，在 $X=\pm1$ 为8，故两阱底的频率均为 $\sqrt{8V_*/(m\ell^2)}$。
\item 零阶 $z=ut-gt^2/2$，一阶 $\ddot x=-2\Omega\dot z$。利用 $\dot x(0)=0$，得 $\dot x=-2\Omega z$、$x=-\Omega ut^2+\Omega gt^3/3$。在 $t=2u/g$，$x=-4\Omega u^3/(3g^2)$，向西。

上升、下落时 $\ddot x$ 反向，返回初始高度时 $\dot x$ 恰好恢复为零，但中间始终有 $z>0$、$\dot x<0$，所以两段位移同向累加。
\item $\ddot w+2if\dot w+\omega_0^2w=0$，替换后 $\ddot z+(\omega_0^2+f^2)z=0$。在线性模型内快速摆动的外包方向以 $-f$ 转动；北半球从上方看顺时针。$T=2\pi/|f|$；物理上采用小摆角且 $|f|\ll\omega_0$。

实际代入时 $\dot w=e^{-ift}(\dot z-ifz)$、$\ddot w=e^{-ift}(\ddot z-2if\dot z-f^2z)$，一阶导数项相消。所给周期按方位转满 $2\pi$ 计，若只辨认无向直线，则转过 $\pi$ 即重合。赤道 $f=0$，该模型中没有进动。
\item 近、远端相对地心的加速度分别沿地月连线向外，大小均近似 $2GM_mR_E/d^3$，矢量之差大小为 $4GM_mR_E/d^3$。背月端绝对引力仍指向月球，只是比地心受到的引力小。

设 $\boldsymbol n$ 从地心指向月球。对外场 $-GM_m\boldsymbol s/|\boldsymbol s|^3$ 关于质点位置求导，在地心 $\boldsymbol s=-d\boldsymbol n$，得到 $GM_m(3\boldsymbol n\boldsymbol n^{\mathsf T}-\mathsf I)/d^3$。分别乘以 $\pm R_E\boldsymbol n$ 就得到两端的相对加速度。
\item (a) 球面环带质量为常数乘 $\dd z$，$z\in[0,R]$，故 $z_G=R/2$。(b) $m_0/m_f=e^{8/3.5}\simeq9.83$。能量来自推进剂，且喷气携带能量；火箭本体是开放系统。

以 $z=R\cos\theta$ 参数化球面，环带面积 $2\pi R^2\sin\theta\,|\dd\theta|=2\pi R|\dd z|$。因此
\[
z_G=\frac{\int_0^R z\,\dd z}{\int_0^R\dd z}=\frac R2.
\]火箭方程 $\dd v=-v_e\dd m/m$ 从 $m_0$ 积分到 $m_f$ 即给出质量比。
\item 球内质量 $M(r)=M(r/R)^3$，$\dd M=3Mr^2\dd r/R^3$，因此 
\[V=-\int_0^R GM(r)\dd M/r=-3GM^2/(5R)\]
代入位力关系 $2T+V=0$ 及 $T=3M\sigma_{\rm los}^2/2$ 得 $M=5R\sigma_{\rm los}^2/G$。结论依赖均匀球和位力平衡等条件，缺少可见质量与误差时不能作所问的观测判断。
\item 沿入射方向取正。动量与恢复条件给 $v_1+v_2=u$、$v_2-v_1=u/2$，所以 $v_1=u/4,v_2=3u/4$。动能比为 $(v_1^2+v_2^2)/u^2=5/8$。在质心系先减去 $u/2$，初速为 $(u/2,-u/2)$，末速为 $(-u/4,u/4)$。质心系相对动能从 $mu^2/4$ 变为 $mu^2/16$，损失 $3mu^2/16$，占实验室系初动能 $mu^2/2$ 的 $3/8$。
\item 机器人与环心相对总质心的距离分别为 $MR/(M+m)$ 与 $mR/(M+m)$。总角动量为 $MR^2\dot\varphi+[Mm/(M+m)]R^2\dot\alpha=0$，配合 $\alpha=\varphi+\theta$ 得 $\Delta\varphi=-2\pi m/(M+2m)$。$M=m$ 时为 $-2\pi/3$，$M\gg m$ 时约为 $-2\pi m/M$。

代入 $\dot\alpha=\dot\varphi+\dot\theta$ 并约去 $MR^2/(M+m)$，有 $(M+2m)\dot\varphi=-m\dot\theta$。从 $\Delta\theta=2\pi$ 直接积分。
\item 取 $r_1=1837.4\,\mathrm{km}$、$r_2=2737.4\,\mathrm{km}$，$\mu=4902.800118\,\mathrm{km^3/s^2}$，转移椭圆的半长轴 $a=(r_1+r_2)/2=2287.4\,\mathrm{km}$。由单位质量能量 $v^2/2-\mu/r=-\mu/(2a)$，得
\begin{align*}
\Delta v_1&=\sqrt{\mu(2/r_1-1/a)}-\sqrt{\mu/r_1}
            =153.47\,\mathrm{m/s},\\
\Delta v_2&=\sqrt{\mu/r_2}-\sqrt{\mu(2/r_2-1/a)}
            =138.84\,\mathrm{m/s},\\
t_{\rm tr}&=\pi\sqrt{a^3/\mu}=4908.4\,\mathrm{s}.
\end{align*}
$r_2\to r_1$ 时两个速度差趋零，形式上的转移时间趋于圆轨道半周期。两次点火通过喷气改变航天器动量和机械能，中间的弹道段则保持机械能和关于力心的角动量。
\item $\dot{\bJ}=0$，$\dot\bp=-k\er/r^2$，且 $\er\times\bJ=mr(\dot r\er-\bv)$，所以 
\[\dot\bp\times\bJ=(mk/r)(\bv-\dot r\er)=mk\dot{\er}\]
与另一项相消。力心运动或附加非中心力时此证明失效。
\item 沿用正文初值，用初始能量、角动量的绝对值归一化每个节点的误差，再分别取最大值。为恰好到达十个理论周期，取 $N=\lceil 10T/h\rceil$，实际步长为 $10T/N$。
\[
\begin{array}{c|r|ccc}
h & N & \max|\Delta E/E_0| & \max|\Delta J/J_0| & \min r\\ \hline
0.01 &3962&2.12975\times10^{-4}&3.89\times10^{-15}&0.470678\\
0.005&7924&5.32759\times10^{-5}&9.02\times10^{-15}&0.470611
\end{array}
\]
能量误差之比为 $3.998$，符合二阶收敛。中心力速度更新不产生角冲量，位置更新沿动量方向，故角动量误差在舍入量级。最小半径可用来检查近心段是否被合理解析，还可与解析近心距比较。即使守恒量误差很小，轨道相位仍可能累积误差。
\end{bookexercises}

\section*{从牛顿力学走向变分法}
本章展示了矢量受力方程、守恒律与几何方法怎样相互补充。但对复杂约束系统，逐个求约束反力仍然繁琐。第\ref{ch:variation}章先研究“把整条路径作为自变量”的变分法：什么叫相邻路径，固定端点意味着什么，驻值为何不必为最小值。第\ref{sec:action-first}节在固定时空端点的一维保守模型中重现牛顿方程，再由第\ref{ch:lagrange}章推广。新方法必须重现本章的运动方程与守恒结论，而不会取消它们。

\section*{进一步阅读与来源}
链接复核：2026-11-08。月球轨道例题采用公开引力参数及平均半径；其余未标明观测来源的例题数值均为教学设定。
\begin{bookreading}
\item MIT 8.01SC，\textit{Chapter 15: Collision Theory}，含质心系与碰撞能量分解：\href{https://ocw.mit.edu/courses/8-01sc-classical-mechanics-fall-2016/mit8_01scs22_chapter15.pdf}{MIT公开碰撞讲义}。
\item NASA/JPL，\textit{Astrodynamic Parameters}：\href{https://ssd.jpl.nasa.gov/astro_par.html}{JPL天体动力学参数表}；\textit{Planetary Satellite Physical Parameters}：\href{https://ssd.jpl.nasa.gov/sats/phys_par/sep.html}{JPL卫星物理参数表}。本章采用月球 $GM=4902.800118\,\mathrm{km^3/s^2}$、平均半径 $1737.4\,\mathrm{km}$，并在数值计算中转换到SI单位。
\item NASA，\textit{Electric Propulsion Hits Its Stride in Space}：\href{https://www.nasa.gov/organizations/osi/electric-propulsion-hits-its-stride-in-space/}{NASA电推进介绍}，用于区分单粒子漂移模型与实际霍尔推进器的电离、输运等过程。
\item NASA，\textit{Webb's Orbit at Sun-Earth Lagrange Point 2 (L2)}：\href{https://science.nasa.gov/asset/webb/webbs-orbit-at-sun-earth-lagrange-point-2-l2/}{NASA韦布轨道示意}，说明航天器运行于拉格朗日点附近的轨道，而非静止在该点。
\item Richard Fitzpatrick，得克萨斯大学公开天体力学讲义 \textit{Motion in nearly circular orbit}：\href{https://farside.ph.utexas.edu/teaching/celestial/Celestial/node43.html}{近圆轨道与拱点角}，用于进一步学习本章的近圆分析。
\item NASA/JPL，\textit{Fundamentals of Electric Propulsion, Second Edition}：\href{https://descanso.jpl.nasa.gov/SciTechBook/SciTechBook.html}{电推进专著}，进一步区分霍尔漂移、电离与阴极中和。
\end{bookreading}

\chapter{变分法}\label{ch:variation}

\begin{guidebox}
从两条连接同一对端点的路线开始。你可以比较它们的长度，也可以比较走完全程需要多久。每条路线都算出一个数，这种用整条函数作输入的规则叫泛函。把其中一条路线稍微挪一挪，比较所得数值怎样改变，就进入了变分法。

先用折射练习只有一个待选点的情形，再把这套做法推广到整条曲线。每次算完一阶变分，记得检查两件事。端点是否允许移动？算出的路径能否进一步证明是极小？第三章会用同一方法从作用量求运动方程。
\end{guidebox}

\section{变分思想：从费马原理开始}
两点之间可以画很多条路。先从光在不同介质中传播的时间开始，看看怎样把选择整条路的问题化成可计算的问题。

\subsection{先说清楚比较哪些路径}
可以给每条路线算一个总长度。改变路线时，先检查两端是否仍接在指定地点，途中是否还满足给定条件。

普通函数 $f(z)$ 的自变量是数，泛函（functional）的自变量是一整条函数。例如固定两端的曲线 $y(x)$，其长度为
\[
\mathcal L[y]=\int_a^b\sqrt{1+y'^2}\,\dd x.
\]
不同曲线给出不同长度，即使它们在端点取相同数值。要问哪一条最短，必须先规定\textbf{容许路径}：定义区间、边界值、光滑性、是否允许自交，以及必须满足的约束。不同的容许集合可能产生不同的答案。

本章先讨论标量图形 $y=y(x)$，通常取 $y\in C^2[a,b]$，扰动取满足端点条件的光滑函数；被积函数在路径邻域具有所需的连续偏导数。含尖点、无穷斜率或分段介质的问题，要单独说明分段求解和极限步骤。参数化曲线适用于不能由单个 $y(x)$ 覆盖的情形。

\subsection{费马原理及其适用范围}
光从空气进入另一种介质时，两段传播速度可能不同。即使把路线拉直会缩短距离，也还要比较总用时。费马原理要求小改路径时，传播时间的一阶变化为零。

在静止、各向同性介质的几何光学近似下，设折射率为正函数 $n(\br)$，局部光速为 $c/n$。沿路径的传播时间与光程为
\begin{equation}\label{eq:fermat}
T[\gamma]=\frac1c\int_\gamma n(\br)\,\dd s,\qquad
\mathcal L_{\rm opt}[\gamma]=\int_\gamma n\,\dd s.
\end{equation}
对固定空间端点的邻近容许路径，费马原理给出
\begin{equation}\label{eq:fermat-variation}
\delta T=0\quad\Longleftrightarrow\quad\delta\mathcal L_{\rm opt}=0.
\end{equation}
这里暂把 $\delta$ 理解为“一阶路径变化”；严格定义将在下一节给出。驻值不预先保证全局最短时间，存在多条光线、焦点或反射约束时尤其如此。
\begin{physicsbox}
比较路径是理论中的操作，不是光线预先知道终点。在短波长近似中，邻近光学路径的相位由光程决定；远离驻相区域的贡献容易相消。这个图像说明几何光学与波动光学的联系，但不能用一句“自然总在优化”代替对容许路径和适用极限的规定。
\end{physicsbox}

\subsection{折射定律：先做一个普通极值问题}
固定光源和终点，先假定每种均匀介质内的光线都是直线。此时唯一要选的是过界面的位置，可以直接对该位置求导。

\begin{example}[固定两端的平面折射]
界面为 $y=0$，$A=(0,h_1)$、$B=(d,-h_2)$，$h_1,h_2,d>0$，两侧均匀折射率为 $n_1,n_2>0$。先在各均匀介质中取直线段，令折射点 $P=(\xi,0)$，则
\[
\mathcal L(\xi)=n_1\sqrt{\xi^2+h_1^2}+n_2\sqrt{(d-\xi)^2+h_2^2}.
\]
求导得到
\[
\mathcal L'(\xi)=n_1\frac{\xi}{\sqrt{\xi^2+h_1^2}}
-n_2\frac{d-\xi}{\sqrt{(d-\xi)^2+h_2^2}}.
\]
定义两侧光线与法线的夹角为 $\theta_1,\theta_2$，驻值条件为
\[
\boxed{n_1\sin\theta_1=n_2\sin\theta_2.}
\]
而
\[
\mathcal L''(\xi)=\frac{n_1h_1^2}{(\xi^2+h_1^2)^{3/2}}
+\frac{n_2h_2^2}{((d-\xi)^2+h_2^2)^{3/2}}>0.
\]
因此本例在折射点这一变量上严格凸，驻点确为唯一最小值。并且 $\mathcal L'(0)<0,\mathcal L'(d)>0$，可知它在两端垂足之间。$n_1=n_2$ 时两段接成同一直线，是结果的直接检验。
\end{example}
\begin{figure}[htbp]
\centering
\begin{tikzpicture}[>=Stealth]
\fill[blue!4] (-2.5,0) rectangle (2.7,2.2);
\fill[blue!12] (-2.5,0) rectangle (2.7,-2.2);
\draw[thick] (-2.5,0)--(2.7,0);
\coordinate (A) at (-1.5,2);\coordinate (P) at (0,0);\coordinate (B) at (1,-2);
\draw[very thick,red!70!black] (A)--(P)--(B);
\draw[dashed] (0,-2.3)--(0,2.3);
\draw (0,.65) arc (90:126.869898:.65);
\draw (0,-.65) arc (-90:-63.434949:.65);
\node at (-.45,.8) {$\theta_1$};\node at (.4,-.85) {$\theta_2$};
\fill (A) circle (2pt) node[above]{$A$};\fill (P) circle (2pt) node[below left]{$P$};
\fill (B) circle (2pt) node[below]{$B$};
\node at (2,1) {$n_1=1$};\node at (2,-1) {$n_2=3/\sqrt5$};
\end{tikzpicture}
\caption{按斯涅尔关系选取的折射示例。两角都从同一界面法线测量；图形斜率与标出的折射率相容。}
\label{fig:var-refraction}
\end{figure}

\subsection{从有限个折射点走向一条未知曲线}
若有许多薄层，每层界面的穿过位置都可以调。层数越来越多时，这一串可调的点便用一条连续曲线来描述。

若把连续变化的介质切成许多薄层，未知量变成每个界面的折射点；层数趋于无穷时，未知量成为整个函数。最速降线也有相同结构：取 $y$ 向下、出发点为原点，均匀重力下无摩擦滑动且初速为零，则由第一章能量守恒，$v=\sqrt{2gy}$。在 $y>0$ 的曲线段上，
\begin{equation}\label{eq:brachistochrone-functional}
T[y]=\int_0^{x_B}\frac{\sqrt{1+y'^2}}{\sqrt{2gy}}\,\dd x.
\end{equation}
最短几何长度与最短时间权衡不同。出发点 $y=0$ 使被积函数奇异，因此不能不加说明地把下面的光滑定理直接用于端点；本章稍后将在内部推导，再检查时间积分的可积性。

\section{泛函、一阶变分与欧拉--拉格朗日方程}
让原曲线稍微鼓起一点，看看整条路径的评分怎样变。把变化的一次项写清楚，就能分别得到曲线内部和端点应满足的条件。

\subsection{有限差、变分与时间位移的区别}
把鼓起的形状固定，只用一个小系数控制鼓起的程度。有限差比较两条已经画好的曲线，一阶变分则看这个系数从零开始改变时，评分变得多快。

考虑
\begin{equation}\label{eq:functional}
I[y]=\int_a^b F(x,y,y')\,\dd x,\qquad y(a)=y_A,\quad y(b)=y_B.
\end{equation}
$F$ 称为被积函数；这里的 $x$ 是独立变量，未必是位置，也可以是时间。取一条试探函数 $y$，选定扰动形状 $\eta(a)=\eta(b)=0$，组成路径族
\[
y_\epsilon(x)=y(x)+\epsilon\eta(x).
\]
$\epsilon$ 是标记路径的小参数。严格定义\textbf{一阶变分}为对参数的方向导数：
\begin{equation}\label{eq:var-definition}
\delta I[y;\eta]=\left.\frac{\dd}{\dd\epsilon}I[y_\epsilon]\right|_{\epsilon=0},
\qquad \delta y=\eta,\quad \delta y'=\eta'.
\end{equation}
本章约定 $\delta y$ 表示单位参数的变分方向，不把 $\epsilon$ 重复包含其中。因此有限变化满足
\[
I[y+\epsilon\eta]-I[y]=\epsilon\delta I[y;\eta]
+\tfrac12\epsilon^2\delta^2I[y;\eta]+o(\epsilon^2).
\]
\textbf{驻值}意味着对所有容许方向都有 $\delta I[y;\eta]=0$；并非要求所有邻近曲线都给出完全相等的泛函值。
\begin{figure}[htbp]
\centering
\begin{tikzpicture}[>=Stealth]
\draw[->] (-.3,0)--(5.5,0) node[right]{$x$};\draw[->] (0,-.2)--(0,2.7) node[above]{$y$};
\draw[very thick,blue!70!black,domain=0:5,samples=100] plot (\x,{.6+.2*\x+.25*sin(36*\x)});
\draw[dashed,red!70!black,domain=0:5,samples=100] plot (\x,{.6+.2*\x+.85*sin(36*\x)});
\draw[<->] (2.5,1.35)--(2.5,1.95);
\draw[thin] (2.5,1.95)--(2.5,2.2);
\node[above] at (2.5,2.2) {$\epsilon\eta(x_*)$};
\fill (0,.6) circle (2pt);\fill (5,1.6) circle (2pt);
\draw[dotted] (2.5,0)--(2.5,1.3);\node[below] at (2.5,0) {$x_*$};
\node[blue!70!black] at (4.4,1.15) {$y$};\node[red!70!black] at (3.8,2.3) {$y_\epsilon$};
\end{tikzpicture}
\caption{在同一个独立变量 $x_*$ 处比较两条路径，竖直差为 $\epsilon\eta(x_*)$；两端相同。变分不是沿原路径从 $x_*$ 走到 $x_*+\dd x$ 的真实位移。}
\label{fig:var-neighbor-paths}
\end{figure}

\subsection{从链式法则到边界项}
曲线的高度和斜率都可能影响积分。把斜率变化那一项分部积分后，一项留在端点，另一项和高度变化一起留在积分里。

对参数 $\epsilon$ 求导，并在光滑条件下交换微分与积分：
\[
\delta I=\int_a^b\left(F_y\eta+F_{y'}\eta'\right)\dd x.
\]
其中各偏导沿未扰动路径 $y(x)$ 取值。$y,y'$ 在求偏导时作为独立槽位；之后才代入 $y'=\dd y/\dd x$。关键的一步是分部积分：
\[
\int_a^b F_{y'}\eta'\dd x
=[F_{y'}\eta]_a^b-\int_a^b\frac{\dd F_{y'}}{\dd x}\eta\dd x.
\]
因此完整的一阶变分为
\begin{equation}\label{eq:variation-boundary}
\delta I=[F_{y'}\eta]_a^b+
\int_a^b\left(F_y-\frac{\dd F_{y'}}{\dd x}\right)\eta\dd x.
\end{equation}
先保留边界项，再根据端点条件处理。固定端点给出 $\eta(a)=\eta(b)=0$，从而
\begin{equation}\label{eq:first-variation}
\delta I=\int_a^b\left(F_y-\frac{\dd F_{y'}}{\dd x}\right)\eta\dd x.
\end{equation}
若括号是连续函数且对所有内部光滑紧支扰动积分均为零，\textbf{变分法基本引理}保证括号逐点为零。理由是：若某点处括号为正，则附近仍为正，取支撑在该邻域的非负扰动即可得到正积分，与假设矛盾；负值同理。故驻值路径满足
\begin{equation}\label{eq:euler-lagrange}
\boxed{F_y-\frac{\dd}{\dd x}F_{y'}=0.}
\end{equation}
这就是欧拉--拉格朗日（Euler--Lagrange）方程，是内部驻值的必要条件，还须满足原边界条件及任何约束。

\begin{example}[一个可直接检验的驻值函数]
令独立变量与函数值均为无量纲量，求
\[
I[y]=\int_0^1(y'^2+2xy)\dd x,\qquad y(0)=0,\quad y(1)=1.
\]
有 $F_y=2x,F_{y'}=2y'$，所以 $y''=x$。两次积分得 $y=x^3/6+C_1x+C_2$，端点给出 $C_2=0,C_1=5/6$。对任意固定端点扰动 $h$，精确展开为
\[
I[y+h]-I[y]=\int_0^1\bigl[2y'h'+2xh+h'^2\bigr]\dd x
=\int_0^1h'^2\dd x>0\quad(h\ne0).
\]
中间用 $y''=x$ 及 $h(0)=h(1)=0$ 消去一次项。因此这个驻点还是整个图形容许类中的唯一全局极小，结论比仅解出方程更强。
\end{example}

\subsection{离散化：把路径看成很多个坐标}
先把曲线画成许多小线段。固定两头，只调中间顶点的高度，就能用熟悉的多变量求极值来近似这个问题。

令 $h=(b-a)/N$、$x_k=a+kh$，固定 $y_0,y_N$，用
\begin{equation}\label{eq:discrete-functional}
I_N=\sum_{k=1}^NhF\left(x_k,y_k,\frac{y_k-y_{k-1}}h\right)
\end{equation}
近似积分。内部 $y_k$ 只影响第 $k$ 和 $k+1$ 项，普通偏导数给出
\begin{equation}\label{eq:discrete-el}
\frac{\partial I_N}{\partial y_k}
=hF_y|_k+F_{y'}|_k-F_{y'}|_{k+1}=0.
\end{equation}
除以 $h$，差分形式趋近 $F_y-\dd F_{y'}/\dd x=0$。这解释连续方程与多元极值条件的联系，但不是一般的数值收敛定理；离散驻点是否逼近指定连续解，还依赖正则性、离散格式和解的稳定性。

\subsection{首次积分与多个未知函数}
有些问题在曲线的不同位置遵守同一条局部规则。这种不显含某个变量的结构有时能直接给一个常量，少解一阶微分方程。

沿驻值路径应用乘积法则，有
\[
\frac{\dd}{\dd x}(F-y'F_{y'})
=F_x+y'\left(F_y-\frac{\dd F_{y'}}{\dd x}\right)=F_x.
\]
若 $F$ 不显含 $x$，则有\textbf{贝尔特拉米恒等式}
\begin{equation}\label{eq:beltrami}
F-y'F_{y'}=C.
\end{equation}
这是沿解成立的首次积分，不是任意曲线上的恒等式；也不保证凡满足它的退化解均满足原欧拉--拉格朗日方程，应代回核验。

若 $I$ 依赖多个函数 $y_i(x)$，在无额外约束时各方向独立，故
\begin{equation}\label{eq:el-multi}
F_{y_i}-\frac{\dd}{\dd x}F_{y_i'}=0,\qquad i=1,\ldots,n.
\end{equation}
若 $F$ 不显含某个 $y_i$，则
\begin{equation}\label{eq:cyclic}
F_{y_i'}=\text{常量}.
\end{equation}
第三章会把这解释为循环坐标与守恒动量。若再作 $F\to F+\dd G(x,y)/\dd x$，泛函只改变 $[G]_a^b$；固定端点时它不改变内部方程，但自由端点时边界条件可能改变。

\section{边界条件和约束：容许变分决定答案}
端点钉在墙上与端点沿导轨滑动，是两种不同的题。写完内部方程以后，还要回来看端点或总长度究竟允许怎样变化。

\subsection{自然边界条件}
如果右端可以上下动，就不能直接把右端变分置零。为了让任意允许的小移动都不产生一次变化，必须要求它所乘的系数为零。

若 $a,b$ 固定，而 $y(b)$ 自由，则先取内部扰动得到欧拉--拉格朗日方程，再任取终点扰动，式 \eqref{eq:variation-boundary} 要求
\begin{equation}\label{eq:natural-bc}
F_{y'}(b)=0.
\end{equation}
起点自由时同理；仅自由的端点才有这条条件，不能在固定端点上额外强加它。若泛函另含端点代价 $B(y(b))$，其变分为 $B_y\eta(b)$。与积分产生的 $F_{y'}(b)\eta(b)$ 相加，再令任意的 $\eta(b)$ 的系数为零，便有 $F_{y'}(b)+B_y=0$。
\begin{example}[到一条竖直直线的最短距离]
取 $I=\int_0^L\sqrt{1+y'^2}\dd x$，$y(0)=0$，终点在 $x=L>0$ 上自由滑动。内部方程给出直线 $y=kx$，边界条件给出 $k/\sqrt{1+k^2}=0$，所以 $y=0$。它垂直目标直线；并且所有容许图形长度至少为 $\int_0^L1\dd x=L$，故确实达到全局最短。

内部方程为 $\dd[y'/\sqrt{1+y'^2}]/\dd x=0$，其括号内函数对 $y'$ 严格单调，故 $y'$ 为常数。积分分部的右端项为 $[y'/\sqrt{1+y'^2}]_{x=L}\eta(L)$，端点可竖直滑动意味着 $\eta(L)$ 任意，因此斜率必须为零。
\end{example}

\subsection{移动端点与横截条件的推导}
端点沿斜线移动时，横坐标和纵坐标会一起改变。先把允许的移动方向写出来，再看边界项沿这个方向的分量是否为零。

现在令终点横坐标也可变，$b_\epsilon=b+\epsilon\Delta b$，终点总纵向变化为 $\Delta y_B$。若 $y_\epsilon(x)=y(x)+\epsilon\eta(x)$，则
\[
y_\epsilon(b_\epsilon)=y(b)+\epsilon[\eta(b)+y'(b)\Delta b]+O(\epsilon^2),
\]
故 $\eta(b)=\Delta y_B-y'(b)\Delta b$。积分上限移动还产生 $F(b)\Delta b$，于是内部方程满足后，终点贡献为
\begin{equation}\label{eq:var-transversality}
\delta I\big|_B=F_{y'}\Delta y_B+(F-y'F_{y'})\Delta b.
\end{equation}
若端点沿 $y_B=\psi(b)$ 移动，$\Delta y_B=\psi'(b)\Delta b$，由 $\Delta b$ 的任意性得到\textbf{横截条件}
\[
\boxed{F+(\psi'-y')F_{y'}=0\quad\text{在终点成立}.}
\]
对长度泛函，化为 $1+y'\psi'=0$，即两切线正交。目标竖直线应采用前面的固定 $b$ 情形，不应把无穷大的 $\psi'$ 硬代入这个公式。
\begin{figure}[htbp]
\centering
\begin{tikzpicture}[>=Stealth]
\draw[thick,gray] (-.5,0)--(4,2.25) node[right]{目标直线};
\coordinate (B) at (2,1.25);\coordinate (A) at (3,-.75);
\draw[very thick,blue!70!black] (A)--(B);
\draw[->,thick,red!70!black] (B)--(3,1.75) node[above]{$\Delta\br_B$};
\fill (A) circle (2pt) node[below]{$A$};\fill (B) circle (2pt) node[above left]{$B$};
\draw (2.16,.93)--(2.48,1.09)--(2.32,1.41);
\end{tikzpicture}
\caption{移动端点的容许位移沿目标线。最短连接段垂直目标线，正交条件来自终点边界项，而非内部微分方程。}
\label{fig:var-free-end}
\end{figure}

\subsection{积分约束与乘子的条件}
规定绳长后，可以改一小段的形状，但其他段可能需要同时调整，才能保持总长。用一个乘子把总长条件和原目标放到同一个变分式中。

设还要求 $J[y]=\int_a^bG(x,y,y')\dd x=\ell$。在约束正则，即一阶线性泛函 $\delta J$ 不恒为零的条件下，受约束驻点满足存在\textbf{常数} $\lambda$ 使
\begin{equation}\label{eq:lagrange-multiplier}
\delta(I-\lambda J)=0,
\end{equation}
也就是对 $\widetilde F=F-\lambda G$ 求
\begin{equation}\label{eq:constrained-el}
\widetilde F_y-\frac{\dd}{\dd x}\widetilde F_{y'}=0.
\end{equation}
乘子与积分上限无关，由边界值和积分约束共同确定。取固定端点扰动 $\zeta$ 满足 $\delta J[\zeta]\ne0$。对任意满足 $\delta J[\eta]=0$ 的固定端点扰动，用隐函数定理选择 $u(\epsilon)$，使 $J[y+\epsilon\eta+u(\epsilon)\zeta]=\ell$。对约束求导得 $u'(0)=0$，所以受约束驻点满足 $\delta I[\eta]=0$。对一般扰动 $\eta$，减去 $\zeta\,\delta J[\eta]/\delta J[\zeta]$ 就得到这样的核方向。由线性性，$\delta I=\lambda\delta J$，其中 $\lambda=\delta I[\zeta]/\delta J[\zeta]$。若 $\delta J$ 处处为零，这个常规推理失效，需分析异常约束，不能直接套用归一化乘子形式。

点态约束 $g(x,y_1,\ldots,y_n)=0$ 则在每个 $x$ 都施加条件，通常引入函数乘子 $\lambda(x)$。积分约束的常数乘子和点态约束的函数乘子不要混淆。第三章还会说明：对非完整速度约束，简单把约束加入作用量未必给出物理运动。

\begin{example}[圆弧面积问题：限制图形类后再求解]
固定端点 $(0,0),(d,0)$，要求 $y(x)\ge0$ 且弧长为 $\ell$。本例先限制为光滑图形，并取 $d<\ell<\pi d/2$。最大化面积 $A=\int_0^dy\dd x$。令
\[
\widetilde F=y-\lambda\sqrt{1+y'^2},
\qquad
1+\lambda\frac{y''}{(1+y'^2)^{3/2}}=0.
\]
曲率为常数，取上凸圆弧对应 $\lambda=R>0$。由对称性，圆心横坐标为 $d/2$。用半圆心角 $0<\alpha<\pi/2$ 表示，
\[
d=2R\sin\alpha,\qquad \ell=2R\alpha,\qquad
\frac\ell d=\frac\alpha{\sin\alpha},
\]
\[
y_*(x)=\sqrt{R^2-(x-d/2)^2}-R\cos\alpha.
\]
上述比例在此范围单调，端点与长度唯一确定 $R,\alpha$。还可以证明极大性：$-R\sqrt{1+p^2}$ 对斜率 $p$ 严格凹，而 $y$ 项线性，故对任意同端点图形
\[
(A-R\mathcal L)[y]-(A-R\mathcal L)[y_*]
\le\delta(A-R\mathcal L)[y_*;y-y_*]=0.
\]
若弧长相同，便有 $A[y]\le A[y_*]$，等号只对应同一图形。

不能把这个图形证明无条件用于任意大弧长：$\ell\to\pi d/2$ 时圆弧趋于半圆，端点斜率发散；再大的圆弧会有回折，不能由一个光滑 $y(x)$ 表示。闭合曲线等周定理需要参数化曲线的容许类，本例没有假装已证明它的全部版本。
\end{example}

\section{二阶变分：驻值为什么不一定是极小？}
一个数的导数为零可能对应谷底、山顶或更平的地方。对整条曲线也要做第二步检查，看看往各个允许方向稍作改变时，原量是增还是减。

\subsection{二次型、必要条件与充分条件}
同一条曲线可以有很多种鼓起方式。二阶变分要检查的是这些方式各自引起的二次变化，其中只要有一个方向使它减小，就排除了极小。

固定端点时，对同一族 $y+\epsilon\eta$ 再求一次参数导数，定义
\[
\delta^2 I[y;\eta]=\left.\frac{\dd^2}{\dd\epsilon^2}I[y+\epsilon\eta]\right|_0.
\]
注意在泰勒展开中它前面有 $1/2$，这只是约定，不影响符号。直接求导两次：
\begin{align}\label{eq:second-variation}
\delta^2 I
&=\int_a^b(F_{yy}\eta^2+2F_{yy'}\eta\eta'+F_{y'y'}\eta'^2)\dd x\nonumber\\
&=\int_a^b\left(P\eta'^2+Q\eta^2\right)\dd x,\qquad
P=F_{y'y'},\quad Q=F_{yy}-\frac{\dd F_{yy'}}{\dd x}.
\end{align}
第二步对 $2F_{yy'}\eta\eta'=F_{yy'}(\eta^2)'$ 分部积分，端点项因 $\eta=0$ 消失。各系数沿待检验的驻值路径取值。

局部极小要求 $\delta^2I\ge0$ 对所有容许方向成立，这是必要条件；若存在正、负两种方向，则为鞍点。对正则标量问题，局部极小还必须满足勒让德条件 $F_{y'y'}\ge0$。严格的 $F_{y'y'}>0$ 常作为加强条件，但它\textbf{单独并不保证极小}，因为 $Q\eta^2$ 可能贡献负值。

在无限维函数空间里，也不能不说明拓扑和余项控制，就将“每个方向二阶导数均正”直接当作统一邻域内的充分条件。一个实用的充分条件是：$F$ 在候选路径附近足够光滑，二次型有统一正下界
\[
\delta^2I[y;\eta]\ge c\int_a^b\eta'^2\dd x,\qquad c>0,
\]
且泰勒余项相对于该范数可一致控制。例如连续系数满足 $P\ge p_0>0,Q\ge0$ 时，配合固定端点的庞加莱不等式即可得到所需估计，因而在足够小的 $C^1$ 邻域内有严格弱极小。这里“弱”指函数值和一阶导数均接近，不是说极小结论不严格。

\begin{example}[直线不仅驻值，而且全局最短]
对 $F=\sqrt{1+y'^2}$，有 $P=(1+y'^2)^{-3/2}>0,Q=0$。固定两端的驻值路径为直线，因此二阶变分为正。还可不依赖局部判据：对任意参数化路径，$\int|\dd\br|\ge|\br_B-\br_A|$，等号对应同向直线段。这个三角不等式给出全局最短，且说明同一直线若往返折返，虽仍落在一条直线上，却不达到等号。
\end{example}

\subsection{作用量随时间区间改变性质的例子}
谐振子的零路径在短区间和长区间都能满足运动方程。不过端点相隔越远，中间可容纳的扰动越多，作用量的极小性可能随之改变。

取谐振子作用量
\[
S[q]=\frac m2\int_0^T(\dot q^2-\omega^2q^2)\dd t,
\qquad q(0)=q(T)=0,\quad m,\omega>0.
\]
$q_*=0$ 对所有 $T$ 都是驻值路径。对固定端点扰动 $\eta$，
\[
\delta^2S=m\int_0^T(\dot\eta^2-\omega^2\eta^2)\dd t.
\]
取 $\eta=A\sin(n\pi t/T)$，利用正弦与余弦平方的积分均为 $T/2$，得
\begin{equation}\label{eq:var-osc-mode}
\delta^2S=\frac{mA^2T}{2}\left[\left(\frac{n\pi}{T}\right)^2-\omega^2\right].
\end{equation}
若 $T<\pi/\omega$，固定端点的庞加莱不等式
$\int\dot\eta^2\dd t\ge(\pi/T)^2\int\eta^2\dd t$
保证所有非零方向均为正；本例是精确二次泛函，所以 $q=0$ 还是全局严格极小。若 $T=\pi/\omega$，$\eta=A\sin\omega t$ 给出零二阶变分，而且这些路径的作用量确实全都为零，极小不严格。若 $T>\pi/\omega$，$n=1$ 给负方向，足够大的 $n$ 给正方向，故是鞍点。
\begin{physicsbox}
“最小作用量”是传统名称。即使 $F_{\dot q\dot q}=m>0$，谐振子在较长时间区间的作用量仍有负变分方向。作用量极小性也不等同于动力学稳定性：谐振子的轨道稳定，并不意味着任意两时刻之间的作用量都是最小值。
\end{physicsbox}
\begin{figure}[htbp]
\centering
\begin{tikzpicture}
\begin{axis}[width=.78\textwidth,height=5.7cm,axis lines=left,
 xlabel={$\tau=\omega T/\pi$},ylabel={$\Lambda_n$},xmin=.55,xmax=2.8,ymin=-1,ymax=4,
 xtick={1,2},ytick={-1,0,1,2,3,4},legend style={at={(.97,.97)},anchor=north east}]
\addplot[thick,blue!70!black,domain=.55:2.8,samples=120] {1/x^2-1};
\addplot[thick,red!70!black,domain=.9:2.8,samples=120] {4/x^2-1};
\addplot[dashed,gray,domain=.55:2.8] {0};
\legend{$n=1$,$n=2$}
\end{axis}
\end{tikzpicture}
\caption{无量纲二阶变分系数 $\Lambda_n=n^2/\tau^2-1$。第一个模在 $T=\pi/\omega$ 变号；一个负模已足以否定极小性。图线由式 \eqref{eq:var-osc-mode} 直接绘出。}
\label{fig:var-second-sign}
\end{figure}

\subsection{进一步探索：雅可比方程与共轭点}
有时存在一个小扰动，从起点出发后恰好又在某个时刻回到零。找出这个时刻，有助于判断原路径的极小性质何时可能改变。

对二次型 $\mathcal Q[\eta]=\int(P\eta'^2+Q\eta^2)\dd x$ 作变分，内部驻值条件为
\begin{equation}\label{eq:jacobi}
-\frac{\dd}{\dd x}(P\eta')+Q\eta=0.
\end{equation}
它是\textbf{零模方程}；若要在 $\int\eta^2\dd x=1$ 下最小化该二次型，则右边是 $\lambda\eta$，不能漏掉归一化乘子。总存在零解，讨论共轭点时寻找的是非平凡解。

在 $P>0$ 且系数正则的标量问题中，固定起点 $a$，若存在非平凡解满足 $\eta(a)=\eta(c)=0$，则称 $c>a$ 与 $a$ 共轭。它不是要求区间内完全不存在非零解，而是要求不存在具有指定两个零点的解。加强勒让德条件与区间 $(a,b]$ 上无共轭点相结合，给出通常的固定端点严格弱极小判据；证明涉及二次型分解，作为进阶阅读，不用它替代前面已给出的直接检查。

谐振子中 $P=m,Q=-m\omega^2$，雅可比方程是 $\ddot\eta+\omega^2\eta=0$，起点为零的非平凡解正比于 $\sin\omega t$，首个共轭时刻正是 $\pi/\omega$，与上面的模态计算一致。学习这个拓展需要二阶线性常微分方程和边值问题。

\section{几何与力学中的变分问题}
找最短路、算绳子的形状和比较下滑时间，分别要最小化或取驻值的是长度、势能和时间。先写明每题比较什么，再套用变分步骤。

\subsection{平面最短路径与曲面测地线}
平面上能直接连直线，球面上则只能沿表面走。即使一段大圆弧处处符合测地线方程，也应比较同一对端点间另一侧的弧长。

长度泛函的被积函数不显含 $y$，由式 \eqref{eq:cyclic} 得 $y'/\sqrt{1+y'^2}=C$，所以 $y'$ 为常数，端点确定直线。

在曲面上，使用局部坐标 $q^1,q^2$，定义正定度规
\[
\dd s^2=g_{ij}(q)\dd q^i\dd q^j,
\]
重复指标求和，$g_{ij}=g_{ji}$，逆矩阵为 $g^{ij}$。用任意正则参数 $u$ 表示曲线，令
\[
N=\sqrt{g_{ij}q'^iq'^j},\qquad \mathcal L=\int N\dd u,
\]
此处撇号表示 $\dd/\dd u$，$u$ 不必是物理时间。先求偏导：
\[
N_{q'^k}=\frac{g_{kj}q'^j}{N},\qquad
N_{q^k}=\frac{\partial_kg_{ij}q'^iq'^j}{2N}.
\]
将它们代入欧拉--拉格朗日方程，乘以 $N$ 并展开：
\[
g_{kj}q''^j+\partial_i g_{kj}q'^iq'^j
-\frac12\partial_kg_{ij}q'^iq'^j
=\frac{N'}N g_{kj}q'^j.
\]
把二次速度项对称化，乘以 $g^{mk}$，得到
\[
q''^m+\Gamma^m_{ij}q'^iq'^j=\frac{N'}Nq'^m,
\]
\begin{equation}\label{eq:christoffel}
\Gamma^m_{ij}=\frac12g^{mk}(\partial_i g_{kj}+\partial_j g_{ki}-\partial_k g_{ij}).
\end{equation}
若在导出方程后选择弧长 $s$ 或其仿射变换为参数，$N$ 为常数，右侧消失：
\begin{equation}\label{eq:geodesic}
\boxed{\frac{\dd^2q^m}{\dd s^2}+\Gamma^m_{ij}\frac{\dd q^i}{\dd s}\frac{\dd q^j}{\dd s}=0.}
\end{equation}
不能先把 $N=1$ 代入被积函数再变分，那样会错误地把长度泛函变成常数。测地线是弧长驻值的曲线，只在足够短、未跨过相应共轭点的弧段上保证局部最短，未必全局最短。

\begin{example}[球面：大圆长弧为何不是最短？]
半径为 $R$ 的球面，取 $q^1=\theta,q^2=\varphi$，有
\[
g_{\theta\theta}=R^2,\quad g_{\varphi\varphi}=R^2\sin^2\theta.
\]
唯一非零的度规偏导是 $\partial_\theta g_{\varphi\varphi}=2R^2\sin\theta\cos\theta$。因此
\[
\Gamma^\theta_{\varphi\varphi}=-\sin\theta\cos\theta,\qquad
\Gamma^\varphi_{\theta\varphi}=\Gamma^\varphi_{\varphi\theta}=\cot\theta,
\]
\[
\theta''-\sin\theta\cos\theta\,\varphi'^2=0,\qquad
\varphi''+2\cot\theta\,\theta'\varphi'=0,
\]
撇号现表示对弧长求导。赤道解为 $\theta=\pi/2,\varphi=\varphi_0\pm s/R$；系数须满足单位速率条件，不能任取。球面对称性及初值唯一性说明一般正则测地线是大圆的一段。

两个非对跖点间存在同一大圆上的短弧和长弧，长度分别为 $R\beta$ 与 $R(2\pi-\beta)$，$0<\beta<\pi$。两者都满足测地线方程，只有短弧是全局最短。对跖点间则有无穷多等长的大圆半圆，最短路径不唯一。球坐标的极点奇异需换坐标，不是曲面本身有奇点。
\end{example}
\begin{figure}[htbp]
\centering
\begin{tikzpicture}[>=Stealth]
\draw[very thick,gray!65] (0,0) circle (1.8);
\draw[very thick,blue!70!black] (1.8,0) arc (0:110:1.8);
\fill (1.8,0) circle (2pt) node[right]{$A$};
\fill ({1.8*cos(110)},{1.8*sin(110)}) circle (2pt) node[above]{$B$};
\draw[dashed] (0,0)--(1.8,0);\draw[dashed] (0,0)--({1.8*cos(110)},{1.8*sin(110)});
\draw (.6,0) arc (0:110:.6);\node at (.6,.6) {$\beta$};
\node[blue!70!black] at (1.6,1.6) {短弧};\node[gray!65] at (-1.8,-1.5) {长弧};
\end{tikzpicture}
\caption{在过球心的大圆平面中观察两端点。短、长两弧均是球面测地线；满足驻值方程不等于在所有连接路径中长度最小。}
\label{fig:var-great-circle}
\end{figure}

\subsection{悬链线：坐标方向、长度约束与张力}
均匀绳子下垂时，靠近支点的张力要托住更多绳段。由于载荷按绳的弧长分布，坡越陡，每单位水平距离要托住的绳子也越多。

\begin{example}[均匀柔绳的平衡形状]
柔绳线密度为 $\rho_\ell$，两端固定，绳长为 $\ell$，不可伸长，忽略弯曲刚度，均匀重力竖直向下。本例取 $y$ 向上，故势能应为
\[
V[y]=\rho_\ell g\int y\dd s
=\rho_\ell g\int y\sqrt{1+y'^2}\dd x,
\]
不能写成与向上坐标相反的负号。固定长度 $\mathcal L=\int\sqrt{1+y'^2}\dd x=\ell$，将常数 $\rho_\ell g$ 提出，并用长度量纲的乘子 $\lambda$，得到
\[
\widetilde F=(y-\lambda)\sqrt{1+y'^2}.
\]
由于不显含 $x$，式 \eqref{eq:beltrami} 给出
\[
\frac{y-\lambda}{\sqrt{1+y'^2}}=a.
\]
选择实际下垂、张力为正的分支 $a>0$。由平方关系
\[
y'^2=\frac{(y-\lambda)^2-a^2}{a^2}
\]
分离变量。在最低点右侧取正斜率，左侧取负斜率，两支合起来为
\begin{equation}\label{eq:var-catenary}
y=\lambda+a\cosh\frac{x-x_0}{a}
=y_{\min}+a\left[\cosh\frac{x-x_0}{a}-1\right],
\qquad y_{\min}=\lambda+a.
\end{equation}
$x_0$、$a$ 和高度常数须由两端位置及绳长确定，不能只通过“吸收常数”忽略物理支点高度。

对等高支点 $(-d,H),(d,H)$，有 $x_0=0$，
\[
\ell=\int_{-d}^d\cosh(x/a)\dd x=2a\sinh(d/a),\qquad
H-y_{\min}=a[\cosh(d/a)-1].
\]
当 $\ell>2d$，长度方程给出唯一 $a>0$；$\ell\downarrow2d$ 时 $a\to\infty$，绳近乎拉直，有限重量下张力趋大。若 $\ell<2d$，不可伸长绳根本无法连接支点。

从第一章受力分析独立核验：设水平张力分量 $H_T$ 恒定，以最低点为弧长原点，带符号弧长为
\[s=a\sinh(x/a).\]
竖直张力分量为 $\rho_\ell gs$，切线斜率为
\[
y'=\frac{\rho_\ell gs}{H_T}=\sinh(x/a)
\quad\Longrightarrow\quad H_T=\rho_\ell ga.
\]
这既确认了坐标符号，也说明 $a$ 的力学意义。变分方程给出平衡候选，实际悬挂支取正张力；这里没有仅凭一阶变分就宣称所有可能曲线类中的全局唯一稳定性。
\end{example}
\begin{figure}[htbp]
\centering
\begin{tikzpicture}[>=Stealth]
\draw[->] (-2.6,0)--(2.7,0) node[right]{$x$};\draw[->] (0,-.2)--(0,3) node[above]{$y$};
\draw[very thick,blue!70!black,domain=-2:2,samples=120] plot (\x,{1.5*(cosh(\x/1.5)-1)+.4});
\fill (-2,{1.5*(cosh(2/1.5)-1)+.4}) circle (2pt) node[above left]{$(-d,H)$};
\fill (2,{1.5*(cosh(2/1.5)-1)+.4}) circle (2pt) node[above right]{$(d,H)$};
\fill (0,.4) circle (2pt) node[below right]{$y_{\min}$};
\draw[->,thick,red!70!black] (1.25,1)--(1.25,.25) node[right]{$\bg$};
\end{tikzpicture}
\caption{由解析双曲余弦函数绘制的下垂绳。$y$ 向上而重力向下；绳长是沿弧线的积分，不能误用水平跨度。}
\label{fig:var-catenary}
\end{figure}
\begin{physicsbox}
均匀绳的重量按\emph{弧长}均匀分布，产生悬链线；若悬索负载主要按\emph{水平长度}均匀分布，受力方程不同，可得到抛物线。曲线看起来接近不代表模型相同。悬链线在 $|x-x_0|\ll a$ 时也有局部展开 $y\simeq y_{\min}+(x-x_0)^2/(2a)$。
\end{physicsbox}

\subsection{最速降线：解、端点与时间验证}
陡一点的路线能更早获得速度，却也可能绕得更远。最速降线正是比较沿途速度与路程这两方面后得到的路径。

本例重新取 $y$ 向下。内部 $y>0$ 时，令
\[
F=\frac{\sqrt{1+y'^2}}{\sqrt{2gy}},
\qquad F-y'F_{y'}=\frac1{\sqrt{2gy}\sqrt{1+y'^2}}=C>0.
\]
定义 $a=1/(4gC^2)>0$，得到
\begin{equation}\label{eq:brachistochrone-ode}
y'^2=\frac{2a-y}{y}.
\end{equation}
在首拱令 $y=a(1-\cos\theta)$，并令 $\theta$ 随 $x$ 增大。式 \eqref{eq:brachistochrone-ode} 给出 $y'^2=\cot^2(\theta/2)$。下行与回升段分别取 $y'$ 的正根和负根，两段合写为 $y'=\cot(\theta/2)$。由链式法则
\[
\frac{\dd x}{\dd\theta}
=\frac{a\sin\theta}{\cot(\theta/2)}
=a(1-\cos\theta),
\]
在 $\theta=\pi$ 处按连续极限取值。用 $x(0)=0$ 积分，得
\begin{equation}\label{eq:cycloid}
\boxed{x=a(\theta-\sin\theta),\qquad y=a(1-\cos\theta).}
\end{equation}
当 $0<\theta<\pi$，$x_\theta=a(1-\cos\theta)>0$、$y_\theta=a\sin\theta>0$，曲线下行；$\pi<\theta<2\pi$ 时 $x$ 仍增大而 $y$ 减小，曲线回升。故分离变量写 $\dd x=\sqrt{y/(2a-y)}\dd y$ 只适用于前一支，不能用正根覆盖整条摆线。

由端点 $B=(X,Y)$，$X,Y>0$，要求
\[
\frac XY=\frac{\theta_B-\sin\theta_B}{1-\cos\theta_B},\qquad
 a=\frac Y{1-\cos\theta_B}.
\]
设 $u=\theta_B/2$。上述比值对 $\theta_B$ 的导数分子为 $2(1-\cos\theta_B)-\theta_B\sin\theta_B=4\sin u(\sin u-u\cos u)$，分母为 $(1-\cos\theta_B)^2$。由于 $(\sin u-u\cos u)'=u\sin u>0$，且该括号在 $u=0$ 为零，导数在首拱为正。比值两端的极限分别为零与正无穷，故正端点唯一确定首拱参数解；若 $X/Y<\pi/2$，终点在下行段，若 $X/Y>\pi/2$，最速降线候选先降到终点以下再回升，不能始终描述成“先陡后平”。

直接计算弧元和速度：
\[
\frac{\dd s}{\dd\theta}=2a\sin(\theta/2),\qquad
v=\sqrt{2gy}=2\sqrt{ga}\sin(\theta/2),
\]
于是
\begin{equation}\label{eq:var-brach-time}
\dd t=\sqrt{a/g}\dd\theta,\qquad T_B=\theta_B\sqrt{a/g}.
\end{equation}
尽管起点速度为零、切线竖直，二者的零因子相消后时间积分有限。这是对起点奇性的检查；曲线不是在端点具有有限导数的 $C^2$ 图形，前面的正则推导仅用于内部。

\begin{example}[同一终点下与直线比较]
取教学参数 $X=Y=1\,\mathrm m$、$g=9.8\,\mathrm{m/s^2}$。端点方程数值求得 $\theta_B\simeq2.4120$、$a\simeq0.5729\,\mathrm m$，故 $T_B\simeq0.583\,\mathrm s$。直线长为 $D=\sqrt{X^2+Y^2}$，沿线的恒定加速度为 $gY/D$，所以
\[
T_{\rm line}=\sqrt{\frac{2D}{gY/D}}
=\sqrt{\frac{2(X^2+Y^2)}{gY}}\simeq0.639\,\mathrm s.
\]
两条路径使用相同起点、终点、初速及重力模型，比较才有意义。该数值不是实际实验测量。

直线的弧长为 $\sqrt{X^2+Y^2}$，以沿线弧长 $s$ 参数化时 $y=Ys/\sqrt{X^2+Y^2}$，从静止出发有 $v=\sqrt{2gy}$。积分 $\int\dd s/v$ 得 $t_{\rm line}=\sqrt{2(X^2+Y^2)/(gY)}$，再与摆线时间 $\sqrt{a/g}\,\theta_B$ 比较。
\end{example}
\begin{figure}[htbp]
\centering
\begin{tikzpicture}[>=Stealth,x=2.7cm,y=-1.8cm]
\draw[->] (0,0)--(2.15,0) node[right]{$x/a$};\draw[->] (0,0)--(0,2.25) node[below]{$y/a$};
\draw[dashed,gray] (0,0)--({2.4-sin(2.4 r)},{1-cos(2.4 r)});
\draw[very thick,blue!70!black,domain=0:2.4,samples=150]
 plot ({\x-sin(\x r)},{1-cos(\x r)});
\fill (0,0) circle (1pt) node[above left]{$A$};
\fill ({2.4-sin(2.4 r)},{1-cos(2.4 r)}) circle (1pt) node[right]{$B$};
\node[gray] at (.9,.6) {直线};\node[blue!70!black] at (.6,1.5) {摆线};
\end{tikzpicture}
\caption{参数为 $0\le\theta\le2.4$ 的精确摆线及同端点直线。$y$ 轴向下，与能量关系 $v^2=2gy$ 的约定一致。}
\label{fig:var-brachistochrone}
\end{figure}

\paragraph{何时已证明“最速”，而非只求到驻值？}
对下行段 $0<\theta_B<\pi$，可在单调深度的路径类中以 $y$ 为自变量，写 $x=x(y)$。时间泛函变为
\[
T[x]=\int_0^Y f(y,w)\dd y,\quad w=\frac{\dd x}{\dd y},\qquad
f(y,w)=\frac{\sqrt{1+w^2}}{\sqrt{2gy}}.
\]
$f_{ww}>0$，且摆线解上 $f_w$ 是常数。因此凸性给出
\[
T[x]-T[x_*]\ge\int_0^Y f_w(y,x_*')(x'-x_*')\dd y
=C_x[x-x_*]_0^Y=0.
\]
在端点极限可积、同端点且有限时间的这类图形中，摆线达到严格最小。这里利用了一个比欧拉--拉格朗日方程更强的论证。对于允许深度回升的更大曲线类，或 $\theta_B>\pi$ 的情形，本段单调深度证明不适用；完整最速降线理论需另作参数化比较，本章不将其冒充为已证明的结论。

\subsection{作用量初步：从变分返回牛顿方程}\label{sec:action-first}
比较的是同一时间区间上连接同一对端点的路径。先给每条路径算动能减势能的时间积分，再检查小改路径时这个积分怎样变化。

\begin{example}[固定时间端点的一维保守粒子]
取质量恒定、势能 $V(q)$ 光滑的粒子，定义
\[
L(q,\dot q)=\tfrac12m\dot q^2-V(q),\qquad
S[q]=\int_{t_1}^{t_2}L\dd t.
\]
固定 $t_1,t_2$ 及两端位置，取 $q_\epsilon=q+\epsilon\eta$、$\eta(t_1)=\eta(t_2)=0$。逐项变分并分部积分：
\[
\delta S=\int_{t_1}^{t_2}(m\dot q\dot\eta-V'(q)\eta)\dd t
=[m\dot q\eta]_{t_1}^{t_2}-\int_{t_1}^{t_2}(m\ddot q+V')\eta\dd t.
\]
由扰动的任意性得 $m\ddot q=-V'(q)$。在此模型与边界条件下，牛顿方程与哈密顿驻作用量条件相互对应；不是对所有力学系统、非完整约束或耗散系统都无条件等价。
\end{example}
作用量的量纲是能量乘时间；它不是能量，也不要求沿真实轨道每时刻最小。给定 $L$ 是指定动力学模型，$\delta S=0$ 才是求其运动的规则。第\ref{sec:ham-principle}节将用独立广义坐标写出相应的多自由度方程。

\section{进一步探索：从几何光学到受约束优化}
同样的思想还能处理光线在连续介质中弯曲、整根弦一起运动，以及怎样安排随时间变化的控制量。

本节各项是拓展应用，不要求初学者掌握所有推导。它们分别展示多分量路径、空间场和动力学约束；不能把任意优化问题都笼统称为同一个物理原理。
\subsection{梯度折射率中的光线方程}
折射率连续变化时，可以把介质想成许多性质略有差别的薄层。光线的方向在每一小段上微调，整体看便是弯曲的。

设 $\br(u)$ 为正则路径，$I=\int n(\br)|\br'|\dd u$。对各坐标分量取偏导，
\[
\frac{\partial F}{\partial r_i'}=n\frac{r_i'}{|\br'|},\qquad
\frac{\partial F}{\partial r_i}=|\br'|\partial_i n.
\]
欧拉--拉格朗日方程给出
\[
\frac{\dd}{\dd u}\left(n\frac{\br'}{|\br'|}\right)=|\br'|\nabla n.
\]
导出后再取弧长 $s$ 为参数，得到
\begin{equation}\label{eq:ray-equation}
\boxed{\frac{\dd}{\dd s}\left(n\frac{\dd\br}{\dd s}\right)=\nabla n.}
\end{equation}
令 $\be_t=\dd\br/\dd s$，用 $\dd n/\dd s=\be_t\cdot\nabla n$ 分离出
\[
n\frac{\dd\be_t}{\dd s}=\nabla n-\be_t(\be_t\cdot\nabla n).
\]
只有折射率梯度的横向部分使光线弯曲；沿切向的梯度改变 $n\be_t$ 的大小。均匀介质中方向不变。

\begin{example}[近轴梯度折射率模型]
理想折射率在轴附近满足 $n^2(x,y)=n_0^2[1-\alpha^2(x^2+y^2)]$，$\alpha$ 的量纲为逆长度。仅在 $\alpha r\ll1$ 且 $|\dd x/\dd z|,|\dd y/\dd z|\ll1$ 的近轴范围使用此模型，则 $n\simeq n_0(1-\alpha^2r^2/2)$、$\dd s\simeq\dd z$。保留线性横向项，
\[
n_0\frac{\dd^2x}{\dd z^2}\simeq\partial_xn=-n_0\alpha^2x,
\qquad \frac{\dd^2y}{\dd z^2}\simeq-\alpha^2y.
\]
给定 $x(0)=x_0,x'(0)=\beta_0$，
\[
x(z)=x_0\cos\alpha z+\frac{\beta_0}{\alpha}\sin\alpha z.
\]
空间周期为 $2\pi/\alpha$，不是时间周期；$\alpha\to0$ 恢复直线 $x=x_0+\beta_0z$。对一般不同横向相位，三维光线不必落在一张平面中，不能一概称为平面正弦线。模型只处理给定光线的几何传播，不据此计算完整波导模或通信色散性能。
\end{example}
\begin{figure}[htbp]
\centering
\begin{tikzpicture}[>=Stealth]
\fill[blue!5] (0,-1) rectangle (6.6,1);
\draw[dashed,gray] (0,0)--(7,0) node[right]{$z$};
\draw[thick,red!70!black,domain=0:6.6,samples=150] plot (\x,{.65*cos(\x r)});
\draw[<->] (0,1.3)--(6.283185,1.3) node[midway,above]{$2\pi/\alpha$};
\node[left] at (0,.65) {$x_0$};
\end{tikzpicture}
\caption{近轴平面光线的解析示意，$\beta_0=0$，横轴按 $\alpha z$ 缩放。背景仅标示可用的轴附近区域，不表示已指定实际光纤的包层结构。}
\label{fig:var-grin}
\end{figure}

\begin{example}[折射率为 $n(y)=y_0/y$ 的理想介质]
限定 $y>0,y_0>0$，局部用 $x=x(y)$ 表示光线。$F=y_0\sqrt{1+x'^2}/y$ 不显含 $x$，故
\[
\frac{y_0x'}{y\sqrt{1+x'^2}}=C,\qquad k=C/y_0.
\]
若 $k\ne0$，在 $|ky|<1$ 的单调 $y$ 分支上，保留 $k$ 的符号得
\[
x'=\frac{ky}{\sqrt{1-k^2y^2}},\qquad
x-x_0=-\frac1k\sqrt{1-k^2y^2}.
\]
因此 $(x-x_0)^2+y^2=1/k^2$。完整光线可用参数拼接为上半平面的圆弧；另一侧分支需改变参数图，不是关于 $x$ 轴反射到禁用的 $y<0$ 区域。若 $C=0$，还存在 $x=\mathrm{const}$ 的竖直光线，除以 $k$ 时会漏掉它。$y\to0$ 处折射率发散，该理想模型不能无条件外推到实际材料边界。
\end{example}

\subsection{空间场：振动弦的欧拉--拉格朗日方程}
弦的每个位置都可以在不同的时刻有不同位移。现在小改的是一个同时依赖位置和时间的函数，分部积分也要分别对这两个变量做。

若未知量是 $u(x,t)$，而非单一自变量的函数，可对时空积分作同样的分部积分。理想小振幅弦的线密度为 $\rho_\ell$、恒定张力为 $\mathcal T$，作用量为
\[
S[u]=\int_{t_1}^{t_2}\int_0^L
\left(\tfrac12\rho_\ell u_t^2-\tfrac12\mathcal T u_x^2\right)\dd x\dd t.
\]
在时间端点和固定弦端令 $\eta=0$。对两项分别按 $t$、$x$ 分部积分：
\[
\delta S=\iint(\rho_\ell u_t\eta_t-\mathcal T u_x\eta_x)\dd x\dd t
=\iint(-\rho_\ell u_{tt}+\mathcal T u_{xx})\eta\dd x\dd t.
\]
边界项消失，故 $\rho_\ell u_{tt}-\mathcal T u_{xx}=0$，波速为 $\sqrt{\mathcal T/\rho_\ell}$。自由弦端则从空间边界项得到 $\mathcal T u_x=0$，体现自然边界条件。这一拓展需要偏导数和二重积分，是进一步学习经典场论的入口。

\subsection{动力学约束与最优控制的最简单例子}
若规定总时间和总位移，可以用很多种速度安排走完同一段路。加入最大速度和代价之后，就能问哪一种安排在这些约束下最好。

令状态 $x(t)$ 满足 $\dot x=u$，端点 $x(0)=0,x(T)=D$，时间 $T>0$ 固定。选择使 $J=\frac12\int_0^Tu^2\dd t$ 最小的控制。引入函数乘子 $\lambda(t)$，写
\[
\mathcal J=\int_0^T\left[\tfrac12u^2+\lambda(\dot x-u)\right]\dd t.
\]
分别对 $u,\lambda,x$ 变分，固定状态端点后得到
\[
u-\lambda=0,\quad \dot x=u,\quad \dot\lambda=0.
\]
所以 $u=D/T$，$x=Dt/T$。用柯西--施瓦茨不等式独立验证：
\[
D^2=\left(\int_0^Tu\dd t\right)^2\le T\int_0^Tu^2\dd t,
\]
故 $J\ge D^2/(2T)$，等号正是恒定控制。本例没有加速度、推力限幅或障碍，不能直接当成真实飞行器最优轨道。更一般的问题还需要控制约束和端点条件，详见章末 MIT 最优控制讲义。

\section{计算与模拟：离散极值与连续解对照}
用等距节点代替连续曲线，把每个内点高度当作未知数。除了比较算出的曲线，还要看网格加密后误差怎样减小。

回到 $I[y]=\int_0^1(y'^2+2xy)\dd x$。令 $h=1/N$，固定 $y_0=0,y_N=1$，取一致的离散泛函
\[
I_N=\sum_{k=1}^N\frac{(y_k-y_{k-1})^2}{h}
+2h\sum_{k=1}^{N-1}x_ky_k.
\]
第二项的端点梯形修正只差一个固定常数，不改变驻点。逐个求偏导得三对角线性方程
\[
-y_{k-1}+2y_k-y_{k+1}=-h^2x_k,\qquad k=1,\ldots,N-1.
\]
矩阵正定，所以离散解是唯一极小。对本例，连续解是三次多项式，而中心差分二阶导数对三次多项式恰好精确，故节点值会直接重现解析解，不能用这个特例声称算法对一般问题没有截断误差。

为检验一般二阶收敛，可另取 $I=\frac12\int_0^1(y'^2+y^2)\dd x$，$y(0)=0,y(1)=1$。连续解为 $y=\sinh x/\sinh1$，离散方程为
\[
-y_{k-1}+(2+h^2)y_k-y_{k+1}=0.
\]
使用 $N=20,40,80$，检查节点最大误差是否随步长减半约缩小4倍。随附数值脚本实现这两组检查。脚本位置为
\begin{center}\texttt{review/verify\_ch02.py}\end{center}
它还核验悬链线、最速降线和二阶变分模型；计算检验不能代替对容许路径与边界条件的判断。

\section*{关键概念与常见误区}
\begin{itemize}[leftmargin=1.8em]
\item 固定一条扰动形状，慢慢增大扰动幅度。一阶变分只看幅度从零开始时的变化率；两条曲线的评分差还可能含二次项。
\item 路径变分比较同一时刻或同一横坐标上的两条曲线。真实位移则沿一条给定轨迹向前走。
\item 端点钉住时，对应的端点变分为零。端点可动时，保留边界项，从中求允许方向上的条件。
\item 解出欧拉--拉格朗日方程，只确定了小改动没有一次影响。再看二阶变化，才能进一步辨认极小、极大或鞍点。
\item 限制整条曲线的总长度只需一个乘子；若每一点各有一条约束，乘子一般也随位置改变。
\item 用曲线长度作目标时，先在容许曲线间作变分，得到方程后再用弧长标记曲线上的点。
\item 图形若出现竖直切线或回折，可以换一个参数继续描述。坐标写法失效时，先检查曲线本身是否仍光滑。
\end{itemize}

\section*{本章小结}
\addcontentsline{toc}{section}{本章小结}
变分法把问题组织成四步：规定容许函数及边界条件；对路径族求一阶导数；分部积分以分离内部方程和边界项；最后用二阶变分、凸性、几何不等式或独立物理方法判断候选的性质。费马光程、最速降线时间、绳的势能和粒子作用量有相似数学结构，但物理假设、量纲与容许路径不同。

本章为下一章准备的是一套可复现的方法，而非一句“自然选择最优路径”。在第三章中，独立路径变量将换成广义坐标，固定端点变分将导出拉格朗日方程；本章的边界项、约束正则性和驻值判别仍然适用。

\section*{分层习题}
\addcontentsline{toc}{section}{分层习题}
先写清哪些端点固定、哪些量必须保持不变，再取一个满足这些条件的小扰动。求出驻值路径后，另用二阶变化或直接比较检查极小性。

\begin{bookexercises}
\item \exlevel{基础} \textbf{基础概念。} 对 $I[y]=\int_0^1y'^2\dd x$、$y=x$、$\eta=x(1-x)$，求 $I[y+\epsilon\eta]-I[y]$、一阶与二阶变分。解释一阶为零为何不等于有限变化为零。
\item \exlevel{基础} \textbf{反射。} 两点在平面镜同一侧，高度 $h_1,h_2>0$，水平间距 $d$。把镜面反射点作为变量，推导反射角等于入射角，并验证其极小性。
\item \exlevel{基础} \textbf{边值问题。} 求 $\int_0^1(y'^2+2xy)\dd x$ 在 $y(0)=0,y(1)=1$ 下的驻值路径和泛函值，证明唯一极小。
\item \exlevel{基础} \textbf{首次积分。} 推导贝尔特拉米恒等式，并对 $F=\tfrac12 y'^2$ 说明：若只解 $F-y'F_{y'}=C$ 而允许不光滑折线，为什么会引入原正则问题不容许的候选？
\item \exlevel{基础} \textbf{自然边界。} 求 $I=\tfrac12\int_0^1(y'^2+y^2)\dd x$ 在 $y(0)=1$、$y(1)$ 自由时的极小路径。再加入端点项 $\kappa y(1)^2/2$，$\kappa\ge0$，写出新边界条件。
\item \exlevel{综合} \textbf{横截。} 从固定点 $A=(a_0,b_0)$ 到直线 $y=kx+c$ 的最短连接，在能用单个图形覆盖且端点横坐标不受限制时，利用横截条件证明正交；再用矢量投影写出垂足，处理水平目标线导致最短连接竖直的情况。
\item \exlevel{综合} \textbf{积分约束。} 令 $I=\tfrac12\int_0^1y'^2\dd x$，$y(0)=y(1)=0$，并固定 $\int_0^1y\dd x=A$。求驻值路径、乘子及最小值。对 $A=0$ 是否仍可使用常规乘子？
\item \exlevel{综合} \textbf{圆弧。} 对 $d<\ell<\pi d/2$ 的图形面积问题，推导面积 $A=R^2(\alpha-\sin\alpha\cos\alpha)$，并说明 $\ell=\pi d/2$ 为何不属于端点斜率有界的光滑类。
\item \exlevel{综合} \textbf{悬链线。} 等高支点跨度 $2d$、绳长 $\ell>2d$。证明 $\sinh z/z=\ell/(2d)$ 对 $z=d/a>0$ 有唯一解，写出最低高度和水平张力。比较按水平距离均匀载荷与按弧长均匀自重的区别。
\item \exlevel{综合} \textbf{最速降线。} 验证摆线在 $0<\theta<2\pi$ 满足驻值方程，求 $\dd t/\dd\theta$，并在 $\theta\to0$ 展开 $x,y$，解释初始竖直切线和有限时间。
\item \exlevel{综合} \textbf{二阶变分。} 对谐振子零路径，分别在 $T=\pi/(2\omega)$、$T=\pi/\omega$、$T=3\pi/(2\omega)$ 判断极小或鞍点；给出最后一种情况中的一个负方向和一个正方向。
\item \exlevel{拓展} \textbf{共轭点。} 求谐振子雅可比方程满足 $\eta(0)=0$ 的非平凡解，列出所有正零点。解释“没有非零解”与“没有共轭点”为什么不同。
\item \exlevel{综合} \textbf{测地线。} 从球面度规计算非零的克里斯托费尔符号；用赤道弧长条件确定 $\varphi(s)$ 的速率。为什么长大圆弧满足方程但不是最短路径？
\item \exlevel{综合} \textbf{光线。} 对 $n=y_0/y$ 写出非竖直圆弧和竖直退化情形。对GRIN模型检查 $\alpha\to0$ 的极限，并说明导数是对空间坐标还是时间。
\item \exlevel{拓展} \textbf{场与控制。} (a) 在弦作用量中保留空间边界项，推导自由端的条件。(b) 固定 $\dot x=u$、$x(0)=0,x(T)=D$，在 $|u|\le u_{\max}$ 下，什么时候存在可行路径，什么时候恒定控制仍是全局最优？
\item \exlevel{拓展} \textbf{数值分析。} 独立推导本章两个离散线性系统，解释三次多项式特例的节点误差为何不能检验一般收敛阶。运行脚本检查双曲正弦模型的最大误差与误差比。
\end{bookexercises}

\section*{习题解答}
先用同一题号找到解答。核对时，从第一条方程开始比较自己的做法，再检查最后的符号、单位和适用条件。
\begin{bookexercises}
\item $\int_0^1\eta'\dd x=0$、$\int_0^1\eta'^2\dd x=1/3$，因此有限变化为 $\epsilon^2/3$，$\delta I=0,\delta^2I=2/3$。这是区分差值与方向导数的直接例子。

逐项展开为 $2\epsilon\int_0^1(1-2x)\dd x+\epsilon^2\int_0^1(1-2x)^2\dd x$。一阶与二阶变分分别是在 $\epsilon=0$ 的一阶与二阶导数。
\item 光程为 $n[\sqrt{\xi^2+h_1^2}+\sqrt{(d-\xi)^2+h_2^2}]$。一次导数为零给出两角正弦相等，在 $[0,\pi/2)$ 内角相等；二阶导数严格为正。也可将一个端点镜像后用直线最短验证。

设 $s_1=\sqrt{\xi^2+h_1^2}$、$s_2=\sqrt{(d-\xi)^2+h_2^2}$。光程导数为 $n[\xi/s_1-(d-\xi)/s_2]$，二阶导数为 $n(h_1^2/s_1^3+h_2^2/s_2^3)>0$。在 $0\le\xi\le d$，驻点条件化为 $\xi/h_1=(d-\xi)/h_2$，故 $\xi=dh_1/(h_1+h_2)$。
\item 欧拉--拉格朗日方程为 $2y''=2x$。积分两次得 $y=x^3/6+C_1x+C_0$，端点给 $C_0=0,C_1=5/6$。代入被积函数得 $7x^4/12+5x^2/2+25/36$，积分为 $74/45$。对任意同端点扰动 $h$，交叉项
\[
2\int_0^1(y'h'+xh)\dd x=2[y'h]_0^1+2\int_0^1(x-y'')h\dd x=0.
\]
因此 $I[y+h]-I[y]=\int_0^1h'^2\dd x$，等号只在 $h=0$ 时成立。
\item 使用乘积法则消去 $y''F_{y'}$。本例首次积分只给 $y'^2=\text{常量}$；允许斜率任意换号会产生折角，而正则欧拉--拉格朗日方程为 $y''=0$。解的光滑性和角点条件不能凭平方关系省略。

沿路径有 $(F-y'F_{y'})'=F_x+y'(F_y-(F_{y'})')$。当 $F_x=0$ 且内部方程成立时左端为零。对本题，在允许分段光滑的变分问题中，角点还要求 $F_{y'}=y'$ 连续，因此也不能任意换号。
\item 内部 $y''=y$，自由端 $y'(1)=0$，故 $y=\cosh(1-x)/\cosh1$。二次型严格为正，驻点是唯一极小。加端点代价后为 $y'(1)+\kappa y(1)=0$；不能仍取零斜率。

积分分部给 $\delta I=[y'\eta]_0^1+\int_0^1(y-y'')\eta\,\dd x$。$\eta(0)=0$，$\eta(1)$ 任意。通解可写成 $A\cosh(1-x)+B\sinh(1-x)$，从右端条件得 $B=0$，左端给 $A=1/\cosh1$。任意非零容许差函数的二次增量 $\frac12\int_0^1(\eta'^2+\eta^2)\dd x$ 严格为正。加端点代价时，边界系数增加 $\kappa y(1)$。
\item $1+ky'=0$（在 $k\ne0$ 的图形情形）。目标线法向为 $\boldsymbol n=(-k,1)$，垂足为 $B=A-\boldsymbol n(\boldsymbol n\cdot A-c)/(1+k^2)$。$k=0$ 时最短连接竖直，原来的 $y(x)$ 图形参数不适合，应使用矢量结果；若 $A$ 已在线上，最短长度为零。

端点沿目标线移动时 $\Delta y=k\Delta x$。边界项为 $(F-y'F_{y'})\Delta x+F_{y'}\Delta y$，取 $F=\sqrt{1+y'^2}$ 即得 $(1+ky')\Delta x/\sqrt{1+y'^2}$。其系数为零等价于路径切向 $(1,y')$ 与目标切向 $(1,k)$ 正交。
\item 用 $\widetilde F=y'^2/2-\lambda y$，得 $y''=-\lambda$，因此 $y=\lambda x(1-x)/2$。积分约束给 $\lambda=12A$，$I_{\min}=6A^2$。对约束差函数 $h$，交叉项为 $\lambda\int h\dd x=0$，余项 $\int h'^2/2$ 为正。$\delta J=\int\eta\dd x$ 不恒为零，与 $A=0$ 无关，正则性仍成立。
\item 圆弓面积为扇形 $R^2\alpha$ 减去三角形 $R^2\sin\alpha\cos\alpha$；$\alpha\to\pi/2$ 时端点切线竖直。再扩大圆弧须更换容许曲线的参数表示。
\item $\dd(\sinh z/z)/\dd z=(z\cosh z-\sinh z)/z^2>0$，因为分子从0起且导数为 $z\sinh z>0$。其值从1增至无穷；$a=d/z$，$y_{\min}=H-a(\cosh z-1)$、$H_T=\rho_\ell ga$。

对等高支点 $y(\pm d)=H$，取 $y=y_{\min}+a(\cosh(x/a)-1)$，沿弧长积分得 $\ell=2a\sinh(d/a)$。两种载荷下，竖直张力分量均为 $H_T y'$。若单位水平长度的向下载荷为 $w$，则 $H_T y''=w$，积分得抛物线。若单位弧长的自重为 $\rho_\ell g$，则 $H_T y''=\rho_\ell g\sqrt{1+y'^2}$，得到悬链线。
\item $y'=\sin\theta/(1-\cos\theta)$，平方满足原式，但在 $\theta=\pi$ 后符号改变。$x\sim a\theta^3/6,y\sim a\theta^2/2$，故 $\dd y/\dd x\sim2/\theta$；$t\sim\sqrt{a/g}\theta$，起点有限时间可离开。

用 $x=a(\theta-\sin\theta)$、$y=a(1-\cos\theta)$，有 $\dd s/\dd\theta=2a\sin(\theta/2)$，而 $v=\sqrt{2gy}=2\sqrt{ga}\sin(\theta/2)$。两者相除得到 $\dd t/\dd\theta=\sqrt{a/g}$，在整个开区间内成立。
\item 前两种分别为严格极小、非严格极小；第三种为鞍点。取 $\eta_1=\sin(\pi t/T)$ 给负方向，$\eta_2=\sin(2\pi t/T)$ 给正方向。

对两端为零的扰动 $\eta$，有 $\int\dot\eta^2\dd t\ge(\pi/T)^2\int\eta^2\dd t$，等号在基频正弦方向取得。作用量增量为 $\tfrac m2\int(\dot\eta^2-\omega^2\eta^2)\dd t$，由此判定前两种极小性。对第 $n$ 个正弦方向，该增量为 $\tfrac{mT}{4}[(n\pi/T)^2-\omega^2]$，在第三种情形代入 $n=1,2$ 即见一负一正。
\item $\eta=C\sin\omega t$，$C\ne0$，正零点为 $n\pi/\omega$。非平凡解始终存在；共轭点要求同一非平凡解除起点外再次为零。
\item $\Gamma^\theta_{\varphi\varphi}=-\sin\theta\cos\theta$，$\Gamma^\varphi_{\theta\varphi}=\Gamma^\varphi_{\varphi\theta}=\cot\theta$；赤道有 $R^2\varphi'^2=1$，故 $\varphi'=\pm1/R$。长弧可被同端点短弧替代，驻值不保证全局最短。

度规的唯一非恒定对角元为 $g_{\varphi\varphi}=R^2\sin^2\theta$。因此
\[
\Gamma^\theta_{\varphi\varphi}=-\frac{\partial_\theta g_{\varphi\varphi}}{2R^2},
\qquad
\Gamma^\varphi_{\theta\varphi}=\frac{\partial_\theta g_{\varphi\varphi}}{2g_{\varphi\varphi}}.
\]
其余由下标对称性或零导数确定。
\item $C\ne0$ 时为中心在 $y=0$ 的上半圆弧，$C=0$ 时 $x$ 恒定。GRIN解的导数对 $z$ 而非 $t$，$\alpha\to0$ 得 $x=x_0+\beta_0z$。

在 $y$ 单调的区段以 $y$ 为参数，$x$ 循环给出 $n x_y/\sqrt{1+x_y^2}=C$。令 $R=y_0/|C|$，当 $C\ne0$ 时 $x_y=\operatorname{sgn}(C)y/\sqrt{R^2-y^2}$。积分得 $(x-x_c)^2+y^2=R^2$，可在转向处换参数延拓。
\item (a) 边界项为 $[-\mathcal T u_x\eta]_0^L$ 在时间上积分，自由端要求 $\mathcal T u_x=0$。(b) 可行性要求 $|D|\le u_{\max}T$；满足时 $u=D/T$ 可行，并由柯西--施瓦茨不等式达到全局最小，若不满足则不存在可行控制。

(b) 积分约束是 $\int_0^T u\,\dd t=D$。控制上界给出所述可行性条件。由 $\int_0^T u^2\dd t\ge D^2/T$，且等号要求 $u$ 为常量，恒定控制达到二次代价的下界。
\item 置 $x_i=ih$、$h=1/N$。对每个内点 $y_i$，左右两条边的平方差项分别贡献 $2(y_i-y_{i-1})/h$、$2(y_i-y_{i+1})/h$。加上势项导数 $2hx_i$ 并置零，得
\[
y_{i-1}-2y_i+y_{i+1}=h^2x_i.
\]
对第二泛函，边项去掉因子2，内点代价导数为 $hy_i$，故 $-y_{i-1}+(2+h^2)y_i-y_{i+1}=0$。两题均将给定端点代入首末行。第一题的四阶导数为零，所以中心差分在解析节点上恰好满足方程。第二题用 $y(x_i)=\sinh x_i/\sinh1$ 比较，$N=20,40,80$ 时的最大节点误差依次为 $1.10469\times10^{-5}$、$2.76237\times10^{-6}$、$6.90896\times10^{-7}$，相邻误差比为 $3.9991,3.9982$，与二阶收敛一致。
\end{bookexercises}

\section*{进一步阅读与来源}
链接复核：2026-11-08。以上算例按本书符号重新推导，例题数值为教学设定。
\begin{bookreading}
\item MIT 16.323，\textit{Calculus of Variations}：\href{https://ocw.mit.edu/courses/16-323-principles-of-optimal-control-spring-2008/resources/lec5/}{官方讲义页面}。可进一步学习可变端点、横截条件及约束优化；课程层次高于本章，不要求照搬全部最优控制内容。
\item MIT 18.086，\textit{7.2 Calculus of Variations}：\href{https://ocw.mit.edu/courses/18-086-mathematical-methods-for-engineers-ii-spring-2006/e94c05947ed036cd6ad0150102087062_am72.pdf}{数学方法讲义}，用于继续理解连续变分与离散方程的联系。
\item David Tong，\textit{Classical Dynamics}：\href{https://davidtong.org/pdfs/teaching/classical-dynamics/clas2.pdf}{剑桥公开讲义}，进一步阅读驻作用量原理及其力学应用。
\item Richard Fitzpatrick，\textit{Conditional Variation}：\href{https://farside.ph.utexas.edu/teaching/sm1/Thermalhtml/node132.html}{得克萨斯大学公开讲义}，对照均匀链的势能及固定长度约束。
\end{bookreading}

\chapter{拉格朗日力学}\label{ch:lagrange}

\begin{guidebox}
描述一个摆，给出摆角就够了。若还用两个直角坐标，就要额外保证摆杆长度不变。拉格朗日方法先挑出刚好够用的独立变量，再用这些变量写动能和势能。这样常能在求轨道时先省去未知的反力。

本章先从摆和滑轮练习选坐标，再用虚功说明哪些力要进入方程。读到守恒量时，试着找出哪个坐标可以改变而不影响力学规律。若题目还要求绳张力或接触力，后面会用乘子把这些力求回来。
\end{guidebox}

\section{约束与广义坐标}

固定长度的摆只需给出一个角度，尽管摆球的位置有两个坐标。先看哪些坐标受限制，再挑出能独立描述运动的少数变量。

\subsection{约束及其分类}

摆杆的长度限制摆球能到哪里，理想滚动则直接限制接触点能怎样移动。遇到约束，先写下它限制的是位置、速度，还是某一侧的接触范围。

在第一章中，我们处理质点或质点系时，往往直接把受力逐个写出。对于由 $n$ 个质点组成的系统，原则上需要 $3n$ 个直角坐标。然而真实力学系统几乎总受到限制：刚体内部两点距离不变，轮子在地面上无滑动滚动，摆球被限制在圆周上，滑块被限制在导轨上。这些限制称为\textbf{约束}。

设第 $i$ 个质点的位矢为 $\br_i=(x_i,y_i,z_i)$，约束通常可以写成
\begin{equation}
f_\ell(\br_1,\br_2,\dots,\br_n,t)=0,
\qquad \ell=1,2,\dots,m.
\label{eq:constraint}
\end{equation}
若约束方程只含坐标而不含速度，称为\textbf{几何约束}；若约束方程中还出现速度，则称为\textbf{运动约束}或\textbf{微分约束}。若微分约束可以积分成式 \eqref{eq:constraint} 的形式，它本质上仍是几何约束，称为\textbf{完整约束}；若不能积分，则称为\textbf{非完整约束}。

\begin{definition}[完整约束与非完整约束]
若约束可以写成只含坐标和时间的等式
\[
f(\br_1,\dots,\br_n,t)=0,
\]
则称为\textbf{完整约束}。若约束必须写成不可积分的微分形式
\[
\sum_{i=1}^{n}\boldsymbol a_i(\br_1,\dots,\br_n,t)\cdot \dd\br_i
+b(\br_1,\dots,\br_n,t)\dd t=0,
\]
则称为\textbf{非完整约束}。
\end{definition}

约束对系统的作用力称为\textbf{约束力}，其余力称为\textbf{主动力}。约束力不是事先给定的力，而是由约束条件和运动状态共同决定的反力；这正是牛顿力学处理受约束系统时最不方便的地方。

例如，质点被限制在半径为 $R$ 的球面上，约束为
\[
x^2+y^2+z^2-R^2=0,
\]
这是完整约束。若球面半径随时间变化，$R=R(t)$，约束仍可写成
\[
x^2+y^2+z^2-R^2(t)=0,
\]
仍是完整约束，只是显含时间。相反，半径为 $R$ 的竖直圆盘在水平面上无滑动滚动时，接触点的速度为零，得到
\[
\dot x-R\dot\varphi\cos\theta=0,\qquad
\dot y-R\dot\varphi\sin\theta=0,
\]
其中 $\theta$ 是圆盘平面与固定方向的夹角，$\varphi$ 是圆盘绕自身对称轴的自转角。这两个微分约束一般不能积分为只含坐标的代数方程，因此是非完整约束。

\begin{figure}[htbp]
\centering
\begin{tikzpicture}[>=Stealth]
\draw[thick,blue!60!black] (0,0) circle (2);
\draw[dashed,gray] (0,0) circle (2.6);
\coordinate (P) at (1.6,1.2);\fill (P) circle (2pt) node[below]{$P$};
\draw[->,thick] (0,0)--(P) node[midway,below]{$R(t)$};
\draw[->,very thick,red!70!black] (P)--(1,2) node[above left]{$\delta\br$};
\draw[->,very thick,green!40!black] (P)--(2.24,1.68) node[right]{$\dd\br$};
\node at (-1.25,-.7) {冻结的约束};\node[gray] at (-1,-2.85) {下一时刻的半径};
\end{tikzpicture}
\caption{半径受驱动增大的圆周约束。冻结时刻的虚位移沿切向；图示实际瞬时角速度为零，真实位移随半径增大而沿径向。两种位移不能互相替代。}
\label{fig:virtual-displacement}
\end{figure}

约束还可以按其他方式分类：
\begin{itemize}[leftmargin=1.8em]
  \item \textbf{定常约束}：约束方程不显含时间，$f(\br)=0$；
  \item \textbf{非定常约束}：约束方程显含时间，$f(\br,t)=0$；
  \item \textbf{双侧约束}：约束是等式，质点不能从两侧脱离；
  \item \textbf{单侧约束}：约束是不等式，例如 $z\ge 0$，质点可以离开约束面；
  \item \textbf{理想约束}：约束力在任意虚位移上所作的总虚功为零。
\end{itemize}

理想约束是本章最重要的结构性假设。光滑曲面、光滑铰链、刚性杆、不可伸长的轻绳都近似为理想约束。对于理想约束，约束力的总虚功为零，因此它不会进入广义坐标下的运动方程；我们只需处理主动力。滑动摩擦通常须作为主动力处理；无滑动滚动中的静摩擦在适当理想接触模型下可不做虚功，不能单凭“有摩擦”就判为非理想。绳的拉紧条件和单侧接触还须检验反力可实现性。

\begin{physicsbox}
“约束力不做虚功”并不意味着约束力不做功。真实运动中，约束力完全可能做功，例如外部机构驱动的光滑导轨可通过法向反力向滑块传递能量。虚功中的位移是同一时刻、只满足约束的假想位移，时间被冻结，所以法向约束力与虚位移正交，虚功才为零。
\end{physicsbox}

\subsection{自由度与广义坐标}

固定长摆杆后，摆球的横、竖坐标不能独立变化。可以用角度代替这两个坐标；有多个约束时，还得检查它们是否真的给出了不同的限制。

设有 $n$ 个质点、$m$ 个完整约束。若在所考察的正则位形邻域中，这些约束对位置坐标的雅可比矩阵具有秩 $m$，则独立坐标数为
\begin{equation}
s=3n-m.
\label{eq:dof}
\end{equation}
$s$ 称为系统的\textbf{自由度}。我们可以选一组 $s$ 个独立参数
\[
q_1,q_2,\dots,q_s,
\]
使得所有质点的位矢都可以表示为
\begin{equation}
\br_i=\br_i(q_1,q_2,\dots,q_s,t),
\qquad i=1,2,\dots,n.
\label{eq:generalized-coordinates}
\end{equation}
这组参数称为\textbf{广义坐标}。它们通常只是位形空间的局部坐标。机械系统中可以是长度或角度；扩展到机电模型时也可以用电荷，但不能把不同模型的自由变量混为一谈。重要的是它们彼此独立，并且自动满足全部完整约束。同一系统可以选多组不同的广义坐标，物理结果不依赖于坐标选择；好的选择通常能让 $T$ 和 $V$ 更简单。

对时间求导，得到\textbf{广义速度}
\[
\dot q_a=\dv{q_a}{t},\qquad a=1,2,\dots,s.
\]
质点的速度可以写成
\begin{equation}
\bv_i=\sum_{a=1}^{s}\pdv{\br_i}{q_a}\dot q_a+\pdv{\br_i}{t}.
\label{eq:velocity-generalized}
\end{equation}
其中第一项是广义速度引起的部分，第二项来自参数化的显式时间依赖，可由移动约束或含时坐标选择引起。

\begin{example}[球坐标与双摆]
质量为 $m$ 的质点被约束在球面上，取球坐标 $(r,\theta,\varphi)$，约束为 $r=R$。自由度 $s=2$，广义坐标可取 $q_1=\theta,q_2=\varphi$。若取 $r$ 也作为广义坐标，则约束没有自动满足，后续必须用乘子处理。

平面双摆由两根长为 $l_1,l_2$ 的轻杆和两个质量 $m_1,m_2$ 组成。四个直角坐标减去两个杆长约束，自由度 $s=2$，广义坐标可取两杆与竖直方向的夹角 $\theta_1,\theta_2$。这两个角度自动满足杆长约束，因此拉格朗日方程中不会出现杆的张力。
\end{example}

\subsection{位形空间}

把系统的每一种姿态看作一个点。双摆的一点就由两个角度标记；摆动时，这个点会在由所有可能姿态组成的空间中画出一条线。

在正则完整约束下，容许位形的集合构成 $s$ 维的\textbf{位形空间}。一组独立广义坐标只给出局部表示。系统的运动对应位形空间中的一条曲线 $q_a=q_a(t)$。局部坐标不必覆盖全部位形，例如摆角按 $2\pi$ 周期识别，不能把坐标图边界当成物理边界。运动由方程和初始数据确定；作用量的固定端点表述则比较连接给定端点的路径。本章将使用上一章的变分法说明二者的关系，固定两端一般不保证解存在或唯一。

\section{虚功原理与达朗贝尔原理}

铁丝的推力可以很大，但它若始终垂直于珠子可移动的方向，就不会出现在沿铁丝的运动方程里。虚功把这种投影推广到多坐标系统。

\subsection{虚位移与虚功}

先把时钟按暂停，再让系统沿约束容许的方向试着挪一点，这就是虚位移。若约束本身正在移动，真实速度还会包含约束的运动。

设系统在时刻 $t$ 处于满足约束的位形。想象在时间 $t$ 冻结的条件下，系统发生一个无穷小位移 $\delta\br_i$，它仍然满足约束的线性化条件。这样的位移称为\textbf{虚位移}。对完整约束，虚位移满足 $\sum_i(\partial f_\ell/\partial\br_i)\cdot\delta\br_i=0$，位于冻结约束的切空间；对本章采用的理想、仿射速度非完整模型，容许虚位移由达朗贝尔规则规定为约束系数的齐次核；它一般不是比较两条都满足速度约束的路径得到的任意变分。详见本章非完整系统一节。虚位移与真实位移 $\dd\br_i$ 的关键区别是：
\begin{equation}
\delta t=0,
\qquad
\dd\br_i=\sum_a\pdv{\br_i}{q_a}\dd q_a+\pdv{\br_i}{t}\dd t,
\qquad
\delta\br_i=\sum_a\pdv{\br_i}{q_a}\delta q_a.
\label{eq:virtual-vs-real}
\end{equation}
真实位移可能包含约束随时间变化产生的项，虚位移没有这一项。若选用无显式时间依赖的定常参数化，$\partial\br_i/\partial t=0$，两者的容许切向形式相同，但真实位移由实际速度确定，虚位移则仍可独立选取；对于非定常约束，两者一般不同。

作用在第 $i$ 个质点上的力为 $\bF_i$，在虚位移上作的功
\[
\delta W=\sum_{i=1}^{n}\bF_i\cdot\delta\br_i
\]
称为\textbf{虚功}。把力分成主动力 $\bF_i^{(a)}$ 和约束力 $\bN_i$，则
\[
\delta W=\sum_i\bF_i^{(a)}\cdot\delta\br_i
+\sum_i\bN_i\cdot\delta\br_i.
\]
\subsection{虚功原理}

在静力问题中，想象系统在每一个允许方向上试着移动一点。如果所有主动力的虚功之和都为零，沿这些方向就没有净推动。

若约束是理想的，第二项为零。于是虚功原理可以表述为：

\begin{keybox}
受理想双侧约束的力学系统，在静止的定常约束上、排除奇异位形时，主动力可由理想约束反力平衡的条件是：对任意满足约束（或其线性化条件）的虚位移，主动力所作的总虚功为零，
\begin{equation}
\delta W=\sum_{i=1}^{n}\bF_i^{(a)}\cdot\delta\br_i=0.
\label{eq:virtual-work}
\end{equation}
\end{keybox}

用广义坐标表示，$\delta\br_i=\sum_a(\partial\br_i/\partial q_a)\delta q_a$，于是
\begin{equation}
\delta W=\sum_{a=1}^{s}Q_a\,\delta q_a,
\qquad
Q_a=\sum_{i=1}^{n}\bF_i^{(a)}\cdot\pdv{\br_i}{q_a}.
\label{eq:generalized-force}
\end{equation}
$Q_a$ 称为与 $q_a$ 对应的\textbf{广义力}。广义力的量纲一般不是力的量纲，因为 $q_a$ 的量纲不一定是长度；但乘积 $Q_a\delta q_a$ 必须具有功的量纲。由于 $\delta q_a$ 彼此独立，平衡条件 $\delta W=0$ 等价于
\begin{equation}
Q_a=0,\qquad a=1,2,\dots,s.
\label{eq:equilibrium-generalized}
\end{equation}
这就是广义坐标下的平衡方程。

\begin{example}[两杆平衡：同一正切值有不同稳定性]
固定铰支点 $O$，等长均匀杆 $OA,AB$ 长均为 $l$、重量为 $P_1,P_2$，铰链理想。取 $x$ 向右、$y$ 向上，两杆绝对角从 $+x$ 逆时针量起，记为 $\alpha,\beta$。在 $B$ 施加恒定向左力 $F>0$，其势能为 $Fx_B$。两杆质心高度分别为 $l\sin\alpha/2$ 与 $l\sin\alpha+l\sin\beta/2$，所以总势为
\[
U=l[(P_1/2+P_2)\sin\alpha+F\cos\alpha]
+l[(P_2/2)\sin\beta+F\cos\beta].
\]
虚功为 $-\delta U$，两坐标独立，平衡要求
\[
(P_1/2+P_2)\cos\alpha-F\sin\alpha=0,
\qquad(P_2/2)\cos\beta-F\sin\beta=0.
\]
只写正切会丢失象限。令 $A_1=P_1/2+P_2,A_2=P_2/2$，稳定下垂支应取
\[
\cos\alpha=-\frac{F}{\sqrt{F^2+A_1^2}},\quad
\sin\alpha=-\frac{A_1}{\sqrt{F^2+A_1^2}},
\]
$\beta$ 将 $A_1$ 换为 $A_2$。两角均在第三象限（或等价的负角区间）。此时 $U_{\alpha\alpha}=l\sqrt{F^2+A_1^2}>0$、$U_{\beta\beta}=l\sqrt{F^2+A_2^2}>0$，交叉导数为零，确为势能极小。两角同时取相反方向也满足相同的正切式，却是极大分支。

$F\to0$ 时稳定支竖直向下；$F$ 很大时向左近乎水平。消去铰链反力简化了计算，但不会替代平衡稳定性和几何分支的检查。
\end{example}

\begin{figure}[htbp]
\centering
\begin{tikzpicture}[>=Stealth]
\coordinate (O) at (0,0);\coordinate (A) at (-1.2,-1.6);\coordinate (B) at (-3.027623097,-2.412276932);
\draw[very thick,blue!60!black] (O)--(A);\draw[very thick,red!60!black] (A)--(B);
\draw[->] (O)--(1.4,0) node[right]{$x$};\draw[->] (O)--(0,1.1) node[above]{$y$};
\fill (O) circle (2pt) node[above left]{$O$};\fill (A) circle (2pt) node[right]{$A$};\fill (B) circle (2pt) node[below]{$B$};
\draw[->,thick,green!40!black] (B)--++(-.9,0) node[left]{$F$};
\draw[->,thick] (-.6,-.8)--++(0,-.8) node[right]{$P_1$};
\draw[->,thick] (-2.113811549,-2.006138466)--++(0,-.8) node[right]{$P_2$};
\node at (-1.5,.65) {稳定支向左下垂};
\end{tikzpicture}
\caption{对应稳定分支的两杆示意，取 $P_1=P_2$ 且 $F=9P_1/8$。两杆长度相等，重力向下、端点力向左，绝对角的象限与正文稳定解一致。}
\label{fig:two-rods}
\end{figure}

\subsection{达朗贝尔原理}

把牛顿方程中的质量乘加速度移到力的一侧，再投影到允许方向。这样就能沿用虚功的写法处理有加速度的物体。

虚功原理原本只适用于静力学。达朗贝尔在 1743 年注意到：若把惯性力 $-m_i\ddot\br_i$ 也看作一种“力”，则动力学问题可以转化为瞬时静力学问题。定义\textbf{达朗贝尔惯性力}
\[
\bF_i^{\rm inertia}=-m_i\ddot\br_i,
\]
则对每个质点有
\[
\bF_i+\bN_i+\bF_i^{\rm inertia}=0,
\]
其中 $\bF_i$ 为主动力，$\bN_i$ 为约束力。在理想约束下，约束力的总虚功为零，因此
\begin{equation}
\sum_{i=1}^{n}\left(\bF_i-m_i\ddot\br_i\right)\cdot\delta\br_i=0.
\label{eq:dAlembert}
\end{equation}
这就是\textbf{达朗贝尔原理}。它不是新的物理定律，而是牛顿第二定律与理想约束假设的结合；其价值在于把约束力消去，并把动力学问题写成变分形式。

\begin{physicsbox}
达朗贝尔原理的“惯性力”不是参考系变换意义下的惯性力，而是一个形式上的记账项。它的作用是把 $m\ddot\br$ 从方程右边移到左边，使动力学方程看起来像平衡方程。真正进入虚功的是主动力与惯性力的合力。
\end{physicsbox}

\section{拉格朗日方程}

一旦选好坐标，只要算出动能、势能和必要的外力投影，就能系统地写方程。这一节给出两条来路，再用具体模型练习。

\subsection{从达朗贝尔原理导出}

对某个广义坐标作一点改动，先问每个质点会往哪个方向移动。把牛顿方程投影到这些方向，再整理成动能的导数。

我们已把位矢写成广义坐标的函数 $\br_i=\br_i(q,t)$。速度
\[
\bv_i=\dot\br_i=\sum_a\pdv{\br_i}{q_a}\dot q_a+\pdv{\br_i}{t}.
\]
由此得到两个关键恒等式：
\begin{equation}
\pdv{\bv_i}{\dot q_a}=\pdv{\br_i}{q_a},
\qquad
\frac{\dd}{\dd t}\left(\pdv{\br_i}{q_a}\right)=\pdv{\bv_i}{q_a}.
\label{eq:kinematic-identities}
\end{equation}
这不是对任意函数都可随意交换两个算子。这里逐项展开为
\[
\frac{\partial\bv_i}{\partial q_a}
=\sum_b\frac{\partial^2\br_i}{\partial q_a\partial q_b}\dot q_b
+\frac{\partial^2\br_i}{\partial q_a\partial t}
=\frac{\dd}{\dd t}\left(\frac{\partial\br_i}{\partial q_a}\right),
\]
偏导时保持 $\dot q$ 不变，并利用 $\br_i(q,t)$ 的混合偏导连续。第一式则直接来自速度对 $\dot q_a$ 的线性依赖。

系统的动能
\[
T=\frac12\sum_i m_i\bv_i^2.
\]
利用式 \eqref{eq:kinematic-identities}，
\[
\pdv{T}{\dot q_a}=\sum_i m_i\bv_i\cdot\pdv{\bv_i}{\dot q_a}
=\sum_i m_i\bv_i\cdot\pdv{\br_i}{q_a},
\]
\[
\pdv{T}{q_a}=\sum_i m_i\bv_i\cdot\pdv{\bv_i}{q_a}.
\]
于是
\[
\frac{\dd}{\dd t}\left(\pdv{T}{\dot q_a}\right)-\pdv{T}{q_a}
=\sum_i m_i\ddot\br_i\cdot\pdv{\br_i}{q_a}.
\]
另一方面，达朗贝尔原理中的惯性力项为
\[
-\sum_i m_i\ddot\br_i\cdot\delta\br_i
=-\sum_a\left[\sum_i m_i\ddot\br_i\cdot\pdv{\br_i}{q_a}\right]\delta q_a.
\]
主动力的虚功为 $\sum_a Q_a\delta q_a$。代入式 \eqref{eq:dAlembert}，并利用 $\delta q_a$ 的独立性，得到
\begin{equation}
\boxed{
\frac{\dd}{\dd t}\left(\pdv{T}{\dot q_a}\right)-\pdv{T}{q_a}=Q_a,
\qquad a=1,2,\dots,s.
}
\label{eq:lagrange-general}
\end{equation}
这就是\textbf{拉格朗日方程}的基本形式。它把 $3n$ 个直角坐标的牛顿方程压缩成 $s$ 个广义坐标方程，约束力自动消失。

若主动力可用位置势 $V(q,t)$ 表示，且没有额外做虚功的非保守力，则
\[
Q_a=-\pdv{V}{q_a}.
\]
定义\textbf{拉格朗日函数}
\begin{equation}
L(q,\dot q,t)=T(q,\dot q,t)-V(q,t),
\label{eq:lagrangian}
\end{equation}
则式 \eqref{eq:lagrange-general} 化为
\begin{equation}
\boxed{
\frac{\dd}{\dd t}\left(\pdv{L}{\dot q_a}\right)-\pdv{L}{q_a}=0,
\qquad a=1,2,\dots,s.
}
\label{eq:lagrange-euler}
\end{equation}
这是分析力学中最常用的形式。

\begin{keybox}
拉格朗日方程的三个要点：
\begin{enumerate}[leftmargin=1.6em]
  \item 只出现广义坐标、广义速度和广义力，不出现理想约束力；
  \item 对完整系统，方程数目等于自由度 $s$，而不是质点数 $3n$；
  \item 对质量恒定、完整理想约束且主动力由位置势描述的模型，写出 $T,V$ 与初始数据即可建立独立坐标方程；但接触保持、摩擦可行性或支承载荷仍可能需要求约束力。
\end{enumerate}
\end{keybox}

\subsection{从哈密顿原理导出}\label{sec:ham-principle}
把第二章的变分步骤用到动能减势能的积分。两端固定后，边界项消失，内部任意的小改动就给出运动方程。

第\ref{sec:action-first}节已给出一维固定端点的分部积分。这里在独立完整坐标中使用同一边界条件；下文的作用量变分应与式 \eqref{eq:virtual-vs-real} 的冻结时间虚位移区分。

拉格朗日方程也可以从整体变分原理得到。定义\textbf{作用量}
\begin{equation}
S[q]=\int_{t_1}^{t_2}L(q,\dot q,t)\,\dd t.
\label{eq:action}
\end{equation}
在所有满足端点条件
\[
q_a(t_1)=q_{a1},\qquad q_a(t_2)=q_{a2}
\]
的曲线中，真实运动使作用量取驻值：
\begin{equation}
\delta S=0.
\label{eq:hamilton-principle-lagrange}
\end{equation}
这是通常的\textbf{哈密顿驻作用量原理}。此处使用独立广义坐标、固定时间和位置端点；$L$ 可显含时间，也可包含电磁速度耦合，不能把它仅限于不显含时间的位置保守势。任意耗散力则不能自动由这个简单作用量给出。把 $\delta S$ 展开，并注意变分与时间导数可交换：$\delta\dot q_a=\frac{\dd}{\dd t}\delta q_a$，
\[
\delta S=\int_{t_1}^{t_2}\sum_a
\left(\pdv{L}{q_a}\delta q_a+\pdv{L}{\dot q_a}\delta\dot q_a\right)\dd t.
\]
对第二项分部积分，并利用端点处 $\delta q_a(t_1)=\delta q_a(t_2)=0$：
\[
\delta S=\int_{t_1}^{t_2}\sum_a
\left[\pdv{L}{q_a}-\frac{\dd}{\dd t}\left(\pdv{L}{\dot q_a}\right)\right]\delta q_a\,\dd t.
\]
由于 $\delta q_a$ 在内部彼此独立且任意，括号必须为零，即得到式 \eqref{eq:lagrange-euler}。在上述完整、理想且力可由所选 $L$ 描述的条件下，两种表述给出同一组方程；前者是“逐时刻”的表述，后者是“整条路径”的表述。

\begin{remark}
哈密顿原理中的端点固定，但端点处的速度不固定。变分只比较在相同时间间隔内、从同一起点到同一终点的路径。若端点时间也可变，则会得到更一般的能量关系；这一点在第六章讨论哈密顿-雅可比理论时还会出现。
\end{remark}

\subsection{非保守广义力与受驱动运动}
摩擦、阻尼和电机力矩常直接按力来给。先算它们沿允许位移做多少功，再把每个坐标的位移前面的系数读作相应的广义力。

若部分力已用 $L$ 表示，另有外力、阻尼或控制力矩，其虚功为 $\sum_a Q_a^{\rm nc}\delta q_a$，则
\begin{equation}\label{eq:lag-nonconservative}
\frac{\dd}{\dd t}L_{\dot q_a}-L_{q_a}=Q_a^{\rm nc}.
\end{equation}
这里不是令非保守作用量驻值，而是采用拉格朗日--达朗贝尔表述 $\delta\int L\dd t+\int\sum_a Q_a^{\rm nc}\delta q_a\dd t=0$。额外力不应既写进势又在右端重复加入。

若黏滞阻尼的广义力可写为 $Q_a^{\rm damp}=-\partial\mathcal D/\partial\dot q_a$，则用耗散函数 $\mathcal D=\frac12\sum_{ab}C_{ab}\dot q_a\dot q_b$ 表示。$C$ 对称半正定时，$\sum_a Q_a^{\rm damp}\dot q_a=-2\mathcal D\le0$；$\mathcal D$ 的量纲为功率，不是势能。它描述这一类黏滞模型，不能自动包括静摩擦阈值或所有耗散。
\begin{example}[有阻尼和关节驱动力矩的单摆]
固定摆长 $l$，角度从向下方向量起，$L=ml^2\dot\theta^2/2+mgl\cos\theta$。若关节黏滞力矩为 $-b\dot\theta$，驱动力矩为 $\tau(t)$，则
\[
ml^2\ddot\theta+b\dot\theta+mgl\sin\theta=\tau(t).
\]
以最低点势能为零的机械能 $E=ml^2\dot\theta^2/2+mgl(1-\cos\theta)$ 满足
\[
\dot E=-b\dot\theta^2+\tau\dot\theta.
\]
无驱动时单调耗散；驱动可输入或取出能量。$b=\tau=0$ 恢复第一章单摆。这既训练广义力矩，也给数值计算提供可检查的能量收支。
\end{example}

\subsection{广义动量和循环坐标的判定}
当位置用角度描述时，与角速度对应的动量常是角动量。把动能和其他耦合项写进拉格朗日量后，可统一用对速度的偏导确定广义动量。

定义\textbf{广义动量}
\begin{equation}
p_a=\pdv{L}{\dot q_a}.
\label{eq:generalized-momentum}
\end{equation}
广义动量的量纲由对应的坐标决定，因为 $q_a$ 的量纲不一定是长度；它与广义坐标共轭，满足 $p_a\delta q_a$ 具有作用量的量纲。若 $L$ 不显含某个广义坐标 $q_a$，即
\[
\pdv{L}{q_a}=0,
\]
则 $q_a$ 称为\textbf{循环坐标}或\textbf{可遗坐标}。若该坐标没有额外广义力，则拉格朗日方程给出
\[
\frac{\dd}{\dd t}\left(\pdv{L}{\dot q_a}\right)=0,
\]
于是
\begin{equation}
p_a=\text{常量}.
\label{eq:cyclic-integral}
\end{equation}
这就是\textbf{循环积分}。

\subsection{典型应用}

每做一题先画坐标方向，写出位置，再求速度和能量。最后检查所得方程在某个简单极限下是否回到熟悉的模型。

本节依次训练三个步骤：先用坐标自动满足约束，再从惯性速度构造动能，最后以守恒量或牛顿方程独立验算。所有绳、杆均按题目理想化，工程背景只用于说明模型用途，不代表具体设备的实测参数。

\begin{example}[阿特伍德机]
质量分别为 $m_1,m_2$ 的两物体用不可伸长的轻绳跨过定滑轮，绳与滑轮间无滑动，滑轮质量和转动惯量均忽略不计。取 $q$ 为 $m_1$ 下落的距离，则 $m_2$ 上升 $q$。动能和势能分别为
\[
T=\frac12(m_1+m_2)\dot q^2,
\qquad
V=-m_1gq+m_2gq.
\]
拉格朗日函数
\[
L=T-V=\frac12(m_1+m_2)\dot q^2+(m_1-m_2)gq.
\]
代入式 \eqref{eq:lagrange-euler}：
\[
(m_1+m_2)\ddot q=(m_1-m_2)g,
\]
即
\[
\ddot q=\frac{m_1-m_2}{m_1+m_2}g.
\]
绳的张力没有出现在方程中；若要它，需回到约束力或使用乘子法。

为独立复核，取 $m_1$ 向下为正，写 $m_1g-\mathcal T=m_1\ddot q$、$\mathcal T-m_2g=m_2\ddot q$。相加给同一 $\ddot q$，再回代得 $\mathcal T=2m_1m_2g/(m_1+m_2)$。
\end{example}

\begin{figure}[htbp]
\centering
\begin{tikzpicture}[>=Stealth]
\draw[thick] (-1.4,3.8)--(1.4,3.8);\draw[thick] (0,3.8)--(0,3);
\draw[thick,blue!60!black] (0,3) circle (.5);
\draw[thick] (-.5,3)--(-.5,1.2);\draw[thick] (.5,3)--(.5,1.8);
\draw[fill=red!12] (-.85,.6) rectangle (-.15,1.2);\node at (-.5,.9) {$m_1$};
\draw[fill=green!12] (.15,1.2) rectangle (.85,1.8);\node at (.5,1.5) {$m_2$};
\draw[->,thick] (-1.1,1.4)--(-1.1,.5) node[below]{$q$};
\draw[->,thick] (1.1,.8)--(1.1,1.7) node[above]{$q$};
\end{tikzpicture}
\caption{质量与转动惯量均忽略的定滑轮模型。两段竖直绳与轮相切，一质量向下时另一质量等量向上，坐标符号与势能定义相同。}
\label{fig:atwood}
\end{figure}

\begin{example}[双滑轮三质量系统]
一个定滑轮 $O$ 固定，一根轻绳跨过 $O$，左端挂质量为 $m_1$ 的物体，右端悬挂一个动滑轮 $B$；第二根轻绳跨过 $B$，两端分别挂质量为 $m_2,m_3$ 的物体。绳与滑轮间无滑动，绳不可伸长，滑轮质量和转动惯量均忽略。取 $q_1$ 为 $m_1$ 向下的位移，取 $q_2$ 为 $m_2$ 相对动滑轮 $B$ 向下的位移。求系统的运动方程。

\begin{center}
\begin{tikzpicture}[>=Stealth]
\draw[thick] (-1.5,4.7)--(1.7,4.7);\draw[thick] (0,4.7)--(0,4);
\draw[thick,blue!60!black] (0,4) circle (.45);
\draw[thick] (-.45,4)--(-.45,1.8);\draw[thick] (.45,4)--(.45,2.6);
\draw[thick,blue!60!black] (.45,2.6) circle (.35);\fill (.45,2.6) circle (1.5pt) node[right=12pt]{$B$};
\draw[thick] (.1,2.6)--(.1,1.3);\draw[thick] (.8,2.6)--(.8,.8);
\draw[fill=red!12] (-.7,1.1) rectangle (-.2,1.8);\node at (-.45,1.45) {$m_1$};
\draw[fill=green!12] (-.15,.6) rectangle (.35,1.3);\node at (.1,.95) {$m_2$};
\draw[fill=orange!12] (.55,.1) rectangle (1.05,.8);\node at (.8,.45) {$m_3$};
\draw[->,thick] (-1,1.8)--(-1,.8) node[below]{$q_1$};
\node at (.1,4.85) {$O$};
\end{tikzpicture}
\end{center}

由于第一根绳长不变，$B$ 的速度为 $-\dot q_1$（取向下为正）。第二根绳长不变，因此
\[
v_2=-\dot q_1+\dot q_2,
\qquad
v_3=-\dot q_1-\dot q_2.
\]
系统的动能和势能分别为
\[
T=\frac12m_1\dot q_1^2+\frac12m_2(-\dot q_1+\dot q_2)^2
+\frac12m_3(-\dot q_1-\dot q_2)^2,
\]
\[
V=-m_1gq_1-m_2g(-q_1+q_2)-m_3g(-q_1-q_2).
\]
拉格朗日函数 $L=T-V$。代入拉格朗日方程，得到
\[
(m_1+m_2+m_3)\ddot q_1+(m_3-m_2)\ddot q_2
=(m_1-m_2-m_3)g,
\]
\[
(m_3-m_2)\ddot q_1+(m_2+m_3)\ddot q_2
=(m_2-m_3)g.
\]
写成矩阵形式：
\[
\begin{pmatrix}
m_1+m_2+m_3 & m_3-m_2\\
m_3-m_2 & m_2+m_3
\end{pmatrix}
\begin{pmatrix}\ddot q_1\\ \ddot q_2\end{pmatrix}
=g\begin{pmatrix}m_1-m_2-m_3\\ m_2-m_3\end{pmatrix}.
\]
两个广义坐标自动满足两根绳的约束，绳中张力不必出现在方程中。

\begin{modernbox}
这类“多绳—多滑轮”系统是电梯配重、塔吊起升、缆索驱动并联机器人和空间吊装系统的基本模型。工程上关心的是各段绳的张力、负载的加速度和系统的固有频率；拉格朗日方法先消去绳张力，再用矩阵方程统一求解，比逐段列牛顿方程更适合多自由度系统。
\end{modernbox}
\end{example}

\begin{example}[球坐标中的质点]
质量为 $m$ 的质点在中心力势 $V(r)$ 中运动，取球坐标 $(r,\theta,\varphi)$ 为广义坐标，并把极轴取为 $z$ 轴。速度为
\[
\bv=\dot r\,\er+r\dot\theta\,\etheta+r\sin\theta\,\dot\varphi\,\ephi,
\]
动能为
\[
T=\frac12m\left(\dot r^2+r^2\dot\theta^2+r^2\sin^2\theta\,\dot\varphi^2\right).
\]
拉格朗日函数
\[
L=\frac12m\left(\dot r^2+r^2\dot\theta^2+r^2\sin^2\theta\,\dot\varphi^2\right)-V(r).
\]
$\varphi$ 是循环坐标，因此
\[
p_\varphi=\pdv{L}{\dot\varphi}=mr^2\sin^2\theta\,\dot\varphi
\]
守恒，这就是角动量的 $z$ 分量。对 $\theta$ 的方程为
\[
\frac{\dd}{\dd t}\left(mr^2\dot\theta\right)
-mr^2\sin\theta\cos\theta\,\dot\varphi^2=0.
\]
在 $r>0$ 且方程正则的区域，若初始 $\theta=\pi/2,\dot\theta=0$，赤道平面是方程的不变子集；径向距离不必固定。对 $r$ 的方程给出
\[
m\ddot r-mr\dot\theta^2-mr\sin^2\theta\,\dot\varphi^2
=-\pdv{V}{r}.
\]
这正是中心力场中的径向方程。
\end{example}

\begin{example}[双摆：把耦合项逐项算出来]
两根无质量刚杆长为 $l_1,l_2$，端点质量为 $m_1,m_2>0$，固定支点，运动限于竖直平面。两绝对角 $\theta_1,\theta_2$ 均从向下竖直方向量起，坐标为
\[
\br_1=(l_1\sin\theta_1,-l_1\cos\theta_1),\qquad
\br_2=\br_1+(l_2\sin\theta_2,-l_2\cos\theta_2).
\]
令 $\Delta=\theta_1-\theta_2$，微分后由点积 $\bv_1\cdot(\bv_2-\bv_1)=l_1l_2\dot\theta_1\dot\theta_2\cos\Delta$ 得
\begin{align}
T&=\tfrac12(m_1+m_2)l_1^2\dot\theta_1^2+\tfrac12m_2l_2^2\dot\theta_2^2
+m_2l_1l_2\cos\Delta\,\dot\theta_1\dot\theta_2,\\
V&=-(m_1+m_2)gl_1\cos\theta_1-m_2gl_2\cos\theta_2.
\end{align}
取 $D=m_2l_1l_2$，第一方程的动量和坐标导数是
\[
p_1=(m_1+m_2)l_1^2\dot\theta_1+D\cos\Delta\dot\theta_2,
\quad L_{\theta_1}=-D\sin\Delta\dot\theta_1\dot\theta_2-(m_1+m_2)gl_1\sin\theta_1.
\]
对 $p_1$ 求时间导数时还须微分 $\cos\Delta$。与 $L_{\theta_1}$ 相减后，两个含 $\dot\theta_1\dot\theta_2$ 的项抵消，得到
\begin{align}\label{eq:lag-double-pendulum}
(m_1+m_2)l_1^2\ddot\theta_1+D\cos\Delta\ddot\theta_2
+D\sin\Delta\dot\theta_2^2+(m_1+m_2)gl_1\sin\theta_1&=0,\nonumber\\
m_2l_2^2\ddot\theta_2+D\cos\Delta\ddot\theta_1
-D\sin\Delta\dot\theta_1^2+m_2gl_2\sin\theta_2&=0.
\end{align}
第二式由相同方法得到，注意两项正负号。写成 $M(\theta)\ddot\theta+b(\theta,\dot\theta)=0$，质量矩阵行列式为
\[
\det M=m_2l_1^2l_2^2[m_1+m_2\sin^2\Delta]>0.
\]
因此物理参数为正时可解出两角加速度；不能把速度耦合项略去后仍声称是精确双摆。无驱动、无阻尼下 $T+V$ 守恒；小角展开将导向第四章耦合小振动。加入关节力矩后，右侧为各角对应的广义力矩，是平面机械臂模型的一个出发点，但本例的质点杆端质量假设不同于一般有质量机械臂。
\end{example}
\begin{figure}[htbp]
\centering
\begin{tikzpicture}[>=Stealth]
\coordinate (O) at (0,3.7);\coordinate (A) at (1,1.967949);\coordinate (B) at (.25,.668911);
\draw[very thick,blue!70!black] (O)--(A);\draw[very thick,red!70!black] (A)--(B);
\draw[dashed] (O)--(0,1.4);\draw[dashed] (A)--(1,.1);
\fill (O) circle (2pt) node[above]{$O$};\fill (A) circle (3pt) node[right]{$m_1$};\fill (B) circle (3pt) node[below]{$m_2$};
\draw (0,2.9) arc (-90:-60:.8);\node at (.3,2.7) {$\theta_1$};
\draw[->] (1,1.317949) arc (-90:-120:.65);\node at (.7,1.1) {$\theta_2$};
\node[right] at (.55,2.85) {$l_1$};\node[left] at (.55,1.3) {$l_2$};
\end{tikzpicture}
\caption{双摆使用两根杆的绝对角，而非两杆之间的夹角。图取 $\theta_1=30^\circ,\theta_2=-30^\circ$；杆长与坐标采用同一几何关系。}
\label{fig:lag-double-pendulum}
\end{figure}

\begin{example}[旋转圆环上的小珠：驱动约束改变平衡]
半径 $R$ 的竖直光滑圆环绕竖直直径以给定 $\Omega$ 匀速旋转。小珠质量 $m$，用从最低点量起的角 $\theta$ 描述位置：
\[
\br=(R\sin\theta\cos\Omega t,R\sin\theta\sin\Omega t,-R\cos\theta).
\]
直接微分给出 $v^2=R^2\dot\theta^2+R^2\Omega^2\sin^2\theta$，所以
\[
L=\tfrac12mR^2\dot\theta^2+\tfrac12mR^2\Omega^2\sin^2\theta+mgR\cos\theta,
\]
\[
\ddot\theta=\sin\theta(\Omega^2\cos\theta-g/R).
\]
平衡为 $\theta=0,\pi$，并在 $\Omega^2>g/R$ 时出现两支 $\theta=\pm\arccos[g/(\Omega^2R)]$。最低点线性方程为 $\ddot\eta+(g/R-\Omega^2)\eta=0$，故临界转速前稳定、之后不稳定；新分支的平方频率为 $\Omega^2\sin^2\theta_0>0$。临界点线性恢复项消失，须看高阶，不能仅凭零频率判为不稳定。

本例 $E_L=mR^2\dot\theta^2/2-mR^2\Omega^2\sin^2\theta/2-mgR\cos\theta$ 守恒；电机维持给定转速，珠的惯性机械能不一定守恒。第1章给出了有效势方法，本章从惯性系动能重建同一结果；第四章再系统讨论平衡附近的展开。
\end{example}

\begin{example}[弹簧摆]
一根自然长度为 $l$、劲度系数为 $k$ 的轻弹簧，一端固定，另一端连接质量为 $m$ 的质点。取弹簧伸长量为 $x$、摆角为 $\theta$，则摆长为 $l+x$。动能和势能分别为
\[
T=\frac12m\left[\dot x^2+(l+x)^2\dot\theta^2\right],
\]
\[
V=-mg(l+x)\cos\theta+\frac12kx^2.
\]
拉格朗日函数 $L=T-V$。对 $x$ 和 $\theta$ 分别求导，得到
\[
m\ddot x-m(l+x)\dot\theta^2-mg\cos\theta+kx=0,
\]
\[
m(l+x)^2\ddot\theta+2m(l+x)\dot x\dot\theta
+mg(l+x)\sin\theta=0.
\]
这两个方程是非线性的，且彼此耦合。平衡点为 $\theta=0,x=mg/k$；小振动需围绕这个伸长后的平衡展开，而非默认 $x=0$。第四章将系统处理。

由 $\partial L/\partial\dot\theta=m(l+x)^2\dot\theta$，其导数含 $2m(l+x)\dot x\dot\theta$。角度偏导为 $-mg(l+x)\sin\theta$。径向偏导同时包含 $m(l+x)\dot\theta^2$、弹簧项 $-kx$ 和重力项 $mg\cos\theta$。令速度和加速度为零，先取下垂角 $\theta=0$，径向平衡即化为 $kx=mg$。
\end{example}

\begin{example}[旋转杆上的滑块与弹簧]
质量为 $m_1$、长为 $l$ 的均匀杆 $OA$ 绕水平轴 $O$ 在竖直平面内转动。质量为 $m_2$ 的滑块 $B$ 套在杆上，可在杆上无摩擦滑动；一根劲度系数为 $k$、自然长度为 $r_0$ 的轻弹簧沿杆放置，一端固定在 $O$，另一端连接 $B$。取杆与竖直向下方向的夹角为 $\theta$，$r=OB$ 为滑块到 $O$ 的距离。求系统的运动方程。

\begin{center}
\begin{tikzpicture}[>=Stealth,scale=0.95]
  \coordinate (O) at (0,3.2);
  \coordinate (A) at (2.9,0.9);
  \coordinate (B) at (1.45,2.05);
  \draw[very thick,blue!60!black] (O)--(A);
  \draw[dashed,gray] (O)--(0,0.5);
  \draw[decorate,decoration={coil,aspect=0.6,segment length=5pt,amplitude=4pt},thick,green!50!black] (0.18,3.05)--(1.30,2.12);
  \fill[red!70!black] (B) circle (0.12);
  \node[above right] at (B) {$B$};
  \node[left] at (O) {$O$};
  \node[below right] at (2.7,0.85) {$A$};
  \draw[->,thick,red!70!black] (O)--(1.45,2.05) node[midway,above]{$r$};
  \draw (0.0,2.45) arc (-90:-38.416:0.75);
  \node at (0.28,2.52) {$\theta$};
\end{tikzpicture}
\end{center}

杆绕 $O$ 的转动惯量为
\[
I=\int_0^l s^2\frac{m_1}{l}\,\dd s=\frac13m_1l^2.
\]
系统的动能为
\[
T=\frac12I\dot\theta^2+\frac12m_2\left(\dot r^2+r^2\dot\theta^2\right),
\]
势能为
\[
V=-m_1g\frac l2\cos\theta-m_2gr\cos\theta+\frac12k(r-r_0)^2.
\]
拉格朗日函数 $L=T-V$。对 $\theta$ 的方程为
\[
\left(I+m_2r^2\right)\ddot\theta+2m_2r\dot r\dot\theta
+\left(\frac{m_1l}{2}+m_2r\right)g\sin\theta=0,
\]
对 $r$ 的方程为
\[
m_2\left(\ddot r-r\dot\theta^2\right)-m_2g\cos\theta
+k(r-r_0)=0.
\]
第一式是转动方程，第二式是滑块沿杆的径向方程；$2m_2r\dot r\dot\theta$ 是由转动惯量随 $r$ 改变产生的耦合项。两条式子均由惯性系动能导出，不能再额外加一次非惯性系的惯性力。

\begin{modernbox}
这个刚性伸缩机构模型把径向运动与转动耦合写成两个广义坐标方程。滑块移动改变转动惯量，角速度又影响径向加速度，可据此练习角动量与能量的交叉检查。若研究实际柔性伸缩机构，还须增加杆的弹性自由度、驱动和阻尼；这里的刚杆模型本身不预测其变形或工程性能。
\end{modernbox}
\end{example}

\begin{example}[两体问题的约化]
两个质量 $m_1,m_2$ 的质点只受彼此相互作用 $V(|\br_1-\br_2|)$。引入质心 $\bR$ 和相对位矢 $\br=\br_1-\br_2$：
\[
\bR=\frac{m_1\br_1+m_2\br_2}{m_1+m_2},
\qquad
\br=\br_1-\br_2,
\qquad
\mu=\frac{m_1m_2}{m_1+m_2}.
\]
反解出
\[
\br_1=\bR+\frac{m_2}{m_1+m_2}\br,
\qquad
\br_2=\bR-\frac{m_1}{m_1+m_2}\br.
\]
代入总动能并利用一次项相消，得到
\[
T=\frac12(m_1+m_2)\dot{\bR}^2+\frac12\mu\dot{\br}^2.
\]
于是拉格朗日函数分解为
\[
L=\frac12(m_1+m_2)\dot{\bR}^2
+\frac12\mu\dot{\br}^2-V(r).
\]
质心运动是自由的，相对运动等价于一个质量为 $\mu$ 的质点在中心势 $V(r)$ 中运动。这个分解在第一章已从牛顿方程得到；拉格朗日方法把它变成动能项的代数分解，更加直接。
\end{example}

\begin{example}[圆锥面上的圆轨道]
一个顶点在下、开口向上的光滑圆锥，对称轴竖直向上，半顶角为 $\alpha$。质量为 $m$ 的质点被约束在圆锥面上运动。取 $r$ 为质点沿母线到顶点的距离，$\varphi$ 为绕轴的方位角。求质点的运动方程，并讨论圆轨道的可能性。

\begin{center}
\begin{tikzpicture}[>=Stealth]
\draw[thick] (0,0)--(-1.732051,3);\draw[thick] (0,0)--(1.732051,3);
\draw[dashed] (0,0)--(0,3.4) node[above]{$z$};
\draw[thick,blue!60!black] (0,1.905256) ellipse (1.1 and .23);
\draw[->,thick,red!70!black] (0,0)--(1.1,1.905256) node[midway,right]{$r$};
\fill (1.1,1.905256) circle (2pt) node[right]{$m$};
\draw[->] (1.1,1.905256)--(1.1,2.4) node[right]{$\dot\varphi$};
\draw (0,.8) arc (90:60:.8);\node at (.3,1.05) {$\alpha$};
\node at (-1.1,-.35) {顶点};
\end{tikzpicture}
\end{center}

由几何关系，质点到轴的距离为 $\rho=r\sin\alpha$，相对顶点的竖直高度为 $z=r\cos\alpha$。因此
\[
T=\frac12m\left(\dot r^2+r^2\sin^2\alpha\,\dot\varphi^2\right),
\qquad
V=mgr\cos\alpha.
\]
拉格朗日函数
\[
L=\frac12m\left(\dot r^2+r^2\sin^2\alpha\,\dot\varphi^2\right)
-mgr\cos\alpha.
\]
$\varphi$ 是循环坐标，故
\[
p_\varphi=mr^2\sin^2\alpha\,\dot\varphi=\text{常量}.
\]
对 $r$ 的方程为
\[
m\ddot r-mr\sin^2\alpha\,\dot\varphi^2+mg\cos\alpha=0.
\]
若质点做半径为 $r_0$ 的圆轨道，则 $\dot r=0$，上式给出
\[
\dot\varphi^2=\frac{g\cos\alpha}{r_0\sin^2\alpha}.
\]
这里 $r_0$ 是沿母线量取的距离，实际圆半径为 $r_0\sin\alpha$。取 $0<\alpha<\pi/2$、$r_0>0$。考察同一守恒动量 $p_\varphi$ 下的径向扰动，将 $\dot\varphi=p_\varphi/(mr^2\sin^2\alpha)$ 代入已经得到的径向方程，得 $m\ddot r=-V_{\rm eff}'(r)$，其中
\[
V_{\rm eff}(r)=\frac{p_\varphi^2}{2mr^2\sin^2\alpha}+mgr\cos\alpha
\]
圆轨道满足 $p_\varphi^2/(mr_0^3\sin^2\alpha)=mg\cos\alpha$。因此 $V_{\rm eff}''(r_0)=3p_\varphi^2/(mr_0^4\sin^2\alpha)=3mg\cos\alpha/r_0$。令 $r=r_0+\eta$，保留一阶项得 $m\ddot\eta=-V_{\rm eff}''(r_0)\eta$，故 $\omega_r^2=3g\cos\alpha/r_0>0$。

\begin{modernbox}
锥面上的圆轨道是锥形离心分离器、旋转料斗、颗粒筛选和锥形离心机的简化模型。本例仅说明约束面上的轨道。该单质点方程中质量约去，不能单独推出按粒径或密度分选；真实分离还需流体阻力、接触等机制。
\end{modernbox}
\end{example}

\begin{example}[旋转系中的两中心势：有效势不能代替全部拉格朗日量]
两个主星绕共同质心以恒定角速度 $n$ 作圆周运动。取随其旋转、原点在质心的平面坐标 $(x,y)$，惯性速度的旋转基分量是 $(\dot x-ny,\dot y+nx)$。令 $V(x,y)=-GM_1m/r_1-GM_2m/r_2$，其中 $r_1,r_2$ 为到两主星的距离。由惯性系动能直接构造
\[
L=\tfrac12m(\dot x^2+\dot y^2)+mn(x\dot y-y\dot x)
+\tfrac12mn^2(x^2+y^2)-V.
\]
对 $x$，$p_x=m\dot x-mny$，$L_x=mn\dot y+mn^2x-V_x$，故
\[
\ddot x-2n\dot y=-\partial_x\Phi_{\rm eff},\qquad
\ddot y+2n\dot x=-\partial_y\Phi_{\rm eff},
\quad \Phi_{\rm eff}=V/m-\tfrac12n^2(x^2+y^2).
\]
速度一次项产生科里奥利项，不能把 $L$ 错写成仅有转动系相对动能减去有效势，否则会漏掉它。原来的动能已用惯性速度计算，也不能再额外加入惯性力重复计算。

由于这个 $L$ 不显含时间，它的能量函数为
\[
E_L=\tfrac12m(\dot x^2+\dot y^2)+m\Phi_{\rm eff},\qquad
C_J=-2E_L/m.
\]
$C_J$ 是常用的雅可比积分，不是粒子在惯性系中的普通机械能。平衡点仍由 $\nabla\Phi_{\rm eff}=0$ 给出，位置及稳定性已在第1章讨论。本例的目的，是看出速度耦合、离心势及守恒量如何同时从一个 $L$ 得到。
\end{example}

\subsection{冲量与广义冲量}

突然的一下推力可能太短，来不及明显改变位置，却足以改变动量。把作用期间的广义力积分，就可直接求速度所对应的动量跳变。

若系统在极短时间内受到冲力，位移可以忽略，但广义速度发生突变。对拉格朗日方程在 $[t_-,t_+]$ 上积分，仅忽略有界项 $\partial L/\partial q_a$ 及有限普通力的积分，保留冲力的 $Q_a$ 积分，得到
\begin{equation}
\left.\pdv{L}{\dot q_a}\right|_{t_+}
-\left.\pdv{L}{\dot q_a}\right|_{t_-}
=\widehat Q_a,
\label{eq:impulsive-lagrange}
\end{equation}
其中 $\widehat Q_a=\int_{t_-}^{t_+}Q_a\,\dd t$ 称为\textbf{广义冲量}。这里假设冲击期间位形及约束几何不变，约束本身不突变；若存在非理想冲量或接触模式切换，还须给出冲击约束和恢复定律，不能仅靠该式决定所有末速度。

\section{对称性与守恒律}

如果改变某个角度或平移整个系统后，力学规律依旧一样，方程往往能给出一个沿运动保持不变的量。

\subsection{循环坐标与广义动量}

角度本身若不影响拉格朗日量，它对应的广义动量就没有来自这一项的变化。要判断它是否守恒，还需看有无额外广义力。

式 \eqref{eq:generalized-momentum} 定义的广义动量在满足循环条件时给出守恒量 \eqref{eq:cyclic-integral}。本节进一步说明如何从对称性识别这些量。

本小节的循环动量结论假设独立完整坐标且对应的额外广义力为零；非完整约束的反力将另行讨论。循环坐标不一定“看不见”。在中心力场中，极角 $\varphi$ 是循环坐标，对应的广义动量正是角动量。在均匀重力场中，水平坐标 $x$ 是循环坐标，对应水平动量守恒。在完整、理想约束下，循环坐标的本质是系统在某个连续变换下保持不变。

\begin{example}[中心力场的角动量]
质量为 $m$ 的质点在势 $V(r)$ 中运动。在极坐标下
\[
L=\frac12m\left(\dot r^2+r^2\dot\varphi^2\right)-V(r).
\]
$\varphi$ 不出现，故
\[
p_\varphi=mr^2\dot\varphi=\text{常量}.
\]
这正是角动量守恒。若用直角坐标，$L$ 同时含有 $x,y$，两个坐标都不是循环坐标，但组合 $xp_y-yp_x$ 守恒。拉格朗日方法自动把对称性对应的守恒量识别出来。
\end{example}

\subsection{能量函数与能量收支}
旋转的导轨可以通过反力给小珠输入机械能，但选用适当坐标后，方程仍可能有另一个守恒量。先从拉格朗日量计算能量函数，再核对它和机械能是否相同。

定义
\begin{equation}\label{eq:energy-function}
E_L(q,\dot q,t)=\sum_a\dot q_ap_a-L,\qquad p_a=L_{\dot q_a}.
\end{equation}
若有未包含于 $L$ 的广义力 $Q_a^{\rm nc}$，方程为 $\dot p_a-L_{q_a}=Q_a^{\rm nc}$。对 $E_L$ 逐项求导，含 $p_a\ddot q_a$ 的两项抵消，留下
\begin{align}\label{eq:energy-function-derivative}
\dot E_L
&=\sum_a[\dot q_a\dot p_a+p_a\ddot q_a]
-\sum_a[L_{q_a}\dot q_a+L_{\dot q_a}\ddot q_a]-L_t\nonumber\\
&=\sum_a Q_a^{\rm nc}\dot q_a-L_t.
\end{align}
因此无额外广义力且 $L_t=0$ 时
\begin{equation}\label{eq:energy-integral}
E_L=\text{常量}.
\end{equation}
若参数化 $\br_i(q,t)$ 含时间，动能通常分为速度二次、一次、零次项 $T=T_2+T_1+T_0$。对位置势 $V(q,t)$，欧拉齐次函数公式给出
\[
\sum_a\dot q_aT_{\dot q_a}=2T_2+T_1,\qquad
E_L=T_2-T_0+V.
\]
它与惯性系机械能 $T+V$ 一般不同。在惯性系的定常参数化下，若 $T=T_2$ 且势能与速度无关，则 $E_L=T+V$。这是常用充分条件。拉格朗日量的速度一次项在勒让德组合中抵消，例如本章的电磁耦合项。$L_t=0$ 是当前表示中 $E_L$ 守恒的充分条件，不能反过来断言任意显含时间的 $L$ 都没有任何能量守恒量；全导数或坐标变化也会改变显式时间依赖。

\begin{example}[外部驱动旋转杆：守恒的 $E_L$ 与增长的机械能]
水平光滑直杆由电机以恒定角速度 $\Omega$ 绕固定端转动。质量 $m$ 的滑块在杆上径向运动，弹簧势为 $kr^2/2$；忽略重力对水平运动的影响，取 $r>0$，不碰端点。惯性速度由 $\br=r(\cos\Omega t,\sin\Omega t)$ 求导，给出
\[
T=\tfrac12m(\dot r^2+\Omega^2r^2),\qquad
L=\tfrac12m\dot r^2-\tfrac12(k-m\Omega^2)r^2.
\]
径向方程与能量函数是
\[
m\ddot r+(k-m\Omega^2)r=0,\qquad
E_L=\tfrac12m\dot r^2+\tfrac12(k-m\Omega^2)r^2=\text{常量}.
\]
但物理机械能为 $E_{\rm mech}=m\dot r^2/2+(k+m\Omega^2)r^2/2$。由第一章极坐标加速度，杆的切向反力为 $N_\theta=2m\Omega\dot r$，故
\[
\dot E_{\rm mech}=2m\Omega^2r\dot r
=N_\theta(\Omega r)=\Omega\frac{\dd}{\dd t}(mr^2\Omega).
\]
电机通过移动的理想约束交换能量；零虚功不等于零真实功。$\Omega\to0$ 时两个能量一致。$k=m\Omega^2$ 是径向恢复项退化的情况，不能继续套用非零振动频率。
\end{example}
\begin{figure}[htbp]
\centering
\begin{tikzpicture}[>=Stealth]
\draw[->] (-.3,0)--(4,0) node[right]{$x$};\draw[->] (0,-.3)--(0,3) node[above]{$y$};
\draw[very thick,blue!60!black] (0,0)--(3.2,1.84752);
\fill (2,1.15470) circle (3pt) node[below right]{$m$};
\draw[->,thick,red!70!black] (2,1.15470)--(1.5,2.020725) node[above]{$N_\theta$};
\draw[->,thick,green!40!black] (2,1.15470)--(2.8,1.61658) node[right]{$\dot r\be_r$};
\draw (.9,0) arc (0:30:.9);\node at (1.2,.3) {$\Omega t$};
\draw[->] (1.6,0) arc (0:55:1.6);\node at (1.55,1.4) {$\Omega$};
\end{tikzpicture}
\caption{受驱动旋转杆的惯性系示意。对 $\dot r>0,\Omega>0$，切向反力沿正切向；它垂直冻结时刻的径向虚位移，却与实际速度的切向部分做正功。}
\label{fig:lag-driven-rod}
\end{figure}

\subsection{诺特定理：把变换和边界项一同写清楚}
有的变换让拉格朗日量本身不变，有的只让作用量增加端点上的一项。两种情况都可以给出守恒量，后一种必须把端点项一起算进去。

本节限于具有独立坐标的完整系统、无额外非保守力。采用一阶点变换
\begin{equation}\label{eq:one-parameter-transform}
t_\epsilon=t+\epsilon\tau(q,t),\qquad
q_{a,\epsilon}(t_\epsilon)=q_a(t)+\epsilon\xi_a(q,t)+O(\epsilon^2).
\end{equation}
总坐标变化的系数为 $\xi_a$，而在同一时间比较两条曲线的变化为
\begin{equation}\label{eq:infinitesimal-variation}
\eta_a=\xi_a-\dot q_a\tau.
\end{equation}
点表示沿原路径的全导数。由 $\dd t_\epsilon/\dd t=1+\epsilon\dot\tau$ 和 $\dd q_{a,\epsilon}/\dd t=\dot q_a+\epsilon\dot\xi_a$，速度按商法则展开：
\[
\frac{\dd q_{a,\epsilon}}{\dd t_\epsilon}
=\frac{\dot q_a+\epsilon\dot\xi_a}{1+\epsilon\dot\tau}
=\dot q_a+\epsilon(\dot\xi_a-\dot q_a\dot\tau)+O(\epsilon^2).
\]
在同一时间比较坐标时，还须扣掉原轨道随时间平移的 $\epsilon\dot q_a\tau$，这就给出式 \eqref{eq:infinitesimal-variation}。

若对\textbf{任意邻近路径}而不只对运动方程的解，有恒等关系
\begin{equation}\label{eq:symmetry-total-derivative}
\sum_a\left[L_{q_a}\xi_a+p_a(\dot\xi_a-\dot q_a\dot\tau)\right]
+L_t\tau+L\dot\tau=\frac{\dd B(q,t)}{\dd t},
\end{equation}
则作用量只改变端点项，称为准对称性；$B=0$ 是严格不变的特例。必须比较 $L\dd t$，时间也变时不能只比较 $L$。

定义欧拉方程左侧 $\mathcal E_a=\dot p_a-L_{q_a}$。利用能量函数的恒等关系
$\dot E_L=\sum_a\dot q_a\mathcal E_a-L_t$
以及上式，直接对以下组合求导：
\begin{equation}\label{eq:noether-identity}
\frac{\dd}{\dd t}\left(\sum_ap_a\xi_a-E_L\tau-B\right)
=\sum_a\mathcal E_a(\xi_a-\dot q_a\tau).
\end{equation}
具体可写为
\[
\dot C=\sum_a(\dot p_a-L_{q_a})\xi_a-(\dot E_L+L_t)\tau
=\sum_a\mathcal E_a(\xi_a-\dot q_a\tau).
\]
右边还没有使用运动方程；沿真实解 $\mathcal E_a=0$，故
\begin{equation}\label{eq:noether-charge}
\boxed{C=\sum_ap_a\xi_a-E_L\tau-B=\text{常量}.}
\end{equation}
这就是本章采用的诺特定理。固定时间的坐标对称性只需取 $\tau=0$；一维时间平移取 $\tau=1,\xi_a=0,B=0$，得到 $C=-E_L$，守恒量整体差一个负号不影响其意义。

\begin{example}[三个对称性及一个边界项]
\textbf{平移。} 若 $q_a$ 为循环坐标，取 $\xi_a=1$、其他生成元为零、$\tau=B=0$，则 $C=p_a$。但完整系统中仍有 $\dot p_a=L_{q_a}+Q_a^{\rm nc}$，存在额外广义力时须另检验。

\textbf{转动。} 对 $L=m(\dot r^2+r^2\dot\varphi^2)/2-V(r)$，取 $\xi_\varphi=1$，得到 $C=mr^2\dot\varphi$。这与第一章零力矩证明是同一事实的两种表述，不能否定其中任何一种。

\textbf{时间平移。} 自治系统给出 $E_L$ 守恒；在定常自然系统中才进一步识别为机械能。

\textbf{伽利略匀速变换。} 一维自由粒子 $L=m\dot x^2/2$，取 $\xi=t,\tau=0$，则 $\delta L/\epsilon=m\dot x=\dd(mx)/\dd t$，所以 $B=mx$，守恒量为 $pt-mx$。直接用 $p=\text{常量}$、$x=x_0+pt/m$ 得 $pt-mx=-mx_0$。若漏掉边界项，会误将 $pt$ 当作守恒量。
\end{example}
\begin{physicsbox}
诺特定理中的对称变换不是固定端点变分原理中的任意试探变分：对称变换可以移动端点。守恒律来自“作用量的变化是边界项”与运动方程的结合。先猜出一个量沿解为常数，也不等于已经证明对应变换在任意路径上是作用量对称性。
\end{physicsbox}

\subsection{全导数等价与它的边界}
如果积分中多出某个函数的时间全导数，整个积分只增加该函数的终值与初值之差。固定端点的小改动就不受这一项影响。

设 $G(q,t)$ 足够光滑，$L'=L+\dd G/\dd t$。固定时间和位置端点时，$S'-S=[G]_{t_1}^{t_2}$ 为固定数，不改变驻值方程。其动量和能量函数变为
\[
p_a'=p_a+G_{q_a},\qquad E_L'=E_L-G_t.
\]
因此它们的坐标表达未必不变。若边界位置自由，$[G]$ 会改变自然边界条件；若取依赖 $\dot q$ 的 $G$，固定位置并不固定该端点项，且 $\dd G/\dd t$ 引入高阶导数，不能直接冒充这里的同类一阶变分问题。

\section{约束力恢复与非完整系统}
选择最少的坐标能省去反力，但绳子会不会松、接触会不会失效仍要用反力来判断。这时要把约束条件和待求反力一起带回方程。

\subsection{完整约束的乘子：求回被消去的反力}
为了找回摆绳的拉力，可以暂时用两个直角坐标描述摆球，并另加长度不变的方程。乘子帮助写出拉力的方向和大小。

若希望保留冗余坐标 $q$ 并求反力，设独立完整约束为 $f_\ell(q,t)=0$。对
\[
S_{\rm aug}=\int\left[L(q,\dot q,t)+\sum_\ell\lambda_\ell(t)f_\ell(q,t)\right]\dd t
\]
分别变分 $q$ 与乘子，得到
\begin{equation}\label{eq:lag-holonomic-multipliers}
\frac{\dd}{\dd t}L_{\dot q_a}-L_{q_a}
=\sum_\ell\lambda_\ell f_{\ell,q_a}+Q_a^{\rm nc},\qquad f_\ell=0.
\end{equation}
反力的广义分量为 $R_a=\sum\lambda_\ell f_{\ell,q_a}$。乘子的单位取决于 $f$ 的归一化：把约束乘2会把同一反力对应的乘子除2，故不能不作换算就将乘子本身称为张力。

\begin{example}[用直角坐标恢复单摆张力]
取 $y$ 向上，$L=m(\dot x^2+\dot y^2)/2-mgy$，约束 $f=(x^2+y^2-l^2)/2=0$。上式给出
\[
m\ddot x=\lambda x,\qquad m\ddot y+mg=\lambda y.
\]
对约束二次求导，$x\ddot x+y\ddot y=-(\dot x^2+\dot y^2)=-v^2$，故
\[
\lambda=\frac{-mv^2+mgy}{l^2},\qquad
\mathcal T=-\lambda l=m\left(\frac{v^2}{l}-\frac{gy}{l}\right).
\]
用 $y=-l\cos\theta$ 得 $\mathcal T=m(l\dot\theta^2+g\cos\theta)$，与第一章法向方程一致。轻绳只能拉，需 $\mathcal T\ge0$；若算出负值，圆周约束解对绳失效，但刚杆可承受轴向推力。
\end{example}
对显含时间的理想完整约束，真实速度满足 $\sum_a f_{\ell,q_a}\dot q_a=-f_{\ell,t}$，约束力功率为 $\sum_aR_a\dot q_a=-\sum_\ell\lambda_\ell f_{\ell,t}$。这再次说明零虚功不等于零真实功。

\subsection{仿射速度约束与达朗贝尔乘子方程}
某些约束直接告诉我们速度在某个方向上的分量。理想滚动的反力不会沿允许的虚位移做功，因此可用这些速度条件的系数来写反力。

完整约束已由坐标消去后，若还存在满秩的仿射速度约束
\begin{equation}\label{eq:nonholonomic-constraint}
\sum_a A_{\ell a}(q,t)\dot q_a+B_\ell(q,t)=0,
\qquad\ell=1,\ldots,k,
\end{equation}
本章采用理想达朗贝尔模型，其\emph{容许虚位移规则}是
\begin{equation}\label{eq:nonholonomic-virtual}
\sum_a A_{\ell a}\delta q_a=0.
\end{equation}
这是对约束反力方向的物理建模，不能不加论证地推广到任意非线性速度约束。对上式核空间中任意向量，达朗贝尔方程左端与之正交，因此其残量位于 $A$ 的行空间：
\begin{equation}\label{eq:lagrange-multipliers}
\boxed{\frac{\dd}{\dd t}L_{\dot q_a}-L_{q_a}
=Q_a^{\rm nc}+\sum_\ell\lambda_\ell A_{\ell a}.}
\end{equation}
连同速度约束、相容初始条件求解。$s$ 仍是位形维数，$s-k$ 是当前容许速度方向的维数，不代表总能消去成 $s-k$ 个独立位置坐标。

齐次约束 $B_\ell=0$ 时反力功率为零；仿射约束则为 $-\sum\lambda_\ell B_\ell$。即使 $L$ 自治，后者也可改变 $E_L$。循环坐标的动量满足 $\dot p_a=Q_a^{\rm nc}+\sum\lambda_\ell A_{\ell a}$，一般不守恒；若对称生成元还满足 $A\xi=0$，相应约束反力投影才消失。

\subsection{滚动圆盘：完整的求解与守恒检查}
直立圆盘可以改变朝向，但无滑动要求盘心的速度沿盘面。先沿盘面投影加速度，再看垂直盘面的反力怎样让轨迹转弯。

\begin{example}[保持竖直的理想滚动盘]
盘质量 $m$、半径 $R$，模型额外限制盘始终竖直、中心高度为 $R$，只允许平面位移 $(x,y)$、朝向角 $\theta$ 与自转角 $\varphi$。可将保持竖直的支承视为理想约束；这不是允许自由倾斜的真实圆盘的一般运动。取转动惯量 $I,J>0$ 分别对应自转与朝向转动，
\[
L=T=\tfrac12m(\dot x^2+\dot y^2)+\tfrac12I\dot\varphi^2+\tfrac12J\dot\theta^2,
\]
\[
\dot x=R\dot\varphi\cos\theta,\qquad
\dot y=R\dot\varphi\sin\theta.
\]
乘子方程逐项为
\begin{align}
m\ddot x&=\lambda_1,& m\ddot y&=\lambda_2,\\
I\ddot\varphi&=-R(\lambda_1\cos\theta+\lambda_2\sin\theta),&
J\ddot\theta&=0.
\end{align}
对速度约束求导后，沿朝向投影得 $\lambda_1\cos\theta+\lambda_2\sin\theta=mR\ddot\varphi$，于是
\[
(I+mR^2)\ddot\varphi=0,\qquad \ddot\theta=0.
\]
令 $\dot\varphi=\omega,\dot\theta=\Omega$ 为常数，则
\[
(\lambda_1,\lambda_2)=mR\omega\Omega(-\sin\theta,\cos\theta).
\]
$\Omega=0$ 时中心走直线，$\Omega\ne0$ 时走半径 $|R\omega/\Omega|$ 的圆。这说明虽然 $L$ 不含 $x,y$，$m\dot x,m\dot y$ 仍可改变；不是所有循环坐标都给独立守恒动量。

反力功率为 $\lambda_1\dot x+\lambda_2\dot y-R(\lambda_1\cos\theta+\lambda_2\sin\theta)\dot\varphi=0$，故总动能守恒。还必须确认静摩擦等真实支承能提供所需反力；本题未给摩擦系数，不能无条件断言任意 $\omega,\Omega$ 均可实现。
\end{example}
\begin{figure}[htbp]
\centering
\begin{tikzpicture}[>=Stealth]
\draw[->] (-.2,0)--(4.3,0) node[right]{$x$};\draw[->] (0,-.3)--(0,3.2) node[above]{$y$};
\coordinate (C) at (1.6,1.1);
\draw[very thick,blue!70!black] (.8,.63812)--(2.4,1.56188);
\fill (C) circle (2pt) node[below]{$C$};
\draw[dashed] (C)--(3,1.1);
\draw[->,thick,green!40!black] (C)--(3,1.90829) node[right]{$\bv_C=R\omega\be$};
\draw[->,thick,red!70!black] (C)--(1,2.13923) node[above]{$\boldsymbol\lambda$};
\draw (2.25,1.1) arc (0:30:.65);\node at (2.5,1.3) {$\theta$};
\end{tikzpicture}
\caption{竖直滚动盘的俯视运动学：粗线为盘平面与地面的交线方向，$\be=(\cos\theta,\sin\theta)$ 为朝向单位矢量；当 $\omega\Omega>0$ 时，合水平反力垂直朝向向左。此图不把俯视平面误画成圆盘的侧面。}
\label{fig:rolling-disk}
\end{figure}

\subsection{为什么不能直接把速度约束加进作用量？}
限制瞬时可移动的方向，与在整条试探路径上强制满足速度条件，会改变可作的变分。把两种模型各写一遍，就能比较它们给出的方程。

若把上述约束作为 $L_{\rm aug}=L+\mu_1(\dot x-R\dot\varphi\cos\theta)+\mu_2(\dot y-R\dot\varphi\sin\theta)$，再对所有路径作普通固定端点变分，$\theta$ 方程变为
\[
J\ddot\theta=R\dot\varphi(\mu_1\sin\theta-\mu_2\cos\theta),
\]
一般不同于达朗贝尔模型的 $J\ddot\theta=0$。这种对受限路径作用量求驻值的方法称为变分约束方法（vakonomic）；它可对应某些最优控制问题，但不能自动代替真实理想滚动的动力学。

差别来自允许的变分：若一族路径都满足 $A(q,t)\dot q+B(q,t)=0$，其线性化包含 $A\dot\eta$、$A_{,q}\eta\dot q$ 和 $B_{,q}\eta$，并不等价于每个时刻仅要求 $A\eta=0$。而达朗贝尔虚位移正是后一类力学投影。该区分是处理非完整系统的核心，不是乘子符号的选择问题。

\begin{modernbox}[title={进一步探索：不可积性与工程多体方法}]
滚动盘的容许速度由 $X_1=\partial_\theta$ 与 $X_2=R\cos\theta\partial_x+R\sin\theta\partial_y+\partial_\varphi$ 张成。先改变朝向再滚动与反序操作并不等效；其交换子为 $[X_1,X_2]=-R\sin\theta\partial_x+R\cos\theta\partial_y$，不在两者张成的空间内。由弗罗贝尼乌斯判据，该速度分布不能积分为二维位形曲面。

阿佩尔和凯恩方法也沿独立运动方向投影动力学。对独立完整坐标，$\mathcal S=\sum_i m_i|\ba_i|^2/2$ 给出 $\partial\mathcal S/\partial\ddot q_a=\sum_i m_i\ba_i\cdot\partial\br_i/\partial q_a=Q_a$。非完整模型须先选相容独立速度。
\end{modernbox}

\section{带电粒子、广义动量与规范等价}
同一电磁场可以由不同的势来描述。先用势求粒子怎样受力，再看换一种势时哪些量会改变、哪些仍有相同的物理含义。

\subsection{从势耦合逐分量导出洛伦兹力}
磁力侧向推转粒子，因此速度的方向可能改变。拉格朗日量通过一个与速度有关的项，产生这种仅用位置势能无法描述的力。

设非相对论粒子质量 $m$ 恒定、电荷为 $e$（可正可负），外场给定，不计辐射反作用。定义
\[
\bE=-\nabla\phi-\partial_t\bA,\qquad \bB=\nabla\times\bA,
\]
\begin{equation}\label{eq:em-lagrangian}
L=\tfrac12m\bv^2-e\phi+e\bA\cdot\bv.
\end{equation}
广义动量与机械动量不同：
\begin{equation}\label{eq:canonical-momentum-em}
\bp=\frac{\partial L}{\partial\bv}=m\bv+e\bA.
\end{equation}
对第 $i$ 个分量，取坐标偏导时保持速度独立，
\[
\frac{\dd}{\dd t}(mv_i+eA_i)=-e\partial_i\phi+e\sum_j v_j\partial_i A_j.
\]
用全导数 $\dd A_i/\dd t=\partial_t A_i+\sum_jv_j\partial_jA_i$ 展开，得
\[
m\dot v_i=e(-\partial_i\phi-\partial_tA_i)
+e\sum_jv_j(\partial_iA_j-\partial_jA_i).
\]
最后一项正是 $e(\bv\times\bB)_i$，于是
\begin{equation}\label{eq:lag-lorentz-force}
m\dot\bv=e(\bE+\bv\times\bB).
\end{equation}
对速度点乘，磁力不做功，$\dot T=e\bE\cdot\bv$。能量函数为 $E_L=T+e\phi$；若所选势不显含时间，它守恒。外场若随时间变化则可能向粒子传能，完整“粒子加场及场源”的能量平衡不应只对测试粒子闭合计算。

\subsection{规范变换为什么不改变轨道？}
两个势若给出同一组电场和磁场，就可描述同一粒子轨道。计算时检查拉格朗日量的差是否为一个时间全导数。

取单值光滑函数 $\chi(\br,t)$，势变为
\begin{equation}\label{eq:gauge-transform}
\bA'=\bA+\nabla\chi,\qquad \phi'=\phi-\partial_t\chi.
\end{equation}
梯度的旋度为零、混合偏导相消，故场强不变，而
\[
L'=L+e\frac{\dd\chi}{\dd t},\qquad
\bp'=\bp+e\nabla\chi,\qquad E_L'=E_L-e\partial_t\chi.
\]
这正是 $G=e\chi$ 的全导数等价。固定时空端点时轨道不变，动量和能量函数的表示可以改变。不可把“改变规范”说成一种产生新轨道或新力的物理操作；它是同一外场的描述冗余。

\begin{example}[均匀磁场的同一平移守恒量，两种规范]
取 $\bB=B\kk$、$\phi=0$。在 $\bA=(-By,0,0)$ 下，
\[
L=\tfrac12m(\dot x^2+\dot y^2+\dot z^2)-eBy\dot x,
\]
$x$ 为循环坐标，所以 $C=m\dot x-eBy$ 守恒。作 $\chi=Bxy$，新势为 $\bA'=(0,Bx,0)$，新 $L'$ 显含 $x$，但 $x\to x+\epsilon$ 给出
\[
\delta L'/\epsilon=eB\dot y=\frac{\dd(eBy)}{\dd t}.
\]
因此诺特边界项为 $B_{\rm N}=eBy$，守恒量为 $p'_x-B_{\rm N}=m\dot x-eBy$，与先前相同。机械动量 $m\dot x$ 一般不守恒；广义动量是否直接显为守恒，依赖所用规范和边界项。洛伦兹方程也直接给出 $m\ddot x-eB\dot y=0$，形成独立核验。
\end{example}
\begin{modernbox}[title={进一步探索：从势到场，而非把所有场都叫规范场}]
场是在时空各点定义的量。电磁势具有规范冗余，场强由其导数给出；温度场等一般标量场不能因此称为规范联络。局域规范冗余与普通机械坐标的整体诺特对称性不同。进一步讨论需电磁场本身的作用量与边界条件，超出本章。
\end{modernbox}

\section{通向小振动与哈密顿力学}
第三章的方程有两个常用的下一步。研究轻微偏离平衡时可保留最低阶的项；研究完整状态时，可用广义动量替换速度。

对定常自然系统，$L=\dot q^{\mathsf T}M(q)\dot q/2-V(q)$。在平衡 $q_0$ 附近令 $q=q_0+\eta$，$\nabla V(q_0)=0$，保留二阶拉格朗日量为
\[
L^{(2)}=\tfrac12\dot\eta^{\mathsf T}M(q_0)\dot\eta
-\tfrac12\eta^{\mathsf T}K(q_0)\eta,
\qquad K_{ab}=\left.\frac{\partial^2V}{\partial q_a\partial q_b}\right|_{q_0}.
\]
它直接给出 $M\ddot\eta+K\eta=0$。第\ref{ch:oscillation}章将以等质量等长双摆（式 \eqref{eq:osc-double-pend}）说明稳定性、简正模和质量加权正交性，而非只把复杂摆都当作单个简谐振子。

几何上，广义坐标是位形空间的局部参数，$(q,\dot q)$ 包含位置与速度信息。第六章在速度海森矩阵可逆等条件下用勒让德变换把速度换成 $p=L_{\dot q}$；这里的 $E_L$ 才转成相空间上的哈密顿函数。不要提前把所有约束模型、奇异速度势或非完整投影都直接当作普通可逆的哈密顿系统。

\section{计算与模拟：检查方程，而非只看动画}
数值求解后，先用方程预测该有怎样的能量变化，再和实际输出比较。不同模型允许的能量收支不同，不能对每一题都要求机械能恒定。

对无驱动双摆，先从式 \eqref{eq:lag-double-pendulum} 求解 $2\times2$ 线性系统得到 $\ddot\theta$，再将状态写成 $(\theta_1,\theta_2,\dot\theta_1,\dot\theta_2)$ 积分。用同一组初值和物理参数，比较步长 $h$ 与 $h/2$；记录能量 $E=T+V$ 的误差，并确认质量矩阵在沿途可逆。四阶Runge--Kutta方法的短时收敛可作检验，但它不是严格能量保持算法，不能由一段看似合理的动画推断长期保真。

对阻尼摆应检查 $E(t)-E(0)=\int_0^t[-b\dot\theta^2+\tau\dot\theta]\dd t$，不能仍要求 $E$ 恒定；对受驱动旋转杆则分别记录 $E_L$ 与 $E_{\rm mech}$。随附数值脚本为
\begin{center}\texttt{review/verify\_ch03.py}\end{center}
它对本章解析模型作独立检查，输出误差和收敛记录。所有模拟参数为教学设定，不代表实际机器人或航天器。

\section*{关键概念与常见误区}
\begin{itemize}[leftmargin=1.8em]
\item 选坐标时，可试着逐个改变它们，看看能否独立变化。用坐标数减去约束数前，先检查这些约束是否独立。
\item 算虚位移时先固定时刻。若导轨自己在移动，它的反力可能对真实运动做功，即使对允许的虚位移不做功。
\item 用惯性系速度写过动能后，坐标变换造成的效应已包含在方程里。再加一遍惯性力会重复计算。
\item 重力和弹簧常可用 $L=T-V$。带电粒子还可能有随速度变化的耦合项，此时广义动量要从 $L$ 重新求。
\item 对转动导轨一类题，分别写出能量函数、物体机械能和电机功率，再用方程检查各自怎样变。
\item 用诺特定理时，要检查整个作用量的变化，包含时间积分的测度与可能出现的端点项。
\item 滚动约束能让循环坐标的动量受到反力改变。先确定采用的是达朗贝尔投影还是受限路径变分，再写方程。
\end{itemize}

\section*{本章小结}
\addcontentsline{toc}{section}{本章小结}
本章从约束几何选坐标，以达朗贝尔原理投影牛顿方程，得到包含非保守广义力的拉格朗日方程；在合适的完整系统中，第二章的固定端点变分给出等价的作用量表述。乘子法用于恢复反力，不能把它与独立坐标消元混淆。对称性及其边界项组织守恒律，能量收支则检验受驱动和受阻尼模型。第四章将在这一框架内系统展开稳定平衡附近的小振动。

\section*{分层习题}
\addcontentsline{toc}{section}{分层习题}
先数独立坐标，并把各质点的位置写成这些坐标的函数。求出运动后，若还需张力或支持力，再把相应的约束方程带回来。

\begin{bookexercises}
\item \exlevel{基础} \textbf{约束与秩。} 三维双摆的两个杆长约束在杆长非零时有多少独立约束和位形自由度？若另限制两球在同一固定竖直平面中，如何恢复平面双摆的两个自由度？
\item \exlevel{基础} \textbf{虚位移。} 质点在 $x^2+y^2=R(t)^2$ 的受驱动圆周上运动。写出虚位移和真实速度必须满足的两个关系，并说明为何径向反力可零虚功而非零功率。
\item \exlevel{基础} \textbf{推导训练。} 对 $\br_i(q,t)$ 逐项验证式 \eqref{eq:kinematic-identities}，再从动能导出拉格朗日惯性项，不把“偏导和全导总可交换”作为理由。
\item \exlevel{综合} \textbf{平衡分支。} 两杆例题中取 $P_1=P_2=P,F=9P/8$。写出稳定支的两组正弦余弦，并分析只把一根杆反向会得到极小、极大还是鞍点。
\item \exlevel{综合} \textbf{阿特伍德机。} 在正文忽略滑轮惯量的假设下求绳张力；再令滑轮有给定转动惯量 $I$、半径 $a$，且无滑动，求加速度并说明两侧张力是否仍相等。
\item \exlevel{综合} \textbf{双摆。} 独立导出双摆第二方程，核对 $-D\sin\Delta\dot\theta_1^2$ 的符号；在最低平衡附近写出 $M_0,K_0$，而不直接把两角看成独立单摆。
\item \exlevel{综合} \textbf{弹簧摆。} 在 $x_0=mg/k,\theta_0=0$ 附近线性化，求径向和角向频率，说明为什么应使用伸长后的摆长。
\item \exlevel{综合} \textbf{转动杆滑块。} 对正文带有质量杆和弹簧的两坐标系统，用角动量定理核验 $2m_2r\dot r\dot\theta$ 项；说明令 $\theta=\Omega t$ 给定时为什么必须另有外部驱动力矩。
\item \exlevel{综合} \textbf{耗散。} 对无驱动阻尼摆，证明 $E$ 单调不增；若希望保持非零匀速 $\dot\theta=\omega_c$，写出所需力矩随角度的函数，不可只令平均驱动补偿阻尼。
\item \exlevel{综合} \textbf{驱动能量。} 对恒速旋转水平杆上的弹簧滑块，验证 $\dot E_L=0$、$\dot E_{\rm mech}=\Omega\dot J_z$；讨论 $k-m\Omega^2$ 为正、零、负时的径向行为及有限杆端点的限制。
\item \exlevel{综合} \textbf{乘子归一化。} 用 $f=x^2+y^2-l^2$ 而非 $f/2$ 重算单摆反力，比较乘子数值和物理张力；判断什么时候轻绳约束会失效。
\item \exlevel{综合} \textbf{诺特边界项。} 对自由粒子匀速变换 $\xi=t$，证明守恒的是 $pt-mx$，并解释为何 $pt$ 本身不守恒。再取时间平移 $\tau=1,\xi=0$ 核对诺特荷的符号。
\item \exlevel{综合} \textbf{全导数。} 取 $L'=L+\dd G(q,t)/\dd t$，独立算出 $p'$、$E_L'$。若 $G=atq$（$a$ 为常数且单位使 $G$ 为作用量），显式时间依赖是否必然意味着原物理系统没有能量守恒量？
\item \exlevel{综合} \textbf{非完整滚动。} 用正文约束与乘子方程推出 $\ddot\varphi=\ddot\theta=0$，给出 $\Omega\ne0$ 时圆心轨迹，并核验反力功率；说明圆盘 $x$ 坐标循环却动量不守恒的原因。
\item \exlevel{拓展} \textbf{两种受限变分。} 将正文滚动盘的速度约束直接加入作用量，计算增广模型的 $\theta$ 方程，再与达朗贝尔模型比较。仅仅给乘子改名能消掉差异吗？
\item \exlevel{综合} \textbf{规范与守恒量。} 在均匀磁场两种规范下，分别由诺特边界项和洛伦兹方程验证同一个 $m\dot x-eBy$ 守恒；机械动量和广义动量各有什么不同？
\item \exlevel{综合} \textbf{小振动与旋转系。} (a) 对旋转圆环的小珠推导非竖直平衡支的小振动频率。(b) 对两中心旋转系求 $E_L$，说明为何仅用有效势构造 $L$ 会漏掉科里奥利项。
\item \exlevel{拓展} \textbf{计算检查。} 运行双摆两步长模拟，比较相同短时区间的能量误差；对受驱动杆同时输出两种能量。说明本章各模型应检验的量为何不同，不把不守恒的物理量误当成算法误差。
\end{bookexercises}

\section*{习题解答}
先用同一题号找到解答。核对时，从第一条方程开始比较自己的做法，再检查最后的符号、单位和适用条件。
\begin{bookexercises}
\item 六个位置坐标、两个独立杆长约束，位形维数为4。再各加一个离开指定平面的坐标约束，剩2；若约束奇异或重复，不能简单照数相减。

约束可取 $|\br_1|^2-l_1^2=0$、$|\br_2-\br_1|^2-l_2^2=0$。在关于 $\br_2$ 的三列中，第一行全零，第二行有非零向量 $2(\br_2-\br_1)$，而第一行关于 $\br_1$ 的梯度也非零，故秩为2。选平面为 $xz$ 平面，另两条 $y_i=0$ 的梯度只在法向列非零，与杆长约束独立。
\item 冻结时间给 $\br\cdot\delta\br=0$，真实速度给 $\br\cdot\bv=R\dot R$。若反力为 $N\br/R$，功率为 $N\dot R$，虚功为零。

对 $f=\br^2-R(t)^2$，固定 $t$ 变分得 $\delta f=2\br\cdot\delta\br=0$。沿实际路径则有 $\dot f=2\br\cdot\bv-2R\dot R=0$。在 $R>0$ 时，第二式给出真实速度的径向分量 $\dot R$。
\item 速度是 $\bv_i=\sum_b\br_{i,q_b}\dot q_b+\br_{i,t}$。第一式从 $\dot q_a$ 的系数直接读出；第二式两边分别为
\[
\sum_b\br_{i,q_bq_a}\dot q_b+\br_{i,tq_a},\qquad
\sum_b\br_{i,q_aq_b}\dot q_b+\br_{i,q_at}.
\]
在 $\br_i$ 二阶偏导连续时二者相等。由 $\partial T/\partial\dot q_a=\sum_i m_i\bv_i\cdot\br_{i,q_a}$ 得
\[
\frac{\dd}{\dd t}\frac{\partial T}{\partial\dot q_a}
=\sum_i m_i\ba_i\cdot\br_{i,q_a}
 +\sum_i m_i\bv_i\cdot\bv_{i,q_a}.
\]
第二项恰是 $\partial T/\partial q_a$，相减即得到惯性投影。
\item 第一杆 $\cos\alpha=-3/5,\sin\alpha=-4/5$；第二杆 $\cos\beta=-9/\sqrt{97},\sin\beta=-4/\sqrt{97}$。只反向一根杆使势能海森矩阵一正一负，是鞍点。

两个角各自的势能为 $l(A_j\sin\psi_j+F\cos\psi_j)$，其中 $A_1=3P/2,A_2=P/2$。驻点单位向量 $(\cos\psi_j,\sin\psi_j)$ 必与 $(F,A_j)$ 共线。取反向使势能极小，曲率为 $l\sqrt{F^2+A_j^2}$。将任一杆反向只翻转对应曲率的符号，混合导数仍为零。
\item 无惯量时 $\mathcal T=2m_1m_2g/(m_1+m_2)$。有惯量时有效质量增为 $m_1+m_2+I/a^2$，加速度为 $(m_1-m_2)g/(m_1+m_2+I/a^2)$；两侧张力差提供转矩，$\mathcal T_1-\mathcal T_2=I\ddot q/a^2$。

取 $m_1$ 向下、$m_2$ 向上为正。两质点方程为 $m_1g-\mathcal T_1=m_1\ddot q$、$\mathcal T_2-m_2g=m_2\ddot q$。滑轮角加速度是 $\ddot q/a$，转矩平衡给 $(\mathcal T_1-\mathcal T_2)a=I\ddot q/a$。将三式相加消元，再令 $I=0$ 便可求共同张力。
\item 对 $\cos\Delta\dot\theta_1$ 求导后与坐标偏导相减，混合速度项抵消。最低点处
\[
M_0=\begin{pmatrix}(m_1+m_2)l_1^2&m_2l_1l_2\\m_2l_1l_2&m_2l_2^2\end{pmatrix},\quad
K_0=\begin{pmatrix}(m_1+m_2)gl_1&0\\0&m_2gl_2\end{pmatrix}.
\]

令 $\Delta=\theta_1-\theta_2$。对 $\partial L/\partial\dot\theta_2=m_2l_2^2\dot\theta_2+D\cos\Delta\,\dot\theta_1$ 求导，再减去 $\partial L/\partial\theta_2=D\sin\Delta\,\dot\theta_1\dot\theta_2-m_2gl_2\sin\theta_2$，得
\[
m_2l_2^2\ddot\theta_2+D\cos\Delta\,\ddot\theta_1
-D\sin\Delta\,\dot\theta_1^2+m_2gl_2\sin\theta_2=0.
\]
在最低点取 $\cos\Delta=1$、$\sin\theta_j\simeq\theta_j$，速度二次项高于线性阶。
\item 径向频率 $\omega_r=\sqrt{k/m}$，角向频率 $\omega_\theta=\sqrt{g/(l+mg/k)}$。径向和角向在这一平衡的线性阶解耦，有限振幅仍耦合。

取 $r=l+x$，则 $V=k(r-l)^2/2-mgr\cos\theta$。平衡半径 $r_*=l+mg/k$ 由 $k(r_*-l)=mg$ 给出。令 $\rho=r-r_*$，二次动能为 $m(\dot\rho^2+r_*^2\dot\theta^2)/2$，二次势能为 $(k\rho^2+mgr_*\theta^2)/2$，于是 $m\ddot\rho+k\rho=0$、$mr_*^2\ddot\theta+mgr_*\theta=0$。
\item 关于固定点的角动量为 $(I+m_2r^2)\dot\theta$，求导产生该耦合项；重力矩为
\[-(m_1l/2+m_2r)g\sin\theta.\]
给定 $\dot\theta=\Omega$ 时，驱动力矩应补足耦合项与重力矩，不是把自由转动方程直接置零。

具体地，$\dot J=(I+m_2r^2)\ddot\theta+2m_2r\dot r\dot\theta$。弹簧力沿杆，对铰点的力矩为零。因此维持恒速所需电机力矩为 $2m_2r\dot r\Omega+(m_1l/2+m_2r)g\sin\theta$。
\item $\dot E=-b\dot\theta^2\le0$。精确匀速要求 $\tau(\theta)=b\omega_c+mgl\sin\theta$，不仅是常量 $b\omega_c$。

由 $ml^2\ddot\theta=-mgl\sin\theta-b\dot\theta+\tau$，乘以 $\dot\theta$ 得 $\dd[ml^2\dot\theta^2/2+mgl(1-\cos\theta)]/\dd t=-b\dot\theta^2+\tau\dot\theta$。求匀速驱动力矩时应在原方程置 $\ddot\theta=0$。
\item 用径向方程代入能量导数；$J_z=mr^2\Omega$。有效刚度为正时数学解振荡，为零时径向匀速，为负时有指数模。仅在不穿过原点、不碰杆端、驱动可提供所需力矩的区段使用这些解。

记 $k_*=k-m\Omega^2$，径向方程为 $m\ddot r=-k_*r$。故 $\dot E_L=(m\ddot r+k_*r)\dot r=0$。又 $E_{\rm mech}=E_L+m\Omega^2r^2=E_L+\Omega J_z$，恒定 $\Omega$ 下直接求导即给出电机功率。$k_*$ 的三种符号分别对应特征根 $\pm i\sqrt{k_*/m}$、重根0和 $\pm\sqrt{-k_*/m}$。
\item 新乘子是正文乘子的一半，反力仍为原来的 $\lambda\br$，张力不变。轻绳要求张力非负；它失去拉紧时必须更换运动模型。

用新约束 $f$ 时，方程为 $m\ba=-mg\hat{\boldsymbol y}+2\lambda_{\rm new}\br$。由 $\br\cdot\ba=-v^2$ 得 $2\lambda_{\rm new}l^2=-mv^2+mgy$，因此 $\mathcal T=-2l\lambda_{\rm new}=m(l\dot\theta^2+g\cos\theta)$。
\item $B=mx$，$C=pt-mx=-mx_0$；直接求导 $\dd(pt)/\dd t=p$ 一般不为零。时间平移荷为 $-E_L$，其相反数同样守恒。

自由粒子 $L=m\dot x^2/2$ 在 $\delta x=\epsilon t$ 下有 $\delta L/\epsilon=m\dot x$，正好为 $\dd(mx)/\dd t$。时间平移取 $B=0$，代入 $C=p\xi-E_L\tau-B$ 得 $-E_L$。
\item $p'=p+G_q,E_L'=E_L-G_t$。$G=atq$ 时 $E_L'=E_L-aq$，所以原守恒量在新表示为 $E_L'+aq$；不能只看 $L'$ 是否显含时间作物理判断。

展开 $L'=L+G_q\dot q+G_t$。速度偏导增加 $G_q$，代入 $p'\dot q-L'$ 后两个 $G_q\dot q$ 项相消。多坐标情形用 $\sum_a G_{q_a}\dot q_a$ 作同样计算。
\item 令 $\theta=\theta_0+\Omega t$、$\dot\varphi=\omega$，则
\[
x=x_0+\frac{R\omega}{\Omega}(\sin\theta-\sin\theta_0),\quad
y=y_0-\frac{R\omega}{\Omega}(\cos\theta-\cos\theta_0).
\]
反力与容许速度配对总功率为零，但 $x,y$ 方向单独不在约束核内，故相应动量可变。

先对两个速度约束求导，再乘 $\cos\theta,\sin\theta$ 相加，得 $\ddot x\cos\theta+\ddot y\sin\theta=R\ddot\varphi$。乘子方程给 $I\ddot\varphi=-mR^2\ddot\varphi$、$J\ddot\theta=0$，从而两个角速度均为常量。反力功率为
\[
\lambda_1(\dot x-R\dot\varphi\cos\theta)
+\lambda_2(\dot y-R\dot\varphi\sin\theta)=0.
\]
\item 变分约束给 $J\ddot\theta=R\dot\varphi(\mu_1\sin\theta-\mu_2\cos\theta)$，物理投影给零右端。差异来自变分集合，不是乘子记号。

增广项中含 $\theta$ 的部分可合写成
\[
-R\dot\varphi(\mu_1\cos\theta+\mu_2\sin\theta).
\]
对它求角度偏导即得上式右端，而 $\partial L_{\rm aug}/\partial\dot\theta=J\dot\theta$。达朗贝尔约束的虚位移系数在 $\delta\theta$ 列为零，因而没有这个广义反力矩。
\item 旧规范直接循环，新规范 $x$ 平移只差 $\dd(eBy)/\dd t$，均给同一个荷；洛伦兹式 $m\ddot x=eB\dot y$ 直接核验。$m\bv$ 规范不变，$\bp=m\bv+e\bA$ 随规范变化。

对 $\bA=(-By,0,0)$，有 $p_x=m\dot x-eBy$；对 $\bA'=(0,Bx,0)$，有 $p_x'=m\dot x$。后一规范平移时的边界项为 $B_{\rm N}=eBy$，由 $p_x'-B_{\rm N}$ 恢复同一荷。从 $m\ddot x=eB\dot y$ 也直接得到 $\dd(m\dot x-eBy)/\dd t=0$。
\item (a) 对右端 $\sin\theta(\Omega^2\cos\theta-g/R)$ 求导，在新平衡点得 $-\Omega^2\sin^2\theta_0$，所以 $\omega_{\rm small}=|\Omega\sin\theta_0|$。(b) $E_L=m(\dot x^2+\dot y^2)/2+m\Phi_{\rm eff}$，速度一次项不出现在该积分中，但其变分产生科里奥利项，不能从 $L$ 中删去。

(b) 正文的速度一次项为 $mn(x\dot y-y\dot x)$。它对速度偏导后再与坐标偏导相减，在 $x,y$ 的方程左端分别贡献 $-2mn\dot y$ 与 $2mn\dot x$。求能量函数时，同一项又从 $\dot xL_{\dot x}+\dot yL_{\dot y}$ 和 $L$ 中等量抵消。因此能量式里不出现它，运动方程中仍须保留它。
\item 双摆用相同初值和终止时刻，按 $\max_t|E(t)-E(0)|/|E(0)|$ 统计，$h=0.01$ 时为 $1.54\times10^{-11}$，$h=0.005$ 时为 $1.04\times10^{-12}$，比约 $14.84$。受驱动杆每步可同时计算 $E_L$ 和 $E_{\rm mech}=E_L+m\Omega^2r^2$；后者的差值应与 $\Omega[J_z(t)-J_z(0)]$ 一致。伴随脚本的瞬时功率测试给出 $3.36$（所用单位制），其计算式为 $2m\Omega^2r\dot r$。阻尼摆则比较能量变化与 $\int(-b\dot\theta^2+\tau\dot\theta)\dd t$，这些检查分别对应各模型的能量收支。
\end{bookexercises}

\section*{进一步阅读与来源}
链接复核：2026-11-08。算例按本书坐标与符号重建；无具体任务数据的模型不视为真实装置预测。
\begin{bookreading}
\item David Tong，\textit{The Lagrangian Formalism}：\href{https://davidtong.org/pdfs/teaching/classical-dynamics/clas2.pdf}{剑桥公开讲义}，进一步学习广义坐标、对称性、双摆和电磁耦合。
\item MIT 8.09，\href{https://ocw.mit.edu/courses/8-09-classical-mechanics-iii-fall-2014/pages/syllabus/}{Classical Mechanics III官方大纲}：完整与非完整约束、达朗贝尔与哈密顿原理、守恒律是相关本科进阶教学主题。本书以本章明确的模型假设为准，不从课程简介推断定理细节。
\item A. D. Lewis、R. M. Murray，\textit{Variational Principles for Constrained Systems: Theory and Experiment}：\href{https://authors.library.caltech.edu/records/2pdbe-k1104}{Caltech原始研究记录及全文入口}，用于继续区分非完整力学与变分约束理论。所需基础是本章约束投影、变分法及常微分方程。
\end{bookreading}

\chapter{小振动}\label{ch:oscillation}

\begin{guidebox}
轻推一只摆，它会在原来的平衡位置附近摆动。若把两只摆连起来，有些动作能保持固定比例，有些动作则会出现时强时弱的拍。先理解单摆如何回来，再找两只摆能一起完成的简单动作。

本章先找平衡，再只保留小位移下最先起作用的恢复力和惯性。主要模型采用定常完整约束，动能正定，恢复力由位置势能给出。阻尼、外部驱动、磁力或随时间改变的参数都可能改变结论，后面会按模型分别处理。
\end{guidebox}

\section{平衡、稳定与近似的边界}
把静止的摆轻轻推开，再看它是摆回来、越离越远，还是缓慢漂开。先分辨这三种结果，才知道何时可以谈小振动频率。

\subsection{平衡点与势能曲率}
找平衡像找势能图上切线水平的地方。接着看曲线是向上弯还是向下弯，便能读出微小位移后恢复力的方向。

设 $q_0$ 是坐标图内部的静止平衡，$M(q)$ 在邻域内光滑、对称且正定。因 $\dot q=\ddot q=0$，第三章的方程给出
\begin{equation}
V_{q_a}(q_0)=0,\qquad a=1,\ldots,s.
\label{eq:equilibrium}
\end{equation}
这个条件适用于当前自然系统；含外驱、接触边界或速度势时应重写平衡条件。令 $\eta=q-q_0$，泰勒展开为
\begin{equation}
V(q_0+\eta)=V(q_0)+\frac12\eta^{\mathsf T}K\eta+O(\|\eta\|^3),
\quad K_{ab}=V_{q_aq_b}(q_0)=K_{ba}.
\label{eq:potential-quadratic}
\end{equation}
一次项因平衡消失，常数项不影响方程。$K$ 是\textbf{刚度矩阵}；广义坐标可以是长度或角度，其矩阵元不必全有弹簧劲度的单位。对频率有意义的是 $K$ 与质量矩阵的配对，而非孤立地读取 $K$ 的本征值。

\begin{definition}[相空间中的稳定性]
平衡状态是 $(q_0,0)$。若对它的每一个给定小邻域，都存在更小的初始位置和速度邻域，使其中所有初态的未来运动始终留在给定邻域，则称平衡为李雅普诺夫稳定。若不满足此条件，则不稳定；不稳定不必表现为指数增长。若邻近状态还随时间趋近 $(q_0,0)$，则称渐近稳定。
\end{definition}
位置与速度的距离须先用固定尺度无量纲化。只扰动位置而不允许初速度变化，不能检验完整的稳定性。

\subsection{能量判据与零曲率}
势能图的谷底能限制小能量运动的范围。如果二次曲率恰好是零，还得继续看更高次项，或直接求运动。

若 $K$ 正定，由泰勒展开和 $M(q)$ 的连续正定性，在足够小邻域中存在正常数，使
\[
E-V(q_0)=\frac12\dot\eta^{\mathsf T}M(q_0+\eta)\dot\eta
+V(q_0+\eta)-V(q_0)
\ge c_1\|\dot\eta\|^2+c_2\|\eta\|^2.
\]
在邻域边界上能量有正下界，选择初始能量小于该下界，能量守恒就阻止轨道穿出边界。这给出稳定性证明，也说明保守振子一般不渐近稳定：非零的微小能量不会自动耗散。

后文将证明 $M_0^{-1/2}KM_0^{-1/2}$ 为实对称矩阵，其中 $M_0=M(q_0)$。若它有负本征值 $\lambda=-\sigma^2$，相应线性方程含 $e^{\sigma t}$，平衡不稳定。若有零本征值，仅凭二次势能无法作出非线性稳定结论；零模的速度还可能引起长期漂移。

\begin{figure}[htbp]
\centering
\begin{tikzpicture}
\begin{axis}[width=0.83\textwidth,height=5.1cm,axis lines=middle,xlabel={$q/q_*$},ylabel={$V/V_*$},xmin=-1.15,xmax=1.15,ymin=-1.3,ymax=1.5,legend style={at={(1.02,1)},anchor=north west,font=\small,draw=none},xtick={-1,0,1},ytick={-1,1}]
\addplot[blue!65!black,very thick,domain=-1.05:1.05,samples=100]{x^2};\addlegendentry{$q^2$：正曲率}
\addplot[green!45!black,thick,dashed,domain=-1.05:1.05,samples=100]{x^4};\addlegendentry{$q^4$：零曲率}
\addplot[red!70!black,thick,dashdotted,domain=-1.05:1.05,samples=100]{-x^4};\addlegendentry{$-q^4$：零曲率}
\end{axis}
\end{tikzpicture}
\caption{以固定尺度无量纲化后的三个势能，均在原点驻定。$q^4$ 和 $-q^4$ 的二阶导数都为零，稳定性却不同；二阶判据没有覆盖所有平衡。}
\label{fig:potential-stability}
\end{figure}

\begin{example}[单摆：稳定点与倒立点]
取从竖直向下量起的摆角，$T=ml^2\dot\theta^2/2$，$V=-mgl\cos\theta$。平衡为 $\theta=0,\pi$（模 $2\pi$）。在最低点
\[
V''(0)=mgl>0,\qquad \omega_0^2=\frac{mgl}{ml^2}=\frac gl.
\]
在倒立点令 $\theta=\pi+\eta$，线性式为 $\ddot\eta-(g/l)\eta=0$，有指数增长解。惯性系的固定支点、无驱动假设不可省去；快速竖直振动支点是另一个时变系统。
\end{example}

\begin{example}[同样零曲率，三种不同运动]
令 $T=m\dot q^2/2$。对 $V=aq^4$、$a>0$，原点为严格极小，能量边界仍证明稳定，但振动不是固定频率的简谐运动。转折振幅为 $A$ 时
\[
T(A)=4\sqrt{\frac{m}{2a}}\frac1A
\int_0^1\frac{\dd u}{\sqrt{1-u^4}},
\]
故频率随 $A$ 趋零而趋零。对 $V=-aq^4$，初态 $q(0)=\varepsilon>0,\dot q(0)=0$ 满足 $\ddot q=4aq^3/m>0$，必离开固定小邻域，原点不稳定。对 $V=0$，$q=q(0)+\dot q(0)t$；任意小非零初速度都可漂离原点，单个平衡也不稳定。这不排斥把整体平移约化后，内部振动保持稳定。
\end{example}

\begin{physicsbox}
“很小的力”不等于“能把质点拉回”。正曲率给出与位移同阶的恢复力；四次势阱给出更弱的三次恢复力；平坦方向根本没有恢复力。零频模应保留为自由漂移方程，不能把它当作周期无穷大的普通余弦解而漏掉初速度。
\end{physicsbox}

\section{单自由度：自由衰减与受迫响应}
先想象把弹簧拉开后放手，再比较有摩擦时的衰减和持续推它时的响应。三种模型中，测量的频率和能量收支各有不同。

\subsection{自由振动与初值}
知道释放时偏离了多少、推了一下有多快，就给了两个初值。解中的两个常数应由这两个初值确定。

把一维平衡移到 $q=0$。线性近似下，恢复力是 $-kq$，所以运动方程为
\begin{equation}
m\ddot q+kq=0,\quad \omega_0=\sqrt{k/m},\quad m,k>0.
\label{eq:sho}
\end{equation}
$m$ 是该坐标的等效惯性。例如 $q$ 是单摆的摆角时，$m$ 对应 $m_{\rm b}l^2$，其中 $m_{\rm b}$ 是摆球质量。将方程除以 $m$，试代 $\cos\omega_0t$ 和 $\sin\omega_0t$，二者的二阶导数均为自身乘 $-\omega_0^2$。这两个独立解张成二阶线性方程的解空间，故
\[
q(t)=C\cos\omega_0t+D\sin\omega_0t,\qquad
\dot q(t)=-\omega_0 C\sin\omega_0t+\omega_0 D\cos\omega_0t .
\]
给定 $q(0)=q_i,\dot q(0)=v_i$，就有 $C=q_i$、$\omega_0D=v_i$。若把解写成 $A\cos(\omega_0t+\phi)$，展开余弦得
\[
A\cos\phi=q_i,\qquad -A\sin\phi=v_i/\omega_0,\qquad
A^2=q_i^2+v_i^2/\omega_0^2 .
\]
两式共同确定非零振幅时的相位。频率 $\nu=\omega_0/(2\pi)$ 的单位为 Hz，角频率 $\omega_0$ 通常用 $\mathrm{rad/s}$ 表示。

定义机械能 $E=m\dot q^2/2+kq^2/2$。对时间求导，再用运动方程，得到
\[
\dot E=\dot q(m\ddot q+kq)=0,\qquad
E=\tfrac12mv_i^2+\tfrac12kq_i^2=\tfrac12kA^2 .
\]
最后一步使用了 $k=m\omega_0^2$。接下来加入阻力，仍可用同一乘积检验能量怎样变化。

\subsection{阻尼与功率}
阻尼就像把振子的一部分机械能不断传给外界。用阻尼力与速度的乘积可直接看出能量损失的快慢。

线性黏性阻力为 $-c\dot q$，其中 $c\ge0$ 是阻尼系数。它也可由瑞利耗散函数 $\mathcal R=c\dot q^2/2$ 写成 $-\partial\mathcal R/\partial\dot q$。用 $m\ddot q+kq=-c\dot q$ 代入上一小节的能量导数，并令 $\gamma=c/(2m)$，得到
\begin{equation}
\ddot q+2\gamma\dot q+\omega_0^2q=0,\qquad
\dot E=\dot q(m\ddot q+kq)=-c\dot q^2\le0.
\label{eq:damped-osc}
\end{equation}
因此 $c\dot q^2$ 是瞬时耗散功率。阻尼力总与速度方向相反，即使质点此刻正向平衡点运动，也会损失机械能。

\paragraph{先把衰减因子分离出来。}
令 $q=e^{-\gamma t}y$，则
\[
\dot q=e^{-\gamma t}(\dot y-\gamma y),\qquad
\ddot q=e^{-\gamma t}(\ddot y-2\gamma\dot y+\gamma^2y).
\]
代入式 \eqref{eq:damped-osc} 后，含 $\dot y$ 的两项抵消，剩下
\[
\ddot y+(\omega_0^2-\gamma^2)y=0,\qquad
y(0)=q_i,\quad \dot y(0)=v_i+\gamma q_i.
\]
于是只要看 $\omega_0^2-\gamma^2$ 的符号，就能由同一个方程得到三种解。

\begin{itemize}
\item \textbf{欠阻尼}（$0\le\gamma<\omega_0$）。取 $\omega_d=\sqrt{\omega_0^2-\gamma^2}$。$y$ 是角频率为 $\omega_d$ 的简谐振动，代入两个初值便得
\[
q=e^{-\gamma t}
\left[q_i\cos\omega_dt+\frac{v_i+\gamma q_i}{\omega_d}\sin\omega_dt\right].
\]
当 $\gamma>0$ 时，振幅包络按 $e^{-\gamma t}$ 衰减。$\gamma=0$ 时恢复无阻尼解。
\item \textbf{临界阻尼}（$\gamma=\omega_0$）。此时 $\ddot y=0$，所以 $y=q_i+(v_i+\omega_0q_i)t$，即
\[
q=e^{-\omega_0t}\big[q_i+(v_i+\omega_0q_i)t\big].
\]
因子 $t$ 来自 $y$ 的线性解，也可从欠阻尼解中取 $\sin(\omega_dt)/\omega_d\to t$ 得到。
\item \textbf{过阻尼}（$\gamma>\omega_0$）。令 $\alpha=\sqrt{\gamma^2-\omega_0^2}$，则 $y$ 由 $e^{\alpha t}$ 和 $e^{-\alpha t}$ 叠加。乘回衰减因子得 $q=Ae^{s_+t}+Be^{s_-t}$，其中 $s_\pm=-\gamma\pm\alpha$ 均为负。由 $A+B=q_i$、$s_+A+s_-B=v_i$ 解得
\[
A=\frac{v_i-s_-q_i}{s_+-s_-},\qquad
B=\frac{s_+q_i-v_i}{s_+-s_-}.
\]
\end{itemize}
一般初值下，过阻尼的长期运动由较慢的 $s_+$ 项控制。其衰减率为 $-s_+=\omega_0^2/(\gamma+\alpha)<\omega_0$。临界阻尼的渐近衰减指数为 $\omega_0$，但是否越过平衡点、何时进入给定误差带，还须结合初值判断。

\paragraph{用一周期能量损失理解品质因数。}
对 $\gamma>0$ 定义 $Q=\omega_0/(2\gamma)$。欠阻尼时取 $T_d=2\pi/\omega_d$。括号内的正弦和余弦一周期后复原，故 $q$ 与 $\dot q$ 都乘以 $e^{-\gamma T_d}$，能量因而乘以 $e^{-2\gamma T_d}$。将正的损失量记为 $\Delta E=E(t)-E(t+T_d)$，则
\[
\frac{\Delta E}{E(t)}=1-e^{-2\gamma T_d}
\simeq 2\gamma T_d\simeq\frac{4\pi\gamma}{\omega_0},\qquad
Q\simeq\frac{2\pi E(t)}{\Delta E}.
\]
第一个近似要求一周期内衰减很小，第二个近似用 $\omega_d\simeq\omega_0$。因此用单周期损失估计 $Q$ 的条件是 $Q\gg1$。

\subsection{稳态、相位和平均吸收功率}
持续用同一节拍推振子时，经过自由振动逐渐衰减的阶段，剩下的运动会跟着驱动频率走，但峰值通常比驱动力来得晚。位移最大的频率和平均吸收功率最大的频率要分别求。

\paragraph{从运动方程求复振幅。}
设外力为 $F_0\cos\omega t$，其中 $F_0>0,\omega>0,\gamma>0$。运动方程为
\[
\ddot q+2\gamma\dot q+\omega_0^2q=(F_0/m)\cos\omega t .
\]
先把右端写成 $(F_0/m)e^{i\omega t}$，试取复特解 $Xe^{i\omega t}$，最后取实部。这样做是因为左端是实系数线性算子，对复数取实部与微分可交换。一阶和二阶微分分别乘以 $i\omega$ 和 $-\omega^2$，于是
\[
(\omega_0^2-\omega^2+2i\gamma\omega)Xe^{i\omega t}
=(F_0/m)e^{i\omega t}.
\]
约去共同因子，令 $X=Ae^{-i\delta}$，得到
\begin{align}
q_{\rm st}(t)&=\operatorname{Re}[Xe^{i\omega t}]=A\cos(\omega t-\delta),\\
X&=\frac{F_0/m}{\omega_0^2-\omega^2+2i\gamma\omega},\quad
A=\frac{F_0/m}{\sqrt{(\omega_0^2-\omega^2)^2+4\gamma^2\omega^2}}.
\label{eq:forced}
\end{align}
若令 $\Delta=\omega_0^2-\omega^2$、$B=2\gamma\omega$、$D=\Delta^2+B^2$，则 $1/(\Delta+iB)=(\Delta-iB)/D$，所以
\[
\cos\delta=\Delta/\sqrt D,\qquad \sin\delta=B/\sqrt D .
\]
由于 $B>0$，滞后角在 $(0,\pi)$ 内。低频时 $\delta$ 接近 $0$，$\omega=\omega_0$ 时为 $\pi/2$，高频时趋于 $\pi$。只用正切会丢掉余弦的符号。给定初值时，完整解为 $q=q_{\rm st}+q_{\rm h}$，其中齐次解的初值是 $q_i-q_{\rm st}(0)$、$v_i-\dot q_{\rm st}(0)$。等齐次解衰减后，观测到的才是这里求出的周期稳态。

\paragraph{位移峰在哪里？}
振幅 $A$ 最大等价于 $D$ 最小。直接微分得
\[
D'(\omega)=4\omega(\omega^2-\omega_0^2+2\gamma^2).
\]
若 $\gamma<\omega_0/\sqrt2$，导数在正频率根的左侧为负、右侧为正，故
\[
\omega_R=\sqrt{\omega_0^2-2\gamma^2}.
\]
若 $\gamma\ge\omega_0/\sqrt2$，$D$ 在 $\omega>0$ 上递增，振幅的最大值出现在静态极限 $A(0)=F_0/k$。

\paragraph{功率为什么在另一频率取峰值？}
外力的瞬时功率是 $F\dot q$。对稳态，展开速度
\[
\dot q_{\rm st}=A\omega\big[\sin\delta\cos\omega t-\cos\delta\sin\omega t\big].
\]
在一个驱动周期内，$\overline{\cos^2\omega t}=1/2$、$\overline{\sin\omega t\cos\omega t}=0$。乘以 $F_0\cos\omega t$ 并平均，再代入 $\sin\delta=B/\sqrt D$，得
\begin{equation}
\overline P=\frac{F_0A\omega}{2}\sin\delta
=m\gamma\omega^2A^2
=\frac{\gamma F_0^2\omega^2/m}{(\omega_0^2-\omega^2)^2+4\gamma^2\omega^2}.
\label{eq:osc-power}
\end{equation}
也可直接用能量收支 $\dot E=F\dot q-c\dot q^2$ 检查。稳态的 $E$ 一周期后复原，故 $\overline{\dot E}=0$，输入平均功率等于 $c\overline{\dot q^2}=cA^2\omega^2/2$。只有速度与驱动力同相的分量在整周期内贡献净输入功。

\paragraph{半功率带宽怎样得到？}
将功率式分子、分母同除以 $\omega^2$，可写成
\[
\overline P=\frac{\gamma F_0^2/m}
 {(\omega_0^2/\omega-\omega)^2+4\gamma^2}.
\]
分母的平方项在 $\omega=\omega_0$ 恰为零，所以功率峰值为 $P_{\max}=F_0^2/(4m\gamma)$。功率降到一半时分母加倍，故 $|\omega_0^2-\omega^2|=2\gamma\omega$。在峰左侧，去绝对值后得 $\omega^2+2\gamma\omega-\omega_0^2=0$。在峰右侧则得 $\omega^2-2\gamma\omega-\omega_0^2=0$。分别取正根，得到
\[
\omega_\pm=\sqrt{\omega_0^2+\gamma^2}\pm\gamma,\qquad
\omega_+-\omega_-=2\gamma .
\]
在这个线性黏性模型中，功率半高宽是精确结果。两个半功率点的算术中点为 $\sqrt{\omega_0^2+\gamma^2}$，弱阻尼时才近似等于 $\omega_0$。

\begin{example}[同一个振子中的三种频率]
取 $m=1\,\mathrm{kg}$、$k=100\,\mathrm{N/m}$、$c=4\,\mathrm{N\,s/m}$。先求出 $\omega_0=10\,\mathrm{rad/s}$、$\gamma=2\,\mathrm{s^{-1}}$。撤去驱动后，欠阻尼自由振动的角频率是 $\omega_d=\sqrt{96}\,\mathrm{rad/s}\approx9.80\,\mathrm{rad/s}$。若改变驱动频率测稳态位移，最大振幅出现在 $\omega_R=\sqrt{92}\,\mathrm{rad/s}\approx9.59\,\mathrm{rad/s}$。测平均功率时，峰在 $10\,\mathrm{rad/s}$，半功率点约为 $8.20$ 与 $12.20\,\mathrm{rad/s}$。先区分测量量，再决定应使用哪一个频率。

三个频率来自不同方程。自由衰减的根为 $-\gamma\pm i\sqrt{\omega_0^2-\gamma^2}$，而受迫位移峰由 $\omega_0^2-2\gamma^2$ 给出。半功率两根为 $\sqrt{100+4}\mp2$（以 $\mathrm{rad/s}$ 为单位），间距恰为 $4\,\mathrm{rad/s}$。
\end{example}

\begin{figure}[htbp]
\centering
\begin{tikzpicture}
\begin{axis}[width=0.8\textwidth,height=5.5cm,axis lines=left,xlabel={$\omega/\omega_0$},ylabel={$A/(F_0/k)$},xmin=0,xmax=2,ymin=0,ymax=5.5,legend style={draw=none,at={(0.98,0.96)},anchor=north east},samples=400]
\addplot[blue!65!black,very thick,domain=0:2]{1/sqrt((1-x^2)^2+.04*x^2)};\addlegendentry{$\gamma/\omega_0=0.1$}
\addplot[red!70!black,thick,dashed,domain=0:2]{1/sqrt((1-x^2)^2+.16*x^2)};\addlegendentry{$0.2$}
\addplot[green!45!black,thick,dashdotted,domain=0:2]{1/sqrt((1-x^2)^2+.64*x^2)};\addlegendentry{$0.4$}
\end{axis}
\end{tikzpicture}
\caption{无量纲位移幅频曲线。图的纵轴由静态位移 $F_0/k$ 归一化；位移峰随阻尼偏移，不能直接拿它代替吸收功率峰。}
\label{fig:resonance-osc}
\end{figure}

\begin{example}[有限阻尼响应与无阻尼共振]
令驱动频率为 $\omega_0$。若 $\gamma>0$，稳态振幅为 $F_0/(2m\gamma\omega_0)$，相位滞后 $\pi/2$，并不随时间线性发散。若先令 $\gamma=0$，正弦和余弦已是齐次解，代回左端都给零。为产生与外力同相的余弦项，试取 $q_p=Ct\sin\omega_0t$。乘积求导给出 $\dot q_p=C(\sin\omega_0t+\omega_0t\cos\omega_0t)$，再求导并加上 $\omega_0^2q_p$，两个含 $t\sin\omega_0t$ 的项抵消：
\[
\ddot q_p+\omega_0^2q_p=2C\omega_0\cos\omega_0t,
\qquad C=\frac{F_0}{2m\omega_0}.
\]
零位移、零初速度下该特解就是全解。由此可见，先取长期稳态再令阻尼趋零，与先取无阻尼再观察有限时间，是不同极限。线性发散预告近似终会失效，不证明任何真实系统都按某种固定方式饱和。
\end{example}


\section{从精确方程到简正坐标}
两个摆连在一起时，轻推一只常会带动另一只。若能找到两只摆按固定比例一起动的方式，复杂运动就可以拆成几种简单振动。

\subsection{小参数如何进入}
把所有偏离平衡的量一起缩小，按它们缩小的快慢分组。保留决定首要恢复力和惯性的项，就会得到线性方程。

在本章自然系统的假设下，动能精确为
\begin{equation}
T=\frac12\sum_{a,b}M_{ab}(q)\dot q_a\dot q_b.
\label{eq:exact-T-multi}
\end{equation}
逐项求导得到
\[
\frac{\dd}{\dd t}T_{\dot q_a}=\sum_b M_{ab}\ddot q_b+
\sum_{b,c}M_{ab,q_c}\dot q_b\dot q_c,\qquad
T_{q_a}=\frac12\sum_{b,c}M_{bc,q_a}\dot q_b\dot q_c.
\]
把第二项对 $b,c$ 对称化，精确方程成为
\begin{equation}
\sum_bM_{ab}(q)\ddot q_b+
\sum_{b,c}\Gamma_{abc}(q)\dot q_b\dot q_c+V_{q_a}(q)=0,
\label{eq:exact-multi-eq}
\end{equation}
其中
\[
\Gamma_{abc}=\tfrac12(M_{ac,q_b}+M_{ab,q_c}-M_{bc,q_a}).
\]
$\Gamma$ 项来自坐标相关惯性，不是另加的外力。

令 $\eta=q-q_0$，定义 $M_0=M(q_0)$ 和精确余项
\begin{align*}
\mathcal N_a={}&-\sum_b[M_{ab}(q_0+\eta)-M_{ab}(q_0)]\ddot\eta_b\\
&-\sum_{b,c}\Gamma_{abc}(q_0+\eta)\dot\eta_b\dot\eta_c
-\left[V_{q_a}(q_0+\eta)-\sum_bK_{ab}\eta_b\right].
\end{align*}
于是 $M_0\ddot\eta+K\eta=\mathcal N$ 是精确重写。取固定长度、角度和时间尺度后，令 $\eta=\epsilon u$，并要求在考察的固定时间区间上 $u,\dot u,\ddot u$ 一致有界。三组余项分别是 $O(\epsilon^2)$、$O(\epsilon^2)$、$O(\epsilon^2)$，除以 $\epsilon$ 后均趋零。因此最低阶方程是
\begin{equation}
M_0\ddot{\boldsymbol\eta}+K\boldsymbol\eta=0,
\qquad
L_2=\tfrac12\dot{\boldsymbol\eta}^{\mathsf T}M_0\dot{\boldsymbol\eta}
-\tfrac12\boldsymbol\eta^{\mathsf T}K\boldsymbol\eta.
\label{eq:matrix-equation}
\end{equation}
也可直接在 $L$ 中数总阶次：$\eta\dot\eta^2$ 是三阶，$\dot\eta^2$ 是二阶。位移小而速度大不满足这一标度；小非线性频移在很长时间内还会积累相位差，不能把固定区间近似称为任意长时间精确。

\begin{physicsbox}
$M_0$ 是“移动这些坐标需要多少惯性”的二次型，$K$ 是“偏离平衡需要多少势能”的二次型。求频率要同时比较二者。弧度虽无量纲，摆角的惯性仍是转动惯量；把它误写为质点质量，会直接给出错误频率。
\end{physicsbox}

\subsection{广义本征值与质量加权}
假设各坐标以同一频率振动，振幅保持固定比例。代入后，刚度矩阵告诉我们恢复力，质量矩阵告诉我们各方向加速有多难。

以后简记 $M=M_0$。试解 $\eta=\bA e^{i\omega t}$，得到
\begin{equation}
K\bA=\lambda M\bA,\qquad \lambda=\omega^2,\qquad
\det(K-\lambda M)=0.
\label{eq:secular}
\end{equation}
行列式方程称为久期方程。有 $s$ 个实本征值（计重数），但它们为零或负时应分别求自由漂移或指数解，而非一律叫作实振动频率。

为何根一定实、能否总找齐振型？先把各坐标用固定尺度表示为数值，再取正定矩阵 $M$ 的对称正平方根。令
\[
y=M^{1/2}\eta,\qquad D=M^{-1/2}KM^{-1/2}=D^{\mathsf T}.
\]
因 $M$ 固定在平衡点，$\ddot\eta=M^{-1/2}\ddot y$。将它与 $\eta=M^{-1/2}y$ 代入，再左乘 $M^{-1/2}$，惯性项变成 $\ddot y$，故 $\ddot y+Dy=0$。实对称矩阵的谱定理给出正交矩阵 $O$，使
\[
O^{\mathsf T}DO=\Lambda=\operatorname{diag}(\lambda_1,\ldots,\lambda_s).
\]
取 $G=M^{-1/2}O$，则
\begin{equation}
G^{\mathsf T}MG=I,\quad G^{\mathsf T}KG=\Lambda,\quad
G^{-1}=G^{\mathsf T}M.
\label{eq:osc-simultaneous}
\end{equation}
例如 $G^{\mathsf T}MG=O^{\mathsf T}O=I$，而 $G^{\mathsf T}KG=O^{\mathsf T}DO=\Lambda$。第一式说明 $G^{\mathsf T}M$ 是 $G$ 的左逆，因 $G$ 是方阵，它也是逆矩阵。$DO=O\Lambda$ 左乘 $M^{1/2}$ 还给出 $KG=MG\Lambda$。于是 $G$ 的各列是完备、质量正交的振型。$M^{-1}K$ 通常不对称，不能直接把它交给只适用于实对称矩阵的数值程序。

设 $G$ 的第 $r$ 列为 $g_r$，定义质量归一化简正坐标
\begin{equation}
\eta=G\xi,\qquad \xi=G^{\mathsf T}M\eta.
\label{eq:normal-coordinates}
\end{equation}
代入二次型，得到
\[
T_2=\tfrac12\sum_r\dot\xi_r^2,\quad V_2=\tfrac12\sum_r\lambda_r\xi_r^2,
\qquad \ddot\xi_r+\lambda_r\xi_r=0.
\]
对 $\lambda_r=\omega_r^2>0$，初值解为
\begin{equation}
\xi_r(t)=(g_r^{\mathsf T}M\eta_i)\cos\omega_rt+
\frac{g_r^{\mathsf T}Mv_i}{\omega_r}\sin\omega_rt.
\label{eq:general-solution}
\end{equation}
$\lambda_r=0$ 时改为 $\xi_r(0)+\dot\xi_r(0)t$；$\lambda_r<0$ 时使用双曲正弦、余弦。若用未归一化振型 $a_r$ 展开 $\eta=\sum_r a_r c_r$，则
\[
c_r=\frac{a_r^{\mathsf T}M\eta}{a_r^{\mathsf T}Ma_r},\qquad
T_2=\frac12\sum_r(a_r^{\mathsf T}Ma_r)\dot c_r^2.
\]
这与质量归一化坐标的公式不同；改变振型长度会改变模态坐标的标度，物理位移保持不变。长度坐标下，$\xi$ 的单位一般为质量平方根乘长度。

\subsection{正交性、简并、坐标变换与频率估计}
若两种振动的频率不同，它们的正交性要按质量来衡量。频率相同则允许在同一个子空间里重选振型；用试探振型还能估计最低频率。

若 $Ka_r=\lambda_rMa_r$，利用对称性，有
\[
a_s^{\mathsf T}Ka_r=\lambda_r a_s^{\mathsf T}Ma_r
=\lambda_s a_s^{\mathsf T}Ma_r.
\]
故不同本征值的振型满足 $a_s^{\mathsf T}Ma_r=0$。重根时该等式本身不强迫正交，但谱定理允许在简并子空间内选质量正交基。重复频率不等于缺少振型；简并空间内的基旋转也不改变运动。

对非零试探方向 $a$，Rayleigh商为
\[
R(a)=\frac{a^{\mathsf T}Ka}{a^{\mathsf T}Ma}.
\]
把 $a=\sum_r b_rg_r$ 代入，得
\[
R(a)=\frac{\sum_r\lambda_rb_r^2}{\sum_r b_r^2},\qquad
\lambda_{\min}\le R(a)\le\lambda_{\max}.
\]
因此任一试探方向给最低本征值的上界，取遍方向才得到精确最小值；若有平移零模而只估计内部振动，必须先限制在其质量正交补。

换局部坐标 $\eta=Jz$，$J$ 可逆时，$M_z=J^{\mathsf T}MJ,K_z=J^{\mathsf T}KJ$，从而
\[
\det(K_z-\lambda M_z)=\det(J)^2\det(K-\lambda M).
\]
坐标单位或标度改变不改变频率。这也解释了为何只对 $K$ 做普通对角化通常不够。

\begin{keybox}
求解顺序：先核验平衡和坐标独立性，再求 $M,K$；检查对称性及 $M$ 正定；解广义本征值问题；按质量内积归一化；投影位置和速度初值；最后还原位移，并用原方程残差和能量核验。零模、负模和简并应在求根后分别处理。
\end{keybox}

\begin{modernbox}[title={进一步探索：陀螺项为什么另需处理}]
若二次拉格朗日量含 $\dot\eta^{\mathsf T}A\eta$，常矩阵 $A$ 的对称部分是全导数，反对称部分却给出
\[
M\ddot\eta+(A-A^{\mathsf T})\dot\eta+K\eta=0.
\]
令扰动按 $e^{i\omega t}$ 变化，得到频率方程 $\det[K-\omega^2M+i\omega(A-A^{\mathsf T})]=0$。第\ref{sec:sleeping-top}节给出含这种速度项的稳定例子。例如旋转参考系中的科里奥利项属于此类，不能只查 $K$ 的符号就套用前文判据。理解这一拓展需要第三章的速度耦合及复线性代数；本章核心仍限于无此项的自然系统。
\end{modernbox}


\section{耦合振动的典型模型}
以下几个模型都先回答三个问题。谁在动，谁把谁拉回去，怎样的集体动作能保持固定形状。

\subsection{双振子、耦合摆与拍}
把两只相同的摆同时往一边拉，或分别往相反方向拉，弹簧的伸长不同。只放开一只摆时，两种运动会叠在一起，并出现一快一慢的节奏。

\begin{example}[交叉项的符号决定哪个振型频率低]
设两个长度坐标的二次拉格朗日量为
\[
L_2=\frac m2(\dot x^2+\dot y^2)
-\frac{m\omega_0^2}{2}(x^2+y^2)+m\kappa xy,
\quad 0<\kappa<\omega_0^2.
\]
$\kappa$ 的单位为时间的负二次方。方程为
\[
\ddot x+\omega_0^2x-\kappa y=0,\qquad
\ddot y+\omega_0^2y-\kappa x=0.
\]
相加、相减分别得到
\[
\ddot c_++(\omega_0^2-\kappa)c_+=0,\quad
\ddot c_-+(\omega_0^2+\kappa)c_-=0,\quad
c_\pm=\frac{x\pm y}{\sqrt2}.
\]
故低频为同相 $(1,1)$，高频为反相 $(1,-1)$。直接代回两条原方程是最简单的符号检查。这里 $T_2=m(\dot c_+^2+\dot c_-^2)/2$；质量归一化坐标为 $\xi_\pm=\sqrt m\,c_\pm$。当 $\kappa=\omega_0^2$ 出现零模，超过该值势能有下降方向。
\end{example}

\begin{example}[无预紧耦合摆：先确定弹簧伸长]
两个固定支点等高，相距 $d>0$；两条轻刚杆长均为 $l$，摆锤质量均为 $m$。轻弹簧连接两摆锤，自然长度恰为 $d$、劲度 $k>0$。用同一向右为正的摆角 $\theta_1,\theta_2$，两端相对位置差为
\[
(d+l\sin\theta_2-l\sin\theta_1,\;l\cos\theta_1-l\cos\theta_2).
\]
弹簧长度减 $d$ 的一阶项为 $l(\theta_2-\theta_1)$，所以其二次势能是 $kl^2(\theta_2-\theta_1)^2/2$。这一步使用了无预紧假设；若弹簧已有张力，几何长度的二阶项也可能贡献刚度。
\[
M=ml^2I,\qquad
K=\begin{pmatrix}mgl+kl^2&-kl^2\\-kl^2&mgl+kl^2\end{pmatrix}.
\]
令 $d_\lambda=mgl+kl^2-ml^2\lambda$，久期式为 $(d_\lambda-kl^2)(d_\lambda+kl^2)=0$，得到
\begin{equation}
\omega_+^2=g/l,\quad a_+=(1,1)^{\mathsf T};\qquad
\omega_-^2=g/l+2k/m,\quad a_-=(1,-1)^{\mathsf T}.
\label{eq:coupled-pend-freq}
\end{equation}
同相运动中弹簧在该理想几何下长度不变；反相运动增加恢复力。两极限 $k\to0$ 给两个简并单摆，$g\to0$ 给整体横向零模与内部伸缩模。
\end{example}

\begin{figure}[htbp]
\centering
\begin{tikzpicture}[>=Stealth]
\foreach \shift/\sign/\lab in {0/1/同相低频,6/-1/反相高频}{
\begin{scope}[xshift=\shift cm]
\draw[thick](-.4,0)--(4,0);
\coordinate (P) at (0,0);\coordinate (Q) at (3,0);
\coordinate (B) at ({2*sin(12)},{-2*cos(12)});
\coordinate (C) at ({3+\sign*2*sin(12)},{-2*cos(12)});
\draw[dashed,gray](P)--(0,-2.2) (Q)--(3,-2.2);
\draw[blue!65!black,very thick](P)--(B) (Q)--(C);
\draw[green!45!black,thick,decorate,decoration={coil,aspect=.5,segment length=4pt,amplitude=3pt}](B)--(C);
\fill (P) circle(1.8pt);\fill (Q) circle(1.8pt);
\fill[blue!65!black](B)circle(3pt);\fill[blue!65!black](C)circle(3pt);
\draw[->,red!70!black,thick](0,-2.45)--(.65,-2.45);
\draw[->,red!70!black,thick](3,-2.45)--({3+\sign*.65},-2.45);
\node at (1.6,.35){\lab};
\end{scope}}
\end{tikzpicture}
\caption{耦合摆的两种振型快照。虚线为平衡位置，红箭头表示该快照的位移方向；弹簧严格连接两摆锤，而非连接平衡位置。左右图使用相同摆长和角度大小。}
\label{fig:coupled-pendulum}
\end{figure}

\begin{example}[只释放一只摆：从投影得到拍]
在上一模型中令 $\theta_1(0)=\Theta,\theta_2(0)=0$，初速度均为零。分解为两模各一半，得到
\begin{align*}
\theta_1&=\tfrac\Theta2(\cos\omega_+t+\cos\omega_-t)
=\Theta\cos(\bar\omega t)\cos(\Delta\omega t/2),\\
\theta_2&=\tfrac\Theta2(\cos\omega_+t-\cos\omega_-t)
=\Theta\sin(\bar\omega t)\sin(\Delta\omega t/2),
\end{align*}
其中 $\bar\omega=(\omega_-+\omega_+)/2$、$\Delta\omega=\omega_--\omega_+>0$。弱耦合 $k/m\ll g/l$ 时，快振动由慢包络调制，第一次近乎完全转移发生在 $t\simeq\pi/\Delta\omega$。取 $k/m=0.1g/l$，则 $\omega_-/\omega_+=\sqrt{1.2}$，首次转移约在 $\omega_+t=32.9$。

每个简正模的能量恒定。说“能量在两摆间交换”时，局部能量怎样分配耦合弹簧有约定；弱耦合下，分别用两只裸单摆的能量观察交换最直观。$k\to0$ 时差频趋零，转移时间发散，不发生无耦合能量传递。
\end{example}

\begin{figure}[htbp]
\centering
\begin{tikzpicture}
\begin{axis}[width=.84\textwidth,height=5.2cm,xlabel={$\omega_+t$},ylabel={$\theta_1/\Theta$},xmin=0,xmax=70,ymin=-1.15,ymax=1.15,axis lines=left,samples=1000]
\addplot[blue!65!black,thick,domain=0:70]{.5*(cos(deg(x))+cos(deg(sqrt(1.2)*x)))};
\addplot[red!70!black,dashed,domain=0:70]{cos(deg((sqrt(1.2)-1)*x/2))};
\addplot[red!70!black,dashed,domain=0:70]{-cos(deg((sqrt(1.2)-1)*x/2))};
\end{axis}
\end{tikzpicture}
\caption{耦合摆的拍：蓝线严格取正文的线性叠加解，红虚线为其有符号包络的正负两支。包络消失并不代表两个简正模的能量消失。}
\label{fig:osc-beats}
\end{figure}

\subsection{不等质量与双摆}
两只质量不同时，同样的位移并不意味着同样的惯性。双摆还要注意，下球速度受到两根杆的转动共同影响。

\begin{example}[欧氏不正交，质量内积仍正交]
两质点质量为 $m,2m$，左端与墙连接劲度 $k$ 的弹簧，两质点间连接劲度 $2k$ 的弹簧，第二质点与右墙也连接劲度 $2k$ 的弹簧。沿直线的平衡位移为 $x_1,x_2$，质量矩阵为 $M=m\operatorname{diag}(1,2)$，刚度矩阵为
\[
K=k\begin{pmatrix}3&-2\\-2&4\end{pmatrix}.
\]
记 $z=m\lambda/k$，久期式为
\[
(3-z)(4-2z)-4=2(z-1)(z-4)=0.
\]
故 $\omega_1^2=k/m,\omega_2^2=4k/m$，未归一化振型为 $(1,1)$ 与 $(2,-1)$。其普通点积为1，而
\[
(1,1)M(2,-1)^{\mathsf T}=0.
\]
质量归一化矩阵及投影为
\begin{equation}
G=\left(\frac{(1,1)^{\mathsf T}}{\sqrt{3m}},
\frac{(2,-1)^{\mathsf T}}{\sqrt{6m}}\right),\quad
\xi_1=\sqrt{\frac m3}(x_1+2x_2),\quad
\xi_2=\sqrt{\frac{2m}3}(x_1-x_2).
\label{eq:osc-unequal-G}
\end{equation}
代回验证 $G^{\mathsf T}MG=I$、$G^{\mathsf T}KG=\operatorname{diag}(k/m,4k/m)$。若 $x(0)=(a,0),\dot x(0)=0$，
\[
x(t)=\frac a3(1,1)^{\mathsf T}\cos\sqrt{k/m}\,t
+\frac a3(2,-1)^{\mathsf T}\cos(2\sqrt{k/m}\,t).
\]
初始总能量为 $3ka^2/2$，两模能量分别为 $ka^2/6$ 与 $4ka^2/3$，相加一致。
\end{example}

\begin{example}[双摆：同样两根杆，惯性有交叉项]
在式 \eqref{eq:lag-double-pendulum} 中取 $m_1=m_2=m,\ l_1=l_2=l$，两角都从向下竖直量起，在 $(0,0)$ 附近线性化。
沿用第三章的两个绝对摆角，取 $m_1=m_2=m,l_1=l_2=l$，在最低平衡展开。二次型为
\[
T_2=\tfrac12ml^2(2\dot\theta_1^2+2\dot\theta_1\dot\theta_2+\dot\theta_2^2),\quad
V_2=\tfrac12mgl(2\theta_1^2+\theta_2^2).
\]
这里两杆之间没有弹簧，耦合来自下球速度包含上球运动。
\[
M=ml^2\begin{pmatrix}2&1\\1&1\end{pmatrix},\quad
K=mgl\begin{pmatrix}2&0\\0&1\end{pmatrix}.
\]
令 $z=\omega^2l/g$，行列式给出
\[
(2-2z)(1-z)-z^2=z^2-4z+2=0.
\]
因此
\begin{equation}
\omega_{\rm low}^2=(2-\sqrt2)g/l,\quad a_{\rm low}=(1,\sqrt2)^{\mathsf T};\qquad
\omega_{\rm high}^2=(2+\sqrt2)g/l,\quad a_{\rm high}=(1,-\sqrt2)^{\mathsf T}.
\label{eq:osc-double-pend}
\end{equation}
以第一行求角度比 $a_2/a_1=2(1-z)/z$，即可独立检查正负。两振型质量范数平方为 $ml^2(4\pm2\sqrt2)$，且彼此质量正交。给定小角初值后应按式 \eqref{eq:general-solution} 投影，不能各自用单摆频率 $\sqrt{g/l}$。
\end{example}

\begin{figure}[htbp]
\centering
\begin{tikzpicture}[>=Stealth]
\foreach \shift/\sgn/\lab in {0/1/低频同向,5.8/-1/高频反向}{
\begin{scope}[xshift=\shift cm]
\coordinate(O)at(0,0);
\coordinate(A)at({2*sin(8)},{-2*cos(8)});
\coordinate(B)at({2*sin(8)+2*sin(\sgn*sqrt(2)*8)},{-2*cos(8)-2*cos(\sgn*sqrt(2)*8)});
\draw[dashed,gray](O)--(0,-4.2);
\draw[dashed,gray](A)--++(0,-2.1);
\draw[very thick,blue!65!black](O)--(A)--(B);
\fill(O)circle(2pt);\fill(A)circle(3pt)node[right]{$m$};\fill(B)circle(3pt)node[right]{$m$};
\draw[->,red!70!black](0,-1.8)--(.55,-1.8);
\draw[->,red!70!black](0,-4.4)--({.55*(1+\sgn*sqrt(2))},-4.4);
\node at (.3,.4){\lab};
\end{scope}}
\end{tikzpicture}
\caption{双摆两振型，取 $\theta_1=8^\circ,\theta_2=\pm\sqrt2\theta_1$ 画几何快照。红箭头表示质点横向位移的线性近似：第二球位移是 $l(\theta_1+\theta_2)$，不只由第二杆角度决定。}
\label{fig:osc-double-modes}
\end{figure}

\subsection{纵向三原子模型与零模}
把三颗质点一起平移，各段弹簧的长度都不变。剩下的两种独立动作才会伸缩弹簧并产生恢复力。

\begin{example}[伸缩频率与质心分离]
三个质点沿轴运动，质量为 $m,\mu,m$，两条相同轻弹簧劲度为 $k$，用 $x_i$ 表示相对平衡的轴向位移。模型只保留纵向自由度：
\[
M=\operatorname{diag}(m,\mu,m),\quad
V_2=\tfrac k2(x_2-x_1)^2+\tfrac k2(x_3-x_2)^2,\quad
K=k\begin{pmatrix}1&-1&0\\-1&2&-1\\0&-1&1\end{pmatrix}.
\]
先直接作用于 $(1,1,1)$，得零本征值。其坐标是质心平移，$X_{\rm cm}=(mx_1+\mu x_2+mx_3)/(2m+\mu)$ 满足 $\ddot X_{\rm cm}=0$。

取端点反向 $a_1=(1,0,-1)$，有 $Ka_1=(k/m)Ma_1$，故 $\omega_1^2=k/m$。再取端点同向 $(1,b,1)$，非零模必须与平移模质量正交，因而 $2m+\mu b=0$，$b=-2m/\mu$。第一行方程给出
\[
\omega_2^2=\frac{k}{m}(1-b)=k\left(\frac1m+\frac2\mu\right),\qquad
 a_2=(1,-2m/\mu,1).
\]
三组范数平方分别为 $2m+\mu,2m,2m+4m^2/\mu$，除以各自范数就得到 $G$。每个非零模满足总动量为零；例如 $\mu\to\infty$ 时中心近乎固定，两伸缩频率趋于简并。

为防止“同相”命名歧义，直接看两段键长变化 $\delta d_1=x_2-x_1,\delta d_2=x_3-x_2$：第一伸缩模两段同时伸长/缩短，是对称伸缩；第二伸缩模一长一短，是反对称伸缩。原子位移的对称性还须考虑空间反演使矢量反向。
\end{example}

\begin{figure}[htbp]
\centering
\begin{tikzpicture}[>=Stealth]
\foreach \yy/\sgn/\lab in {0/-1/两键同伸缩,-2.2/1/两键反向伸缩}{
\node[anchor=east]at(-.7,\yy){\lab};
\foreach \xx/\mass in {0/m,2/\mu,4/m}{\draw[fill=blue!10](\xx,\yy)circle(.19);\node[below]at(\xx,\yy-.25){$\mass$};}
\draw[decorate,decoration={coil,aspect=.5,segment length=4pt,amplitude=3pt},green!45!black,thick](.19,\yy)--(1.81,\yy) (2.19,\yy)--(3.81,\yy);
\draw[red!70!black,->,thick](0,\yy+.5)--(.6,\yy+.5);
\draw[red!70!black,->,thick](4,\yy+.5)--({4+\sgn*.6},\yy+.5);
}
\node at(2,.52){$0$};
\draw[red!70!black,->,thick](2,-1.7)--(.8,-1.7);
\end{tikzpicture}
\caption{三质点模型的两个伸缩模，箭头表示位移振幅。本图取 $\mu=m$，下排幅度比严格为 $1:-2:1$；上排为 $1:0:-1$。整体平移零模未画。}
\label{fig:triatomic}
\end{figure}

\begin{modernbox}[title={进一步探索：从振型到光谱还缺什么？}]
这个模型说明质量、刚度与伸缩频率的关系；同位素改变质量而近似保留势能曲率，是应用广义本征值的一个切入点。但真实分子还有横向、转动自由度和电子结构。一个模能否被红外或拉曼观测到，还分别取决于偶极矩或极化率随该模是否变化，不能把所有数学振型都说成吸收峰。继续学习需三维振动、对称性及电磁耦合；本例不拟合具体分子的实验谱线。
\end{modernbox}


\section{模态受迫响应与阻尼的限制}
外力若顺着某个振型推，通常更容易把这个模激发出来。阻尼器若只装在一个位置，则可能同时感受到几个模的速度。

\subsection{广义力如何激发某个模}
把外力在各振型上的分量算出来，就能判断直接推到了哪些模。若模态阻尼矩阵有非对角项，一个模还可能被另一个模带动。

若有小外力向量 $f(t)$ 和线性黏性矩阵 $C=C^{\mathsf T}\ge0$，则
\begin{equation}
M\ddot\eta+C\dot\eta+K\eta=f(t),\qquad
\dot E_2=\dot\eta^{\mathsf T}f-\dot\eta^{\mathsf T}C\dot\eta.
\label{eq:osc-multi-driven}
\end{equation}
采用无阻尼振型 $\eta=G\xi$ 并左乘 $G^{\mathsf T}$，得到
\[
\ddot\xi+C_N\dot\xi+\Lambda\xi=f_N(t),\qquad
C_N=G^{\mathsf T}CG,\quad f_N=G^{\mathsf T}f.
\]
模态力 $f_{N,r}=g_r^{\mathsf T}f$ 不是 $g_r^{\mathsf T}Mf$；质量已经在惯性归一化中处理。对 $\Lambda>0$ 且 $C_N$ 为对角阵的情形，每个模有独立响应
\[
\xi_r^{\rm st}=\operatorname{Re}\left[
\frac{g_r^{\mathsf T}F}{\omega_r^2-\omega^2+i c_r\omega}e^{i\omega t}\right],\qquad
c_r=(C_N)_{rr}.
\]
若模态力为零，即使驱动频率靠近该本征频率，也不激发这个模。阻尼为零时，该式只适用于非共振特解，并不意味着初始自由响应会消失。

一个常用而受限的选择是比例阻尼 $C=\alpha M+\beta K$，此时
\[
C_N=\alpha I+\beta\Lambda,\quad c_r=\alpha+\beta\omega_r^2,
\]
$\alpha$ 的单位为时间的负一次方，$\beta$ 为时间。只要求模型总 $C$ 半正定；$\alpha,\beta\ge0$ 是保证它的一种充分选择。

\begin{example}[一个局部阻尼器一般耦合无阻尼模]
取式 \eqref{eq:osc-unequal-G} 的两质量模型，只在第一质点与墙之间加阻尼器，$C=\operatorname{diag}(c,0)$。则
\[
C_N=\frac c m\begin{pmatrix}1/3&\sqrt2/3\\\sqrt2/3&2/3\end{pmatrix}.
\]
非对角元不为零，两个无阻尼简正坐标通过速度耦合。尽管 $\dot E_2=-c\dot x_1^2\le0$，不能分别往两个独立简谐方程随意加衰减率；弱阻尼下忽略非对角项也需额外论证。一般谐波响应应直接解
\[
(K-\omega^2M+i\omega C)X=F.
\]
能量不增本身不保证任一模型的所有模都会衰减到零；可能存在阻尼看不到的运动方向。

因为阻尼只作用于第一坐标，变换后有 $(C_N)_{rs}=c(g_r)_1(g_s)_1$。第一行的两个振型分量是 $1/\sqrt{3m}$、$2/\sqrt{6m}$，所以交叉元为 $\sqrt2c/(3m)$。这也可从耗散功率 $-c[(g_1)_1\dot\xi_1+(g_2)_1\dot\xi_2]^2$ 的交叉项核对。
\end{example}

\begin{example}[调谐吸振器：共振之外还有反共振]
主质量 $m$ 由劲度 $k$ 的弹簧连墙，附加质量 $\mu$ 通过劲度 $\kappa$ 的弹簧连主质量，均沿直线无摩擦运动。对主质量施加 $F_0\cos\omega t$，两个平衡位移为 $x_1,x_2$，有
\[
M=\begin{pmatrix}m&0\\0&\mu\end{pmatrix},\quad
K=\begin{pmatrix}k+\kappa&-\kappa\\-\kappa&\kappa\end{pmatrix},\quad
F=(F_0,0)^{\mathsf T}.
\]
无阻尼、避开本征频率时，解两条代数方程得
\begin{align*}
X_1&=\frac{(\kappa-\mu\omega^2)F_0}{D(\omega)},\quad
X_2=\frac{\kappa F_0}{D(\omega)},\\
D(\omega)&=(k+\kappa-m\omega^2)(\kappa-\mu\omega^2)-\kappa^2.
\end{align*}
当 $\omega_a^2=\kappa/\mu$ 时，$D=-\kappa^2\ne0$，因此
\[
X_1=0,\qquad X_2=-F_0/\kappa.
\]
附加弹簧力恰好抵消驱动，主质量的受迫分量为零，这是反共振，不是整个系统没有运动。取调谐条件 $\kappa/\mu=k/m$，即可在原主振子频率消除理想特解的主质量位移；两个新本征频率则由 $D=0$ 给出。

该结论针对指定频率的无阻尼特解。任意初值会叠加不衰减的自由振动；阻尼、失谐和有限行程也会改变零点附近响应。因此不能据此宣称任意扰动或任意频率都被消除。
\end{example}

\begin{figure}[htbp]
\centering
\begin{tikzpicture}[>=Stealth]
\draw[thick](-.2,-.5)--(-.2,.7);
\draw[decorate,decoration={coil,aspect=.5,segment length=5pt,amplitude=4pt},thick](-.2,.1)--(1.6,.1)node[midway,above=.18cm]{$k$};
\draw[fill=blue!10](1.6,-.35)rectangle(2.7,.55);\node at(2.15,.1){$m$};
\draw[decorate,decoration={coil,aspect=.5,segment length=5pt,amplitude=4pt},thick](2.7,.1)--(4.7,.1)node[midway,above=.18cm]{$\kappa$};
\draw[fill=green!10](4.7,-.35)rectangle(5.6,.55);\node at(5.15,.1){$\mu$};
\draw[->,red!70!black,thick](2.15,.9)--(3.6,.9)node[right]{$F_0\cos\omega t$};
\draw[->](1.8,-.7)--(2.8,-.7)node[right]{$x_1$};\draw[->](4.7,-.7)--(5.7,-.7)node[right]{$x_2$};
\end{tikzpicture}
\caption{理想调谐吸振器。附加质量的运动产生补偿主驱动的弹簧力；主质量反共振频率与两自由度系统的本征频率不同。}
\label{fig:osc-absorber}
\end{figure}


\section{从有限链到长波连续极限}
许多质点排成一列时，一个质点的运动可沿弹簧传给邻居。先看一串有限质点允许哪些波，再看波长远大于间距时如何近似为连续介质。

\subsection{周期边界、色散与模态计数}
把链的两端接成环，绕一圈后的位移必须接得上。这个要求使波数只能取离散值；固定两端时则要满足另一套条件。

把 $N\ge3$ 个质量为 $m$ 的质点放在一维周期单元中，相邻间隔为 $a$，纵向位移为 $u_n$，周期边界 $u_{n+N}=u_n$。这里只把周期边界当作理想无限链的有限近似，不引入实际圆环弯曲的几何力。相邻弹簧劲度均为 $k$，则
\begin{equation}
m\ddot u_n=k(u_{n+1}+u_{n-1}-2u_n).
\label{eq:chain-eq}
\end{equation}
试解 $u_n=Ae^{i(pna-\omega t)}$，周期边界要求 $p=2\pi j/(Na)$，$j$ 为整数。代入并消去公共因子，得到
\[
-m\omega^2=k(e^{ipa}+e^{-ipa}-2)=-4k\sin^2(pa/2),
\]
所以
\begin{equation}
\omega(p)=2\sqrt{k/m}\,|\sin(pa/2)|.
\label{eq:dispersion-chain}
\end{equation}
$p$ 与 $p+2\pi/a$ 在所有格点给同一位形，因此取第一布里渊区 $-\pi<pa\le\pi$ 就足够；并非区外完全没有数学表示。$p=0$ 是整体平移零模。

对 $0<pa<\pi$，$\pm p$ 是同频的复共轭对，实基可选 $\cos pna$ 与 $\sin pna$。它们合起来只有两个实位形自由度，不能把共轭波数再各算两个。若 $N$ 偶数，$pa=\pi$ 的格点正弦恒零，只剩一个交替位移模；这加上零模和其余波数对，恰给 $N$ 个实坐标。

\begin{figure}[htbp]
\centering
\begin{tikzpicture}
\begin{axis}[width=.82\textwidth,height=5.3cm,axis lines=left,xlabel={$pa/\pi$},ylabel={$\omega/\sqrt{k/m}$},xmin=-1.05,xmax=1.05,ymin=0,ymax=3.4,legend style={at={(.5,.98)},anchor=north,draw=none,font=\small},samples=240]
\addplot[blue!65!black,very thick,domain=-1:1]{2*abs(sin(deg(pi*x/2)))};\addlegendentry{离散链}
\addplot[red!70!black,dashed,domain=-1:1]{pi*abs(x)};\addlegendentry{长波近似}
\addplot[only marks,mark=*,mark size=1.7pt,blue!65!black]coordinates{(-.75,1.847759)(-.5,1.414214)(-.25,.765367)(0,0)(.25,.765367)(.5,1.414214)(.75,1.847759)(1,2)};
\end{axis}
\end{tikzpicture}
\caption{第一布里渊区内的准确色散与 $\omega\simeq a\sqrt{k/m}|p|$。在当前横轴下，长波直线斜率为 $\pi$；圆点为 $N=8$ 周期链的允许波数，$-\pi$ 与 $\pi$ 是同一格点模。}
\label{fig:dispersion}
\end{figure}

\subsection{相速度、群速度与边界}
波峰移动的速度和一群相近波叠成的包络速度可以不同。看清题目追踪的是波峰还是包络，再选相速度或群速度。

对正波数，定义
\[
v_{\rm ph}=\omega/p,\qquad
v_g=\frac{\dd\omega}{\dd p}=a\sqrt{k/m}\cos(pa/2).
\]
群速度描述在无限或足够长的无耗散链中，窄波数包络的传播；有限环会有回返，单个驻波也没有可以追踪的行进包络。长波 $|pa|\ll1$ 时两速度均趋 $c_s=a\sqrt{k/m}$；区边界相邻格点反相、$v_g=0$，但 $\omega=2\sqrt{k/m}$ 并不为零。

\begin{example}[固定端与周期端不可混用]
另取 $N$ 个活动质点，左右固定端 $u_0=u_{N+1}=0$。令振型为 $u_n=A\sin(pna)$，边界给
\[
p_j a=\frac{j\pi}{N+1},\quad
\omega_j=2\sqrt{k/m}\sin\frac{j\pi}{2(N+1)},\quad j=1,\ldots,N.
\]
全部频率正，无整体平移模。最小频率在大 $N$ 时约为 $\pi c_s/[(N+1)a]$。同样的局部弹簧方程，不同边界改变允许波数和模态形状；不能只记色散式而丢掉边界条件。
\end{example}

\subsection{长波展开与连续能量}
波长远大于质点间距时，相邻两点的位移差可用斜率乘间距近似。把每一小段的质量和弹性加起来，就得到连续模型。

用平滑场 $u(x,t)$ 插值格点，$x=na$，泰勒展开使奇数空间导数相消：
\[
u(x+a)+u(x-a)-2u(x)
=a^2u_{xx}+\frac{a^4}{12}u_{xxxx}+\cdots.
\]
保留最低阶得
\begin{equation}
u_{tt}=c_s^2u_{xx},\qquad c_s^2=ka^2/m.
\label{eq:wave-equation}
\end{equation}
忽略项相对大小约为 $(pa)^2/12$。这个近似只用于长波，不能把截断高阶方程任意外推到区边界。

连续质量密度 $\rho=m/a$，纵向弹性系数 $\mathcal C=ka$，离散能量相应趋于
\[
T\simeq\frac12\int\rho u_t^2\,\dd x,\qquad
V\simeq\frac12\int\mathcal C u_x^2\,\dd x,\qquad
c_s^2=\mathcal C/\rho.
\]
变分给出波动方程，同时边界项为 $\mathcal C u_x\delta u$：固定端取 $\delta u=0$，自由端取 $\mathcal C u_x=0$。本例是纵向弹性链，不能把 $k$ 直接当作弦张力。若真的取 $a\to0$ 且保持有限连续介质，须按 $m=\rho a,k=\mathcal C/a$ 缩放；固定 $m,k$ 而只令 $a\to0$ 会使声速趋零。

\begin{modernbox}[title={进一步探索：声子与有限模态}]
每个正频简正坐标都是谐振子。若进一步量子化，其能级为 $E_n=\hbar\omega(n+1/2)$，振动量子称为声子；这里需要量子谐振子基础。零频整体平移不能照搬这一束缚振子能级式。色散给频率与波数关系，但有限温度热输运还需模占据、边界及散射，不是由本章无耗散链单独决定。
\end{modernbox}


\section{进一步探索：线性理论之外}
位移增大或外部机构开始周期地改变刚度时，线性定常模型会漏掉一些现象。下面用几个简单模型看看最早出现的修正。

本节不把弱非线性结果混入前文的精确线性叠加。它回答两个问题：被略去的项首先改变什么？周期改变参数为何不同于施加外力？

\subsection{幅频移动与组合频率}
理想弹簧振得大些和小些，周期相同。加入更高次恢复力后，周期可能随振幅变化，几个模一起运动时也可能产生新的组合频率。

\begin{example}[对称四次势能的首阶频移]
取
\[
V(q)=\tfrac12m\omega_0^2q^2+\tfrac14\beta q^4,\qquad
\ddot q+\omega_0^2q+\frac\beta m q^3=0.
\]
$\beta$ 在势能和方程中的含义保持一致。令转折振幅为 $A$，小参数为 $|\beta|A^2/(m\omega_0^2)\ll1$。最低阶取 $q=A\cos\omega t$，用 $\cos^3z=(3\cos z+\cos3z)/4$，基频项系数为
\[
A\left[-\omega^2+\omega_0^2+\frac{3\beta A^2}{4m}\right].
\]
消除这一阶的基频残差，得到
\begin{equation}
\omega=\omega_0+\frac{3\beta A^2}{8m\omega_0}
+O\left(\omega_0\left[\frac{\beta A^2}{m\omega_0^2}\right]^2\right).
\label{eq:osc-duffing-shift}
\end{equation}
还存在小三次谐波，不能说解仍是精确正弦。$\beta>0$ 为硬化，$\beta<0$ 为局部软化；负四次项在大位移不提供全局束缚，本近似只在原点附近小能量势阱内使用。若势能另有三次项，其二阶效应也可贡献频移，不能照搬这个对称势公式。
\end{example}

\begin{example}[从明确的耦合势产生和频、差频]
取三个等效质量 $m$ 的坐标，势能含 $\epsilon c q_1q_2q_3$，$c$ 的单位使该项为能量；$\epsilon$ 为无量纲小参数。在有限时间的最低阶，令
\[
q_1^{(0)}=A\cos\omega_1t,\qquad q_2^{(0)}=B\cos\omega_2t,
\]
则第三模的一阶受迫式含
\[
\ddot q_3+\omega_3^2q_3
=-\frac{\epsilon c AB}{2m}
[\cos(\omega_1+\omega_2)t+\cos(\omega_1-\omega_2)t].
\]
非共振时，一个频率 $\nu$ 的力 $f\cos\nu t$ 给特解 $f\cos\nu t/(\omega_3^2-\nu^2)$；若 $\nu=\omega_3$，首阶特解为 $ft\sin\omega_3t/(2\omega_3)$。这表示普通展开将出现长期失效，须跟踪各模振幅反馈，不能把线性增长当成封闭非线性系统的无限能量来源。此三次势也只作为局部模型，远处可需高阶稳定项。
\end{example}

\subsection{主参数共振的慢变振幅推导}
荡秋千时按时改变身体姿势，等于周期地改变系统参数。这里先找调制节拍与自由振动的关系，再看输入能量能否抵消阻尼。

\begin{example}[刚度调制如何克服阻尼]
取
\begin{equation}
\ddot q+2\gamma\dot q+\omega_0^2[1+\epsilon\cos\Omega t]q=0,
\quad 0<\epsilon\ll1,\quad
\delta=\Omega-2\omega_0=O(\epsilon\omega_0),\quad
\gamma=O(\epsilon\omega_0).
\label{eq:mathieu}
\end{equation}
当 $\gamma=0$ 时，这是马蒂厄方程的一种参数化。取试探振荡的载波角频率 $\omega_{\rm c}=\Omega/2$，
\[
q=a(t)\cos\omega_{\rm c} t+b(t)\sin\omega_{\rm c} t,
\]
$a,b$ 在一个振动周期内只作 $O(\epsilon)$ 的相对尺度变化。一阶计算中，$\ddot q+\omega_{\rm c}^2q$ 的慢项是 $2\omega_{\rm c}(\dot b\cos\omega_{\rm c} t-\dot a\sin\omega_{\rm c} t)$，含 $\ddot a,\ddot b$ 的项为二阶。阻尼只需乘 $\dot q$ 的快项，因为 $\gamma\dot a,\gamma\dot b$ 也是二阶。调制项用
\[
\cos2\omega_{\rm c} t\,\cos\omega_{\rm c} t=\tfrac12(\cos3\omega_{\rm c} t+\cos\omega_{\rm c} t),\quad
\cos2\omega_{\rm c} t\,\sin\omega_{\rm c} t=\tfrac12(\sin3\omega_{\rm c} t-\sin\omega_{\rm c} t).
\]
在一周期投影到 $\cos\omega_{\rm c} t,\sin\omega_{\rm c} t$ 时，$3\omega_{\rm c}$ 项与二者正交。两条投影方程为
\begin{align*}
2\omega_{\rm c}\dot b+2\gamma\omega_{\rm c} b+(\omega_0^2-\omega_{\rm c}^2+\epsilon\omega_0^2/2)a&=0,\\
-2\omega_{\rm c}\dot a-2\gamma\omega_{\rm c} a+(\omega_0^2-\omega_{\rm c}^2-\epsilon\omega_0^2/2)b&=0.
\end{align*}
因为 $\omega_{\rm c}=\omega_0+\delta/2$，有 $(\omega_0^2-\omega_{\rm c}^2)/(2\omega_{\rm c})=-\delta/2+O(\epsilon^2\omega_0)$。调制项的系数 $\epsilon\omega_0^2/(4\omega_{\rm c})$ 可在同一精度下取 $\epsilon\omega_0/4$。记 $h=\epsilon\omega_0/4$，化成
\[
\frac{\dd}{\dd t}\begin{pmatrix}a\\b\end{pmatrix}
=\begin{pmatrix}-\gamma&-\delta/2-h\\\delta/2-h&-\gamma\end{pmatrix}
\begin{pmatrix}a\\b\end{pmatrix}.
\]
设慢振幅与 $e^{st}$ 成正比。系数矩阵减去 $sI$ 后的行列式为 $(s+\gamma)^2-(h^2-\delta^2/4)$，令它为零得两个慢指数
\begin{equation}
s_\pm=-\gamma\pm\sqrt{h^2-\delta^2/4}.
\label{eq:osc-param-growth}
\end{equation}
因此无阻尼主不稳定区在首阶为 $|\delta|<\epsilon\omega_0/2$。有阻尼时需 $h^2>\gamma^2+\delta^2/4$；正调谐 $\delta=0$ 的阈值为 $\epsilon>4\gamma/\omega_0$。边界附近高阶修正会影响判断，这不是任意调制深度下的精确稳定图。

调制中的瞬时机械能取
\[
E(t)=\tfrac m2\dot q^2+\tfrac{m\omega_0^2}2[1+\epsilon\cos\Omega t]q^2,
\quad
\dot E=-2m\gamma\dot q^2-\frac{m\omega_0^2\epsilon\Omega}{2}\sin\Omega t\,q^2.
\]
第二项是改变刚度的机构输入/取出的功率。$q\equiv0$ 仍为精确解；“不稳定”表示存在增长的微扰，不表示零初态自行产生非零运动，也不是每个特选初值都会增长。
\end{example}

\begin{figure}[htbp]
\centering
\begin{tikzpicture}
\begin{axis}[width=.78\textwidth,height=5.2cm,axis lines=left,xlabel={$\Omega/\omega_0$},ylabel={$\epsilon$},xmin=1.8,xmax=2.2,ymin=0,ymax=.32,xtick={1.8,1.9,2,2.1,2.2},ytick={0,.1,.2,.3}]
\fill[red!12](axis cs:2,0)--(axis cs:1.85,.3)--(axis cs:2.15,.3)--cycle;
\addplot[red!70!black,thick,domain=0:.3]({2-x/2},x);
\addplot[red!70!black,thick,domain=0:.3]({2+x/2},x);
\node at(axis cs:2,.23){首阶不稳定区};
\end{axis}
\end{tikzpicture}
\caption{无阻尼、弱调制的主舌形区首阶边界 $\Omega/\omega_0=2\pm\epsilon/2$。舌尖在零调制处收拢，宽度随调制增大；本图不画高阶窄舌，也不冒充完整精确的马蒂厄稳定图。}
\label{fig:parametric}
\end{figure}

周期系数问题可用弗洛凯（Floquet）方法。取状态 $(q,\dot q)^{\mathsf T}$，分别以 $(1,0)^{\mathsf T}$、$(0,1)^{\mathsf T}$ 为初值，推进一个调制周期 $T_{\rm m}=2\pi/\Omega$。将所得两列并排，得到单周期矩阵 $U$。系数每周期复原，所以第 $n$ 个周期末的状态由 $U^n$ 给出，$U$ 的本征值称为弗洛凯乘子。若某乘子模大于1，相应模随周期数指数增长。无阻尼时，两个解的行列式 $W=q_1\dot q_2-\dot q_1q_2$ 满足 $\dot W=q_1\ddot q_2-\ddot q_1q_2=0$，故 $\det U=1$。若模为1的重根只有一个本征方向，$U^n$ 中还会出现正比于 $n$ 的项。NIST的马蒂厄稳定图使用标准参数 $(a,q)$，若要比较，先作 $z=\Omega t/2$，得到
\[
\frac{\dd^2q_{\rm osc}}{\dd z^2}+[a-2q_{\rm M}\cos2z]q_{\rm osc}=0,\quad
a=\frac{4\omega_0^2}{\Omega^2},\quad q_{\rm M}=-\frac{2\epsilon\omega_0^2}{\Omega^2}.
\]
这里 $q_{\rm osc}$ 是位移，$q_{\rm M}$ 是马蒂厄参数，二者不同；这一映射限于无阻尼式。

\subsection{自激、分岔与混沌：不由二次型包办}
外力一直按节拍推、系统自己补充振动能量、以及初值很接近却逐渐分离，是不同的现象。下面用不同的判据来区分。

范德波尔模型可写成无量纲形式
\[
x''-\mu(1-x^2)x'+x=0,\quad \mu>0,
\]
撇号表示对无量纲时间求导，$x$ 也无量纲。能量 $e=(x'^2+x^2)/2$ 满足 $e'=\mu(1-x^2)x'^2$，小位移区可注入能量，大位移区可耗散。弱 $\mu$ 下令 $x\simeq A\cos(\tau+\phi)$，用
\[
\overline{x'^2}=A^2/2,\quad \overline{x^2x'^2}=A^4/8,
\]
得到 $A'=\mu A(1-A^2/4)/2$。$A=2$ 是这一平均方程的稳定振幅，不是任意 $\mu$ 下精确正弦轨道；相位还与初始条件有关。原点精确初态仍停在原点，极限环也不是保守系统的能量等高线。

分岔研究的是改变参数后解的结构如何改变。例如 $V(q)=rq^2/2+bq^4/4$、$b>0$，平衡为 $q=0$ 及 $r<0$ 时的 $q=\pm\sqrt{-r/b}$；$r=0$ 处二次曲率消失，必须保留四次项。混沌则需更丰富的动力学，不能从单个频移、拍或短时轨道分离就断定；二维光滑自治范德波尔系统的平面极限环也不等于混沌。进一步研究需非线性方程、相空间与稳定性理论，超出本章简正模求解。

\section{计算与模拟：用残差和收敛检验模型}
把求出的振型代回原方程，检查两边差多少。再分别缩小初始振幅和积分步长，区分小角近似误差与计算误差。

伴随脚本 \texttt{review/verify\_ch04.py} 使用标准Python库，结果保存在
\begin{center}\texttt{ch04\_verification\_results.json}\end{center}
主要检查包括：
\begin{enumerate}[leftmargin=1.8em]
\item 对每个解析振型计算 $Ka_r-\lambda_rMa_r$，并核验 $G^{\mathsf T}MG$ 和 $G^{\mathsf T}KG$；检查初值重建及两模能量之和。
\item 独立积分等质量双摆精确方程，与本章小角解比较。在固定时间区间逐次减小初角，同时加密步长，区分非线性误差和积分误差；不能要求线性近似任意长期保持相位。
\item 对有限链检查允许波数、零模与固定端谱；数值核验色散的长波相对误差按 $(pa)^2$ 缩小。
\item 用一个调制周期的两个独立解求Floquet乘子，在舌形区内外各取参数点，并与首阶慢指数比较；边界附近须提高精度，不只凭图的颜色判断。
\end{enumerate}
所有数值参数为教学设定，不代表实验材料参数或工程性能。

\section*{关键概念与常见误区}
\begin{itemize}[leftmargin=1.8em]
\item 先找平衡点，再看微小偏离后的恢复力。势能的二次项若消失，就检查更高次项或直接研究运动。
\item 简正坐标使本章的线性保守模型拆成独立振动。加上磁力、陀螺项、非比例阻尼或周期驱动后，要重新检查各模是否仍独立。
\item 在 $Ka=\lambda Ma$ 中，正根给频率 $\sqrt{\lambda}$，零根允许匀速漂移，负根给指数增长或衰减。先按质量归一化，再用初值确定各模的振幅。
\item 问“哪里共振”之前，先确认是在测位移还是吸收功率。若外力对某个振型没有直接投影，还要看其他模能否通过耦合带动它。
\item 接成环的链和固定两端的链允许不同的波数。用复数写波时，最后应数出独立的实振型。
\item 两种频率接近的振动相加可以产生拍。周期改变刚度则可向振子供能；若恢复力含高次项，还可能出现和频或差频。
\end{itemize}

\section*{本章小结}
\addcontentsline{toc}{section}{本章小结}
第三章的拉格朗日量在平衡附近展开后，以动能和势能二次型组织主要运动；正定惯性使问题归结为实对称广义本征值，质量正交基给出初值投影与模态能量。阻尼、驱动与边界决定实际响应，离散链的长波极限连接连续波动。第\ref{ch:rigid}章把内部距离固定，转而研究刚体的整体平动和转动；刚体转动的惯量张量与本章质量二次型相呼应，但有限转角与姿态几何仍需专门处理。


\section*{分层习题}
\addcontentsline{toc}{section}{分层习题}
先找静止平衡，写出相对于平衡的位移。多自由度题可先找振型，再用初值确定各振型的振幅；有驱动时还要看力在每个振型上的投影。

\begin{bookexercises}
\item \exlevel{基础} \textbf{稳定的定义。} 对 $V=aq^4,-aq^4,0$，$a>0$，分别判断原点平衡的稳定性。说明为什么“没有指数增长”不构成稳定性证明。
\item \exlevel{基础} \textbf{位形相关惯性。} 对 $L=m(1+q^2/l^2)\dot q^2/2-kq^2/2$，写精确方程，取 $q/l=O(\epsilon)$ 判断惯性非线性首次出现的阶次，并求小振动频率。
\item \exlevel{基础} \textbf{阻尼与返回。} 临界阻尼下取 $q(0)=q_i>0,\dot q(0)=v_i$，求 $q(t)$。什么初速度会使其越过原点？计算瞬时能量导数。
\item \exlevel{基础} \textbf{不同的共振峰。} 从受迫响应推导位移幅值峰与平均功率峰，核验半功率带宽 $2\gamma$；解释只给相位正切为何不能确定高频响应的相位。
\item \exlevel{基础} \textbf{交叉项与简并。} 对交叉项例题独立代入两个振型。若 $\kappa=0$，能否选择两个不正交振型？为什么仍可选择质量标准正交基？
\item \exlevel{综合} \textbf{投影与能量。} 对 $M=m\operatorname{diag}(1,2),K=k\bigl(\begin{smallmatrix}3&-2\\-2&4\end{smallmatrix}\bigr)$，取 $x(0)=(0,a),\dot x(0)=0$，求两个质量归一化模态初值、物理位移和总能量。
\item \exlevel{综合} \textbf{坐标标度。} 对上题改用 $z_1=x_1,z_2=2x_2$，求新的 $M_z,K_z$，核验频率不变。
\item \exlevel{综合} \textbf{Rayleigh上界。} 对等质量等长双摆取试探角度比 $(1,1)$，估计最低频率平方的上界，并与精确值 $(2-\sqrt2)g/l$ 比较。给出相应的频率上界。
\item \exlevel{综合} \textbf{双摆初值。} 等质量等长双摆从 $(\theta_1,\theta_2)=(\Theta,0)$ 静止释放，求两角的线性解，并验证初始加速度。
\item \exlevel{综合} \textbf{拍与激发选择。} 耦合摆的两广义力矩若为 $(\tau,\tau)\cos\omega t$，激发哪个模？改为 $(\tau,-\tau)$ 呢？再令外力矩为零，将第一摆从非零小角度静止释放，第二摆静止于平衡位置。由两模叠加推导首次包络转移时间。
\item \exlevel{综合} \textbf{三原子与质量扰动。} 对三质点模型写完整质心运动和两个伸缩模。若中间质量 $\mu$ 增大，哪一频率不变、哪一频率降低？能否仅由频率断言出现红外吸收峰？
\item \exlevel{综合} \textbf{非比例阻尼。} 核验正文单阻尼器的 $C_N$。若改用 $C=\alpha M+\beta K$，求各模阻尼比，并讨论模态外力为零的情况。
\item \exlevel{综合} \textbf{吸振器。} 取 $m=\mu,k=\kappa$，求两个本征频率和主质量反共振频率，说明自由响应未衰减时的限制。
\item \exlevel{综合} \textbf{模态计数。} 列出 $N=8$ 周期链的允许波数与全部 $\omega^2$，标明零模和重数；再写 $N=3$ 固定端链的频率平方。
\item \exlevel{综合} \textbf{连续极限。} 从离散能量推出 $\rho=m/a,\mathcal C=ka$，解释保持声速不变的缩放。计算 $pa\ll1$ 时离散频率相对长波近似的首项误差。
\item \exlevel{拓展} \textbf{弱非线性。} 推导对称四次势的首阶频移，说明负 $\beta$ 的大振幅限制；给出三模耦合的一项和频特解及其共振极限。
\item \exlevel{拓展} \textbf{参数共振。} 取 $\omega_0=1,\epsilon=0.1$，用首阶式求 $\delta=0$、$0.1$ 时无阻尼慢指数；若 $\gamma=0.03$，正调谐还会增长吗？注明近似精度。
\item \exlevel{拓展} \textbf{数值复核。} 运行伴随脚本，检查本征残差、双摆振幅缩小的误差、步长误差及Floquet乘子。哪些检验针对物理近似，哪些针对数值算法？为什么短时能量准确不足以证明长期相位准确？
\end{bookexercises}

\section*{习题解答}
先用同一题号找到解答。核对时，从第一条方程开始比较自己的做法，再检查最后的符号、单位和适用条件。
\begin{bookexercises}
\item 对 $V=aq^4$，若初始能量小于 $\min(aR^4,mU^2/2)$，轨道始终满足 $|q|<R,|\dot q|<U$，故原点在相空间稳定。对 $V=-aq^4$，取任意小的 $q(0)>0$ 且 $\dot q(0)=0$，此后 $\ddot q=4aq^3/m$ 为正并有正下界，轨道终会离开原点邻域。对 $V=0$，任意非零初速度给线性漂移 $q=q(0)+\dot q(0)t$。后两种离开机制均无需线性指数模。
\item 精确式为 $m(1+q^2/l^2)\ddot q+mq\dot q^2/l^2+kq=0$。固定时间尺度上 $q,\dot q,\ddot q=O(\epsilon)$ 时，惯性修正是三阶；$\omega_0=\sqrt{k/m}$。

$\partial L/\partial\dot q=m(1+q^2/l^2)\dot q$，求时间导数产生 $2mq\dot q^2/l^2$；减去 $\partial L/\partial q=mq\dot q^2/l^2-kq$ 后保留一个这样的三阶项。
\item $q=[q_i+(v_i+\omega_0q_i)t]e^{-\omega_0t}$。当 $v_i<-\omega_0q_i$ 时有正时间越零；$\dot E=-2m\omega_0\dot q^2$。无振荡复根不等于任何初态都不越零。

重特征根 $-\omega_0$ 给 $q=(A+Bt)e^{-\omega_0t}$，由 $q(0)=q_i$、$\dot q(0)=B-\omega_0A$ 定系数。满足越零条件时，越零时刻为 $-q_i/(v_i+\omega_0q_i)$。用 $E=m(\dot q^2+\omega_0^2q^2)/2$ 求导并代入运动方程即可得上式功率。
\item 设 $\gamma>0$。令 $s=\omega^2$，响应分母平方为 $D=(\omega_0^2-s)^2+4\gamma^2s$，位移幅值正比于 $D^{-1/2}$。令 $D'=0$ 得 $s=\omega_0^2-2\gamma^2$，只有它为正时才有非零频率的位移峰，否则最大位移幅值在零频。平均功率正比于 $s/D$，其导数的分子为 $D-sD'=\omega_0^4-s^2$，故功率峰在 $\omega_0$。半功率要求 $D=8\gamma^2\omega^2$，即 $|\omega_0^2-\omega^2|=2\gamma\omega$。两个正根为 $\sqrt{\omega_0^2+\gamma^2}\mp\gamma$，带宽为 $2\gamma$。相位应取复分母 $\omega_0^2-\omega^2+2i\gamma\omega$ 的辐角，高频时从第二象限趋于 $\pi$。
\item $(1,1)$ 给 $\omega_0^2-\kappa$，$(1,-1)$ 给 $\omega_0^2+\kappa$。简并时任何独立线性组合都是本征基，但用非正交基不能直接写对角归一化能量；可在质量内积下作Gram--Schmidt正交化。

直接相乘时，第一行作用于两向量分别为 $\omega_0^2-\kappa$、$\omega_0^2+\kappa$，第二行分别与第一行同号、反号。$\kappa=0$ 时频率矩阵为标量矩阵，整个位形空间是同一特征子空间。
\item $\xi_1(0)=2a\sqrt{m/3},\xi_2(0)=-a\sqrt{2m/3}$，故
\[
x=\frac{2a}{3}(1,1)^{\mathsf T}\cos\omega_1t
-\frac a3(2,-1)^{\mathsf T}\cos\omega_2t.
\]
$E_1=2ka^2/3,E_2=4ka^2/3$，和为初始势能 $2ka^2$。

取 $g_1=(1,1)^{\mathsf T}/\sqrt{3m}$、$g_2=(2,-1)^{\mathsf T}/\sqrt{6m}$，直接计算 $g_r^{\mathsf T}Mg_s=\delta_{rs}$。由 $\xi_r(0)=g_r^{\mathsf T}Mx(0)$、$\dot\xi_r(0)=0$ 得上述初值，再用 $E_r=\omega_r^2\xi_r(0)^2/2$ 得各模能量。
\item $x=Jz,J=\operatorname{diag}(1,1/2)$，所以
\[
M_z=m\operatorname{diag}(1,1/2),\quad
K_z=k\begin{pmatrix}3&-1\\-1&1\end{pmatrix}.
\]
久期根仍为 $k/m,4k/m$。

令 $s=m\omega^2/k$，新行列式除以 $k^2$ 后为 $(3-s)(1-s/2)-1=(s-1)(s-4)/2$。
\item $R(1,1)=3g/(5l)=0.6g/l$，略高于 $(2-\sqrt2)g/l\simeq0.5858g/l$。这是最低频率平方的上界。

对 $a=(1,1)^{\mathsf T}$，有 $a^{\mathsf T}Ka=3mgl$、$a^{\mathsf T}Ma=5ml^2$。频率上界为 $\sqrt{3g/(5l)}$，精确最低频率为 $\sqrt{(2-\sqrt2)g/l}$。
\item 两个未归一化振型的系数均为 $\Theta/2$，故
\[
\theta_1=\tfrac\Theta2(\cos\omega_{\rm low}t+\cos\omega_{\rm high}t),\quad
\theta_2=\tfrac\Theta{\sqrt2}(\cos\omega_{\rm low}t-\cos\omega_{\rm high}t).
\]
初加速度为 $\ddot\theta_1(0)=-2g\Theta/l,\ddot\theta_2(0)=2g\Theta/l$，与 $-M^{-1}K\theta(0)$ 一致。

两振型可选 $(1,\pm\sqrt2)^{\mathsf T}$，令系数为 $A,B$，由 $A+B=\Theta$、$\sqrt2(A-B)=0$ 得到初值。两式微分两次后，$\ddot\theta_1(0)$ 取频率平方的和，$\ddot\theta_2(0)$ 取频率平方的差。
\item 同向力矩只激发同相模，反向力矩只激发反相模；在完全对称模型中由 $g_r^{\mathsf T}f$ 直接判断。首次包络转移时间 $\pi/(\omega_--\omega_+)$，裸摆能量几乎全转移需弱耦合解释。

设第一摆从 $\Theta$ 静止释放，第二摆初始静止于平衡，两角分别是 $\frac\Theta2(\cos\omega_+t\pm\cos\omega_-t)$。和差化积后，慢包络分别为 $\Theta|\cos[(\omega_--\omega_+)t/2]|$ 和 $\Theta|\sin[(\omega_--\omega_+)t/2]|$。第一包络首次为零时第二包络达最大。
\item $X_{\rm cm}=X_0+V_0t$，内部模分别沿 $(1,0,-1)$ 和 $(1,-2m/\mu,1)$。$\omega_1^2=k/m$ 不变，$\omega_2^2=k/m+2k/\mu$ 降低。是否红外活性还需偶极矩变化，频率本身不够。

完整位移为
\begin{align*}
x&=(X_0+V_0t)(1,1,1)^{\mathsf T}
+\eta_1(t)(1,0,-1)^{\mathsf T}\\
&\quad+\eta_2(t)(1,-2m/\mu,1)^{\mathsf T},\\
\eta_j(t)&=A_j\cos\omega_jt+B_j\sin\omega_jt.
\end{align*}
在 $M=\operatorname{diag}(m,\mu,m)$ 下，两个内部向量的质量加权分量和均为零。将它们分别代入链的 $K a=\omega^2Ma$，端点恢复力系数依次为 $k$ 与 $k(1+2m/\mu)$，除以端点质量 $m$ 就得到两个频率平方。
\item 由于 $C=\operatorname{diag}(c,0)$，故 $(C_N)_{rs}=c(g_r)_1(g_s)_1$。代入前题归一化振型得
\[
C_N=\frac{c}{3m}\begin{pmatrix}1&\sqrt2\\\sqrt2&2\end{pmatrix}.
\]
若 $C=\alpha M+\beta K$，则 $G^{\mathsf T}CG=\alpha I+\beta\Lambda$，与标准模态方程的 $2\zeta_r\omega_r\dot\xi_r$ 对照，得 $\zeta_r=(\alpha+\beta\omega_r^2)/(2\omega_r)$（$\omega_r>0$）。比例阻尼时，若模态力为零且该模初值全零，则一直静止。有非对角阻尼时，即使 $g_r^{\mathsf T}F=0$，仍可能由其他模的速度激发。
\item $\omega_{1,2}^2=(3\mp\sqrt5)k/(2m)$，$\omega_a^2=k/m$ 位于二者之间。无阻尼任意初态含额外自由解，不能保证主质量总位移为零。

令 $s=m\omega^2/k$，齐次方程的行列式除以 $k^2$ 为 $(2-s)(1-s)-1=s^2-3s+1$。在 $s=1$，受迫方程第二行给 $-kX_1=0$，第一行给 $-kX_2=F_0$，这说明主质量虽无受迫位移，附加质量仍在振动。
\item 周期链 $pa=0,\pm\pi/4,\pm\pi/2,\pm3\pi/4,\pi$。以 $k/m$ 为单位的平方频率依次为 $0,2-\sqrt2$（二重）、$2$（二重）、$2+\sqrt2$（二重）、$4$。三质点固定端为 $(2-\sqrt2,2,2+\sqrt2)k/m$。

周期端条件给 $e^{iNpa}=1$，再代入 $\omega^2=2(k/m)(1-\cos pa)$。共轭的 $\pm p$ 对应两个实独立模，$p=0,\pi/a$ 各只计一个。固定端条件给 $pa=j\pi/(N+1)$，在 $N=3$ 时代入 $j=1,2,3$ 即得所列谱。
\item 把求和换成 $\int\dd x/a$，并用 $u_{n+1}-u_n\simeq au_x$。连续缩放 $m=\rho a,k=\mathcal C/a$。$\omega/(c_s|p|)=1-(pa)^2/24+O((pa)^4)$，故线性长波近似略高估离散频率。

两项能量分别成为 $\tfrac12\int (m/a)\dot u^2\dd x$ 和 $\tfrac12\int ka\,u_x^2\dd x$，故固定 $\rho,\mathcal C$ 即固定 $c_s=\sqrt{\mathcal C/\rho}$。离散频率比为 $\sin(|p|a/2)/(|p|a/2)$，展开即得上式。
\item 基频项平衡给式 \eqref{eq:osc-duffing-shift}，尚有三次谐波。和频 $\nu=\omega_1+\omega_2$ 的非共振特解为
\[q_{3,p}=-\frac{\epsilon cAB}{2m(\omega_3^2-\nu^2)}\cos\nu t.\]
精确共振需改用 $-\epsilon cABt\sin\omega_3t/(4m\omega_3)$，不能直接代入非共振式。

取 $q=A\cos\omega t$ 作零阶解。由 $\cos^3z=(3\cos z+\cos3z)/4$，基频系数为 $A(\omega_0^2-\omega^2)+3\beta A^3/(4m)$。令它为零后对平方根作首阶展开，得 $\omega-\omega_0=3\beta A^2/(8m\omega_0)$。$\beta<0$ 时在 $q^2=m\omega_0^2/|\beta|$ 处存在势垒，小振动展开要求振幅远小于该势垒位置。组合频率项来自 $\cos\omega_1t\cos\omega_2t$ 的和差分解；共振式可由 $(\dd^2/\dd t^2+\omega_3^2)(t\sin\omega_3t)=2\omega_3\cos\omega_3t$ 直接核验。
\item $h=0.025$；$\delta=0$ 时 $s=\pm0.025$；$\delta=0.1$ 时 $s=\pm i\sqrt{0.05^2-0.025^2}$，首阶无增长。正调谐加 $\gamma=0.03$ 后最大慢指数为 $-0.005$，首阶判为衰减。有限 $\epsilon$ 的精确边界须单周期矩阵进一步核验。

此处用 $s^2=h^2-(\delta/2)^2$，$h=\epsilon\omega_0/4$。所以失谐情形 $s\simeq\pm0.04330\,i$，慢包络无首阶指数增长。加阻尼的正调谐情形用 $s_{\rm max}\simeq-\gamma+h$；这些结果只保留 $\epsilon$、$\gamma/\omega_0$ 与 $\delta/\omega_0$ 的首阶效应。
\item 对每个模先算残差 $Kg_r-\omega_r^2Mg_r$。双摆在固定区间内将初角从 $0.08$ 减到 $0.04$、$0.02$，线性解的最大角度偏差依次为 $2.9980\times10^{-3}$、$3.7611\times10^{-4}$、$4.7054\times10^{-5}$，约每次缩小8倍，检验小角近似。独立减半 RK4 步长所得误差比约为 $15.86$，检验积分。把两个独立初态积分一个调制周期，终值作为单周期矩阵的两列。脚本的舌形区内、区外、加阻尼三组最大乘子模分别约为 $1.08168$、$1.00000$、$0.98441$。相位误差可以逐周期累积，因此还需直接比较振动相位或轨道状态。
\end{bookexercises}

\section*{进一步阅读与来源}
链接复核：2026-11-08。正文算例按本书坐标、边界和符号重新推导，数值均为教学设定。
\begin{bookreading}
\item Iain Stewart，MIT 8.09，\href{https://ocw.mit.edu/courses/8-09-classical-mechanics-iii-fall-2014/aa6e7a8b4e0b5eeef24bce949d4fdb27_MIT8_09F14_Chapter_3.pdf}{Chapter 3: Vibrations \& Oscillations}。可对照广义坐标展开、同时对角化与简正坐标；本书另明确零模及含速度项的适用边界。
\item MIT 8.09，\href{https://ocw.mit.edu/courses/8-09-classical-mechanics-iii-fall-2014/4ffa078142eca4a96a4bcb7a90bf0a4c_MIT8_09F14_pset4.pdf}{Problem Set 4}。双摆与弹簧支承杆的题目可作独立训练：先自行建立动能和势能，再画振型，而不是只代入已有矩阵。本章题目与解答按本书模型另行编写。
\item NIST Digital Library of Mathematical Functions，\href{https://dlmf.nist.gov/28.17}{\S28.17: Mathieu稳定区}。进一步学习精确特征曲线；比较前必须使用正文给出的参数映射，不把弱调制的一阶直线当作完整稳定边界。
\end{bookreading}

\chapter{刚体力学}\label{ch:rigid}

\begin{guidebox}
观察一只滚动的轮子，可以先追踪轮心，再看轮子相对轮心转了多少。刚体模型认为物体各点间距离固定，所以描述轮心和朝向就能找回所有点的位置。一般三维刚体用三个量记位置、三个量记朝向。

同样重的物体，质量分布不同，转起来也不同。本章先用惯量表示这种差别，再分析滚动、自由翻转和陀螺。写角动量方程时，先标明矩心，以及各分量是沿地面坐标轴还是沿随物体转动的轴量出的。
\end{guidebox}

\section{刚体运动学}

移动一本书时，可以把它整个平移，也可以一边移动一边转。先学会分别描述这两个动作，再把它们合成书上各点的运动。

\subsection{刚体与自由度}

刚体模型允许物体整体移动和转向，物体内部各点的距离保持不变。选一个参考点描述位置，再数还需几个量才能说明物体朝向哪里。

\begin{definition}[刚体]
若质点系中任意两个质点之间的距离在运动过程中保持不变，则称该质点系为\textbf{刚体}。对任意两个质点 \(i,j\)，有
\[
|\br_i-\br_j|^2=c_{ij},
\]
其中 \(c_{ij}\) 为常数。
\end{definition}

刚体是宏观物体的理想化模型。真实物体在受力时会形变，但当形变远小于整体运动尺度时，把物体看成刚体是很好的近似。若形变本身重要，则必须把它当作弹性体，额外引入振动自由度。

确定一个刚体的位形，只需要确定它上面三个不共线点的位置。三个不共线点构成一个刚性三角形；在选定初始形状与手性、并由连续刚体运动相连的位形分支上，其余物质点相对该三角形的位置也固定。仅凭三点距离的代数关系会有镜像歧义，反射不由连续的正常转动产生。三个点共有 \(9\) 个坐标，但三点之间的三个距离固定，因此独立坐标数为
\[
9-3=6.
\]
这给出刚体最重要的运动学结论：

\begin{keybox}
含不共线物质点的一般自由刚体在三维空间中有 \textbf{6 个自由度}：三个描述质心平动，三个描述绕质心的转动。
\end{keybox}

严格共线且不带额外可辨取向标记的理想点列，绕自身线转动不改变位形，只有5个位形自由度；单个质点只有3个。真实细杆截面通常可辨，仍有6个。

刚体的常见运动模式可以按约束分类：

\begin{center}
\begin{tabular}{lll}
\toprule
运动类型 & 自由度 & 约束特征 \\
\midrule
平动 & 3 & 刚体上所有点速度相同 \\
定轴转动 & 1 & 一条直线上的点固定 \\
平面平行运动 & 3 & 每点始终平行于某一固定平面 \\
定点转动 & 3 & 一个点固定，绕该点任意转动 \\
一般运动 & 6 & 无额外约束 \\
\bottomrule
\end{tabular}
\end{center}

平动时，刚体上每一点的速度和加速度都相同，因此可以用质心代表整个刚体。定轴转动时，只有绕过固定轴的转角一个独立变量。平面平行运动可以分解为质心的平面运动与绕质心的转动。定点转动没有平动，但仍需要三个角度描述转轴的方向和绕轴转过的角度。

\begin{figure}[htbp]
\centering
\begin{tikzpicture}[>=Stealth]
\draw[thick,fill=blue!7](1,0)ellipse(1.7 and .72);
\fill(0,0)circle(2pt)node[below left]{$O$};\fill(2,0)circle(2pt)node[below]{$P$};
\draw[->,thick](0,0)--(2,0)node[midway,below]{$\br_{P/O}$};
\draw[->,red!70!black,very thick](2,0)--(2,1.4)node[right]{$\bv_P=\bom\times\br_{P/O}$};
\node[blue!65!black]at(.3,1.15){$\odot\ \bom=\omega\hat z$};
\draw[->](-.8,-1)--(.3,-1)node[right]{$x$};\draw[->](-.8,-1)--(-.8,.15)node[above]{$y$};
\end{tikzpicture}
\caption{固定点转动的平面截面。角速度沿纸面外，故右侧物质点速度向上；任一点速度由同一个角速度和相对固定点的位矢决定。}
\label{fig:rigid-body}
\end{figure}

\subsection{有限转动与无限小转动}

拿起一本书，先绕竖直方向转再向前翻，和先翻再转，最终朝向通常不同。把转动写成矩阵便能明确记录先后次序。

刚体的转动与平动有一个本质区别：\textbf{绕不同轴的有限转动一般不满足交换律}。把一本书先绕竖直轴转 \(90^\circ\)，再绕水平轴转 \(90^\circ\)，与先绕水平轴再绕竖直轴，最终姿态不同。因此，有限转动不能用一个普通的矢量简单相加。

设刚体绕单位轴 \(\hat{\be}\) 转过有限角度 \(\theta\)，对应的转动矩阵记为 \(R(\hat{\be},\theta)\)。连续两次转动对应矩阵乘积
\[
R_2R_1,
\]
而绕不同轴时一般有
\[
R_2R_1\ne R_1R_2.
\]
绕同一固定轴的有限转动仍然可以交换。

但当转角趋于零时，情况发生变化。取转轴过坐标原点，设无限小转动矢量为
\[
\delta\bth=\delta\theta\,\hat{\be}.
\]
绕该轴的无限小转动矩阵可以写成
\[
R(\hat{\be},\delta\theta)
=\mathbf 1+[\delta\bth]_\times+O(\delta\theta^2),
\]
其中 \([\delta\bth]_\times\) 是由 \(\delta\bth\) 生成的反对称矩阵。作用于位矢 \(\br\)，得到
\[
\delta\br
=R\br-\br
=[\delta\bth]_\times\br+O(\delta\theta^2)
=\delta\bth\times\br+O(\delta\theta^2).
\]
忽略二阶小量，即得刚体上一点的位移
\[
\delta\br=\delta\bth\times\br.
\]
若连续做两次无限小转动 \(\delta\bth_1,\delta\bth_2\)，忽略二阶小量，
\[
\delta\br_{12}
=\delta\bth_2\times\br+\delta\bth_1\times\br
+O(\delta\theta^2),
\]
而
\[
\delta\br_{21}
=\delta\bth_1\times\br+\delta\bth_2\times\br
+O(\delta\theta^2).
\]
两者在最低阶相同。因此：

\begin{keybox}
有限转动不可交换；无限小转动在一阶近似下可以按矢量相加。角速度之所以是矢量，正是因为它是无限小转动的时间变化率。
\end{keybox}

\subsection{角速度与角加速度}

姿态告诉我们物体此刻朝哪里，角速度告诉我们姿态正在怎样变。后者既有大小也有方向，可用空间轴或随物体转的轴来报出分量。

定义刚体的\textbf{角速度}为
\[
\bom=\lim_{\Delta t\to0}\frac{\Delta\bth}{\Delta t},
\]
其中 \(\Delta\bth\) 是时间间隔 \(\Delta t\) 内的无限小转动矢量。它描述刚体瞬时转动的方向和快慢。对绕定点 \(O\) 转动的刚体，任意一点 \(P\) 的速度为
\begin{equation}
\bv_P=\bom\times\br_P.
\label{eq:rigid-velocity}
\end{equation}
若刚体同时有平动，取固连于刚体的物质基点 \(O\)，则
\begin{equation}
\bv_P=\bv_O+\bom\times(\br_P-\br_O).
\label{eq:rigid-velocity-general}
\end{equation}
注意 \(\bom\) 是同一个刚体上所有点共有的角速度，但不同点的线速度一般不同。

对角速度求时间导数，得到\textbf{角加速度}
\[
\boldsymbol\alpha=\dot{\bom}.
\]
对定点转动，将式 \eqref{eq:rigid-velocity} 求导：
\begin{equation}
\ba_P
=\boldsymbol\alpha\times\br_P
+\bom\times(\bom\times\br_P).
\label{eq:rigid-acceleration}
\end{equation}
第二项可以写成
\[
\bom\times(\bom\times\br)
=\bom(\bom\cdot\br)-\omega^2\br.
\]
这两项合起来为 $-\omega^2\br_\perp$，其中 $\br_\perp$ 为到瞬时轴的垂直位矢；不能把其中 $-\omega^2\br$ 单独说成指向轴的完整向心加速度。$\boldsymbol\alpha\times\br$ 来自角速度变化；定轴时它才沿圆周切向。

若基点 \(O\) 也在运动，对
\[
\bv_P=\bv_O+\bom\times\br_{P/O}
\]
求导：
\[
\ba_P=\ba_O+\boldsymbol\alpha\times\br_{P/O}
+\bom\times\dot\br_{P/O}.
\]
而
\[
\dot\br_{P/O}=\bv_P-\bv_O=\bom\times\br_{P/O},
\]
所以有更一般的形式
\begin{equation}
\ba_P
=\ba_O+\boldsymbol\alpha\times\br_{P/O}
+\bom\times(\bom\times\br_{P/O}).
\label{eq:rigid-acceleration-general}
\end{equation}

\begin{remark}
角速度 \(\bom\) 一般不能写成某个矢量 \(\bth\) 的时间导数。有限转动不是矢量，所以不存在一个全局的“角位置矢量”使其导数等于角速度。只有在定轴转动时，才可以选一个固定轴，把 \(\bom=\dot\theta\,\hat{\be}\) 写成标量角速度乘以固定方向。
\end{remark}

\subsection{欧拉角}

若要用三个角度说明一本书的朝向，必须约定每次绕哪根轴转。某些姿态下两根转轴会重合，这时角度的表示会失去唯一性。

描述刚体绕定点转动的三个自由度，最常用的是\textbf{欧拉角}。取固定坐标系 \(Oxyz\) 和固连在刚体上的坐标系 \(O x'y'z'\)。采用 \(z\)-\(x'\)-\(z'\) 约定，转动依次为：

\begin{enumerate}[leftmargin=1.8em]
  \item 绕固定 \(z\) 轴转 \(\phi\)，称为\textbf{进动角}；
  \item 绕第一次转动后的中间 $x$ 轴转 \(\theta\)，称为\textbf{章动角}；这条新的 \(x'\) 轴也称为\textbf{节线}，记其单位矢量为 \(\hat{\boldsymbol N}\)；
  \item 绕刚体自身的 \(z'\) 轴转 \(\psi\)，称为\textbf{自转角}。
\end{enumerate}

把体坐标变到固定坐标的转动矩阵记为 \(R\)。按上述顺序，
\begin{equation}
R(\phi,\theta,\psi)
=R_z(\phi)R_x(\theta)R_z(\psi),
\label{eq:euler-matrix}
\end{equation}
其中
\[
R_z(\alpha)=
\begin{pmatrix}
\cos\alpha & -\sin\alpha & 0\\
\sin\alpha & \cos\alpha & 0\\
0&0&1
\end{pmatrix},
\qquad
R_x(\alpha)=
\begin{pmatrix}
1&0&0\\
0&\cos\alpha&-\sin\alpha\\
0&\sin\alpha&\cos\alpha
\end{pmatrix}.
\]
欧拉角的取值范围通常取
\[
0\le\phi<2\pi,\qquad 0<\theta<\pi,\qquad 0\le\psi<2\pi.
\]
当 \(\sin\theta=0\) 时，第一次和第三次转动绕同一根轴，欧拉角描述出现退化，这就是\textbf{万向节死锁}。在数值姿态控制中，常用四元数避免这一奇点。

\begin{figure}[htbp]
\centering
\begin{tikzpicture}[>=Stealth]
\draw[->] (0,0)--(3,0) node[right]{$x$};
\draw[->] (0,0)--(0,2.5) node[above]{$y$};
\draw[->,thick,blue!65!black] (0,0)--(2.2,1.6) node[right]{$\hat{\boldsymbol N}$};
\draw[->] (0.9,0) arc (0:36:0.9);
\node at (1.05,0.35) {$\phi$};
\node at (1.3,-0.65) {从 $+z$ 向原点看};
\begin{scope}[shift={(6,0)}]
\draw[->] (0,0)--(0,2.7) node[above]{$z$};
\draw[->,thick,blue!65!black] (0,0)--(1.7,2.1) node[right]{$\be_3$};
\draw[->] (0,1.1) arc (90:51:1.1);
\node at (0.55,1.4) {$\theta$};
\node at (0.7,-0.65) {垂直节线的截面};
\end{scope}
\end{tikzpicture}
\caption{欧拉角的两个独立截面：$\phi$ 定位节线，$\theta$ 是固定 $z$ 轴与体轴 $\be_3$ 的夹角。最后的 $\psi$ 绕 $\be_3$ 转动另外两根体轴，不改变 $\be_3$；其完整约定由式 \eqref{eq:euler-matrix} 给出。}
\label{fig:euler-angles}
\end{figure}

把三个角速度贡献分别投影到体轴 \(x_1x_2x_3\) 上。\(\dot\phi\) 沿固定 \(z\) 轴，而固定 \(z\) 轴在体轴中的方向是 \(R^{\mathsf T}\hat{\boldsymbol z}\)，即
\[
(\sin\theta\sin\psi,\ \sin\theta\cos\psi,\ \cos\theta),
\]
故 \(\dot\phi\) 的贡献为
\[
\dot\phi(\sin\theta\sin\psi,\ \sin\theta\cos\psi,\ \cos\theta).
\]
\(\dot\theta\) 沿节线，节线在体轴中的方向是 \(R^{\mathsf T}\hat{\boldsymbol N}\)，即
\[
(\cos\psi,\ -\sin\psi,\ 0),
\]
故 \(\dot\theta\) 的贡献为
\[
\dot\theta(\cos\psi,\ -\sin\psi,\ 0).
\]
\(\dot\psi\) 沿体轴 \(x_3\)，贡献为 \((0,0,\dot\psi)\)。三者在无限小转动意义下相加，得到角速度分量
\begin{equation}
\omega_1=\dot\phi\sin\theta\sin\psi+\dot\theta\cos\psi,
\label{eq:euler-omega1}
\end{equation}
\begin{equation}
\omega_2=\dot\phi\sin\theta\cos\psi-\dot\theta\sin\psi,
\label{eq:euler-omega2}
\end{equation}
\begin{equation}
\omega_3=\dot\phi\cos\theta+\dot\psi.
\label{eq:euler-omega3}
\end{equation}
这三个式子看起来复杂，但含义清楚：\(\dot\phi\) 通过章动角 \(\theta\) 分解到刚体轴上，\(\dot\theta\) 沿节线方向，\(\dot\psi\) 沿刚体自身的 \(z'\) 轴。它们是把欧拉角动力学与欧拉方程连接起来的关键。

\begin{mathbox}
若只关心角速度的大小，可以利用
\[
\omega^2=\dot\theta^2+\dot\phi^2\sin^2\theta
+(\dot\psi+\dot\phi\cos\theta)^2.
\]
这是对称陀螺动能计算的起点。
\end{mathbox}

\subsection{刚体上各点的速度与加速度}

一个点固定在物体上时，它相对参考点的距离不变，方向却随物体转动。把参考点的运动和这段相对运动相加，得到该点的速度和加速度。

将刚体运动分解为基点的平动和绕基点的转动，是处理刚体运动学最常用的方法。取质心 \(C\) 为基点：
\begin{equation}
\bv_P=\bv_C+\bom\times(\br_P-\br_C),
\label{eq:velocity-decomp}
\end{equation}
\begin{equation}
\ba_P=\ba_C+\boldsymbol\alpha\times(\br_P-\br_C)
+\bom\times\bigl[\bom\times(\br_P-\br_C)\bigr].
\label{eq:acceleration-decomp}
\end{equation}
若刚体做平面平行运动，取平面为 \(xy\) 平面，则 \(\bom=\omega\hat{\boldsymbol k}\)，且
\[
\bv_P=\bv_C+\omega\hat{\boldsymbol k}\times(\br_P-\br_C).
\]
当平面运动的 $\omega\ne0$ 时，可找到瞬时速度为零的位置；纯平动没有有限瞬心。瞬心只是一时刻的速度表示，不能将它自动当作固定支点来写加速度或力矩方程。后文滚动将检验这一差别。

\section{刚体的惯量张量}

同样重的两个轮子，质量若更多集中在轮缘，转起来就更费力。对形状不对称的物体，还需同时知道它绕不同方向转动的惯性。

\subsection{角动量与惯量张量}

先算刚体中一个小质量块的速度和角动量，再把所有小块加起来。把总角动量与角速度相联系的矩阵叫惯量张量。

设刚体绕固定点 \(O\) 转动，角速度为 \(\bom\)。第 \(i\) 个质点的速度为
\[
\bv_i=\bom\times\br_i.
\]
刚体对 \(O\) 点的角动量为
\begin{equation}
\bJ_O=\sum_i m_i\br_i\times\bv_i
=\sum_i m_i\br_i\times(\bom\times\br_i).
\label{eq:angular-momentum-def}
\end{equation}
利用矢量恒等式
\[
\boldsymbol{A}\times(\boldsymbol{B}\times\boldsymbol{C})
=\boldsymbol{B}(\boldsymbol{A}\cdot\boldsymbol{C})-\boldsymbol{C}(\boldsymbol{A}\cdot\boldsymbol{B}),
\]
得到
\begin{equation}
\bJ_O=\sum_i m_i\left[r_i^2\bom-\br_i(\br_i\cdot\bom)\right].
\label{eq:angular-momentum-intermediate}
\end{equation}
把 \(\bom\) 提出，定义
\begin{equation}
\bI_O=\sum_i m_i\left(r_i^2\mathbf 1-\br_i\br_i^{\mathsf T}\right),
\label{eq:inertia-tensor-discrete}
\end{equation}
其中 \(\mathbf 1\) 是 \(3\times3\) 单位矩阵。把 \(\bI_O\) 称为刚体对 \(O\) 点的惯量张量，则
\begin{equation}
\boxed{\bJ_O=\bI_O\bom.}
\label{eq:J-equals-I-omega}
\end{equation}
对连续体，
\begin{equation}
\bI_O=\int \rho(\br)\left(r^2\mathbf 1-\br\br^{\mathsf T}\right)\dd V.
\label{eq:inertia-tensor-continuous}
\end{equation}

在直角坐标中，惯量张量的分量为
\begin{equation}
I_{xx}=\int\rho(y^2+z^2)\dd V,\qquad
I_{yy}=\int\rho(x^2+z^2)\dd V,\qquad
I_{zz}=\int\rho(x^2+y^2)\dd V,
\label{eq:inertia-diagonal}
\end{equation}
\begin{equation}
I_{xy}=I_{yx}=-\int\rho xy\dd V,\qquad
I_{xz}=I_{zx}=-\int\rho xz\dd V,\qquad
I_{yz}=I_{zy}=-\int\rho yz\dd V.
\label{eq:inertia-offdiagonal}
\end{equation}
对角元称为\textbf{转动惯量}，非对角元是本书采用负号约定的\textbf{惯量积}。有些书把 $\int xy\,\dd m$ 本身称为惯量积，使用公式前须核对负号。惯量张量是实对称矩阵：
\[
I_{ab}=I_{ba}.
\]

这里必须强调一个容易误解的地方：

\begin{keybox}
角动量 \(\bJ_O\) 一般\textbf{不平行}于角速度 \(\bom\)。对非零角速度，共线条件是 $\bom$ 属于惯量张量的某个本征子空间；在简并子空间中，不必沿选定的一根主轴基矢。惯量张量正是把“转动惯量”从单个标量推广为线性映射的数学对象。
\end{keybox}

例如均匀细杆被机构强迫绕非主轴转动时，角动量通常不沿角速度，即使转轴固定也可能需要横向力矩。反过来，无外力矩时角动量方向固定，角速度仍可变化；后文将分别讨论这两个问题。

\subsection{惯量张量的性质}

绕某条轴转时，离轴远的质量贡献较大。惯量张量的二次型就是把每一小块的质量乘它到该轴的距离平方后求和。

惯量张量有几个基本性质：

\begin{enumerate}[leftmargin=1.8em]
  \item \textbf{对称性}：\(I_{ab}=I_{ba}\)，因此总可以正交对角化。
  \item \textbf{半正定性}：对任意 \(\bom\)，
  \[
  \bom^{\mathsf T}\bI\bom
  =\sum_i m_i|\bom\times\br_i|^2\ge 0.
  \]
  等号成立当且仅当所有质量都位于转轴上。对真正有三维体积的刚体，等号不成立，惯量张量正定；对质量全在某条过参考点直线上的理想模型，该轴惯量为零。非退化薄板的质量不共线，三个主惯量仍都为正。
  \item \textbf{与原点有关}：同一刚体对不同参考点的惯量张量不同，需用平行轴定理换算。
  \item \textbf{与坐标系取向有关}：矩阵分量随坐标转动而变，但本征值不变。
\end{enumerate}

刚体的转动动能也可以写成二次型：
\begin{equation}
T_{\rm rot}
=\frac12\sum_i m_i v_i^2
=\frac12\sum_i m_i(\bom\times\br_i)\cdot(\bom\times\br_i)
=\frac12\bom^{\mathsf T}\bI_O\bom.
\label{eq:rot-kinetic-energy-tensor}
\end{equation}
这一步说明了惯量张量的物理意义：它是转动动能和角动量中的“质量矩阵”。

\begin{physicsbox}
平动惯性的度量是质量 \(m\)，转动惯性的度量是惯量张量 \(\bI\)。质量是标量，惯量张量是二阶张量；这意味着刚体对不同方向的转动“阻力”不同。把一个物体绕长轴和绕短轴旋转，感受到的惯性可以完全不同。
\end{physicsbox}

\subsection{坐标变换与惯量的物理可实现性}
把坐标轴换个方向，惯量矩阵的数字会变，同一物体绕同一根物理轴的惯量仍应相同。真实质量分布还会对主惯量的取值加上限制。

令 $R$ 把体分量变到空间分量，同一个参考点的张量满足
\begin{equation}
\bI_S=R\bI_bR^{\mathsf T},\qquad
\bJ_S=R\bJ_b,\qquad \bom_S=R\bom_b.
\label{eq:rigid-inertia-transform}
\end{equation}
这是由 $\bJ_b=\bI_b\bom_b$ 两边左乘 $R$ 得到的，不同参考点还须另用平行轴定理。在体轴中惯量常数，在空间轴中通常随姿态变化。

主惯量除非负外还满足三角不等式。例如主轴坐标下
\[
I_1+I_2-I_3=2\int x_3^2\,\dd m\ge0,
\]
其余循环置换。故任意三个正数不一定都能作为某个质量分布的主惯量；薄板位于 $x_3=0$ 时恰有 $I_3=I_1+I_2$，并不意味着某个主惯量为零。

\begin{example}[带惯量积的三质点刚架]
三个质量均为 $m$ 的质点，用无质量刚杆固定在
\[
(a,0,0),\quad(0,a,0),\quad(-a,-a,0).
\]
质心为原点。逐点求和得
\[
\bI_C=ma^2\begin{pmatrix}2&-1&0\\-1&2&0\\0&0&4\end{pmatrix}.
\]
例如 $I_{xy}=-m[0+0+a^2]=-ma^2$，$I_{xx}=m(0+a^2+a^2)$。主轴为
\[
\be_1=(1,1,0)/\sqrt2,\quad\be_2=(-1,1,0)/\sqrt2,\quad\be_3=(0,0,1),
\]
对应 $I_1=ma^2,I_2=3ma^2,I_3=4ma^2$。它们满足薄板等式，且全部为正。若瞬时 $\bom=\Omega(1,0,0)$，则
\[
\bJ=ma^2\Omega(2,-1,0),\qquad T=ma^2\Omega^2,
\]
直接显示角动量偏离角速度。用主轴投影计算同一动能，结果不变。
\end{example}


\begin{figure}[htbp]
\centering
\begin{tikzpicture}[>=Stealth]
\draw[thick,blue!50!black](1,0)--(0,1)--(-1,-1)--cycle;
\draw[->,gray](-1.5,0)--(1.7,0)node[right]{$x$};\draw[->,gray](0,-1.4)--(0,1.7)node[above]{$y$};
\draw[->,dashed,red!70!black](0,0)--(1.2,1.2)node[right]{$\be_1$};
\draw[->,dashed,red!70!black](0,0)--(-1.2,1.2)node[left]{$\be_2$};
\foreach \x/\y/\pos in {1/0/below right,0/1/above right,-1/-1/below left}{\fill(\x,\y)circle(3pt);\node[\pos]at(\x,\y){$m$};}
\fill(0,0)circle(1.8pt)node[below right]{$C$};
\end{tikzpicture}
\caption{三质点无质量刚架的坐标截面，以$a$为长度单位。质心在原点，红虚线是面内两个惯量主轴；法向为第三主轴。}
\label{fig:rigid-frame}
\end{figure}

\subsection{平行轴定理与垂直轴定理}

已知绕质心轴的惯量后，若把轴平移到旁边，还要加上由离轴距离增加带来的贡献。薄板在特定三根轴之间还有一个便捷的加法关系。

设 \(C\) 为质心，刚体对质心的惯量张量为 \(\bI_C\)。对任意点 \(O\)，设
\[
\br_i=\bR_C+\boldsymbol\rho_i,
\qquad
\sum_i m_i\boldsymbol\rho_i=0.
\]
代入式 \eqref{eq:inertia-tensor-discrete}，含 \(\boldsymbol\rho_i\) 一次的交叉项为
\[
\sum_i m_i\left[
2(\bR_C\cdot\boldsymbol\rho_i)\mathbf 1
-\bR_C\boldsymbol\rho_i^{\mathsf T}
-\boldsymbol\rho_i\bR_C^{\mathsf T}\right],
\]
它因 \(\sum_i m_i\boldsymbol\rho_i=0\) 而消失，于是得到
\begin{equation}
\boxed{
\bI_O=\bI_C+M\left(R_C^2\mathbf 1-\bR_C\bR_C^{\mathsf T}\right).
}
\label{eq:parallel-axis-tensor}
\end{equation}
其中 \(M\) 为刚体总质量，\(\bR_C\) 是质心相对 \(O\) 的位矢。

若转轴方向为单位矢量 \(\hat{\be}\)，则绕该轴的标量转动惯量为
\begin{equation}
I_O(\hat{\be})=I_C(\hat{\be})+M d_\perp^2,
\label{eq:parallel-axis-scalar}
\end{equation}
其中 \(d_\perp\) 是质心到转轴的垂直距离。这就是通常的\textbf{平行轴定理}。

对薄板，取板面为 \(xy\) 平面，则 \(z=0\)，有
\begin{equation}
I_z=I_x+I_y.
\label{eq:perpendicular-axis}
\end{equation}
这是\textbf{垂直轴定理}，只适用于薄板。

\subsection{惯量主轴与惯量椭球}

总可以找一组互相垂直的方向，使绕这些方向之一转动时角动量沿同一方向。这些方向叫主轴，可用矩阵的本征向量求出。

由实对称矩阵的谱定理，实对称矩阵 \(\bI\) 必可正交对角化。对三维刚体，\(\bI\) 还是正定的；对过参考点的共线质点模型可有零主惯量，非退化薄板没有这一退化。因此存在三个相互正交的单位矢量
\[
\hat{\be}_1,\hat{\be}_2,\hat{\be}_3
\]
和三个非负实数 \(I_1,I_2,I_3\)，使得
\begin{equation}
\bI\hat{\be}_a=I_a\hat{\be}_a.
\label{eq:principal-axis-eq}
\end{equation}
这三个方向称为\textbf{惯量主轴}，\(I_a\) 称为\textbf{主转动惯量}。在以主轴为坐标轴的坐标系中，
\begin{equation}
\bI=\operatorname{diag}(I_1,I_2,I_3).
\label{eq:principal-inertia}
\end{equation}
这些对称性推论都要求惯量参考点位于相应对称面或对称轴上：对称面法线是主轴，至少二重的旋转对称轴是主轴；高于二重的旋转对称轴还使两个横向主惯量相等。

对任意方向 \(\hat{\be}\)，绕该轴的转动惯量为
\begin{equation}
I(\hat{\be})=\hat{\be}^{\mathsf T}\bI\hat{\be}.
\label{eq:inertia-direction}
\end{equation}
若在方向 \(\hat{\be}\) 上取长度
\[
r=\frac{1}{\sqrt{I(\hat{\be})}},
\]
则所有这样的点构成\textbf{惯量椭球}：
\begin{equation}
\br^{\mathsf T}\bI\br=1.
\label{eq:inertia-ellipsoid}
\end{equation}
在主轴坐标系中，它写成
\begin{equation}
I_1x^2+I_2y^2+I_3z^2=1.
\label{eq:inertia-ellipsoid-principal}
\end{equation}
当三个主转动惯量都非零时，三个半轴长度分别为
\[
\frac{1}{\sqrt{I_1}},\qquad
\frac{1}{\sqrt{I_2}},\qquad
\frac{1}{\sqrt{I_3}}.
\]
主轴方向就是椭球的三个主轴方向。转动惯量最大的方向对应椭球最短的半轴，转动惯量最小的方向对应最长的半轴。若某个主转动惯量为零，惯量椭球退化为柱面等退化二次曲面；这时“椭球”应作广义理解。

\begin{figure}[htbp]
\centering
\begin{tikzpicture}[>=Stealth]
\draw[thick,fill=blue!7] (0,0) ellipse (2.4 and 1.35);
\draw[->] (-2.9,0)--(3.1,0) node[right]{$\xi_1$};
\draw[->] (0,-1.8)--(0,2) node[above]{$\xi_2$};
\draw[<->,blue!65!black] (0,-0.3)--(2.4,-0.3);
\node[below,fill=white] at (1.2,-0.32) {$1/\sqrt{I_1}$};
\draw[<->,red!70!black] (-0.3,0)--(-0.3,1.35);
\node[left,fill=white] at (-0.35,0.8) {$1/\sqrt{I_2}$};
\end{tikzpicture}
\caption{惯量椭球在主平面 $\xi_3=0$ 上的截面，取 $I_2>I_1$。较大的惯量对应较短的半轴；第三个半轴垂直本截面。图中 $\xi_i$ 是惯量二次型的辅助坐标，不是物体的真实边界。}
\label{fig:inertia-ellipsoid}
\end{figure}

\begin{remark}
惯量主轴不一定唯一。若两个主转动惯量相等（例如对称陀螺），则垂直于对称轴的任意两个正交方向都可以选为另外两个主轴；若三个主转动惯量都相等（例如均匀球体），则任意正交三轴都是主轴。这种简并来自对称性，不影响物理结果。
\end{remark}

\subsection{典型刚体的转动惯量}

算规则物体的惯量时，先决定绕哪条轴，再选容易分块的质量元。每一小块对惯量的贡献只用它到轴的垂直距离。

下表列出几种常见均匀刚体对质心主轴（或特别注明的固定轴）的转动惯量。表中 \(M\) 为质量，\(R\) 为半径，\(L\) 为长度，\(a,b\) 为矩形边长；薄圆盘可看作 \(L=0\) 的圆柱。

\begin{table}[htbp]
\centering
\small
\begin{tabular}{lll}
\toprule
刚体 & 转轴 & 转动惯量 \\
\midrule
细杆 & 过质心且垂直杆 & \(\dfrac{1}{12}ML^2\) \\
细杆 & 过一端且垂直杆（固定轴） & \(\dfrac{1}{3}ML^2\) \\
圆盘或圆柱 & 对称轴 & \(\dfrac{1}{2}MR^2\) \\
圆柱（长 \(L\)） & 过质心且垂直对称轴 & \(\dfrac{1}{4}MR^2+\dfrac{1}{12}ML^2\) \\
圆环 & 对称轴 & \(MR^2\) \\
实心球 & 任一直径 & \(\dfrac{2}{5}MR^2\) \\
薄球壳 & 任一直径 & \(\dfrac{2}{3}MR^2\) \\
矩形薄板 & 过质心且平行边 \(a\) & \(\dfrac{1}{12}Mb^2\) \\
\bottomrule
\end{tabular}
\caption{常见均匀刚体对质心主轴或特别注明固定轴的转动惯量。}
\label{tab:inertia-moments}
\end{table}

\begin{example}[矩形薄板绕对角线转动]
质量为 \(M\)、边长为 \(a,b\) 的均匀矩形薄板，求绕其对角线转动的转动惯量。

取质心为原点，\(x\) 轴平行于边 \(a\)，\(y\) 轴平行于边 \(b\)。面密度为 $M/(ab)$，例如
\[
I_{xx}=\frac{M}{ab}\int_{-a/2}^{a/2}\dd x\int_{-b/2}^{b/2}y^2\,\dd y=\frac{Mb^2}{12}.
\]
另一对角元交换 $a,b$；$xy$ 项关于任一坐标为奇函数，积分为零。因此
\[
I_{xx}=\frac{1}{12}Mb^2,\qquad
I_{yy}=\frac{1}{12}Ma^2,\qquad
I_{xy}=0.
\]
对角线方向单位矢量为
\[
\hat{\be}=\frac{1}{\sqrt{a^2+b^2}}(a,b,0).
\]
于是
\[
I_{\rm diag}
=I_{xx}\cos^2\alpha+I_{yy}\sin^2\alpha,
\]
其中
\[
\cos\alpha=\frac{a}{\sqrt{a^2+b^2}},
\qquad
\sin\alpha=\frac{b}{\sqrt{a^2+b^2}}.
\]
代入得
\[
I_{\rm diag}
=\frac{1}{12}Mb^2\frac{a^2}{a^2+b^2}
+\frac{1}{12}Ma^2\frac{b^2}{a^2+b^2}
=\frac{Ma^2b^2}{6(a^2+b^2)}.
\]
当 \(a=b\) 时，\(I_{\rm diag}=Ma^2/12\)，与正方形绕任意过质心的面内轴相同，这是正方形对称性的结果。
\end{example}

\begin{physicsbox}
主轴是关于指定参考点的惯性方向，不是凭物体外形随便指定的“长轴”。物体改变内部形状时，惯量也会改变，但它已超出单一刚体的固定惯量模型；应将各刚体部件与内部关节另行建模。
\end{physicsbox}

\section{力系简化与刚体平衡}

同一组外力既能推动整个物体，也能让它转。分别把总力和总力矩记下来，就能判断刚体是否平衡。

\subsection{力系的简化}

把分布在不同位置的力搬到同一个点表示时，要把搬动所改变的力矩补回来。总力决定平动趋势，总力矩还保留力作用位置的信息。

刚体不变形，因此作用在刚体上的力可以沿作用线滑动而不改变外效应。对力系进行简化时，所有力可以向任意一点 \(O\) 平移，但每平移一个力 \(\bF_i\) 都要附加一个力偶
\[
\bM_i=(\br_i-\br_O)\times\bF_i.
\]
于是任意力系等效为
\[
\bF=\sum_i\bF_i,
\qquad
\bM_O=\sum_i(\br_i-\br_O)\times\bF_i.
\]
这就是\textbf{力系向一点简化}的结果：一个合力加一个合力偶。

若 $\bF\ne0$，把合力放在相对 $O$ 的位置 $\ba=(\bF\times\bM_O)/F^2$，则 $\ba\times\bF$ 正好重现垂直于 $\bF$ 的力矩，剩下平行力偶 $(\bF\cdot\bM_O)\bF/F^2$。因此 $\bF\cdot\bM_O=0$ 时可用单个合力替代，否则为力螺旋。若 $\bF=0$，只剩纯力偶或零力系。这个等效针对刚体的整体动力学，不保证局部应力相同。

\subsection{刚体平衡方程}

柜子静止时，既没有平动加速度，也没有转动加速度。写出合力和合力矩为零后，还要检查接触面能否提供所算出的支持力和摩擦力。

这里讨论\textbf{静止平衡}：质心速度和角速度都为零，质心加速度和角加速度也为零。给定非零初速度的运动问题不能只靠这两个代数条件判断，尤其非主轴恒速转动可能需要力矩。由动量定理和角动量定理，
\begin{equation}
\boxed{
\sum_i\bF_i=0,\qquad
\sum_i\bM_O^{(i)}=0.
}
\label{eq:rigid-equilibrium}
\end{equation}
这两式配合静止初态给出平衡的必要动力学条件；还须满足约束几何和单侧反力、静摩擦可行性。它们本身不把任意已有运动自动停止，也不保证平衡稳定。第二式对任意点 \(O\) 都成立；若对某一点成立，则对另一点也成立（在合力为零的前提下）。

在平面问题中，平衡方程通常写成三个标量方程：
\[
\sum F_x=0,\qquad
\sum F_y=0,\qquad
\sum M_O=0.
\]
这三个方程最多能确定三个未知约束力。如果未知力多于独立方程，系统就是\textbf{静不定}的，需要引入变形协调条件；刚体模型本身无法给出唯一解。

\subsection{例题：梯子的平衡}

先给梯子画受力图。以梯脚为矩心可暂时去掉梯脚的未知力，先求墙的推力，再回到平动平衡求摩擦。

\begin{example}[斜靠在墙上的梯子]
长度为 \(2L\)、质量为 \(M\) 的均匀梯子斜靠在光滑竖直墙上，下端放在粗糙水平地面上，与地面夹角为 \(\theta\)。求梯子不滑倒所需的最小静摩擦因数。

设墙面对梯子的法向力为 \(N_w\)，地面对梯子的法向力为 \(N_f\)，静摩擦力为 \(f\)。平衡方程为
\[
N_w=f,\qquad N_f=Mg.
\]
对梯子与地面的接触点取力矩：
\[
N_w(2L\sin\theta)-Mg(L\cos\theta)=0.
\]
因此
\[
N_w=\frac{Mg}{2}\cot\theta.
\]
不滑倒条件为
\[
f=N_w\le \mu N_f=\mu Mg.
\]
代入得
\[
\frac12\cot\theta\le\mu,
\]
即
\[
\tan\theta\ge\frac{1}{2\mu}.
\]
当 \(\theta\) 太小时，梯子越接近水平，墙面对梯子的法向力越大，需要的摩擦力也越大，因此越容易滑倒。
\end{example}

\begin{figure}[htbp]
\centering
\begin{tikzpicture}[>=Stealth]
\draw[thick](-.6,0)--(3,0);\draw[thick](2,0)--(2,3.9);
\draw[very thick,blue!65!black](0,0)--(2,3.464);
\fill(0,0)circle(2pt);\fill(1,1.732)circle(2pt)node[right]{$C$};
\draw[red!70!black,thick,->](1,1.732)--(1,.85)node[right]{$Mg$};
\draw[green!45!black,thick,->](2,3.464)--(.9,3.464)node[left]{$N_w$};
\draw[green!45!black,thick,->](0,0)--(0,1)node[left]{$N_f$};
\draw[green!45!black,thick,->](0,0)--(1,0)node[below]{$f$};
\draw(.55,0)arc(0:60:.55);\node at(.67,.38){$\theta$};
\end{tikzpicture}
\caption{均匀梯子的四个外力，图取$\theta=60^\circ$。光滑墙面只给水平法向力，地面摩擦在此向墙，重力作用于中点；是否平衡还受静摩擦上限限制。}
\label{fig:rigid-ladder}
\end{figure}

这个例题展示了刚体平衡问题的标准步骤：画受力图、写合力方程、写力矩方程、用约束条件（这里是摩擦力上限）判断可行性。注意力矩方程可以取任意点，取在未知力最多的点通常最方便。

\begin{physicsbox}
仅在水平支承、单侧法向接触、重力主导且无额外力偶的静力模型中，质心竖直投影落在支撑区域内才给出抗倾覆条件；还要检查不滑动。侧向外力、加速、抓取或主动驱动会改变条件，不能用一个支撑多边形判据包办所有刚体平衡。
\end{physicsbox}

\section{刚体动力学基本定理}

研究刚体的运动时，通常先求质心怎样走，再求物体绕质心怎样转。必要时可以改用支点或接触点，但要检查这个点是不是在运动。

\subsection{动量与角动量}

把刚体各点的动量加起来，得到质量乘质心速度。再相对于选定的点算角动量，便能分出质心平动和绕质心转动的贡献。

令 \(\br_i=\br_C+\boldsymbol\rho_i\)，其中 \(\sum_i m_i\boldsymbol\rho_i=0\)。刚体的总动量为
\[
\bP=\sum_i m_i\bv_i
=\sum_i m_i(\bv_C+\bom\times\boldsymbol\rho_i)
=M\bv_C+\bom\times\sum_i m_i\boldsymbol\rho_i
=M\bv_C,
\]
即
\begin{equation}
\bP=M\bv_C,
\label{eq:rigid-momentum}
\end{equation}
其中 \(M\) 为总质量，\(\bv_C\) 为质心速度。由质点系动量定理，质心运动方程为
\begin{equation}
M\ba_C=\sum_i\bF_i^{\rm ext}.
\label{eq:rigid-cm-equation}
\end{equation}
这说明刚体的平动部分与一个质量为 \(M\) 的质点完全相同，内部力不影响质心运动。

利用同一分解，速度写成 \(\bv_i=\bv_C+\bom\times\boldsymbol\rho_i\)。代入角动量定义：
\[
\bJ_O=\sum_i m_i(\br_C+\boldsymbol\rho_i)\times(\bv_C+\bom\times\boldsymbol\rho_i).
\]
展开后，含 \(\boldsymbol\rho_i\) 的一次项因 \(\sum_i m_i\boldsymbol\rho_i=0\) 而消失，剩下
\[
\bJ_O=\br_C\times(M\bv_C)
+\sum_i m_i\boldsymbol\rho_i\times(\bom\times\boldsymbol\rho_i)
=\br_C\times\bP+\bJ_C.
\]
因此刚体对固定点 \(O\) 的角动量为
\begin{equation}
\bJ_O=\bJ_C+\br_C\times\bP,
\label{eq:rigid-J-decomp}
\end{equation}
其中 \(\bJ_C=\bI_C\bom\) 是绕质心的自转角动量。由质点系角动量定理，对固定点 \(O\)，
\begin{equation}
\dot\bJ_O=\bM_O^{\rm ext}.
\label{eq:rigid-angular-momentum}
\end{equation}
同样由质点系角动量定理，对质心 \(C\)，
\begin{equation}
\dot\bJ_C=\bM_C^{\rm ext}.
\label{eq:rigid-angular-momentum-cm}
\end{equation}
第二式特别重要：即使质心有加速度，只要取质心为矩心，角动量定理仍然保持简单形式。这就是为什么刚体动力学通常把运动分解为“质心平动 + 绕质心转动”。

\subsection{移动矩心：先说明角动量的定义}
轮子接触地面的那个物质点，可以在这一瞬间速度为零，下一瞬间却已在动。若要绕它写角动量方程，必须先交代采用绝对速度还是相对速度。

取惯性系中的任意移动点 $A$，令 $\rho_i=\br_i-\br_A$。若角动量用质点的\emph{绝对速度}定义，
\[
{\boldsymbol H}_A=\sum_i\rho_i\times m_i\bv_i,
\]
逐项微分得
\[
\dot{\boldsymbol H}_A=\sum_i(\bv_i-\bv_A)\times m_i\bv_i
+\sum_i\rho_i\times\bF_i
=\bM_A^{\rm ext}-\bv_A\times\bP.
\]
其中已用内部中心力矩相消，或将刚体内力视为理想约束。若改用\emph{相对速度}，定义
\[
{\boldsymbol K}_A=\sum_i\rho_i\times m_i(\bv_i-\bv_A),
\]
则第一项的同矢量叉积为零，得到
\begin{equation}
\dot{\boldsymbol K}_A=\bM_A^{\rm ext}-M\rho_C\times\ba_A,
\quad \rho_C=\br_C-\br_A.
\label{eq:rigid-moving-origin}
\end{equation}
${\boldsymbol H}_A$ 与 ${\boldsymbol K}_A$ 不可混称；两者满足 $\boldsymbol H_A-\boldsymbol K_A=M\rho_C\times\bv_A$，上面的时间导数均在惯性系取。若 $A$ 固连于刚体，${\boldsymbol K}_A=\bI_A\bom$，故一般还要保留参考点加速度项。取 $A=C$ 时 $\rho_C=0$；取固定点时 $\bv_A=\ba_A=0$，两种定义都恢复简单角动量定理。这是选质心或固定支点的理由。

\begin{example}[瞬时静止的接触点不是固定支点]
圆轮在静止水平面上向右匀速无滑动滚动，半径 $R$，速度 $v$，角速度矢量为 $-(v/R)\hat z$。当前接触的物质点 $P$ 有 $\bv_P=0$，但
\[
\ba_P=\bom\times[\bom\times(-R\hat y)]
=\frac{v^2}{R}\hat y\ne0.
\]
下一瞬间接触地面的是另一个物质点。当前轮的动能确可写成 $I_P\omega^2/2$，却不能据此把 $P$ 当作一段时间内固定的转轴。若用随轮的这个物质点作矩心，应检验式 \eqref{eq:rigid-moving-origin}；本例修正力矩恰因 $\rho_C\parallel\ba_P$ 为零，但其加速度仍不为零。
\end{example}


\subsection{动能与功}

向前滚的轮子有平动动能，也有自转动能。功率则要看各力作用点的速度，接触点是否滑动尤其重要。

把 \(\bv_i=\bv_C+\bom\times\boldsymbol\rho_i\) 代入动能定义：
\[
T=\frac12\sum_i m_i|\bv_C+\bom\times\boldsymbol\rho_i|^2.
\]
展开后的交叉项为
\[
\bv_C\cdot\left(\bom\times\sum_i m_i\boldsymbol\rho_i\right)=0,
\]
于是只剩两项：
\[
T=\frac12Mv_C^2
+\frac12\sum_i m_i|\bom\times\boldsymbol\rho_i|^2
=\frac12Mv_C^2+\frac12\bom^{\mathsf T}\bI_C\bom.
\]
因此刚体的动能可以分解为质心平动动能和绕质心转动动能：
\begin{equation}
\boxed{
T=\frac12Mv_C^2+\frac12\bom^{\mathsf T}\bI_C\bom.
}
\label{eq:rigid-kinetic-energy}
\end{equation}
这就是刚体版的柯尼希定理。若刚体绕固定点 \(O\) 转动，则
\begin{equation}
T=\frac12\bom^{\mathsf T}\bI_O\bom.
\label{eq:rigid-kinetic-energy-fixed}
\end{equation}
若刚体做定轴转动，\(\bom=\omega\hat{\be}\)，则
\[
T=\frac12 I\omega^2.
\]

对无限小转动 \(\dd\bth\)，刚体上一点 \(\br_i=\br_C+\boldsymbol\rho_i\) 的位移为
\[
\dd\br_i=\dd\br_C+\dd\bth\times\boldsymbol\rho_i.
\]
因此外力做的元功为
\[
\dd W=\sum_i\bF_i\cdot\dd\br_i
=\bF\cdot\dd\br_C
+\dd\bth\cdot\sum_i\boldsymbol\rho_i\times\bF_i
=\bF\cdot\dd\br_C+\bM_C\cdot\dd\bth.
\]
写成方程即
\begin{equation}
\dd W=\sum_i\bF_i\cdot\dd\br_i
=\bF\cdot\dd\br_C+\bM_C\cdot\dd\bth.
\label{eq:rigid-work}
\end{equation}
因此功率为
\begin{equation}
P=\bF\cdot\bv_C+\bM_C\cdot\bom.
\label{eq:rigid-power}
\end{equation}
对保守力系，机械能守恒：
\[
T+V=E.
\]

\subsection{定轴转动与复摆}

钥匙绕固定轴转时，只需一个角度；绕固定轴摆动的硬杆也一样。求小幅摆动的周期时，恢复力矩和整根杆的惯量都要算。

设刚体绕固定轴 \(z\) 转动，转动惯量为 \(I\)。取角动量定理的 \(z\) 分量；理想固定轴的反力对 \(z\) 轴没有力矩，因此只有外力矩进入方程：
\begin{equation}
I\ddot\theta=M_z,
\label{eq:fixed-axis-equation}
\end{equation}
其中 \(M_z\) 是外力对转轴的力矩。若只有重力做功，则能量积分为
\[
\frac12I\dot\theta^2+V(\theta)=E.
\]

\begin{example}[复摆]
刚体绕通过 \(O\) 点的水平固定轴摆动，质心到轴的距离为 \(l\)，总质量为 \(m\)，绕轴转动惯量为 \(I_O\)。求小摆动的周期。

设摆角为 \(\theta\)，重力矩为
\[
M_z=-mgl\sin\theta.
\]
运动方程为
\[
I_O\ddot\theta=-mgl\sin\theta.
\]
小角度下 \(\sin\theta\approx\theta\)，得
\[
\ddot\theta+\frac{mgl}{I_O}\theta=0.
\]
因此
\[
\omega_0=\sqrt{\frac{mgl}{I_O}},
\qquad
T=2\pi\sqrt{\frac{I_O}{mgl}}.
\]
对长为 \(L\) 的均匀细杆绕一端摆动，
\[
I_O=\frac13ML^2,\qquad l=\frac L2,
\]
所以
\[
T=2\pi\sqrt{\frac{2L}{3g}}.
\]
这比同样长度的单摆周期短，因为杆的质量分布更靠近转轴。
\end{example}

\begin{figure}[htbp]
\centering
\begin{tikzpicture}[>=Stealth]
\draw[very thick,blue!65!black](0,0)--(1.5,-2.598);
\fill(0,0)circle(2pt)node[above]{$O$};
\fill(.75,-1.299)circle(3pt)node[right]{$C$};
\draw[dashed,gray](0,0)--(0,-2.8);
\draw(0,-.65)arc(-90:-60:.65);\node at(.24,-.85){$\theta$};
\draw[->,red!70!black,thick](.75,-1.299)--(.75,-2.25)node[right]{$mg$};
\draw[<->](-.25,.05)--(.5,-1.249)node[midway,left]{$l$};
\draw[dashed](0,-1.299)--(.75,-1.299);
\node[below]at(.32,-1.4){$l\sin\theta$};
\end{tikzpicture}
\caption{以均匀杆为示意的复摆，$\theta=30^\circ$。重力作用于质心，$l=OC$，力臂为$l\sin\theta$，角度从竖直向下量起。}
\label{fig:physical-pendulum}
\end{figure}

\subsection{平面平行运动}

在平面上滚动的物体，可以先画质心的速度，再加上各点绕质心转动的速度。检查接触点相对地面的速度是否为零，可判断能否按纯滚动处理。

刚体做平面平行运动时，取质心 \(C\) 为基点。由质心运动方程和绕质心的角动量定理，得到
\begin{equation}
M\ddot x_C=F_x,\qquad
M\ddot y_C=F_y,\qquad
I_C\ddot\theta=M_C,
\label{eq:planar-motion-equations}
\end{equation}
其中 \(F_x,F_y\) 是外力在平面内的分量，\(M_C\) 是外力对过质心且垂直于运动平面的轴的合力矩。
动能写成
\begin{equation}
T=\frac12Mv_C^2+\frac12I_C\omega^2.
\label{eq:planar-kinetic-energy}
\end{equation}
若刚体在静止曲面上无滑动滚动，当前接触的物质点速度为零，这就是滚动约束；曲面运动时应要求双方接触点的相对速度为零。

\begin{figure}[htbp]
\centering
\begin{tikzpicture}[>=Stealth,scale=1.2]
\begin{scope}[rotate=-20]
\draw[thick](-2,0)--(3,0);
\draw[fill=blue!7,thick](0,1)circle(1);
\fill(0,1)circle(2pt)node[above left=3pt]{$C$};\fill(0,0)circle(2pt)node[below]{$P$};
\draw[->,green!45!black,thick](0,0)--(-1.2,0)node[above]{$f$};
\draw[->,green!45!black,thick](0,0)--(0,.65)node[pos=.55,left=3pt]{$N$};
\draw[->,red!70!black,thick](0,1)--(.342,.06)node[pos=.75,right=8pt]{$Mg$};
\draw[->,blue!65!black,thick](0,1)--(1.3,1)node[above]{$v_C$};
\draw[->](0,1.55)arc(90:-15:.55);\node at(.75,1.65){$\omega$};
\end{scope}
\draw[dashed,gray](0,0)--(2.1,0);\draw(.8,0)arc(0:-20:.8);\node at(1.05,-.16){$\beta$};
\end{tikzpicture}
\caption{圆轮截面与斜面严格相切。下坡向右，转动顺时针，静摩擦向上坡、法向力垂直斜面、重力竖直；图取$\beta=20^\circ$，无滑动还须检验静摩擦上限。}
\label{fig:rolling-cylinder}
\end{figure}

\begin{example}[圆柱沿斜面滚下]
质量为 \(M\)、半径为 \(R\)、绕质心转动惯量为 \(I_C\) 的圆柱，沿倾角为 \(\beta\) 的斜面做无滑动滚动。求质心加速度。

取沿斜面向下为质心加速度的正方向，摩擦力 \(f\) 沿斜面向上，并取与该滚动方向相应的角加速度为正。沿斜面方向：
\[
Mg\sin\beta-f=Ma.
\]
绕质心取相应的正方向：
\[
fR=I_C\alpha.
\]
无滑动条件为
\[
a=R\alpha.
\]
由后两式得
\[
f=\frac{I_C}{R^2}a.
\]
代入第一式：
\[
Mg\sin\beta-\frac{I_C}{R^2}a=Ma,
\]
所以
\[
a=\frac{g\sin\beta}{1+I_C/(MR^2)}.
\]
对实心圆柱 \(I_C=\frac12MR^2\)，得
\[
a=\frac23g\sin\beta.
\]
对实心球 \(I_C=\frac25MR^2\)，得
\[
a=\frac57g\sin\beta.
\]
对薄圆环 \(I_C=MR^2\)，得
\[
a=\frac12g\sin\beta.
\]
转动惯量越大，越来越多的重力势能转化为转动动能，质心平动加速度越小。
\end{example}

由 $N=Mg\cos\beta$ 及上述方程，还应检查
\[
f=\frac{\kappa}{1+\kappa}Mg\sin\beta,\qquad
\mu_s\ge\frac{\kappa}{1+\kappa}\tan\beta,\quad
\kappa=\frac{I_C}{MR^2}.
\]
实心圆柱要求 $\mu_s\ge\tan\beta/3$。若不足，应改用滑动摩擦模型，而不是继续强加 $a=R\alpha$。静止斜面上的理想静摩擦总功率为 $\bF_f\cdot\bv_P=0$；它对质心平动的功率 $-fv_C$ 与对转动的功率 $fR\omega$ 抵消。故摩擦传递两种动能，却不在该模型中耗散机械能。



\begin{example}[台球击球点]
质量为 \(M\)、半径为 \(R\) 的均匀实心球静止在水平桌面上。用水平冲量在距球心高度 \(h\) 处击球，求击后瞬间质心速度与角速度，并判断球会先滑动还是直接滚动。

击球瞬间摩擦冲量可忽略。用 \(\mathcal J\) 表示击球冲量。由动量定理和角动量定理，
\[
Mv_0=\mathcal J,\qquad I_C\omega_0=\mathcal J h.
\]
这里取球向右运动，并把顺时针方向的角速度取为正值。因此
\[
v_0=\frac{\mathcal J}{M},\qquad
\omega_0=\frac{\mathcal J h}{I_C}.
\]
球与桌面的接触点速度为
\[
v_P=v_0-R\omega_0
=\mathcal J\left(\frac1M-\frac{Rh}{I_C}\right).
\]
对实心球 \(I_C=\frac25MR^2\)，有
\[
v_P=\frac{\mathcal J}{M}\left(1-\frac{5h}{2R}\right).
\]
当 \(h=\frac25R\) 时，\(v_P=0\)，球击后立即无滑动滚动；当 \(h>\frac25R\) 时，\(v_P<0\)，接触点向后滑，摩擦向前，球会加速；当 \(h<\frac25R\) 时，接触点向前滑，摩擦向后，球会减速。这里的判据依赖均匀实心球、水平冲量、击球期间桌面切向冲量可忽略等假设，不能直接套用于有壳、变形或复杂接触的其他球体。
\end{example}

\begin{example}[击后滑动怎样过渡到滚动]
沿用台球例题的顺时针标量角速度，记 $I_C=\kappa MR^2$、$u_0=v_0-R\omega_0$。击球结束后，假设桌面动摩擦因数 $\mu_k>0$、法向力 $Mg$，且没有其他水平力或滚阻。滑动阶段
\[
\dot v=-\mu_kg\,\operatorname{sgn}u,\qquad
R\dot\omega=\frac{\mu_kg}{\kappa}\operatorname{sgn}u,
\quad \dot u=-\mu_kg(1+1/\kappa)\operatorname{sgn}u.
\]
因此达到 $u=0$ 的时间和末速度为
\[
t_r=\frac{|u_0|}{\mu_kg(1+1/\kappa)},\qquad
v_r=\frac{v_0+\kappa R\omega_0}{1+\kappa},\quad \omega_r=v_r/R.
\]
摩擦的总功率为 $-\mu_kMg|u|$，能量损失
\[
T_0-T_r=\frac12\frac{\kappa M}{1+\kappa}u_0^2.
\]
实心球 $\kappa=2/5$ 时 $v_r=(5v_0+2R\omega_0)/7$。中心击球 $\omega_0=0$ 给 $v_r=5v_0/7$；在 $h=2R/5$ 处击球则 $u_0=0$、无这一滑动损失。进入滚动后，本理想模型无需维持非零摩擦；实际滚阻属于另加效应。
\end{example}


\subsection{非主轴转动：轴承为何承受动态负载？}
轴转得一样快，角动量的方向仍可能随物体一起改变。轴承必须提供相应力矩，才能维持这种转动。

\begin{example}[角速度与角动量不平行的直接后果]
均匀细杆质量 $m$、长度 $l$，质心固定，杆方向为单位矢量 $\bu$。理想细杆惯量张量为
\[
\bI=I_\perp(\mathbf1-\bu\bu^{\mathsf T}),\qquad I_\perp=ml^2/12.
\]
令外部机构迫使杆绕固定 $z$ 轴以恒定角速度 $\bom=\Omega\hat{\boldsymbol z}$ 转动，并保持夹角 $0<\alpha<\pi/2$。于是
\[
\bJ=I_\perp\Omega(\hat{\boldsymbol z}-\cos\alpha\,\bu),
\qquad \dot\bu=\Omega\hat{\boldsymbol z}\times\bu.
\]
对空间角动量求导，得保持这种运动所需的外力矩
\[
\bM=-I_\perp\Omega^2\cos\alpha\,
\hat{\boldsymbol z}\times\bu,
\qquad |\bM|=I_\perp\Omega^2\sin\alpha\cos\alpha.
\]
虽然 $\dot\Omega=0$，外力矩仍不为零，因为角动量方向在改变。$\alpha=\pi/2$ 时固定轴是主轴，动态力矩消失；$\alpha=0$ 时理想细杆绕自身轴惯量为零，这个退化极限不能代表真实转子的纵向惯量。力矩与角速度正交，故功率 $\bM\cdot\bom=0$，与转动动能恒定相容。
\end{example}



\section{欧拉方程与自由刚体}

用固定在刚体上的轴来写方程，惯量矩阵比较简单，但轴自己在转。欧拉方程正是把这个变化算进去后的角动量方程。

\subsection{欧拉方程}

先用角动量定理，再把它改写成随刚体转动的三个分量。即使惯量在这套坐标里不变，也要保留由于坐标轴转动产生的项。

写欧拉方程前，先推导体轴的输运公式。设 $\be_i(t)$ 是固连基矢，任意矢量写为
\[\bA=\sum_i A_i(t)\be_i(t).\]
由 $\dot\be_i=\bom\times\be_i$，乘积法则给出
\begin{equation}
\left(\frac{\dd\bA}{\dd t}\right)_S
=\sum_i\dot A_i\be_i+\sum_i A_i\dot\be_i
=\left(\frac{\dd\bA}{\dd t}\right)_b+\bom\times\bA.
\label{eq:rigid-transport}
\end{equation}
下标 $b$ 表示仅微分体轴分量；它不是另一根独立的物理矢量。取 $\bA=\bom$ 时叉积为零，故角加速度的体轴分量恰为 $\dot\omega_i$；对角动量则一般不能删去叉积。

\begin{physicsbox}[title={物理图像与理解：角动量不动，角速度为何会动？}]
无力矩时 $\bJ$ 在空间固定，但体轴随刚体转动，其体轴分量可以改变。不同主惯量又使 $\bom=\bI^{-1}\bJ$ 通常不沿 $\bJ$。跟踪空间角动量、体轴方向和角速度三个对象，才能解释自由转动，不能把“体分量变化”直接当作外力矩。
\end{physicsbox}


取固连在刚体上的主轴坐标系，原点取在固定点或质心。设主转动惯量为 \(I_1,I_2,I_3\)，角速度分量为 \(\omega_1,\omega_2,\omega_3\)。角动量在体坐标系中为
\[
\bJ=(I_1\omega_1,I_2\omega_2,I_3\omega_3).
\]
在惯性系中，
\[
\dot\bJ_{\rm in}=\bM.
\]
体坐标系随刚体转动，因此同一角动量在惯性系中的变化率等于体坐标系中的变化率加上转动项：
\[
\dot\bJ_{\rm in}
=\dot\bJ_{\rm body}+\bom\times\bJ.
\]
在主轴系中，\(\bJ=\bI\bom\) 的分量为 \(J_a=I_a\omega_a\)，所以
\[
\dot\bJ_{\rm body}
=(I_1\dot\omega_1,\ I_2\dot\omega_2,\ I_3\dot\omega_3).
\]
叉乘项的分量为
\[
\bom\times\bJ=
\begin{pmatrix}
\omega_2 I_3\omega_3-\omega_3 I_2\omega_2\\
\omega_3 I_1\omega_1-\omega_1 I_3\omega_3\\
\omega_1 I_2\omega_2-\omega_2 I_1\omega_1
\end{pmatrix}
=
\begin{pmatrix}
(I_3-I_2)\omega_2\omega_3\\
(I_1-I_3)\omega_3\omega_1\\
(I_2-I_1)\omega_1\omega_2
\end{pmatrix}.
\]
把两式相加并令其等于外力矩 \(\bM=(M_1,M_2,M_3)\)，得到\textbf{欧拉方程}：
\begin{equation}
\boxed{
\begin{aligned}
I_1\dot\omega_1+(I_3-I_2)\omega_2\omega_3&=M_1,\\
I_2\dot\omega_2+(I_1-I_3)\omega_3\omega_1&=M_2,\\
I_3\dot\omega_3+(I_2-I_1)\omega_1\omega_2&=M_3.
\end{aligned}}
\label{eq:euler-equations}
\end{equation}
其中 \(M_a\) 是外力矩在体主轴上的分量。欧拉方程是刚体绕定点或绕质心转动的核心方程；它与质心平动方程
\[
M\dot\bv_C=\bF
\]
合起来，构成刚体动力学的完整方程组。

\begin{remark}
欧拉方程中的非线性项 \((I_3-I_2)\omega_2\omega_3\) 等来自体坐标系的转动，不是“额外的力”。它们描述的是：即使外力矩为零，角动量在体坐标系中也可能看起来在变化。陀螺的进动、章动和自由刚体的复杂翻滚，都来自这些项。
\end{remark}

\subsection{自由刚体}

把外力矩去掉后，空间中的总角动量固定。刚体却还在转，所以从刚体自己的三根轴看，同一个角动量的分量可能变化。

无力矩时，\(\bM=0\)，欧拉方程变为
\begin{equation}
\begin{aligned}
I_1\dot\omega_1+(I_3-I_2)\omega_2\omega_3&=0,\\
I_2\dot\omega_2+(I_1-I_3)\omega_3\omega_1&=0,\\
I_3\dot\omega_3+(I_2-I_1)\omega_1\omega_2&=0.
\end{aligned}
\label{eq:free-euler}
\end{equation}
无力矩时角动量矢量 \(\bJ\) 守恒；又因为没有外力做功，转动动能 \(T\) 也守恒。由此得到两个标量守恒量：角动量大小的平方
\[
J^2=I_1^2\omega_1^2+I_2^2\omega_2^2+I_3^2\omega_3^2
\]
和转动动能
\[
T=\frac12(I_1\omega_1^2+I_2\omega_2^2+I_3\omega_3^2).
\]
这两个守恒量把三维角速度空间中的轨道限制在一条曲线上。

对\textbf{对称自由刚体}，设
\[
I_1=I_2\ne I_3.
\]
由第三式立即得到
\[
I_3\dot\omega_3=0,
\qquad
\omega_3=\Omega=\text{常数}.
\]
前两式化为
\begin{equation}
\dot\omega_1=-n\omega_2,\qquad
\dot\omega_2=n\omega_1,
\qquad
n=\frac{I_3-I_1}{I_1}\Omega.
\label{eq:free-symmetric-omega}
\end{equation}
由第一式求导并代入第二式，得到
\[
\ddot\omega_1+n^2\omega_1=0,
\]
因此 \(\omega_1\) 作简谐振荡，其解为
\begin{equation}
\omega_1=A\cos(nt+\delta),\qquad
\omega_2=A\sin(nt+\delta).
\label{eq:free-symmetric-solution}
\end{equation}
因此 \(\omega_3\) 恒定，横向角速度分量以角频率 \(|n|\) 在体坐标系中旋转。总角速度大小为
\[
\omega=\sqrt{\Omega^2+A^2}=\text{常数}.
\]

在惯性系中，角动量 \(\bJ\) 固定不动。由
\[
\bJ=I_1\bom+(I_3-I_1)\Omega\,\hat{\be}_3
\]
和体轴 \(\hat{\be}_3\) 随刚体转动的关系 \(\dot{\hat{\be}}_3=\bom\times\hat{\be}_3\)，对前式求导：
\[
0=I_1\dot\bom+(I_3-I_1)\Omega\,\bom\times\hat{\be}_3.
\]
又
\[
\bJ\times\bom
=(I_3-I_1)\Omega\,\hat{\be}_3\times\bom
=-(I_3-I_1)\Omega\,\bom\times\hat{\be}_3,
\]
所以
\[
\dot\bom=\frac1{I_1}\bJ\times\bom.
\]
这正是 \(\bom\) 绕固定矢量 \(\bJ\) 的进动。对对称刚体，进动角速度的大小为
\begin{equation}
\Omega_p=\frac{J}{I_1}.
\label{eq:free-precession}
\end{equation}
这就是无力矩对称刚体的\textbf{规则进动}：角速度矢量在空间中绕固定的角动量方向画圆锥，同时在体坐标系中绕对称轴画另一个圆锥。

\begin{figure}[htbp]
\centering
\begin{tikzpicture}[>=Stealth]
\draw[->,black,very thick](0,0)--(0,3.1)node[above]{$\bJ$ 固定};
\draw[->,blue!65!black,very thick](0,0)--({-4/sqrt(13)},{6/sqrt(13)})node[left]{$\be_3$};
\draw[->,red!70!black,very thick](0,0)--({1.6/sqrt(13)},{8/sqrt(13)})node[right]{$\bom$};
\draw[dashed,gray](-1.3,2.55)--(1.3,2.55);
\draw[dashed,gray](0,2.55)ellipse(.6 and .18);
\node at(1.7,2.7){$\Omega_p=J/I_1$};
\end{tikzpicture}
\caption{无力矩对称刚体在包含$\bJ,\bom,\be_3$的瞬时平面内的方向关系，示意参数$I_1=2,I_3=3,\omega_3=1$且横向角速度为1。$\bom$与$\be_3$在$\bJ$两侧，共同绕它进动；箭长只用于显示方向。}
\label{fig:free-top}
\end{figure}

\subsection{主轴稳定性与网球拍定理}

尝试分别绕网球拍的三个主轴轻轻抛转。要解释哪一种最容易翻转，可给每一种匀转加一点横向扰动，再看扰动是否变大。

考虑刚体绕某一主轴的匀速转动，例如绕主轴 3：
\[
\omega_1=\omega_2=0,\qquad \omega_3=\Omega.
\]
给它一个很小的扰动 \(\delta\omega_1,\delta\omega_2\)。把 \(\omega_3=\Omega\) 代入欧拉方程并保留一阶项：
\[
I_1\dot{\delta\omega}_1+(I_3-I_2)\Omega\,\delta\omega_2=0,
\qquad
I_2\dot{\delta\omega}_2+(I_1-I_3)\Omega\,\delta\omega_1=0.
\]
由第一式解出 \(\delta\omega_2\)，再代入第二式，得到
\begin{equation}
\ddot{\delta\omega}_1
+\frac{(I_3-I_2)(I_3-I_1)}{I_1I_2}\Omega^2
\delta\omega_1=0.
\label{eq:principal-stability}
\end{equation}
稳定性取决于括号的符号：
\[
(I_3-I_2)(I_3-I_1)>0.
\]
这等价于 \(I_3\) 是最大或最小的主转动惯量。因此：

\begin{keybox}
绕最大主轴的转动和绕最小主轴的转动是稳定的；绕中间主轴的转动是不稳定的。这就是\textbf{网球拍定理}，也称\textbf{中间轴定理}。
\end{keybox}

把网球拍抛向空中，使它绕中间轴旋转，拍面会突然翻转。这不是因为“手抖”，而是中间轴转动对微小扰动指数敏感。理想自由刚体的翻转可由这个模型解释，主动改变形状的运动员则需要变惯量模型。



\subsection{稳定性结论的范围与守恒量检验}
还要区分转轴附近的小扰动与完整姿态的误差。即使两个物体一直绕同一根主轴转，微小转速差也会让它们逐渐错开相位。

上述稳定指体轴角速度靠近一族主轴匀转，或等价地忽略沿旋转相位的累积偏移，并不表示任意小转速误差都产生永远很小的绝对姿态差。设严格 $I_1<I_2<I_3$，在固定角动量大小的球面上
\[
T=\tfrac12(J_1^2/I_1+J_2^2/I_2+J_3^2/I_3).
\]
绕最小惯量轴时能量取最大，绕最大惯量轴时取最小，中间轴为鞍点。更直接地，两守恒量给出
\[
J^2-2I_1T=I_2(I_2-I_1)\omega_2^2+I_3(I_3-I_1)\omega_3^2,
\]
\[
2I_3T-J^2=I_1(I_3-I_1)\omega_1^2+I_2(I_3-I_2)\omega_2^2.
\]
右端正系数控制相应横向分量，配合能量可证明两极值轴附近的角速度稳定，而不只是依赖线性振荡的猜测。若有重复主惯量，严格中间轴分类退化，须另用对称刚体解；若加入耗散和内部运动，也不再是此无力矩刚体问题。

\subsection{庞索作图：接触平面从哪里来？}
庞索作图把自由转动变成一个几何问题。先画出由动能确定的椭球，再由固定角动量找出与椭球相切的平面。

为几何解释无力矩运动，固定质心作为平动已扣除的原点，令 $T>0$ 为恒定转动动能，并在辅助坐标空间定义
\[
\boldsymbol x=\frac{\bom}{\sqrt{2T}},\qquad \boldsymbol x^{\mathsf T}\bI\boldsymbol x=1.
\]
这里辅助坐标的单位不是长度。该随体转动的椭球在 $\boldsymbol x$ 处的法线为 $\bI\boldsymbol x=\bJ/\sqrt{2T}$，所以切平面满足
\[
\bJ\cdot{\boldsymbol X}=\sqrt{2T}.
\]
$\bJ,T$ 均固定，因此这是空间中的不变平面。刚体在当前接触位置的速度正比于 $\bom\times\boldsymbol x=0$，于是椭球在该平面上无滑动滚动。辅助椭球上的接触迹称极迹，平面上的轨迹称空间极迹。

这幅构造依赖理想无力矩、正定惯量和非零能量；不能把没有能量标度的 $\bom$ 直接说成单位惯量椭球上的点。对称刚体的极迹是圆，三主惯量不等时则由能量面与角动量面共同限制。真正的姿态仍需积分 $\dot R=R[\bom_b]_\times$ 重构。


\section{对称重刚体与陀螺}

转得很快的陀螺受到重力时，轴常沿圆锥面缓慢绕行。我们先确定重力矩怎样改变角动量，再分辨绕行与倾角的起伏。

\subsection{拉格朗日陀螺}

这个模型把支点固定，且陀螺绕自身对称轴的惯量与横向惯量可以不同。先用角度分别记下轴的朝向和绕自身轴的转动。

考虑一个对称刚体，绕其对称轴的质量分布对称，固定点 \(O\) 在对称轴上，质心 \(C\) 到 \(O\) 的距离为 \(l\)，质量为 \(m\)。这就是\textbf{拉格朗日陀螺}。

取欧拉角 \(\phi,\theta,\psi\)，对称轴为体轴 3。质心位置为
\[
\bR_C=l(\sin\theta\sin\phi,-\sin\theta\cos\phi,\cos\theta).
\]
因为固定点 \(O\) 不动，总动能就是绕 \(O\) 的转动动能。这里 \(I_1,I_3\) 是对固定点 \(O\) 的主转动惯量，而不是对质心的主转动惯量。于是
\[
T=\frac12I_1(\omega_1^2+\omega_2^2)+\frac12I_3\omega_3^2.
\]
若改用质心分解，应写成
\[
T=\frac12mv_C^2+\frac12I_{C1}(\omega_1^2+\omega_2^2)+\frac12I_{C3}\omega_3^2,
\]
其中
\[
v_C^2=l^2(\dot\theta^2+\dot\phi^2\sin^2\theta).
\]
由平行轴定理 \(I_1=I_{C1}+ml^2\) 可知两式完全等价。
利用本章的欧拉角运动学，把式 \eqref{eq:euler-omega1} 和 \eqref{eq:euler-omega2} 平方相加：
\[
\begin{aligned}
\omega_1^2+\omega_2^2
&=\dot\phi^2\sin^2\theta(\sin^2\psi+\cos^2\psi)
+\dot\theta^2(\cos^2\psi+\sin^2\psi)\\
&\quad+2\dot\phi\dot\theta\sin\theta
(\sin\psi\cos\psi-\cos\psi\sin\psi)\\
&=\dot\theta^2+\dot\phi^2\sin^2\theta,
\end{aligned}
\]
交叉项正好抵消；同时
\[
\omega_3=\dot\psi+\dot\phi\cos\theta.
\]
因此
\begin{equation}
T=\frac12I_1(\dot\theta^2+\dot\phi^2\sin^2\theta)
+\frac12I_3(\dot\psi+\dot\phi\cos\theta)^2.
\label{eq:lagrange-top-T}
\end{equation}
势能为
\begin{equation}
V=mgl\cos\theta.
\label{eq:lagrange-top-V}
\end{equation}
拉格朗日量为
\begin{equation}
L=\frac12I_1(\dot\theta^2+\dot\phi^2\sin^2\theta)
+\frac12I_3(\dot\psi+\dot\phi\cos\theta)^2
-mgl\cos\theta.
\label{eq:lagrange-top-L}
\end{equation}

\subsection{守恒量与有效势}

利用两个循环角的守恒动量，先求出它们的速度，再把剩下的倾角运动写成一维问题。有效势就是用来判断倾角能到哪里的一条曲线。

\(L\) 不显含 \(\phi\) 和 \(\psi\)，因此有两个循环坐标和两个守恒的广义动量：
\begin{equation}
p_\psi=\frac{\partial L}{\partial\dot\psi}
=I_3(\dot\psi+\dot\phi\cos\theta)=I_3\omega_3,
\label{eq:p-psi}
\end{equation}
\begin{equation}
p_\phi=\frac{\partial L}{\partial\dot\phi}
=I_1\dot\phi\sin^2\theta
+I_3(\dot\psi+\dot\phi\cos\theta)\cos\theta.
\label{eq:p-phi}
\end{equation}
用 \(p_\psi\) 消去 \(\dot\psi\)，得
\begin{equation}
p_\phi=I_1\dot\phi\sin^2\theta+p_\psi\cos\theta.
\label{eq:p-phi-simplified}
\end{equation}
由此可解出进动角速度
\begin{equation}
\dot\phi=\frac{p_\phi-p_\psi\cos\theta}{I_1\sin^2\theta}.
\label{eq:phi-dot}
\end{equation}
由式 \eqref{eq:p-psi} 还可解出
\[
\dot\psi=\frac{p_\psi}{I_3}-\dot\phi\cos\theta.
\]
把 \(\dot\phi,\dot\psi\) 代回转动动能，得到
\[
\begin{aligned}
T
&=\frac12I_1\dot\theta^2
+\frac12I_1\sin^2\theta
\left(\frac{p_\phi-p_\psi\cos\theta}{I_1\sin^2\theta}\right)^2
+\frac{p_\psi^2}{2I_3}\\
&=\frac12I_1\dot\theta^2
+\frac{(p_\phi-p_\psi\cos\theta)^2}{2I_1\sin^2\theta}
+\frac{p_\psi^2}{2I_3}.
\end{aligned}
\]
再加上势能 \(V=mgl\cos\theta\)，能量守恒给出
\begin{equation}
E=\frac12I_1\dot\theta^2
+\frac{(p_\phi-p_\psi\cos\theta)^2}{2I_1\sin^2\theta}
+\frac{p_\psi^2}{2I_3}
+mgl\cos\theta.
\label{eq:top-energy}
\end{equation}
定义\textbf{有效势}
\begin{equation}
V_{\rm eff}(\theta)
=\frac{(p_\phi-p_\psi\cos\theta)^2}{2I_1\sin^2\theta}
+mgl\cos\theta.
\label{eq:top-effective-potential}
\end{equation}
则
\begin{equation}
\frac12I_1\dot\theta^2+V_{\rm eff}(\theta)
=E-\frac{p_\psi^2}{2I_3}
=\text{常数}.
\label{eq:top-1d}
\end{equation}
这样，三维刚体运动被约化为一维的 \(\theta\) 运动。\(\theta(t)\) 描述\textbf{章动}，\(\phi(t)\) 描述\textbf{进动}，\(\psi(t)\) 描述自转。

\begin{figure}[htbp]
\centering
\begin{tikzpicture}
\begin{axis}[width=.8\textwidth,height=5.7cm,axis lines=left,xlabel={$\theta/\mathrm{rad}$},ylabel={$V_{\rm eff}/E_*$},xmin=.15,xmax=2.75,ymin=0,ymax=2,xtick={.5,1,1.5,2,2.5},ytick={0,.5,1,1.5,2},domain=.16:2.74,samples=300]
\addplot[blue!65!black,very thick]{.5*(1.2-2*cos(deg(x)))^2/sin(deg(x))^2+.3*cos(deg(x))};
\addplot[red!70!black,dashed,domain=.2:2.7]{.7};
\node[red!70!black]at(axis cs:2.35,.8){$E'=0.7E_*$};
\addplot[only marks,mark=*,black]coordinates{(0.5563710705352931,.7)(1.5611529346033528,.7)};
\draw[dashed,gray](axis cs:0.5563710705352931,0)--(axis cs:0.5563710705352931,.7);
\draw[dashed,gray](axis cs:1.5611529346033528,0)--(axis cs:1.5611529346033528,.7);
\node[above right]at(axis cs:0.5563710705352931,.75){$\theta_{\min}$};
\node[above left]at(axis cs:1.5611529346033528,.75){$\theta_{\max}$};
\end{axis}
\end{tikzpicture}
\caption{给定教学尺度$I_*,\omega_*,E_*=I_*\omega_*^2$，取$I_1/I_*=1,p_\psi/(I_*\omega_*)=2,p_\phi/(I_*\omega_*)=1.2,mgl/E_*=0.3$。水平能量线与有效势的两个交点给出章动转折点。}
\label{fig:top-effective-potential}
\end{figure}

对有效势的正则极小值，略高能量使 $\theta$ 在两个转折角间章动。进动角速度由式 \eqref{eq:phi-dot} 决定，可能变号，轴的轨迹可出现回绕；不能一律说 $\phi$ 持续增加。一般 $p_\phi\ne\pm p_\psi$ 时两端势垒发散，若恰好相等则一端可相消，应回到非奇异姿态坐标判断极轴运动。

\subsection{约化方程与稳恒进动}
若陀螺轴始终保持同一倾角，倾角方程中的各项必须彼此抵消。满足这个条件的解仍可在水平面方向绕行，称为稳恒进动。

约化时固定两个循环动量后再微分有效势；一般的偏勒让德变换见式 \eqref{eq:routhian}。
能量给出容许区间，但在转折点不应除以 $\dot\theta$ 推导运动方程。固定 $p_\phi,p_\psi$，对循环坐标作罗斯约化。用 $\dot\psi=p_\psi/I_3-\dot\phi\cos\theta$ 改写被减去的两项：
\[
p_\phi\dot\phi+p_\psi\dot\psi
=(p_\phi-p_\psi\cos\theta)\dot\phi+p_\psi^2/I_3
=I_1\sin^2\theta\,\dot\phi^2+p_\psi^2/I_3 .
\]
从原 $L$ 减去它，两个循环速度的动能项都变号。再代入 $\dot\phi$ 的动量表达式，得到
\[
\mathcal R=\tfrac12I_1\dot\theta^2-V_{\rm eff}(\theta)-\frac{p_\psi^2}{2I_3}.
\]
最后常数不影响 $\theta$ 方程，故所有正则点（含章动转折点）均有
\[
I_1\ddot\theta=-V_{\rm eff}'(\theta).
\]
这说明消去\emph{守恒动量}时不能只把速度表达式代回原 $L$ 再当作同一变分问题。

稳恒进动要求 $0<\theta_0<\pi$、$\dot\theta=0$、$\dot\phi=\Omega$，并记恒定的轴向角动量 $h=p_\psi=I_3\omega_3$，$\omega_3=\dot\psi+\Omega\cos\theta_0$。从原拉格朗日量取 $\theta$ 方程，
\[
I_1\ddot\theta=I_1\dot\phi^2\sin\theta\cos\theta-h\dot\phi\sin\theta+mgl\sin\theta,
\]
除以非零 $\sin\theta_0$ 得
\begin{equation}
I_1\Omega^2\cos\theta_0-h\Omega+mgl=0.
\label{eq:rigid-steady}
\end{equation}
当 $c=\cos\theta_0\ne0$ 时
\[
\Omega_\pm=\frac{h\pm\sqrt{h^2-4I_1mgl c}}{2I_1c}.
\]
给定倾角与轴向自转，先检查判别式；$c=0$ 时直接回原式得 $\Omega=mgl/h$（$h\ne0$），不要使用零分母公式。稳恒运动还须选择相容的初始动量 $p_\phi=I_1\Omega\sin^2\theta_0+h\cos\theta_0$；只给自转并不能保证从任意静置倾角释放后无章动。

\subsection{慢支、快支及章动频率}
同样的自转条件下，稳恒进动有时有两个速率。检验它们附近的倾角振动时，应固定各自的守恒动量，再让倾角略微改变。

设 $h>0$ 且 $|4I_1mgl c/h^2|\ll1$。与弱重力连续相接的一支可以写成无数值相消的形式
\begin{equation}
\Omega_{\rm slow}=\frac{2mgl}{h+\sqrt{h^2-4I_1mgl c}}
\simeq\frac{mgl}{h}=\frac{mgl}{I_3\omega_3}.
\label{eq:fast-top-precession}
\end{equation}
另一支在 $c\ne0$ 时约为 $h/(I_1c)$。所谓“快陀螺公式”通常指\emph{轴向自转快、进动慢}的分支，不能用它替代所有根。若角动量近似为 $h\be_3$，重力矩 $mgl\sin\theta$ 也直接给 $\Omega h\sin\theta$，提供独立的力矩核验。这里使用 $\omega_3$ 而非未经说明的 $\dot\psi$；慢支中二者近似相同需要 $|\Omega\cos\theta|\ll|\omega_3|$。

固定两个动量，记 $w(\theta)=(p_\phi-h\cos\theta)/\sin\theta$，则 $V_{\rm eff}=w^2/(2I_1)+mgl\cos\theta$。依次求导得 $w'=h-w\cot\theta$、$w''=-w'\cot\theta+w/\sin^2\theta$。在稳恒角有 $w=I_1\Omega\sin\theta_0$、$w'=h-I_1\Omega c$。代入 $V_{\rm eff}''=[(w')^2+ww'']/I_1-mgl c$，并在第二步用式 \eqref{eq:rigid-steady} 消去 $mgl$，得
\begin{align*}
V_{\rm eff}''(\theta_0)
&=h^2/I_1-3h\Omega c+I_1\Omega^2(1+2c^2)-mgl c\\
&=\frac{(h-2I_1\Omega c)^2}{I_1}+I_1\Omega^2\sin^2\theta_0.
\end{align*}
因此小章动满足
\begin{equation}
\omega_n^2=\frac{V_{\rm eff}''(\theta_0)}{I_1}
=\left(\frac h{I_1}-2\Omega\cos\theta_0\right)^2
+\Omega^2\sin^2\theta_0.
\label{eq:nutation-frequency}
\end{equation}
对 $mgl>0$ 的非竖直稳恒支该值为正，给出约化角度的稳定性；不等同于所有角变量差值都永远不积累。慢进动、自转主导时 $\omega_n\simeq |h|/I_1$。频率高并不自动意味着初始章动幅度小，幅度仍由相对有效势极小的超额能量决定。

\begin{example}[同一自转对应快、慢两支]
取教学单位使 $I_1=1,I_3=0.6,mgl=0.3,h=2$，且 $\theta_0=\pi/3$。轴向自转为 $\omega_3=10/3$，稳恒根是
\[
\Omega_\pm=2\pm\sqrt{3.4}.
\]
慢根约 $0.15609$，快根约 $3.84391$；简单慢支近似 $mgl/h=0.15$。两支分别要求不同的 $p_\phi=1+3\Omega/4$。代入式 \eqref{eq:nutation-frequency}，均为正，但数值不同。这个例子说明“给出转子转速”尚未指定进动状态。

代入 $\cos\theta_0=1/2$ 得二次式 $\Omega^2/2-2\Omega+0.3=0$。每个根对应的 $p_\phi$ 由 $I_1\Omega\sin^2\theta_0+h\cos\theta_0$ 求得，约为 $1.11707$、$3.88293$。分别在各自固定动量下求曲率，$\omega_n^2$ 约为 $3.41827$、$14.48173$。
\end{example}

\begin{example}[高速重陀螺：先检查近似再报数值]
取理想教学模型 $m=0.5\,\mathrm{kg},l=0.05\,\mathrm m,g=9.8\,\mathrm{m\,s^{-2}}$，固定点惯量 $I_1=0.03,I_3=0.02\,\mathrm{kg\,m^2}$，轴向 $\omega_3=3000\,\mathrm{rad\,s^{-1}}$，$\theta_0=\pi/3$，并令初态位于慢稳恒支。此时 $h=60\,\mathrm{kg\,m^2\,s^{-1}}$，
\[
\frac{4I_1mgl\cos\theta_0}{h^2}=4.08\times10^{-6},\qquad
\Omega_{\rm slow}\simeq4.0833\times10^{-3}\,\mathrm{rad\,s^{-1}}.
\]
相应进动周期约 $1.539\times10^3\,\mathrm s$，即25.6分钟；章动频率近似 $h/I_1=2000\,\mathrm{rad\,s^{-1}}$。这些是给定模型与相容初态的计算，不是某个实际陀螺仪的规格或性能。
\end{example}

\subsection{竖直睡眠陀螺须单独线性化}\label{sec:sleeping-top}
陀螺轴竖直时，绕竖直方向的方位角已不能标出一个唯一姿态。改用轴尖在水平面上的两个小偏移，就能直接判断扰动是否增长。

$\theta=0$ 的“睡眠”解是质心在支点上方、对称轴竖直向上的自转。欧拉角在此奇异，不能在约去 $\sin\theta$ 后直接套用。以轴向单位矢量 $\be_3\simeq(u,v,1)$ 代替角坐标，保留横向一阶项。由
\[
\bJ=I_1\be_3\times\dot\be_3+h\be_3,\qquad
\bM=mgl\,\hat z\times\be_3,
\]
其中第一式由 $\dot{\be}_3=\bom\times\be_3$ 解出横向角速度得到。在线性阶，$\bJ_\perp=(-I_1\dot v+hu,I_1\dot u+hv)$，$\bM_\perp=mgl(-v,u)$。令 $\dot{\bJ}_\perp=\bM_\perp$，整理为
\[
I_1\ddot u+h\dot v-mgl\,u=0,\qquad
I_1\ddot v-h\dot u-mgl\,v=0.
\]
令 $w=u+iv$，有 $I_1\ddot w-ih\dot w-mgl\,w=0$。试 $w\propto e^{i\sigma t}$，得到
\[
\sigma_\pm=\frac{h\pm\sqrt{h^2-4I_1mgl}}{2I_1}.
\]
当 $h^2>4I_1mgl$，两个频率实且不同，横向线性扰动有界；低于阈值则有指数增长。等号处重根可产生 $t e^{i\sigma t}$ 项，不能仅因频率实就宣称线性稳定。这里是含陀螺项的动态稳定化，正是第四章简单“只看势能海森矩阵”判据之外的情况。

\begin{modernbox}[title={进一步探索：用什么模型解释陀螺传感器？}]
本节只建立固定支点机械转子的模型。振动式陀螺需用科里奥利耦合，光学陀螺需用光路相位；它们可测同一刚体的角速度，但不共享这里的慢进动公式。进一步学习机械陀螺可参看章末Tong与MIT讲义，需掌握循环坐标、线性稳定性与旋转系输运。
\end{modernbox}


\section{刚体碰撞与冲击}

撞击时间很短，刚体的位置和姿态几乎来不及变，速度却可能突然跳变。先用冲量求这次跳变，再用接触条件决定冲量大小。

\subsection{冲量与冲量矩}

同样大小的冲量打在球心和打在球的上缘，会造成不同的转动变化。要同时写出动量的变化和冲量关于质心的力矩。

当碰撞时间极短、碰撞力极大时，普通力的冲量可以忽略，只保留碰撞冲量。为避免与角动量 \(\bJ\) 混淆，本节用 \(\boldsymbol{\mathcal J}\) 表示冲量、$\boldsymbol{\mathcal L}_C$表示对质心的冲量矩：
\begin{equation}
\boldsymbol{\mathcal J}=\int_{t_1}^{t_2}\bF\dd t,
\qquad
\boldsymbol{\mathcal L}_C=\int_{t_1}^{t_2}\bM_C\dd t.
\label{eq:impulse}
\end{equation}
碰撞前后，刚体的动量发生突变：
\begin{equation}
M(\bv_C'-\bv_C)=\boldsymbol{\mathcal J},
\label{eq:impulse-momentum}
\end{equation}
角动量也发生突变：
\begin{equation}
\bI_C(\bom'-\bom)=\boldsymbol{\mathcal L}_C.
\label{eq:impulse-angular}
\end{equation}
这里假设碰撞瞬间刚体的位形不变，因此惯量张量不变。

\subsection{刚体碰撞方程}

碰撞时恢复系数比较的是两接触点在法线方向的接近、分离速率。每个接触点的速度既有质心平移，也有物体转动的贡献。

考虑自由刚体上一点 \(P\) 与光滑静止固定面单点碰撞。忽略碰撞期间姿态变化及重力等有限力冲量，没有其他支承冲量；接触只给法向冲量，且三维惯量正定。若有铰支、多点接触或摩擦，需联立相应冲量约束。设接触面法线为 \(\hat{\boldsymbol{n}}\)，碰撞前接触点法向速度为 \(v_{Pn}\)，碰撞后为 \(v_{Pn}'\)。恢复系数定义为
\begin{equation}
e_r=-\frac{v_{Pn}'}{v_{Pn}},
\qquad 0\le e_r\le1.
\label{eq:restitution}
\end{equation}
\(e_r=0\) 只使碰后该接触点法向速度为零，不意味着刚体停止或永久黏附；在此无摩擦单冲量模型中，\(e_r=1\) 保持总动能。以下取 \(v_n=|v_{Pn}|>0\) 为接近速度，并取法向冲量 \(J_n\) 沿离开固定面的方向为正。

碰撞过程中，接触点的法向冲量 \(J_n\) 同时改变质心速度和角速度。由式 \eqref{eq:impulse-momentum} 和 \eqref{eq:impulse-angular}，
\[
\bv_C'=\bv_C+\frac{J_n}{M}\hat{\boldsymbol{n}},
\]
\[
\bom'=\bom+\bI_C^{-1}\left[(\br_P-\br_C)\times J_n\hat{\boldsymbol{n}}\right].
\]
接触点速度
\[
\bv_P'=\bv_C'+\bom'\times(\br_P-\br_C).
\]
令 \(\boldsymbol{b}=\br_P-\br_C\)。把 \(\bv_C'\) 和 \(\bom'\) 代入，取法向分量：
\[
v_{Pn}'=v_{Pn}+J_n\left[
\frac1M+\hat{\boldsymbol{n}}\cdot
\left(\bI_C^{-1}(\boldsymbol{b}\times\hat{\boldsymbol{n}})\times\boldsymbol{b}\right)
\right].
\]
利用标量三重积恒等式，方括号中的转动项可写成
\[
(\boldsymbol{b}\times\hat{\boldsymbol{n}})\cdot
\bI_C^{-1}(\boldsymbol{b}\times\hat{\boldsymbol{n}}).
\]
定义有效质量 \(M_{\rm eff}\) 满足
\[
\frac1{M_{\rm eff}}
=\frac1M
+(\boldsymbol{b}\times\hat{\boldsymbol{n}})\cdot
\bI_C^{-1}(\boldsymbol{b}\times\hat{\boldsymbol{n}}),
\]
则
\[
v_{Pn}'=v_{Pn}+\frac{J_n}{M_{\rm eff}}.
\]
代入恢复系数 \(v_{Pn}'=-e_r v_{Pn}\)，得到
\begin{equation}
J_n=-\frac{(1+e_r)v_{Pn}}{1/M
+(\boldsymbol{b}\times\hat{\boldsymbol{n}})\cdot
\bI_C^{-1}(\boldsymbol{b}\times\hat{\boldsymbol{n}})}.
\label{eq:normal-impulse}
\end{equation}
当 \(v_{Pn}<0\)（接触点向固定面接近）时，\(J_n>0\)。对一般刚体，结果依赖于撞击点位置和惯量张量；这正是“击球点”会影响反弹方向的原因。

\subsection{一般单点模型的能量损失}
接触点受冲量后改变速度的难易，由质量、转动惯量和力臂共同决定。把这三方面归入一个等效质量，能把能量损失的计算写得更简洁。

记初接触法向速度 $v_{Pn}<0$，$\boldsymbol b=\br_P-\br_C$。展开碰后平动、转动动能，线性冲量项合为 $J_nv_{Pn}$，二次项合为 $J_n^2/(2M_{\rm eff})$：
\[
T'-T=J_nv_{Pn}+\frac{J_n^2}{2M_{\rm eff}}.
\]
代入 $J_n=-(1+e_r)M_{\rm eff}v_{Pn}$，得
\begin{equation}
T-T'=\frac12M_{\rm eff}(1-e_r^2)v_{Pn}^2.
\label{eq:collision-energy-loss}
\end{equation}
仅在撞击作用线通过质心等无转动效应时，才可把 $M_{\rm eff}$ 换成总质量。固定碰前状态和接触几何时，$e_r=0$ 在本法向模型中损失最大；切向摩擦和多点碰撞的能量不能只由这个恢复系数决定。

\begin{example}[倾斜杆单端碰撞：有效质量不是总质量]
均匀细杆质量 $m$、长度 $l$，在竖直平面内自由下落，杆与水平夹角 $60^\circ$，较低的右端先撞光滑水平面。碰前质心速度为 $-V\hat y$、角速度零，恢复系数 $e_r=1/2$，没有其他瞬时约束冲量。采用平面转动惯量 $I_C=ml^2/12$，接触臂水平分量 $b_x=l/4$，于是
\[
\frac1{m_{\rm eff}}=\frac1m+\frac{(l/4)^2}{ml^2/12}=\frac7{4m},\quad
J_n=(1+e_r)m_{\rm eff}V=\frac67mV.
\]
碰后质心与角速度为
\[
v_{Cy}'=-V+J_n/m=-V/7,\qquad
\omega_z'=\frac{(l/4)J_n}{I_C}=\frac{18V}{7l}.
\]
接触点速度 $v_{Pn}'=-V/7+(18V/7l)(l/4)=V/2$，恰满足恢复系数。损失为 $3mV^2/14$，碰后平动与转动动能之和为 $2mV^2/7$，与初动能相减一致。杆为理想细杆使轴向惯量退化，但这里只反演非零的\emph{平面}惯量，并未使用奇异三维张量的逆。
\end{example}

\begin{figure}[htbp]
\centering
\begin{tikzpicture}[>=Stealth]
\draw[thick](-.5,0)--(3,0);
\draw[very thick,blue!65!black](0,2.598)--(1.5,0);
\fill(.75,1.299)circle(2pt)node[left]{$C$};
\fill(1.5,0)circle(2pt)node[below]{$P$};
\draw[->,red!70!black,thick](.75,1.299)--(.75,.55)node[left]{$V$};
\draw[->,green!45!black,very thick](1.5,0)--(1.5,1.05)node[right]{$J_n$};
\draw[dashed](.75,1.299)--(1.5,1.299);
\draw[<->](.75,1.65)--(1.5,1.65)node[midway,above]{$l/4$};
\node at(2.3,2.2){$60^\circ$ 倾斜杆};
\end{tikzpicture}
\caption{低端单点冲击，法向冲量作用线距质心 $l/4$。图中红箭头为碰前质心速度，绿箭头为冲量；其余杆点尚未接触地面。}
\label{fig:rigid-impact}
\end{figure}


\begin{remark}
刚体碰撞模型把碰撞时间理想化为零，把接触力理想化为冲量。这在高频、短时碰撞中很好用，但当接触时间与波传播时间相当时，必须考虑弹性波和局部变形。工程中的落锤、撞击、机器人足端触地，常需要更精细的接触模型。
\end{remark}

\section{姿态重构与数值表示}
求出角速度后，还要知道物体每个时刻到底朝哪里。下面分别用矩阵和四元数把角速度积分成姿态。

\subsection{转动矩阵的左右乘法}
先确定角速度的三个分量是沿空间轴还是沿物体轴量出的。按这个约定，姿态的微小更新才知道该乘在原矩阵的哪一侧。

前文欧拉方程求出的是体轴角速度，还要由姿态运动学还原刚体方向。令 $R$ 把体分量变成空间分量，$R^{\mathsf T}R=I,\det R=1$。对固定物质坐标 $\rho_b$，$\rho_S=R\rho_b$。微分得
\[
\dot\rho_S=\dot RR^{\mathsf T}\rho_S=\bom_S\times\rho_S.
\]
又因对 $RR^{\mathsf T}=I$ 求导可知 $\dot RR^{\mathsf T}$ 反对称，故
\begin{equation}
\dot R=[\bom_S]_\times R=R[\bom_b]_\times,\qquad \bom_S=R\bom_b.
\label{eq:rotation-matrix-kinematics}
\end{equation}
\begin{equation}
[\bom]_\times=\begin{pmatrix}0&-\omega_3&\omega_2\\\omega_3&0&-\omega_1\\-\omega_2&\omega_1&0\end{pmatrix}.
\label{eq:rotation-matrix-body}
\end{equation}
左右乘法对应不同的分量约定，不能只换下标而不换乘积次序。一般运动另加 $\br_C(t)$，不会改变以上姿态方程。

\subsection{单位四元数：给出乘法与转动约定}
四元数用四个数表示三维转动，其中单位长度条件少掉一个自由度。先学乘法，再看它如何把物体上的方向转到空间中。

为避免欧拉角坐标奇异，取标量在前的单位四元数 $q=(q_0,\boldsymbol q)$。本书采用哈密顿乘法
\[
(a,\bu)\otimes(b,\bv)=(ab-\bu\cdot\bv,\;a\bv+b\bu+\bu\times\bv),
\]
共轭 $q^*=(q_0,-\boldsymbol q)$，$q_0^2+|\boldsymbol q|^2=1$。主动转动规定为
\[
(0,\rho_S)=q\otimes(0,\rho_b)\otimes q^*.
\]
展开乘法即得
\begin{equation}
R(q)=(q_0^2-|\boldsymbol q|^2)I+2\boldsymbol q\boldsymbol q^{\mathsf T}
+2q_0[\boldsymbol q]_\times.
\label{eq:rigid-quaternion-matrix}
\end{equation}
轴角表示为 $q=(\cos(\theta/2),\hat e\sin(\theta/2))$，$q$ 与 $-q$ 给同一个 $R$。单位四元数覆盖所有姿态，但并非一一参数化；没有欧拉角的局部奇异，不代表可以忽略单位长度约束和双覆盖。

在当前乘法约定下，体角速度运动学为
\begin{equation}
\dot q=\tfrac12q\otimes(0,\bom_b),
\qquad
\text{或 }\dot q=\tfrac12(0,\bom_S)\otimes q.
\label{eq:rigid-quaternion-rate}
\end{equation}
可由一小步体轴转动 $q(t+\dd t)=q(t)\otimes(1,\bom_b\dd t/2)+O(\dd t^2)$ 推出；标量、矢量部分相乘也直接给 $\dd(|q|^2)/\dd t=0$。

\begin{example}[恒定体轴自转与双覆盖]
$\bom_b=\Omega\hat e_3,q(0)=(1,0,0,0)$ 时，
\[
q(t)=(\cos(\Omega t/2),0,0,\sin(\Omega t/2)),\quad R(t)=R_z(\Omega t).
\]
$t=2\pi/|\Omega|$ 时四元数变为 $(-1,0,0,0)$，姿态却回到初始。比较两次数值姿态应用 $R_1^{\mathsf T}R_2$ 或考虑四元数正负等价，不应把 $q$ 与 $-q$ 的欧氏差误判为物体翻转。
\end{example}

\subsection{保持姿态约束不等于保持动力学能量}
计算机会把连续转动拆成许多小步。每步都要保持三根姿态轴相互垂直，而能量是否正确，还要另外检查角速度的计算。

显式欧拉法步 $R_{n+1}=R_n(I+hW_n)$、$W_n=[\bom_{b,n}]_\times$，即使 $R_n$ 正交，也有
\[
R_{n+1}^{\mathsf T}R_{n+1}=I+h^2W_n^{\mathsf T}W_n,
\]
所以逐步漂离正交群。可用指数更新
\[
R_{n+1}=R_n\exp(h[\bom_{b,n+1/2}]_\times),
\]
每步保持正常转动；若中点角速度足够准确可获得二阶精度，恒定角速度时该步精确。Rodrigues公式为
\[
\exp(\vartheta[\hat e]_\times)=I+\sin\vartheta[\hat e]_\times
+(1-\cos\vartheta)[\hat e]_\times^2.
\]
四元数也可用对应单位增量复合。标准RK4或归一化修正仍须检查误差；仅保持 $R^{\mathsf T}R=I$ 或 $|q|=1$ 不保证能量、空间角动量或长期相位准确。

\begin{modernbox}[title={进一步探索：几何相位与模型边界}]
无力矩非对称刚体的体角动量完成闭合路径后，空间姿态未必复原，姿态重构可分出与动力学及闭合路径有关的部分。进一步研究需旋转群、微分方程和约化力学；这里不把它直接等同于量子贝里相位。猫翻身还需要内部形状变化：对正定惯量的\emph{单一}刚体，$\bJ=0$ 意味着 $\bom=0$，不能自行翻身。多刚体内部运动则是不同系统。
\end{modernbox}


\section{进一步探索：平台与反作用轮}
航天器能通过内部电机加快一只飞轮，从而让平台向相反方向转。把轮和平台当作一个系统，能同时看清角动量和能量各流向哪里。

\subsection{单轴角动量与内部电机功率}
轮相对于平台的转速与轮相对于空间的转速不同。写总角动量时用绝对转速，算电机在轮与平台之间的功率时用相对转速。

忽略外力矩，平台与轴对称反作用轮绕共同固定方向的主轴转动。平台惯量 $I_b$ 不含轮；轮的轴向惯量为 $I_w$。用 $\Omega$ 表示平台角速度，$\nu$ 表示轮相对平台的角速度，则轮的绝对角速度是 $\Omega+\nu$。总角动量与动能为
\[
J=(I_b+I_w)\Omega+I_w\nu,\qquad
T=\tfrac12I_b\Omega^2+\tfrac12I_w(\Omega+\nu)^2.
\]
因电机只给内力矩，
\begin{equation}
(I_b+I_w)\dot\Omega+I_w\dot\nu=0,\qquad
\Delta\Omega=-\frac{I_w}{I_b+I_w}\Delta\nu.
\label{eq:rigid-wheel}
\end{equation}
若电机给轮的力矩为 $\tau$，轮与平台分别满足
\[
I_w(\dot\Omega+\dot\nu)=\tau,\quad I_b\dot\Omega=-\tau,
\quad \dot T=\tau\nu.
\]
因此零外力矩不等于机械能不变；电源可把内部能转成机械能。$\nu$ 为相对转速也是分母包含 $I_w$ 的原因。

\begin{example}[加速再制动：转速复原，姿态未复原]
取教学参数 $I_b=9,I_w=1\,\mathrm{kg\,m^2}$，初始静止。轮相对转速在2秒内匀增到 $10\,\mathrm{rad\,s^{-1}}$，则平台末角速度为 $-1\,\mathrm{rad\,s^{-1}}$、平台转角为 $-1\,\mathrm{rad}$。所需电机力矩 $\tau=4.5\,\mathrm{N\,m}$，输入机械能
\[
\int_0^2\tau\nu\,\dd t=45\,\mathrm J
=\tfrac12I_b(-1)^2+\tfrac12I_w(9)^2.
\]
随后用相反的恒定力矩，2秒内将相对转速降为零。平台角速度也回零，却又转过 $-1\,\mathrm{rad}$，总转角为 $-2\,\mathrm{rad}$。理想可逆制动可取回45焦耳机械能，实际电气损失需另加。这个两刚体系统不满足“单一刚体零角动量就不能转动”的前提。

由 $J=0$ 得 $\Omega=-\nu/10$，所以第一段平台角加速度为 $-0.5\,\mathrm{rad\,s^{-2}}$，并由 $-\tau=I_b\dot\Omega$ 求得电机力矩。两段平台角速度的平均值都为 $-0.5\,\mathrm{rad\,s^{-1}}$，每段持续2秒，因此两段转角同号。
\end{example}

\begin{figure}[htbp]
\centering
\begin{tikzpicture}
\begin{axis}[width=.78\textwidth,height=4.7cm,axis lines=left,xlabel={$t/\mathrm s$},ylabel={$\Omega/(\mathrm{rad\,s^{-1}})$},xmin=0,xmax=4,ymin=-1.15,ymax=.1,xtick={0,1,2,3,4},ytick={-1,-.5,0}]
\addplot[blue!65!black,very thick]coordinates{(0,0)(2,-1)(4,0)};
\node at(axis cs:2,-.35){有符号面积 $=-2\,\mathrm{rad}$};
\end{axis}
\end{tikzpicture}
\caption{反作用轮加速和制动期间的平台角速度。始末角速度均为零，曲线积分却非零，给出净姿态转角。图中的折点来自理想瞬时切换力矩。}
\label{fig:rigid-wheel}
\end{figure}

\subsection{外力矩、饱和与三维扩展的边界}
如果外力矩一直存在，轮速可能不断增大，最后达到设备上限。先算轮速还有多少余量，再判断能保持姿态多久。

若恒定外力矩 $M_{\rm ext}$ 作用于平台，而控制使 $\Omega=0$，则 $I_w\dot\nu=M_{\rm ext}$。轮在吸收外部输入的角动量，有限轮速意味着不能无限保持这一状态，需要用外部力矩卸载。NASA的航天器系统介绍也以外力矩积累及卸载说明这一限制，见章末来源；本例不使用具体卫星性能数据。

对轮轴固连于平台、轮为轴对称体的三维模型，取组合质心为原点，可令 $\bI_{\rm lock}$ 为各部件相对锁止时的组合惯量，$\boldsymbol h=I_w\nu\be_3$ 为相对自转角动量。体轴中
\[
\bJ_b=\bI_{\rm lock}\bom_b+\boldsymbol h,\qquad
\bI_{\rm lock}\dot\bom_b+\dot{\boldsymbol h}
+\bom_b\times(\bI_{\rm lock}\bom_b+\boldsymbol h)=\bM_{\rm ext,b}.
\]
这里组合惯量可为常数，但总系统不是一个固定形状、各点相对静止的单一刚体。非共轴姿态控制不能逐轴机械套用单轴比例式；需先用总角动量和输运方程建模。


\section{计算与模拟：同时检查动力学和姿态}
数值积分时要同时检查体轴分量和空间姿态。即使能量和角动量大小正确，也可能由于姿态更新用错方向而得到错误的空间角动量。

伴随标准库脚本为 \texttt{review/verify\_ch05.py}，结果文件为
\begin{center}\texttt{ch05\_verification\_results.json}\end{center}
\begin{enumerate}[leftmargin=1.8em]
\item 对离散刚架从质点求和独立构造惯量，检查主轴、本征值、坐标旋转和角动量/能量一致性；数值微分欧拉角矩阵，检查左右角速度约定。
\item 对无力矩欧拉方程同时积分单位四元数，分别记录 $T,J^2,|q|$ 及空间角动量 $R(q)\bI_b\bom_b$。前两个标量正确仍不足以保证姿态重构正确。
\item 用最大轴与中间轴附近的初值比较体轴分量；将步长减半，在相同有限时间内检查收敛。中间轴的早期指数增长只在微扰仍小时与线性公式比较。
\item 以原 $\theta$ 方程、固定动量的有效势二阶有限差分、守恒能量三种方式核验陀螺；滚动、碰撞和反作用轮还分别检验反力可行性、冲量与能量收支。
\end{enumerate}
模拟参数均为教学设定；RK4不保证长期严格保能，归一化四元数也不能代替动力学收敛检查。

\section*{关键概念与常见误区}
\begin{itemize}[leftmargin=1.8em]
\item 一般刚体用六个量描述。若所有质量都在一条线上，绕这条线转动不改变点列，而且轴向惯量为零。
\item 写姿态更新前，先说明矩阵把哪个坐标系的分量变到哪个坐标系。空间角速度和随体角速度在更新式中分别乘在不同侧。
\item 使用 $\bJ=\bI\bom$ 前，先标明矩心和坐标轴。若研究绕质心的自转，就用质心惯量，单独处理质心平动。
\item 轴速恒定时，轴承仍可能要施力矩来改变角动量方向。求静力平衡时，也要核查支承是否真能提供所需反力。
\item 轮子纯滚动时，当前接触点相对地面速度为零，下一刻仍可能移动。摩擦功率要按实际接触点的速度来算。
\item 先问有无外力矩，再分清陀螺轴的绕行、倾角起伏和竖直自转。它们需要各自的近似或稳定性条件。
\item 恢复系数给的是接触点法向相对速度的关系。若题目有摩擦、铰支或多个接触点，需要重新列冲量方程。
\end{itemize}

\section*{本章小结}
\addcontentsline{toc}{section}{本章小结}
刚体运动由质心平动与姿态变化组成。惯量张量把动能二次型和角动量联系起来；选质心或固定点、再用输运定理，才能得到不漏项的欧拉方程。自由转动由守恒量限制，重陀螺由循环动量约化为章动有效势，接触与冲击则需额外检查摩擦、单侧反力和冲量假设。内部电机改变组合体机械能而保持总角动量，说明建模时必须区分单一刚体与多刚体系统。

下一章用共轭动量代替速度，在满足正则性条件时把本章及前文的拉格朗日系统转成哈密顿方程。惯量正定与欧拉角坐标奇异是两件不同的事：真实刚体的姿态动力学可以正则，而某个角坐标图仍会在极轴退化。

\section*{分层习题}
\addcontentsline{toc}{section}{分层习题}
先画质心、转轴和力臂，并标出矩心。用角速度或惯量之前，说明分量所用的坐标轴。

\begin{bookexercises}
\item \exlevel{基础} \textbf{自由度与惯性。} 对一般刚体、严格共线点列及单质点分别计数位形自由度。为什么均匀矩形薄板的惯量正定，而理想细杆绕自身轴惯量为零？
\item \exlevel{基础} \textbf{转动约定。} 分别用 $R_x(\pi/2)R_z(\pi/2)$ 和 $R_z(\pi/2)R_x(\pi/2)$ 作用于 $\hat z$。对ZXZ欧拉角取 $\theta\to0$，说明 $\dot\phi,\dot\psi$ 为何不能由姿态唯一分开。
\item \exlevel{基础} \textbf{角点惯量。} 均匀正方薄板质量 $M$、边长 $a$，以一个角点为原点、两条边为轴，直接积分求完整惯量张量及主惯量，再用质心平行轴公式复算。
\item \exlevel{基础} \textbf{物理可实现性。} 三个正数 $I_1=I_2=1,I_3=3$（同一惯量单位）能否作为真实质量分布对某点的主惯量？证明一般主惯量的三角不等式。
\item \exlevel{基础} \textbf{加载梯子。} 正文长 $2L$、质量 $M$ 的梯子上站一质量 $m$ 的人，距底端沿杆距离为 $s$，$0\le s\le2L$。求不滑所需的最小静摩擦因数，并核验 $m=0$ 极限。
\item \exlevel{综合} \textbf{移动矩心。} 用绝对速度与相对速度分别定义对移动点 $A$ 的角动量，推导各自的输运修正。刚体物质点瞬时速度为零，是否就能当固定点使用？
\item \exlevel{综合} \textbf{非主轴恒速。} 对受驱倾斜细杆，独立算 $\dot\bJ$ 与 $\bom\times\bJ$，核验所需力矩和零功率；解释恒定转速不等于零外力矩。
\item \exlevel{综合} \textbf{复摆最短周期。} 刚体质心惯量 $I_C=m\kappa^2$，转轴平行移动使轴到质心距离为 $l>0$。求周期随 $l$ 的变化及其最小值，说明不允许的轴位置应另受几何约束。
\item \exlevel{综合} \textbf{滚动可行性。} 实心圆柱在倾角 $\beta$ 斜面上从静止滚下，推导加速度、静摩擦和所需最小 $\mu_s$；比较摩擦对平动、转动的功率。
\item \exlevel{综合} \textbf{击球与滑动。} 由正文冲量初态导出实心球进入滚动后的 $v_r$ 对击球高度 $h$ 的关系，计算中心击球的能量损失。假设击球期间切向冲量可忽略，随后按恒定动摩擦滑动。
\item \exlevel{综合} \textbf{两个进动频率。} 无力矩对称刚体取 $I_1=I_2=2,I_3=3,\omega_3=1$，横向角速度振幅 $A=0.3$（教学单位）。求体分量旋转频率 $n$、空间进动率 $J/I_1$、能量，并解释二者不同。
\item \exlevel{综合} \textbf{中间轴与相位。} $I_1<I_2<I_3$ 时推导绕中间轴 $\Omega$ 匀转的线性增长率。若初始角速度仅沿稳定轴增大一个小量，绝对姿态差是否永远小？
\item \exlevel{综合} \textbf{稳恒进动。} 重陀螺取 $I_1=1,I_3=0.6,mgl=0.3,h=2,\theta_0=\pi/3$。求快慢两根及各自 $p_\phi$，以固定动量微分有效势检验章动频率。为什么不能保持 $\Omega$ 不变求二阶导数？
\item \exlevel{综合} \textbf{睡眠阈值。} 从轴向小偏离 $u,v$ 推出竖直向上睡眠解的两个横向频率；说明阈值处重根为何不能仅按“频率为实”判稳。
\item \exlevel{综合} \textbf{单点冲击。} 对与水平夹角为 $0<\theta\le\pi/2$ 的均匀杆，下端撞击光滑水平面，碰前无转动、质心下降速度 $V$，恢复系数 $e_r$，求有效质量、冲量及损失。什么几何或约束变化会使本公式不能直接用？
\item \exlevel{拓展} \textbf{四元数与误差。} 用本书乘法展开 $q\otimes(0,\rho)\otimes q^*$，核验转动矩阵及 $q,-q$ 等价。显式欧拉法姿态步为何漂离正交约束？归一化能否保证保能？
\item \exlevel{拓展} \textbf{轮速与卸载。} 反作用轮两阶段例题中求电机力矩、机械功和净姿态角。若平台固定且有 $0.1\,\mathrm{N\,m}$ 恒定外力矩，$I_w=1$，轮速从10到上限 $50\,\mathrm{rad\,s^{-1}}$，还可维持多久？
\item \exlevel{拓展} \textbf{数值检查。} 运行伴随脚本，对自由刚体同时检查 $T,J^2,|q|$ 与空间 $\bJ$，比较两步长。哪些错误可能保留前两个标量而破坏姿态运动？中间轴增长率应在哪一阶段比较？
\end{bookexercises}

\section*{习题解答}
先用同一题号找到解答。核对时，从第一条方程开始比较自己的做法，再检查最后的符号、单位和适用条件。
\begin{bookexercises}
\item 在正文模型下分别为6、5、3。薄板质量不全在任一过参考点直线上，任一非零转轴都使部分质点有速度，故惯量二次型正；细杆自身线是零速度方向。

一般刚体由参考点的三个平移坐标与三个姿态参数确定。共线点列还需确定轴线方向，只有两个独立方向参数，绕轴旋转不改变点列；单点只需位置。惯量二次型为 $\int|\boldsymbol\omega\times\br|^2\dd m$，为零要求质量几乎处处位于转轴上。
\item 先施加右侧矩阵。第一种次序中 $\hat z$ 经 $R_z$ 保持不变，再经 $R_x$ 变为 $-\hat y$；第二种先得 $-\hat y$，再经 $R_z$ 变为 $\hat x$。$\theta=0$ 时 $R=R_z(\phi+\psi)$，此时姿态只确定两角之和，绕重合轴的角速度也只确定 $\dot\phi+\dot\psi$。若轨迹正离开极点，仍可有与 $\dot\theta$ 有关的横向角速度。
\item $I_{xx}=I_{yy}=Ma^2/3,I_{xy}=-Ma^2/4,I_{zz}=2Ma^2/3$，其余非对角元零。主惯量为 $Ma^2/12,7Ma^2/12,2Ma^2/3$；最小轴沿对角线 $(1,1,0)$。

直接积分用 $\dd m=(M/a^2)\dd x\dd y$，$0\le x,y\le a$。有 $\langle x^2\rangle=\langle y^2\rangle=a^2/3$、$\langle xy\rangle=a^2/4$，因此 $I_{xy}=-M\langle xy\rangle$，平面内本征向量沿 $(1,1)$ 与 $(1,-1)$。独立复算时，质心坐标中的二阶平均为 $a^2/12$，交叉平均为零，所以
\[
I_C=\frac{Ma^2}{12}\operatorname{diag}(1,1,2),\qquad
I_O=I_C+M(|d|^2\mathsf I-dd^{\mathsf T}),\quad
d=(a/2,a/2,0)^{\mathsf T}.
\]
平行轴增量的对角元为 $Ma^2(1/4,1/4,1/2)$，$xy$ 元为 $-Ma^2/4$，与直接积分一致。
\item 不能。$I_1+I_2-I_3=2\int x_3^2\dd m\ge0$，本组违反此式。正定对称矩阵不必就是可实现的惯量张量。

因为在主轴系中 $I_j$ 是另两坐标平方的质量积分，其余两个不等式由指标循环得到。平面质量分布在法向主轴上可达到等号。
\item $N_f=(M+m)g$，$N_w=g(M/2+ms/(2L))\cot\theta$，故
\[\mu_{\min}=\frac{M/2+ms/(2L)}{M+m}\cot\theta.\]
退回 $m=0$ 得正文结果；还假设人保持相对梯子静止、墙面光滑。

以梯脚为矩心，墙面法向力的力臂为 $2L\sin\theta$，两重力的力臂为 $L\cos\theta$、$s\cos\theta$。故
\[
2LN_w\sin\theta=(MgL+mgs)\cos\theta .
\]水平力平衡给 $f=N_w$，再由 $f\le\mu_s N_f$ 得上式。
\item 记 $\rho_i=\br_i-\br_A$，定义 $\boldsymbol H_A=\sum_i\rho_i\times m_i\bv_i$。微分时 $\dot\rho_i=\bv_i-\bv_A$，因此
\[
\dot{\boldsymbol H}_A
=\sum_i(\bv_i-\bv_A)\times m_i\bv_i+\sum_i\rho_i\times m_i\ba_i
=\bM_A^{\rm ext}-\bv_A\times\bP .
\]
若定义 $\boldsymbol K_A=\sum_i\rho_i\times m_i(\bv_i-\bv_A)$，微分的第一项为 $\sum_i\dot\rho_i\times m_i\dot\rho_i=0$，第二项为 $\sum_i\rho_i\times m_i(\ba_i-\ba_A)$。于是 $\dot{\boldsymbol K}_A=\bM_A^{\rm ext}-M\rho_C\times\ba_A$。瞬时 $\bv_A=0$ 只保证绝对定义的速度修正此刻为零；相对定义的加速度修正仍须检查，也不能把该点视为在整个时间段内固定。
\item 两种导数均给 $-I_\perp\Omega^2\cos\alpha\,\hat z\times\bu$；力矩与 $\bom$ 垂直。若用固定转速标量的导数代替角动量矢量导数，就会漏掉此力矩。

对细杆，$\bJ=I_\perp\Omega(\hat z-\cos\alpha\,\bu)$，由 $\dot\bu=\Omega\hat z\times\bu$ 可直接求 $\dot\bJ$。另一算法用 $\bom=\Omega\hat z$，叉乘式中的 $\hat z$ 分量为零。两者均表明 $\bM\cdot\bom=0$。
\item $T=2\pi\sqrt{(\kappa^2+l^2)/(gl)}$，在可实现的 $l=\kappa$ 取最小 $2\pi\sqrt{2\kappa/g}$。$l\to0$ 时没有重力线性恢复矩，不能把周期公式写成有限常数。

由平行轴公式 $I_O=m(\kappa^2+l^2)$ 及线性重力矩 $-mgl\theta$，得 $T^2/(4\pi^2)=(l+\kappa^2/l)/g$。它对 $l$ 的导数为 $(1-\kappa^2/l^2)/g$，先负后正。
\item $a=2g\sin\beta/3,f=Mg\sin\beta/3,\mu_{s,\min}=\tan\beta/3$。平动摩擦功率 $-fv_C$，转动功率 $fR\omega$，无滑动时总和零。

取向下沿坡为正，$Ma=Mg\sin\beta-f$、$I_C\alpha=fR$，无滑动给 $a=R\alpha$。代入 $I_C=MR^2/2$ 得 $f=Ma/2$，解出加速度与摩擦。再除以 $N=Mg\cos\beta$ 得静摩擦下界。
\item $v_r=5\mathcal J(1+h/R)/(7M)$；中心击球 $v_0=\mathcal J/M,\omega_0=0$，损失 $Mv_0^2/7$，最终为 $5v_0/7$。直接滚动高度 $h=2R/5$ 给零滑动损失。

击后 $Mv_0=\mathcal J$、$I_C\omega_0=\mathcal Jh$。令 $f$ 为摩擦力的有符号水平分量，正方向取击球方向，角速度则取顺时针为正。滑动中 $M\dot v=f$、$I_C\dot\omega=-Rf$，所以 $v+I_C\omega/(MR)$ 保持不变。无滑动条件 $\omega_r=v_r/R$ 和 $I_C=2MR^2/5$ 给所求 $v_r$。中心击球时 $T_r=\frac12(M+I_C/R^2)(5v_0/7)^2=5Mv_0^2/14$，从 $T_0=Mv_0^2/2$ 减去即得损失。
\item $n=1/2$，$J=\sqrt{9.36}$，空间进动率 $\sqrt{9.36}/2\simeq1.5297$，$T=1.59$。前者观察体轴分量，后者观察空间角动量周围的方向变化。

具体计算为 $n=(I_3-I_1)\omega_3/I_1$，$J^2=I_1^2 A^2+I_3^2\omega_3^2$，$T=(I_1 A^2+I_3\omega_3^2)/2$。空间中 $\dot\be_3=(\bJ/I_1)\times\be_3$，故空间进动率取 $J/I_1$。
\item $\lambda=|\Omega|\sqrt{(I_2-I_1)(I_3-I_2)/(I_1I_3)}$。稳定轴上的转速误差使转角差按时间积累；角速度/约化稳定不要求同一绝对相位永远贴近。

令 $\omega_2=\Omega$、$\omega_1=u,\omega_3=v$，线性欧拉方程为 $I_1\dot u=(I_2-I_3)\Omega v$、$I_3\dot v=(I_1-I_2)\Omega u$。消去 $v$ 后 $\ddot u=\lambda^2u$，给出增长率。只扰动转速而仍严格沿主轴时，转角差为 $\Delta\Omega\,t$。
\item $\Omega=2\pm\sqrt{3.4}$，$p_\phi=1+3\Omega/4$，且
\[\omega_n^2=(2-\Omega)^2+3\Omega^2/4.\]
改变 $\theta$ 而保持动量时 $\dot\phi$ 随之变化，不等于固定进动率的另一扰动问题。

代入稳恒条件 $I_1\Omega^2\cos\theta_0-h\Omega+mgl=0$，得 $\Omega^2/2-2\Omega+0.3=0$。固定 $p_\phi,h$，记 $w=(p_\phi-h\cos\theta)/\sin\theta$，则 $w'=h-w\cot\theta$、$w''=-w'\cot\theta+w/\sin^2\theta$。将 $w=I_1\Omega\sin\theta_0$ 代入
\[
I_1\omega_n^2=V_{\rm eff}''=\big[(w')^2+ww''\big]/I_1-mgl\cos\theta_0
\]
并用稳恒方程消去 $mgl$，后两项合成 $I_1\Omega^2\sin^2\theta_0$。慢支有 $\omega_n^2\simeq3.41827$，快支有 $\omega_n^2\simeq14.48173$。两者均为正。
\item $\sigma_\pm=[h\pm\sqrt{h^2-4I_1mgl}]/(2I_1)$。等号处二重频率可有线性时间因子，须检验解的有界性；不得只用除过 $\sin\theta$ 的稳恒式。

横向角动量为 $(-I_1\dot v+hu,I_1\dot u+hv)$，重力矩为 $mgl(-v,u)$。逐分量令 $\dot{\bJ}_\perp=\bM_\perp$ 后，用 $w=u+iv$ 合并为 $I_1\ddot w-ih\dot w-mgl\,w=0$。代入 $w=e^{i\sigma t}$ 得 $I_1\sigma^2-h\sigma+mgl=0$。阈值处一般解含 $(A+Bt)e^{iht/(2I_1)}$，所以有增长初值。
\item $m_{\rm eff}=m/(1+3\cos^2\theta)$，$J_n=(1+e_r)m_{\rm eff}V$，损失为 $m_{\rm eff}(1-e_r^2)V^2/2$。多端同时接触、固定铰支、摩擦或表面移动时需重建冲量方程。

若杆长 $l$、倾角从水平量起，冲击点力臂的水平分量绝对值为 $l|\cos\theta|/2$。法向速度每单位冲量的增量为 $1/m+(l\cos\theta/2)^2/I_C=(1+3\cos^2\theta)/m$，其中 $I_C=ml^2/12$。碰前法向速度 $-V$，碰后为 $e_r V$，故 $\Delta v_n=(1+e_r)V$。总动能变化为 $-VJ_n+J_n^2/(2m_{\rm eff})$，代入冲量即得损失。
\item 展开得式 \eqref{eq:rigid-quaternion-matrix}；$q$ 换号使二次项不变。欧拉步产生 $h^2W^{\mathsf T}W$ 缺陷；保持单位四元数仅控制表示约束，不控制角速度方程的能量误差。

令 $q=(s,\boldsymbol a)$，其中 $s^2+|\boldsymbol a|^2=1$。第一次相乘得 $(-\boldsymbol a\cdot\rho,s\rho+\boldsymbol a\times\rho)$，再乘 $(s,-\boldsymbol a)$，标量部为零，矢量部为
\[
(s^2-|\boldsymbol a|^2)\rho+2\boldsymbol a(\boldsymbol a\cdot\rho)
+2s\boldsymbol a\times\rho.
\]
各项均为 $q$ 分量的二次式。对显式欧拉矩阵步，乘积为 $(I+hW)^{\mathsf T}(I+hW)$，由 $W^{\mathsf T}=-W$ 一次项相消，只留上述二次缺陷。
\item 加速、减速力矩分别为 $\pm4.5\,\mathrm{N\,m}$，机械功为 $\pm45\,\mathrm J$，平台总转角 $-2\,\mathrm{rad}$。固定平台时 $\dot\nu=0.1\,\mathrm{rad\,s^{-2}}$，余量40，对应400秒；耗尽前须卸载或改变姿态要求。

由 $J=0$ 有 $\Omega=-\nu/10$。加速段 $\dot\nu=5\,\mathrm{rad\,s^{-2}}$，所以 $\tau=-I_b\dot\Omega=4.5\,\mathrm{N\,m}$。该段 $\int_0^2 \nu\,\dd t=10\,\mathrm{rad}$，于是 $W=\int\tau\nu\,\dd t=45\,\mathrm J$、$\int\Omega\,\dd t=-1\,\mathrm{rad}$。制动段 $\nu$ 仍为正，时间积分相同，力矩反向，故取回相同机械功，平台角位移仍为负。固定平台时，电机必须抵消外力矩，轮的角动量以 $0.1\,\mathrm{N\,m}$ 的速率增长。
\item 把体角速度当空间角速度、左右乘次序或符号错误，都可能使姿态重构错而欧拉方程的两个标量仍正确。指数率仅在横向扰动很小、非线性饱和或翻转尚未发生时比较；步长收敛也限指定有限区间。

以中间轴初值为例，$h=0.005$ 时四个最大误差按题目顺序约为 $3.68\times10^{-13}$、$3.68\times10^{-13}$、$1.09\times10^{-12}$（统计 $|q|^2$）和 $5.13\times10^{-12}$。用三个相邻步长比较末态差，粗步长对细步长的末态差之比约 $16.07$，符合四阶积分的预期。在早期线性区可比较 $\log$ 横向扰动的斜率与理论值 $0.57735$，应避开初始增长模尚未占主导的区段。
\end{bookexercises}

\section*{进一步阅读与来源}
链接复核：2026-11-08。算例按本书惯量参考点、主动转动与ZXZ约定重建，数值参数为教学设定。
\begin{bookreading}
\item David Tong，\href{https://davidtong.org/pdfs/teaching/classical-dynamics/clas3.pdf}{The Motion of Rigid Bodies}，可继续学习自由陀螺、睡眠陀螺和姿态重构；比较公式前先核对转动与惯量约定。
\item Iain Stewart，MIT 8.09，\href{https://ocw.mit.edu/courses/8-09-classical-mechanics-iii-fall-2014/6fe39e8d5ce4ce746ca256dfea665eda_MIT8_09F14_Chapter_2.pdf}{Chapter 2: Rigid Body Dynamics}，对照姿态、惯量与欧拉方程。其坐标变换的取向需与本书体到空间矩阵区分。
\item MIT 8.09，\href{https://ocw.mit.edu/courses/8-09-classical-mechanics-iii-fall-2014/d5fdb8fec0febc6ba799cc86e67af3ee_MIT8_09F14_pset3.pdf}{Problem Set 3}，练习从对称性与质量分布独立建立惯量、并使用旋转坐标系。本章题目与解答另行编写。
\item NASA，\href{https://science.nasa.gov/learn/basics-of-space-flight/chapter11-2/}{Basics of Space Flight, Chapter 11: Onboard Systems}，进一步了解反作用轮的角动量积累与外力矩卸载。这里只借其说明模型边界，不将单轴算例当作实际任务控制系统。
\end{bookreading}


\chapter{哈密顿力学}\label{ch:hamilton}

\begin{guidebox}
给出弹簧振子的位置和速度，就能预测下一刻的运动。也可以记录位置和动量；只要能在所用范围内从动量找回速度，两套数据同样够用。本章用后一套数据描述完整状态，每个状态对应相空间中的一点。

初读时可以一直跟着弹簧振子。先看怎样从拉格朗日量换到哈密顿量，再在相图上观察轨道和守恒量。之后再学怎样换到更好用的成对变量，以及如何用一个函数重建轨道。遇到带电粒子时，留意方程中的正则动量和由实际速度算出的机械动量。
\end{guidebox}

\section{勒让德变换与正则方程}
同一个振子可以用位置和速度描述，也可以用位置和动量描述。把速度换成动量后，运动由一对一阶方程给出。

\subsection{先作变量变换，再使用运动方程}
从速度求动量通常直接，反过来却可能有多个速度给出同一个动量。先判断能否在所用范围内反解，再建立以动量为变量的函数。

设 $L(q,\dot q,t)$ 光滑，$q=(q_1,\ldots,q_s)$ 是独立局部坐标。定义
\begin{equation}
p_a=L_{\dot q_a},\qquad W_{ab}=L_{\dot q_a\dot q_b}.
\label{eq:hamilton-momentum}
\end{equation}
这里的 $p_a$ 是第\ref{ch:lagrange}章定义的广义动量。把它与 $q_a$ 一起用作相空间的正则变量时，称 $p_a$ 为\textbf{正则动量}。
在 $\det W\ne0$ 的一个点附近，隐函数定理保证可局部反解 $\dot q=v(q,p,t)$。定义
\begin{equation}
H(q,p,t)=\sum_a p_av_a(q,p,t)-L(q,v(q,p,t),t).
\label{eq:hamiltonian}
\end{equation}
这个勒让德变换中的速度已被消去，不是额外的独立变量。\emph{不使用运动方程}，先微分：
\begin{align*}
\dd H&=\sum_a(v_a\dd p_a+p_a\dd v_a)
-\sum_a(L_{q_a}\dd q_a+L_{\dot q_a}\dd v_a)-L_t\dd t\\
&=\sum_a(v_a\dd p_a-L_{q_a}\dd q_a)-L_t\dd t.
\end{align*}
故恒等地 $H_{p_a}=v_a,H_{q_a}=-L_{q_a},H_t=-L_t$。最后才在真实路径上用 $\dot p_a=L_{q_a}$，得到
\begin{equation}
\boxed{\dot q_a=H_{p_a},\qquad\dot p_a=-H_{q_a}.}
\label{eq:canonical-equations}
\end{equation}
这里 $H_t$ 保持 $q,p$ 不变，$L_t$ 保持 $q,\dot q$ 不变，再用速度反解比较。反向代入也恢复原拉格朗日方程；两种描述都需要 $2s$ 项初始数据。

若 $L=\dot q^{\mathsf T}M(q,t)\dot q/2+a(q,t)^{\mathsf T}\dot q-V(q,t)$，$M$ 可逆，则
\[
p=M\dot q+a,\qquad
H=\tfrac12(p-a)^{\mathsf T}M^{-1}(p-a)+V.
\]
速度一次项不会使变换失效；它使正则动量偏离通常的机械动量。第五章刚体惯量正定有利于反解，但欧拉角在极轴奇异时仍须换坐标图。

\begin{example}[非退化不等于全局单值]
令 $u=\dot q/v_*$，取正尺度 $E_*,v_*$，
\[
L=E_*(u^4/4-u^2/2)-V(q),\qquad
p=\frac{E_*}{v_*}(u^3-u).
\]
$p=0$ 对应 $u=0,\pm1$，三处海森矩阵分别为 $-E_*/v_*^2,2E_*/v_*^2,2E_*/v_*^2$，均非零，但同一动量对应不同速度分支。各分支局部可用
\[
H=E_*(3u^4/4-u^2/2)+V(q),
\]
并不能据局部非退化就任意选成一个全局函数。严格凸性可保证动量映射在适当凸域中的单射性，覆盖整个动量域还需全局条件。普通正则理论不要求海森矩阵正定，但稳定性和全局性是另外的问题。
\end{example}

若 $L$ 含未消去的乘子，或某个速度根本不出现，海森矩阵可奇异，例如 $L=(\dot x-y)^2/2$ 有 $p_y=0$。不能照常反解所有速度；约束哈密顿理论需另外建立，本章不展开。

\subsection{能量、循环坐标与含时表示}
坐标若不直接出现在拉格朗日量中，常能读出一个守恒动量。哈密顿量能否当作机械能，还要看模型和坐标是否显含时间。

$H$ 是第\ref{ch:lagrange}章能量函数 \eqref{eq:energy-function} 在可逆正则变量中的表达，即 $H(q,p,t)=E_L(q,\dot q(q,p,t),t)$。
沿正则运动
\begin{equation}
\frac{\dd H}{\dd t}=\sum_a(H_{q_a}\dot q_a+H_{p_a}\dot p_a)+H_t=H_t.
\label{eq:hamiltonian-conserved}
\end{equation}
所以 $H_t=0$ 给出 $H$ 守恒。它是否等于实验室系机械能，仍需建模条件：定常自然系统的 $L=T-V$、$T$ 为齐二次型时，$H=T+V$。电磁速度一次项也可抵消而给 $H=T+e\phi$；不能简单断言“凡含速度势，$H$ 必不等于能量”。含时坐标或全导数改变则可能改变这种识别。

若 $H$ 不含 $q_a$，则 $\dot p_a=0$，但守恒的是所选正则动量。若 $L'=L+\dd G(q,t)/\dd t$，则
\[
p'=p+G_q,\qquad
H'(q,p',t)=H(q,p'-G_q,t)-G_t.
\]
显含时间的表示仍可能描述原来守恒的物理量，不能只看 $H'$ 是否显含时间判断整个系统有无能量守恒。

\begin{example}[谐振子：椭圆与时间方向]
对 $L=m\dot q^2/2-m\omega^2q^2/2$，有
\[
p=m\dot q,\quad H=\frac{p^2}{2m}+\frac12m\omega^2q^2,
\quad \dot q=p/m,\quad\dot p=-m\omega^2q.
\]
消元得 $\ddot q+\omega^2q=0$。令固定长度尺度为 $a$，用 $Q=q/a,P=p/(m\omega a)$，等能线为 $Q^2+P^2=2E/(m\omega^2a^2)$，在 $(Q,P)$ 图上沿顺时针演化；右端点速度朝下，最高点速度向右。这两个无量纲量用于画图，它们的括号为$1/(m\omega a^2)$；作无量纲正则化时需连同作用量、时间尺度一致处理。
\end{example}

\begin{figure}[htbp]
\centering
\begin{tikzpicture}[>=Stealth]
\begin{axis}[width=.72\textwidth,height=5.6cm,axis equal image,axis lines=middle,xlabel={$q/a$},ylabel={$p/(m\omega a)$},xlabel style={at={(axis description cs:1.04,.5)},anchor=west},ylabel style={at={(axis description cs:.5,1.04)},anchor=south,rotate=0},xmin=-2.25,xmax=2.25,ymin=-2.25,ymax=2.25,xtick={-2,-1,1,2},ytick={-2,-1,1,2}]
\addplot[blue!65!black,thick,domain=0:360,samples=160]({2*cos(x)},{2*sin(x)});
\addplot[red!70!black,thick,dashed,domain=0:360,samples=160]({cos(x)},{sin(x)});
\draw[->,blue!65!black,thick](axis cs:2,0)--(axis cs:2,-.55);
\draw[->,red!70!black,thick](axis cs:0,1)--(axis cs:.5,1);
\end{axis}
\end{tikzpicture}
\caption{谐振子按同一尺度归一化的相轨道，内外能量比为1:4。两曲线采用相同轴比，箭头与正则方程的时间方向一致。}
\label{fig:phase-oscillator}
\end{figure}

\subsection{相空间作用量与端点条件}
现在把位置和动量暂时当作可以各自变动的路径。积分分部留下的端点项，告诉我们固定哪一种端点数据时可以直接得到正则方程。

取固定端点时间，定义
\begin{equation}
\mathcal A[q,p]=\int_{t_1}^{t_2}\left[\sum_ap_a\dot q_a-H(q,p,t)\right]\dd t.
\label{eq:ham-phase-action}
\end{equation}
$q,p$ 是独立的试探路径。对 $p\delta\dot q$ 分部积分，得到
\begin{align*}
\delta\mathcal A={}&\left[\sum_ap_a\delta q_a\right]_{t_1}^{t_2}\\
&+\int\sum_a[(\dot q_a-H_{p_a})\delta p_a
-(\dot p_a+H_{q_a})\delta q_a]\dd t.
\end{align*}
固定两端 $q$、不固定动量端点，边界项消失，任意内部变分就给两组正则方程。沿驻值路径有 $p=L_{\dot q}$，作用量恢复第三章的 $\int L\dd t$；任意试探相路径却不需预先满足该关系。

若固定动量端点，应改用 $\mathcal A-[\sum p_aq_a]_{t_1}^{t_2}$，边界变分成为 $-[\sum q_a\delta p_a]$。任意再加固定两端全部动量的求解条件，通常会使边值问题过度约束；固定两个构型端点也不保证任意时间间隔都存在唯一连接轨道。


\section{从模型辨认正则动量}
先看两个熟悉的模型，中心力与磁场。它们能说明为何用角坐标或加入电磁耦合后，正则动量的表达会变。

\subsection{中心势中的球坐标}
用球坐标时，改变角度要沿圆弧移动，圆弧长度还取决于半径。角速度对应的动量因此含有几何因子，不能一律写成质量乘速度。

\begin{example}[中心势的哈密顿方程与径向核验]
非相对论质点质量 $m$，势能 $V(r)=-\kappa/r$，$\kappa>0$。只把它当作经典中心势模型，不包含辐射反作用、原子稳定性或量子跃迁。球坐标正则域取 $r>0,\sin\theta\ne0$，动能给
\[
p_r=m\dot r,\quad p_\theta=mr^2\dot\theta,\quad
p_\varphi=mr^2\sin^2\theta\dot\varphi.
\]
逐项勒让德变换得到
\[
H=\frac{p_r^2}{2m}+\frac{p_\theta^2}{2mr^2}
+\frac{p_\varphi^2}{2mr^2\sin^2\theta}-\frac\kappa r.
\]
于是
\begin{align*}
\dot r&=p_r/m,&\dot p_r&=\frac{p_\theta^2}{mr^3}+\frac{p_\varphi^2}{mr^3\sin^2\theta}-\frac\kappa{r^2},\\
\dot\theta&=p_\theta/(mr^2),&\dot p_\theta&=\frac{p_\varphi^2\cos\theta}{mr^2\sin^3\theta},\\
\dot\varphi&=p_\varphi/(mr^2\sin^2\theta),&\dot p_\varphi&=0.
\end{align*}
$p_\varphi$ 是角动量的 $z$ 分量。若初态 $\theta=\pi/2,p_\theta=0$，方程保持这两个条件，化为
\[
\dot p_r=\frac{\ell^2}{mr^3}-\frac\kappa{r^2},\quad \ell=p_\varphi.
\]
与第一章 $m(\ddot r-r\dot\varphi^2)=-\kappa/r^2$ 一致。坐标轴极点处的分母奇异是球坐标图问题，不代表质点动力学必然奇异；$r=0$ 的库仑势奇异则是另一回事。
\end{example}

\subsection{带电粒子与规范变换}
磁场中，决定粒子实际速度的是机械动量。哈密顿方程使用的正则动量还含矢势，换规范后要同时更新这部分。

带电量记为 $e$（可正可负），与广义坐标 $q$ 区分。外电磁势给定，质量 $m$ 恒定且不计辐射反作用：
\[
L=\tfrac12m\bv^2+e\bA(\br,t)\cdot\bv-e\phi(\br,t),\qquad
\bp=m\bv+e\bA.
\]
故
\begin{equation}
H(\br,\bp,t)=\frac{[\bp-e\bA(\br,t)]^2}{2m}+e\phi(\br,t).
\label{eq:em-hamiltonian}
\end{equation}
定义机械动量 $\boldsymbol\pi=\bp-e\bA$，正则方程是
\[
\dot x_i=\pi_i/m=v_i,\qquad
\dot p_i=e\sum_jv_j\partial_i A_j-e\partial_i\phi.
\]
再用 $\dot\pi_i=\dot p_i-e\partial_tA_i-e\sum_jv_j\partial_jA_i$，得到
\[
\dot\pi_i=e[-\partial_i\phi-\partial_tA_i]
+e\sum_jv_j(\partial_iA_j-\partial_jA_i)
=e[\bE+\bv\times\bB]_i.
\]
这与第三章洛伦兹力推导吻合，但 $\dot p_i$ 本身一般不等于洛伦兹力。

\begin{example}[匀强磁场：机械动量不构成正则对]
取 $\bB=B\hat z,\phi=0,\bA=(0,Bx,0)$。则
\[
H=\frac{p_x^2+(p_y-eBx)^2+p_z^2}{2m},\quad
p_y,p_z\text{ 守恒}.
\]
$x$ 方程给 $\dot p_x=eB(p_y-eBx)/m=eBv_y$，$\dot p_y=0$ 却不意味着 $mv_y$ 守恒，因为 $\pi_y=p_y-eBx$。代入即得平面回旋方程 $m\dot v_x=eBv_y,m\dot v_y=-eBv_x$，机械动能守恒。后文将算出 $\{\pi_x,\pi_y\}=eB$，所以不能把 $(x,y,\pi_x,\pi_y)$ 当成普通正则变量而仍用原括号公式。
\end{example}

规范变换 $\bA'=\bA+\nabla\chi,\phi'=\phi-\partial_t\chi$ 使
\[
\bp'=\bp+e\nabla\chi,\qquad
H'(\br,\bp',t)=H(\br,\bp'-e\nabla\chi,t)-e\partial_t\chi.
\]
它是一种含时正则表示变化，后文的第二类生成函数可取 $F_2(\br,\bp',t)=\br\cdot\bp'-e\chi$。相同轨道、相同机械动量并不要求不同规范下的正则动量和哈密顿函数数值相同。

\subsection{罗斯约化的符号与循环条件}
若某个角度的动量已守恒，可以只写剩下的径向方程。消去角速度前要先按固定动量改写拉格朗日量，才能得到正确的离心项。

第\ref{ch:rigid}章约化重陀螺使用 $\mathcal R=L-p_\phi\dot\phi-p_\psi\dot\psi$。这里保持同一\textbf{拉格朗日型罗斯函数}约定。设坐标分成 $x_i,y_\alpha$，对选中的速度块有可逆海森矩阵，可局部由 $\mu_\alpha=L_{\dot y_\alpha}$ 消去 $\dot y$，定义
\begin{equation}
\mathcal R(x,y,\dot x,\mu,t)=L-\sum_\alpha\mu_\alpha\dot y_\alpha.
\label{eq:routhian}
\end{equation}
把动量关系代入微分，速度 $\dd\dot y$ 项相消，得到
\[
\mathcal R_{\dot x_i}=L_{\dot x_i},\quad \mathcal R_{x_i}=L_{x_i},\quad
\mathcal R_{y_\alpha}=L_{y_\alpha},\quad
\mathcal R_{\mu_\alpha}=-\dot y_\alpha.
\]
真实路径因而满足
\[
\frac{\dd}{\dd t}\mathcal R_{\dot x_i}-\mathcal R_{x_i}=0,\quad
\dot\mu_\alpha=\mathcal R_{y_\alpha},\quad
\dot y_\alpha=-\mathcal R_{\mu_\alpha}.
\]
若 $y$ 循环，$\mu$ 才可固定为常数，剩下 $x$ 的约化拉格朗日问题。另一些文献取哈密顿型相反号 $R_H=\sum\mu\dot y-L=-\mathcal R$，其混合正则方程符号随之改变；不能把两种约定无说明地混用。完整海森矩阵可逆也不自动保证任意选中的子块可逆。

\begin{example}[中心力约化为什么要加离心有效势？]
平面中心势 $L=m(\dot r^2+r^2\dot\varphi^2)/2-V(r)$，循环动量 $\ell=mr^2\dot\varphi$。先取 $L-\ell\dot\varphi$ 再代入，得
\[
\mathcal R=\frac m2\dot r^2-\left[V(r)+\frac{\ell^2}{2mr^2}\right].
\]
所以 $m\ddot r=-V'(r)+\ell^2/(mr^3)$。若只代回 $L$，会得到离心项的错误符号，因为改变的是固定动量的变分问题。用 $\dot\varphi=-\mathcal R_\ell=\ell/(mr^2)$ 可重构轨道方位。
\end{example}


\section{相空间流与刘维尔定理}
同时画出位置与动量，每个完整状态就成了图上的一个点。一组初态随时间一起移动时，可以观察它们占据的相空间区域怎样改变。

\subsection{完整状态与轨道唯一性}
同一位置可以有朝左或朝右的运动。把动量也记下来，才能区分这两个初态并唯一追踪它们的后续轨道。

在正则局部图中，$z=(q,p)$ 是 $2s$ 维相点。记
\[
\Omega=\begin{pmatrix}0&I\\-I&0\end{pmatrix},\qquad
\dot z=\Omega\nabla_zH.
\]
满足局部唯一性条件时，同一时刻的一个相点给唯一的未来。自治系统的不同相轨道在正则点不能相交；含时系统若把时间投影掉，两条图上曲线可以相交，不能据投影否定唯一性。位形曲线在同一位置交叉更不奇怪，因为动量可以不同。

\begin{example}[单摆：闭合摆动、分离曲线与持续转动]
取角度 $q$ 从竖直向下量起，$I=ml^2,K_0=mgl$，
\[
H=\frac{p^2}{2I}+K_0(1-\cos q),\quad
\dot q=p/I,\quad \dot p=-K_0\sin q.
\]
能量给 $p=\pm\sqrt{2I[E-K_0(1-\cos q)]}$。$0<E<2K_0$ 为有两转折点的摆动；$E>2K_0$ 动量不变号，是持续转动；$E=2K_0$ 为分离曲线。$q$ 按 $2\pi$ 识别，完整相空间是圆柱，转动轨道在展开平面上不闭合，却在圆柱上闭合。顶点 $(\pi,0)$ 为鞍点，线性式 $\ddot\eta-(K_0/I)\eta=0$ 含增长与衰减方向。
\end{example}

\begin{figure}[htbp]
\centering
\begin{tikzpicture}[>=Stealth]
\begin{axis}[width=.84\textwidth,height=5.6cm,axis lines=middle,xlabel={$q$},ylabel={$p/\sqrt{IK_0}$},xlabel style={at={(axis description cs:1.04,.5)},anchor=west},ylabel style={at={(axis description cs:.5,1.06)},anchor=south,rotate=0},xmin=-3.5,xmax=3.5,ymin=-2.7,ymax=2.7,xtick={-3.14159,0,3.14159},xticklabels={$-\pi$,$0$,$\pi$},ytick={-2,2},samples=220]
\addplot[blue!65!black,thick,domain=-pi/2:pi/2]{sqrt(max(0,2*cos(deg(x))))};
\addplot[blue!65!black,thick,domain=-pi/2:pi/2]{-sqrt(max(0,2*cos(deg(x))))};
\addplot[red!70!black,thick,dashed,domain=-pi:pi]{2*cos(deg(x)/2)};
\addplot[red!70!black,thick,dashed,domain=-pi:pi]{-2*cos(deg(x)/2)};
\addplot[green!45!black,thick,dashdotted,domain=-pi:pi]{sqrt(2*(2+cos(deg(x))))};
\addplot[green!45!black,thick,dashdotted,domain=-pi:pi]{-sqrt(2*(2+cos(deg(x))))};
\draw[->,red!70!black](axis cs:0,2)--(axis cs:.5,2);
\draw[->,red!70!black](axis cs:0,-2)--(axis cs:-.5,-2);
\end{axis}
\end{tikzpicture}
\caption{单摆相图的展开表示：蓝线$E/K_0=1$，红虚线为分离曲线$E/K_0=2$，绿点划线为$E/K_0=3$。上下半平面时间方向相反；左右边按周期识别。}
\label{fig:ham-pendulum}
\end{figure}

\subsection{散度、流映射与密度}
把一群初态想成相空间里的一团墨点。墨团可以拉长、倾斜，投影到位置轴上也可以变宽，但哈密顿演化保持它的相空间体积。

设 $H$ 二阶偏导连续，$\mathcal D_t$ 是随流移动的相区域。速度场散度为
\[
\nabla_z\cdot\dot z=\sum_a(\partial_{q_a}H_{p_a}-\partial_{p_a}H_{q_a})=0.
\]
由输运定理，
\begin{equation}
\frac{\dd}{\dd t}\int_{\mathcal D_t}\dd^sq\,\dd^sp=0.
\label{eq:liouville}
\end{equation}
等价地，流映射 $z(t)=\Phi_{t,t_0}(z_0)$ 的雅可比行列式满足
\[
\frac{\dd}{\dd t}\det D\Phi=(\nabla_z\cdot\dot z)\det D\Phi=0,
\quad \det D\Phi_{t_0,t_0}=1.
\]
这是刘维尔定理；不要求自治，体积保持也不要求形状保持。

\begin{figure}[htbp]
\centering
\begin{tikzpicture}[>=Stealth]
\foreach \shift/\lab in {0/初始区域,5.2/自由流后的区域}{
\begin{scope}[xshift=\shift cm]
\draw[->](0,0)--(3.5,0)node[right]{$Q$};\draw[->](0,0)--(0,2)node[above]{$P$};
\node at(1.65,2.4){\lab};
\end{scope}}
\draw[thick,fill=blue!12](.5,.5)rectangle(1.5,1.5);
\draw[thick,fill=red!12](6.2,.5)--(7.2,.5)--(8.2,1.5)--(7.2,1.5)--cycle;
\end{tikzpicture}
\caption{取无量纲自由粒子$H=P^2/2$，时间步长为1，准确映射为$(Q,P)\mapsto(Q+P,P)$。左方单位面积矩形变为右方单位面积平行四边形，顶点均由同一映射计算。}
\label{fig:liouville}
\end{figure}

对随流输运的概率密度 $\rho(q,p,t)$，连续性方程为
\[
\partial_t\rho+\nabla_z\cdot(\rho\dot z)=0,
\quad \partial_t\rho+\dot z\cdot\nabla_z\rho=0.
\]
因此精细密度沿轨道不变，即使分布显含时间也如此。对自治 $H$，适当的 $\rho=f(H)$ 是静态分布，但体积保持本身既不保证遍历性，也不推出能量面上某个特定均匀概率；微正则平衡仍需统计假设。相体积元素是正则坐标的 $\dd q\dd p$，在一般非正则变量中不能丢掉测度雅可比。


\section{泊松括号、守恒量与生成元}
若想知道某个量沿运动会不会改变，可以把它与哈密顿量的各个偏导按固定方式配对。这种运算叫泊松括号，也可用来写出小的状态变换。

\subsection{定义及恒等式的来由}
计算泊松括号时，把每一对位置和动量的偏导交叉相乘再相减。先在一对变量上练习，再把所有变量对的贡献相加。

对相空间函数定义
\begin{equation}
\{f,g\}=\sum_a(f_{q_a}g_{p_a}-f_{p_a}g_{q_a})
=(\nabla f)^{\mathsf T}\Omega\nabla g.
\label{eq:poisson-bracket}
\end{equation}
用正则方程代入总导数，得到
\begin{equation}
\dot f=\{f,H\}+\partial_tf.
\label{eq:poisson-evolution}
\end{equation}
所以基本括号为 $\{q_a,q_b\}=\{p_a,p_b\}=0,\{q_a,p_b\}=\delta_{ab}$。括号对实常数线性、反对称，且由导数乘积法则有
\[
\{f,gh\}=\{f,g\}h+g\{f,h\}.
\]
若系数依赖相点，就不能把它随意当常数移出括号。

雅可比恒等式为
\begin{equation}
\{f,\{g,h\}\}+\{g,\{h,f\}\}+\{h,\{f,g\}\}=0.
\label{eq:pb-jacobi}
\end{equation}
它不是额外的运动假设。以 $z_i$ 记全部正则变量，$f_i=\partial_if$，常矩阵 $\Omega$ 给
\[
\{f,\{g,h\}\}=f_i\Omega_{ij}g_{kj}\Omega_{k\ell}h_\ell
+f_i\Omega_{ij}g_k\Omega_{k\ell}h_{\ell j}.
\]
加上两个循环置换后，含 $f_{ij}$、$g_{ij}$、$h_{ij}$ 的项分别因二阶导数对称与 $\Omega$ 反对称成对抵消。这也说明它依赖当前括号的微分结构，不能只凭反对称性就认定任意双线性运算满足雅可比式。

\subsection{显含时间的运动积分}
一个量可能一方面直接含有时间，另一方面又随位置和动量改变。两个变化恰好抵消时，它仍可沿实际轨道守恒。

定义 $D_Hf=\partial_tf+\{f,H\}$。$f$ 守恒的条件是 $D_Hf=0$，不要求它不显含时间。利用雅可比式和时间导数乘积规则，
\[
D_H\{f,g\}=\{D_Hf,g\}+\{f,D_Hg\}.
\]
因此两个运动积分的括号仍是运动积分，即泊松定理；该括号可能只是零或已有量，并不保证产生新的独立守恒量。

\begin{example}[角动量代数与空间转动]
取正则Cartesian变量，$J_x=yp_z-zp_y,J_y=zp_x-xp_z$。分别求梯度：
\[
(J_x)_q=(0,p_z,-p_y),\quad(J_x)_p=(0,-z,y),
\]
\[
(J_y)_q=(-p_z,0,p_x),\quad(J_y)_p=(z,0,-x).
\]
两次点积相减得 $\{J_x,J_y\}=xp_y-yp_x=J_z$，循环置换给 $\{J_i,J_j\}=\epsilon_{ijk}J_k$。这是正则轨道角动量的代数；存在磁矢势时，机械轨道角动量与对称性生成元的识别需另查规范和边界项。
\end{example}

\begin{example}[机械动量的磁括号]
$\pi_i=p_i-eA_i(q,t)$ 满足
\[
\{q_i,\pi_j\}=\delta_{ij},\qquad
\{\pi_i,\pi_j\}=e(\partial_iA_j-\partial_jA_i)=e\epsilon_{ijk}B_k.
\]
计算中 $\{p_i,A_j\}=-\partial_iA_j$，$\{A_i,p_j\}=\partial_jA_i$，正负号不可丢。故虽有位置--机械动量的单位括号，完整基本括号仍不正则；两动量之间的非零括号正记录磁场的影响。
\end{example}

\subsection{守恒量如何生成状态变换}
动量可以生成位置的小平移，角动量可以生成方向的小转动。这里用泊松括号把一个量和它能引起的状态变化联系起来。

对光滑 $G(q,p,t)$，在固定时刻令
\[
\delta q_a=\epsilon G_{p_a},\quad\delta p_a=-\epsilon G_{q_a},\quad
\delta f=\epsilon\{f,G\}.
\]
$G=p_x$ 生成 $x$ 平移；$G=J_z=xp_y-yp_x$ 给
\[
\delta x=-\epsilon y,\quad\delta y=\epsilon x,\quad
\delta p_x=-\epsilon p_y,\quad\delta p_y=\epsilon p_x,
\]
即绕 $+z$ 的主动微转动。$G$ 的哈密顿向量场本身生成局部正则流；若 $G$ 自治且 $\{H,G\}=0$，该变换还保持 $H$，因而给动力学对称性。

\begin{example}[伽利略速度变换与含时守恒量]
一维自由粒子 $H=p^2/(2m)$，取 $G=tp-mx$。有
\[
\partial_tG=p,\quad\{G,H\}=-p,\quad\dot G=0;
\qquad\delta x=\epsilon t,\quad\delta p=\epsilon m.
\]
$\epsilon$ 的单位是速度。它生成主动增加速度的状态变换；同第三章含边界项的诺特量一致。若取 $-G$，变换参数方向同时反转，不能只比较守恒量的正负号。
\end{example}

\begin{modernbox}[title={进一步探索：泊松括号与量子对易子的对应范围}]
正则量子化用 $[\hat q_a,\hat p_b]=i\hbar\delta_{ab}$，角动量也有相应对易代数。这是结构联系，不意味着对任意经典函数都存在保持所有乘积和泊松括号的简单量子映射；高次乘积有算符次序等问题。继续学习需线性算符与量子力学，本章不以替换符号代替量子化的完整构造。
\end{modernbox}


\section{正则变换：保持的结构与局部生成函数}
有些新变量能让方程更好算，但要保留位置与动量成对的关系。先判断变换是否正则，再用合适的生成函数来构造它。

\subsection{辛条件比保持体积更强}
四维相空间中，可以拉大一对位置与动量所张成的面积，同时按相反比例缩小另一对的面积，总体积仍不变。判断正则性还要检查全部基本括号，包括不同变量对之间的括号。

令 $Z=(Q,P)=Z(z,t)$ 是可逆局部变换，$\mathcal J=\partial Z/\partial z$。固定时刻的新变量括号矩阵是
\[
\{Z_i,Z_j\}=(\mathcal J\Omega\mathcal J^{\mathsf T})_{ij}.
\]
正则变换要求基本括号保持，即
\begin{equation}
\mathcal J\Omega\mathcal J^{\mathsf T}=\Omega
\quad\Longleftrightarrow\quad
\mathcal J^{\mathsf T}\Omega\mathcal J=\Omega.
\label{eq:symplectic-condition}
\end{equation}
对左边的矩阵等式取逆，再用 $\Omega^{-1}=-\Omega$，得到 $\mathcal J^{-\mathsf T}\Omega\mathcal J^{-1}=\Omega$。左乘 $\mathcal J^{\mathsf T}$、右乘 $\mathcal J$，即得右边的等式。反向交换 $\mathcal J$ 与 $\mathcal J^{\mathsf T}$ 即可。此条件还使任意 $f,g$ 的泊松括号经链式法则保持，不仅是某个给定 $H$ 的一条轨道偶然能写成相似方程。

对不显含时间的正则变换，$K(Z)=H(z(Z))$，于是
\[
\dot Z=\mathcal J\Omega\nabla_zH
=\mathcal J\Omega\mathcal J^{\mathsf T}\nabla_ZK=\Omega\nabla_ZK.
\]
显含时间时还须加 $\partial_tZ$，并相应改变 $K$；单独验证某时刻辛性不允许把这一项删去。

几何上，辛条件保持二形式 $\sum_a\dd q_a\wedge\dd p_a$。将它作$s$次外积便得到正则体积（差一个固定符号和阶乘），故正则变换保持定向体积、$\det\mathcal J=1$。在一自由度的二维空间，行列式为1也足够；多自由度时仅有体积保持远远不够。

\begin{example}[四维体积保持的反例]
取两自由度变换
\[
Q_1=2q_1,\ P_1=p_1,\quad Q_2=q_2,\ P_2=p_2/2.
\]
总雅可比行列式为1，但 $\{Q_1,P_1\}=2,\{Q_2,P_2\}=1/2$，不是正则变换。相反，$Q=aq,P=p/a$、$a\ne0$ 的一对互逆缩放是正则的。体积守恒不能代替逐对括号检验。

按 $(q_1,p_1,q_2,p_2)$ 排列，雅可比矩阵对角元为 $2,1,1,1/2$。基本括号分别是 $2\{q_1,p_1\}$ 和 $\{q_2,p_2\}/2$。这两个因子在总行列式中互抵，但在各自的正则平面上并未抵消。
\end{example}

\begin{figure}[htbp]
\centering
\begin{tikzpicture}[>=Stealth]
\foreach \shift/\lab/\w/\h in {0/第一正则平面/2/1,5/第二正则平面/1/.5}{
\begin{scope}[xshift=\shift cm]
\draw[->](0,0)--(3,0);\draw[->](0,0)--(0,2.2);
\draw[thick,fill=blue!10](.2,.2)rectangle({.2+\w},{.2+\h});
\draw[dashed,gray](.2,.2)rectangle(1.2,1.2);
\node at(1.3,2.45){\lab};
\end{scope}}
\node at(1.3,-.45){面积乘2};\node at(6.3,-.45){面积乘$1/2$};
\end{tikzpicture}
\caption{四维反例在两个共轭平面的面积变化，虚线表示原单位方块。两个面积因子的乘积为1，却各自不保持，因而整体体积保持仍不辛。}
\label{fig:ham-volume-not-symplectic}
\end{figure}

\subsection{从作用量边界项推导四类公式}
将旧变量和新变量写进作用量，比较它们在端点上相差什么。选择不同的独立变量后，同一个变换就有不同形式的生成函数。

局部正则变换可由相空间作用量的一形式关系构造：
\[
\sum_ap_a\dd q_a-H\dd t
=\sum_aP_a\dd Q_a-K\dd t+\dd F_1.
\]
把 $F_1$ 取为 $(q,Q,t)$ 的函数，比较独立微分系数，就有 $p=(F_1)_q,P=-(F_1)_Q,K=H+(F_1)_t$。进一步令
\[
F_2=F_1+\sum_aP_aQ_a,\quad
F_3=F_1-\sum_ap_aq_a,\quad
F_4=F_1-\sum_ap_aq_a+\sum_aP_aQ_a,
\]
并在各自可反解的图中改用相应自变量。例如 $\dd F_2=\dd F_1+\sum_a(P_a\dd Q_a+Q_a\dd P_a)$，其中 $P_a\dd Q_a$ 抵消了 $\dd F_1$ 中的 $-P_a\dd Q_a$。留下 $p_a\dd q_a$ 与 $Q_a\dd P_a$，所以 $p=(F_2)_q$、$Q=(F_2)_P$。其余两类用同样的乘积微分整理，得到
\begin{align}
F_1(q,Q,t):&\quad p=(F_1)_q,\quad P=-(F_1)_Q,\label{eq:F1}\\
F_2(q,P,t):&\quad p=(F_2)_q,\quad Q=(F_2)_P,\label{eq:F2}\\
F_3(p,Q,t):&\quad q=-(F_3)_p,\quad P=-(F_3)_Q,\label{eq:F3}\\
F_4(p,P,t):&\quad q=-(F_4)_p,\quad Q=(F_4)_P.\label{eq:F4}
\end{align}
每一类都有 $K=H+(F_i)_t$，偏时间导数保持该类列出的自变量固定。用旧变量写出右端后，还需代入反变换，才能成为 $K(Q,P,t)$。

这些是局部构造，需相应混合海森矩阵可逆。例如 $F_2$ 需要 $\det(F_{2,qP})\ne0$，才能由 $p=F_{2,q}$ 解出 $P$。恒等变换有 $F_2=q\cdot P$，但不能把依赖的 $q,Q=q$ 强作独立变量构成常规 $F_1(q,Q)$ 图。存在正则变换不意味着四类生成函数在所有点或全局都同时可用。

\begin{example}[可逆点变换及其动量]
设 $Q=f(q)$，$Df$ 可逆，取 $F_2(q,P)=f(q)\cdot P$。则
\[
Q=f(q),\qquad p=(Df)^{\mathsf T}P.
\]
这是第三章位形坐标变换的正则提升。对 $Q=aq$，$P=p/a$，新量纲随坐标尺度变化，$p\dd q=P\dd Q$；若只改 $q$ 而不相应改 $p$，一般不再正则。
\end{example}

\begin{example}[移动坐标：新哈密顿量要包含时间项]
自由粒子改用 $Q=x-Vt$，取 $F_2(x,P,t)=(x-Vt)P$，故
\[
p=P,\quad Q=x-Vt,\quad K=H-VP=\frac{P^2}{2m}-VP.
\]
正则方程 $\dot Q=P/m-V,\dot P=0$ 与坐标定义一致。也可从 $L'=m(\dot Q+V)^2/2$ 直接勒让德变换复算；$K$ 不等于原系 $p^2/(2m)$。

若还希望新动量是移动系机械动量，另取
\[
\widehat F_2=(x-Vt)(\widehat P+mV)+\tfrac12mV^2t.
\]
得到 $p=\widehat P+mV,\widehat Q=x-Vt,\widehat K=\widehat P^2/(2m)$。两个描述的坐标相同、动量和边界项不同，不能将其中一套的$P$与另一套的$K$混用。

第二种生成函数的显时项为 $-V(\widehat P+mV)+mV^2/2$。将它加到 $(\widehat P+mV)^2/(2m)$ 后，一次项与常数项分别相消，因此只剩 $\widehat P^2/(2m)$。
\end{example}

\begin{example}[规范变换作为含时正则变换]
对电磁例取 $F_2(\br,\bp',t)=\br\cdot\bp'-e\chi(\br,t)$，则
\[
\bp=\bp'-e\nabla\chi,\quad\br'=\br,\quad
K=H-e\partial_t\chi.
\]
代入旧 $H$，正好成为势 $\bA'=\bA+\nabla\chi,\phi'=\phi-\chi_t$ 下的同一最小耦合形式。机械动量 $\bp'-e\bA'=\bp-e\bA$ 不变，独立验证了轨道不变；哈密顿量改变不是额外施加了一个电场。
\end{example}


\section{哈密顿--雅可比理论：从完全积分重构运动}
若能找到在运动中保持不变的一整组新坐标，就可以用它们标记初态。哈密顿--雅可比方法先求一个函数，再从它的导数找回原来的位置和动量。

\subsection{把运动转成常变量的条件}
先用含有足够多参数的主函数表示可能的轨道，再用参数和它们对应的常量来重建初值。参数必须能彼此独立地改变轨道。

希望构造第二类生成函数 $S(q,\alpha,t)$，使新动量为 $\alpha$，新坐标为 $\beta$，新哈密顿量为零。由 $p=S_q$、$K=H+S_t$，得到
\begin{equation}
\boxed{S_t+H(q,S_q,t)=0.}
\label{eq:HJ}
\end{equation}
这是一阶偏微分方程，称哈密顿--雅可比（HJ）方程，$S$称哈密顿主函数，具有作用量量纲。求到\emph{一个}无参数的 $S(q,t)$，只给一族 $p=S_q$，通常不足以重建全部初值。所需完全积分含 $s$ 个参数 $\alpha_a$，并在工作图中满足
\[
\det S_{q_a\alpha_b}\ne0.
\]
于是
\[
p_a=S_{q_a},\quad \beta_a=S_{\alpha_a},\quad
\dot\alpha_a=\dot\beta_a=0.
\]
先由 $\beta=S_\alpha$ 反解 $q(t;\alpha,\beta)$，再求 $p$，共给出$2s$个轨道常量。可任加一个不影响导数的常数，但它不算新的独立轨道参数。

沿一支光滑的真实轨道，把固定初端到可变终端的作用量记为 $S$。在固定比较时刻的终端位移为 $\delta q=\dd q-\dot q\,\dd t$，与第\ref{ch:variation}章移动端点的区别相同。内部变分由运动方程消去，终端贡献由构型边界项和积分上限项相加：
\[
\dd S=p\cdot(\dd q-\dot q\,\dd t)+L\,\dd t
=p\cdot\dd q-H\,\dd t .
\]
由独立终端变化读出 $S_q=p,S_t=-H$。多条连接、焦散或转折会使这个端点函数需分支或换图，不能把它不加条件称作全局单值函数。

\begin{example}[自由粒子的完全积分]
$H=p^2/(2m)$，取
\[
S(x,\alpha,t)=\alpha x-\frac{\alpha^2}{2m}t,
\quad S_{x\alpha}=1.
\]
直接代入HJ式为零；$p=\alpha,\beta=x-\alpha t/m$，故 $x=\beta+\alpha t/m$。$\alpha$ 为初动量，$\beta$ 为初位置；量纲 $[S]=$能量$\times$时间，故 $S_x$ 确为动量。
\end{example}

\begin{example}[恒力的完全积分与两个初值]
恒力 $F$ 的势能 $V=-Fx$，取
\[
S=(\alpha+Ft)x-\frac{\alpha^2t}{2m}
-\frac{\alpha Ft^2}{2m}-\frac{F^2t^3}{6m}.
\]
有 $S_x=\alpha+Ft$，而
\[
S_t=Fx-\frac{\alpha^2}{2m}-\frac{\alpha Ft}{m}-\frac{F^2t^2}{2m},
\]
与 $H(x,S_x)=S_x^2/(2m)-Fx$ 相消，且 $S_{x\alpha}=1$。因此
\[
p=\alpha+Ft,\quad\beta=x-\alpha t/m-Ft^2/(2m),
\quad x=\beta+\alpha t/m+Ft^2/(2m).
\]
这独立恢复牛顿恒加速度轨道；$F\to0$ 退回自由粒子完全积分。
\end{example}

\subsection{自治分离与转折点分支}
谐振子到达同一个位置时，可能向左，也可能向右。求主函数时要分开动量的正负支，到转折点后换到另一支继续。

自治 $H$ 可令
\begin{equation}
S=W(q,\alpha)-E(\alpha)t,\qquad H(q,W_q)=E.
\label{eq:HJ-time-independent}
\end{equation}
$W$ 称特征函数。求解它仍可能与原运动一样困难；HJ提供精确的重组方式，而非保证任意系统都能显式分离。

\begin{example}[谐振子的一支HJ解为何不覆盖整周期？]
设 $E>0$，$A=\sqrt{2E/(m\omega^2)}$，在 $|q|<A,p>0$ 的分支上
\[
W_q=\sqrt{2mE-m^2\omega^2q^2},
\]
其原函数可取
\[
W_+(q,E)=\tfrac12q\sqrt{2mE-m^2\omega^2q^2}
+\frac E\omega\arcsin(q/A).
\]
参数 $E$ 在该图中使 $W_{qE}=m/p\ne0$，因此 $S=W_+-Et$ 为局部完全积分。求 $S_E=\beta$ 得
\[
\frac1\omega\arcsin(q/A)-t=\beta,
\quad q=A\sin[\omega(t+\beta)].
\]
但此反正弦主值只覆盖 $p>0$、相位在 $(-\pi/2,\pi/2)$ 的区间。转折后必须改取负动量支并匹配相位；转折点 $p=0$ 时 $W_{qE}$ 发散，该生成函数图失效，而原相空间运动光滑。不能把局部反正弦时间关系直接当作全局周期计时器。
\end{example}

\begin{example}[中心势分离常数就是角动量]
对前文球坐标中心势，取
\[
W=W_r(r)+W_\theta(\theta)+\ell_z\varphi.
\]
代入HJ式并乘 $r^2$，径向与角向分离，定义常数 $\ell^2$：
\[
(W_\theta')^2+\frac{\ell_z^2}{\sin^2\theta}=\ell^2,
\qquad
(W_r')^2=2m[E-V(r)]-\ell^2/r^2.
\]
所以 $\ell^2=J^2$，并要求 $\ell^2\ge\ell_z^2$。各项根号的正负分支对应运动方向，容许区要求根号非负；由参数导数重构角度与时间仍是必需步骤。圆轨道或特殊平面选择可使某参数图退化，不能只数形式上的分离常数就保证全局完全积分。
\end{example}


\section{作用角变量、周期与可积性}
对反复重复的运动，可以用一个量记轨道大小，用另一个量记本周期走到了哪里。这两种量分别是作用变量和角变量。

\subsection{一自由度的面积与频率}
在一自由度的相图中，一条周期轨道围成一块面积。面积除以 $2\pi$ 得到作用变量，沿着这条轨道走过的进度用角变量表示。

对自治一自由度中一族光滑、紧致、无临界点的周期轨道，取真实运动方向的闭合积分
\begin{equation}
J(E)=\frac1{2\pi}\oint p\,\dd q.
\label{eq:action-variable}
\end{equation}
对通常有两个简单转折点的自然系统，$J$ 等于闭合相面积除以 $2\pi$。$q$ 为角坐标的持续转动则要在圆柱的基本周期回路积分，未必是展开平面内的围面积。

在非临界的轨道族上，对 $H(q,p)=E$ 求导得 $H_p\partial_Ep=1$。两支合起来有
\[
\frac{\dd J}{\dd E}=\frac1{2\pi}\oint\frac{\dd q}{\dot q}
=\frac{T(E)}{2\pi}.
\]
简单转折处 $p=0$ 使端点项为零，$1/\dot q$ 的奇异性可积；在分离曲线处周期发散则另论。反解 $E=H(J)$ 后，取局部生成函数 $W(q,J)$，定义
\begin{equation}
\theta=W_J,\qquad\dot\theta=H_J=2\pi/T(E),\qquad\dot J=0.
\label{eq:angle-variable}
\end{equation}
$\theta$ 的周期为 $2\pi$，不是构型角必然相等的另一个名字。若改用不除以 $2\pi$ 的作用，保持正则性就需将其共轭角同时除以 $2\pi$。

\begin{example}[谐振子的全周期相位与局部生成函数]
取
\[
q=\sqrt{\frac{2J}{m\omega}}\sin\theta,\qquad
p=\sqrt{2m\omega J}\cos\theta,\quad J>0.
\]
沿一周积分 $p\dd q$ 得 $J=E/\omega$。反变换是
\[
J=\frac{p^2+m^2\omega^2q^2}{2m\omega},\qquad
\theta=\operatorname{atan2}(m\omega q,p)\pmod{2\pi}.
\]
直接求偏导可验证 $\{\theta,J\}=1$，因此 $H=\omega J$ 给 $\dot\theta=\omega$。角在任一选定切口图上光滑，绕周后加 $2\pi$；原点 $J=0$ 没有可辨相位，不能纳入同一角坐标图。

在 $\sin\theta\ne0$ 的适当分支，也可用第一类生成函数
\[
F_1(q,\theta)=\tfrac12m\omega q^2\cot\theta,
\quad p=(F_1)_q,\quad J=-(F_1)_\theta
=\frac{m\omega q^2}{2\sin^2\theta}.
\]
由此得 $K=\omega J$。$\sin\theta=0$ 时这类投影图失效，并不表示非零动量的物理状态消失；应换图或直接使用上述正反变换。
\end{example}

\begin{figure}[htbp]
\centering
\begin{tikzpicture}[>=Stealth]
\draw[thick,fill=blue!7](0,0)circle(1.6);
\draw[->](-2,0)--(2.3,0)node[right]{$q\sqrt{m\omega}$};
\draw[->](0,-1.9)--(0,2.2)node[above]{$p/\sqrt{m\omega}$};
\draw[dashed](0,0)--(1.13137,1.13137);
\fill(1.13137,1.13137)circle(2pt);
\draw[->](0,.7)arc(90:45:.7);\node at(.4,.85){$\theta$};
\node[fill=blue!7,inner sep=1pt] at(-.55,-.65){面积$=2\pi J$};
\node[above right]at(1.13137,1.13137){$\sqrt{2J}$};
\end{tikzpicture}
\caption{使用互逆缩放后的共轭轴，两坐标均有作用量平方根单位。半径为$\sqrt{2J}$，相面积为$2\pi J$；相位从正动量轴顺时针量起，与正文正弦/余弦约定一致。}
\label{fig:ham-action-angle}
\end{figure}

\subsection{单摆分离曲线附近：有限作用与无限周期}
单摆越接近翻过最高点的门槛，在最高点附近停留得越久。每圈所需时间可趋于无穷，围出的相面积却仍有限。

\begin{example}[周期发散不要求相面积发散]
回到 $H=p^2/(2I)+K_0(1-\cos q)$，$0<E<2K_0$ 时令
\[
k=\sqrt{E/(2K_0)}=\sin(q_{\max}/2)<1.
\]
由 $\dot q=p/I$，四分之一周期为 $\sqrt{I/(2K_0)}\int_0^{q_{\max}}\dd q/\sqrt{E/K_0-1+\cos q}$。作 $\sin(q/2)=k\sin u$，得到
\[
T(E)=4\sqrt{I/K_0}\int_0^{\pi/2}\frac{\dd u}{\sqrt{1-k^2\sin^2u}}.
\]
$k\to0$ 恢复 $2\pi\sqrt{I/K_0}$；$k\to1^-$ 时端点附近被积函数趋 $1/\cos u$，积分对数发散，不能继续用固定小振动频率。

分离曲线的正支为 $p_+(q)=2\sqrt{IK_0}\cos(q/2)$，$-\pi\le q\le\pi$。因此摆动侧的极限作用
\[
J_{\rm lib,s}=\frac1{2\pi}\,2\int_{-\pi}^{\pi}p_+(q)\dd q
=\frac{8\sqrt{IK_0}}\pi
\]
仍有限，而 $\dd J/\dd E=T/(2\pi)$ 发散。正向转动侧只沿一支绕圆柱一圈，其极限为 $4\sqrt{IK_0}/\pi$。两侧采用不同基本回路，不能要求它们的作用值无区别地相等；分离曲线本身不是正则的作用角环面。
\end{example}

\begin{figure}[htbp]
\centering
\begin{tikzpicture}
\begin{axis}[width=.78\textwidth,height=5.1cm,axis lines=left,xlabel={$E/K_0$},ylabel={$T/(2\pi\sqrt{I/K_0})$},xmin=0,xmax=2.03,ymin=.95,ymax=3,xtick={0,.5,1,1.5,2},ytick={1,1.5,2,2.5,3}]
\addplot[blue!65!black,very thick]table[x=energy,y=period,col sep=space]{figures/ch06_pendulum_period.dat};
\addplot[gray,dashed,domain=0:1.99]{1};
\draw[red!70!black,dashed](axis cs:2,.95)--(axis cs:2,3);
\end{axis}
\end{tikzpicture}
\caption{单摆摆动周期的数值积分，与小振动常周期比较。蓝线仅计算$E/K_0<2$的正则轨道，红虚线标出分离能量；不在临界点伪造有限周期。数据由伴随核查脚本生成。}
\label{fig:ham-pendulum-period}
\end{figure}

\subsection{开普勒作用：不同频率何时相同？}\label{sec:kepler-action}
行星一边改变与太阳的距离，一边绕太阳转。用两个作用变量分别记录这两种运动，再看它们的频率为何能配合成闭合轨道。

\begin{example}[有界平面库仑势的径向作用]
取 $V=-\kappa/r,E<0$，并固定正向角动量 $\ell>0$，非圆椭圆的半长轴、偏心率满足
\[
a=-\frac\kappa{2E},\qquad \ell^2=m\kappa a(1-e^2),\quad0<e<1.
\]
角向作用 $J_\varphi=\ell$，径向往复作用为
\[
J_r=\frac1\pi\int_{r_-}^{r_+}\sqrt{2m(E+\kappa/r)-\ell^2/r^2}\,\dd r.
\]
以偏近点角 $u$ 代入 $r=a(1-e\cos u),0\le u\le\pi$，则
\[
p_r=\sqrt{m\kappa/a}\frac{e\sin u}{1-e\cos u},\qquad
J_r=\frac{\sqrt{m\kappa a}}\pi\int_0^\pi
\frac{e^2\sin^2u}{1-e\cos u}\dd u.
\]
用恒等式 $e^2\sin^2u/(1-e\cos u)=1+e\cos u-(1-e^2)/(1-e\cos u)$，再由 $\tan(u/2)$ 代换得 $\int_0^\pi\dd u/(1-e\cos u)=\pi/\sqrt{1-e^2}$，故
\[
J_r=\sqrt{m\kappa a}-\ell
=\kappa\sqrt{\frac{m}{-2E}}-\ell.
\]
反解能量为
\[
H=-\frac{m\kappa^2}{2(J_r+J_\varphi)^2},\quad
\omega_r=\omega_\varphi=\frac{m\kappa^2}{(J_r+J_\varphi)^3}
=\sqrt{\frac\kappa{ma^3}}.
\]
两角频率相同对应闭合的开普勒椭圆，并与式 \eqref{eq:kepler3} 相容（两体相对问题取 $m$ 为约化质量）。第\ref{sec:central}节的平面极角 $\theta$ 在本例记为 $\varphi$，正向平面轨道的角动量 $J$ 在这里记为 $\ell$。圆轨道 $J_r=0$ 是此作用域的边界，零角动量碰撞轨道也不在当前正则环面图内。
\end{example}

\subsection{多自由度可积性与KAM的条件}
几个自由度的系统若有足够多、彼此相容的守恒量，轨道便受到更多限制。加入小扰动后，有些规则轨道可以保留，还须检查频率条件。

对自治 $s$ 自由度系统，刘维尔可积要求包含 $H$ 在内的$s$个函数独立积分 $F_i$，且 $\{F_i,F_j\}=0$。若共同正则水平集紧致、连通，适当邻域中可用环面上的作用角变量，$H=H(J)$，$\dot\theta_a=\partial_{J_a}H,\dot J=0$。一般作用是基本回路上的 $\oint\sum p_a\dd q_a/(2\pi)$，只有在合适分离坐标中才直接变成单项 $(2\pi)^{-1}\oint p_a\dd q_a$。

非平凡运动在环面上周期闭合的条件是存在 $T>0$ 使每个 $\omega_aT$ 都是 $2\pi$ 的整数倍。部分整数共振关系可使轨迹限于较低维环面，但不足以保证所有相位同时复原。非紧致散射、临界能量、分离曲线和坐标覆盖障碍不能被“有$s$个常数”一句话排除。可积也不等于一定能在原坐标中初等分离变量。

\begin{modernbox}[title={进一步探索：近可积不等于任意摄动都稳定}]
对 $H_0(J)+\epsilon H_1(J,\theta)$，科尔莫戈罗夫、阿诺德和莫泽（KAM）定理的一种常见版本要求足够光滑或解析、摄动小，以及频率非退化，例如 $\det\partial_J^2H_0\ne0$；保留频率满足 $|k\cdot\omega|\ge\gamma|k|^{-\tau}$，$k\in\mathbb Z^s\setminus\{0\}$，参数按定理选取。结论是许多环面保留，不保证所有共振环面，也不直接给出天体寿命。谐振子 $H_0=\sum\omega_aJ_a$ 在此海森矩阵条件下退化，需另用合适定理。

自治一自由度鞍点的指数不稳定不等于混沌；线性拍或短时能量误差也不是混沌判据。继续学习需摄动、微分方程及非共振数论条件。
\end{modernbox}


\section{计算与模拟：辛性、能量与误差阶}
数值积分每走一步，也是在相空间做一次变换。除了比较能量和位置，还可以检查这一步是否保留原方程的位置与动量配对结构。

\subsection{显式欧拉法与辛欧拉法的不同结构}
用弹簧振子比较两种更新顺序。一种同时用旧数据更新位置和动量，另一种先更新动量，再用新动量更新位置，长期表现可以很不同。

取无量纲谐振子 $H=(q^2+p^2)/2$，步长 $h$ 固定。显式欧拉法为
\[
q_{n+1}=q_n+hp_n,\quad p_{n+1}=p_n-hq_n,
\quad A_E=\begin{pmatrix}1&h\\-h&1\end{pmatrix}.
\]
它满足 $H_{n+1}=(1+h^2)H_n$、$\det A_E=1+h^2$，任意非零步长都会产生非物理的面积和能量增长。

先踢动量、后漂移位置的辛欧拉法为
\[
p_{n+1}=p_n-hq_n,\quad q_{n+1}=q_n+hp_{n+1},\quad
A_S=\begin{pmatrix}1-h^2&h\\-h&1\end{pmatrix}.
\]
两步是正则剪切的复合，或直接核验 $A_S^{\mathsf T}\Omega A_S=\Omega$。代入还可证明它精确保持二次型
\[
\widetilde H_S=\tfrac12(q^2+p^2-hqp).
\]
仅当 $|h|<2$，这个二次型正定，保证轨道有界；$|h|=2$ 的重本征值一般有Jordan增长。辛性不自动保证任意步长稳定，更不意味着逐步精确保持原能量。

对一般自治 $H(q,p)$，第二类生成函数 $F_2(q,P)=q\cdot P+hH(q,P)$ 给
\[
p=P+hH_q(q,P),\qquad Q=q+hH_p(q,P).
\]
这是一般辛欧拉法的一种形式，通常需隐式解 $P$。不能把只适于可分离模型的显式“踢--漂移”直接套用于任意含磁矢势或坐标相关惯性的哈密顿量。

\subsection{Störmer--Verlet：分裂流的对称复合}
把每个时间步拆成半次受力、一次移动和再半次受力。这样的前后对称安排既便于计算，也使反向走一步能够撤销正向的一步。

若 $H=T(p)+V(q)$，两部分精确流分别是 $q$ 漂移、$p$ 踢动。以半踢--全漂移--半踢复合，得
\begin{equation}
\begin{aligned}
p_{n+1/2}&=p_n-\tfrac h2\nabla V(q_n),\\
q_{n+1}&=q_n+h\nabla T(p_{n+1/2}),\\
p_{n+1}&=p_{n+1/2}-\tfrac h2\nabla V(q_{n+1}).
\end{aligned}
\label{eq:ham-verlet}
\end{equation}
第\ref{sec:orbit-compute}节的速度 Verlet 就是这三步。每一流辛，复合也辛；把 $h$ 变为 $-h$ 并逆序执行恢复原态，因此它还具有时间对称性。对光滑问题，泰勒展开相对精确一步的误差是 $O(h^3)$，在合适固定时间区间内全局误差为 $O(h^2)$。标准逐状态自适应改变步长一般会破坏原有的辛映射论证，需专门构造。

例如对常见的$T=p^{\mathsf T}p/(2m)$，先把分裂步展开：
\[
q_{n+1}=q_n+hp_n/m-\frac{h^2}{2m}\nabla V(q_n),
\]
\[
p_{n+1}=p_n-h\nabla V(q_n)-\frac{h^2}{2m}V^{\prime\prime}(q_n)p_n+O(h^3).
\]
精确方程有$\ddot q=-\nabla V/m,\ddot p=-V^{\prime\prime}p/m$，故两式逐项匹配到$h^2$，直接核验二阶精度。

\begin{example}[保住能量形状，仍会积累相位误差]
对单位谐振子，Verlet矩阵为
\[
A_V=\begin{pmatrix}1-h^2/2&h\\-h(1-h^2/4)&1-h^2/2\end{pmatrix}.
\]
它精确保持 $\widetilde H_V=[(1-h^2/4)q^2+p^2]/2$。在 $|h|<2$ 内，本征值为 $e^{\pm i\vartheta}$，由迹得
\[
\cos\vartheta=1-h^2/2,\quad
\widetilde\omega=\frac2h\arcsin(h/2)=1+\frac{h^2}{24}+O(h^4).
\]
原能量只作小幅振荡，但相位误差约按 $(\widetilde\omega-1)t$ 积累。$h\to0$ 的收敛与固定 $h$ 下任意长时间相位精确不是同一结论。这个模型中的精确保持二次型，也不代表一般非线性系统有同样简单、每步精确保留的修正能量。

单步矩阵的行列式为1，迹为 $2-h^2$，故 $\cos\vartheta=1-h^2/2$。在 $h>0$ 且足够小时取 $\vartheta=2\arcsin(h/2)=h+h^3/24+O(h^5)$，离散频率 $\vartheta/h$ 高于准确频率。经历 $n$ 步后，这一频率差给出 $n(\vartheta-h)$ 的相位增量。
\end{example}

\begin{figure}[htbp]
\centering
\begin{tikzpicture}
\begin{axis}[width=.47\textwidth,height=5cm,ymode=log,xlabel={$t$},ylabel={$H_n/H_0$},xmin=0,xmax=60,ymin=.8,ymax=600,legend style={draw=none,font=\scriptsize,at={(.03,.97)},anchor=north west}]
\addplot[red!70!black,thick]table[x=time,y=euler,col sep=space]{figures/ch06_integrator_energy.dat};\addlegendentry{显式欧拉法}
\addplot[blue!65!black,thick]table[x=time,y=symplectic,col sep=space]{figures/ch06_integrator_energy.dat};\addlegendentry{辛欧拉法}
\addplot[green!45!black,thick,dashed]table[x=time,y=verlet,col sep=space]{figures/ch06_integrator_energy.dat};\addlegendentry{Verlet}
\end{axis}
\begin{axis}[at={(.52\textwidth,0)},width=.47\textwidth,height=5cm,xmode=log,ymode=log,xlabel={$h$},ylabel={末态误差},legend style={draw=none,font=\scriptsize,at={(.97,.03)},anchor=south east}]
\addplot[blue!65!black,thick,mark=*]table[x=step,y=symplectic,col sep=space]{figures/ch06_integrator_convergence.dat};\addlegendentry{辛欧拉法}
\addplot[green!45!black,thick,dashed,mark=square*]table[x=step,y=verlet,col sep=space]{figures/ch06_integrator_convergence.dat};\addlegendentry{Verlet}
\end{axis}
\end{tikzpicture}
\caption{同初态$(q,p)=(1,0)$。左图固定$h=0.1$，纵轴为对数；右图在固定$t=10$比较解析末态，显示一阶与二阶收敛。有限区间比较不替代长期相位核查。}
\label{fig:ham-integrators}
\end{figure}

\subsection{伴随脚本如何独立核验}
脚本分别检验三件事。公式代回方程是否成立，步长减小后误差是否按预期下降，以及许多周期之后能量和相位各自偏了多少。

标准库脚本 \texttt{review/verify\_ch06.py} 输出
\begin{center}\texttt{ch06\_verification\_results.json}\end{center}
并生成上述周期和算法图的数据。检查包括：勒让德偏导、罗斯离心项、磁括号与规范变换、相流面积、正则与非正则反例、含时生成函数、HJ完全积分与重构、谐振子/单摆/开普勒作用，以及积分器辛条件、稳定步长、修正二次型和收敛阶。长期比较还记录能量与相位，避免只凭曲线看似闭合就认定算法正确。

\begin{modernbox}[title={进一步探索：为什么辛方法常有良好的长期能量表现？}]
足够光滑、小步长且轨道位于合适区域时，后向误差分析可将辛算法联系到邻近修正哈密顿系统，解释长期能量表现。这里并未保证任意步长或临界轨道；本章只证明了线性模型的明确二次型。一般理论见章末研究综述，需数值微分方程基础。
\end{modernbox}


\section*{关键概念与常见误区}
\begin{itemize}[leftmargin=1.8em]
\item 从速度算出动量后，要检查能否反解。同一动量在全局可能对应多个速度，即使每一支在局部都可反解。
\item 正则动量用于哈密顿方程，机械动量由实际速度确定。遇到含时的坐标变换，还要重算新变量中的哈密顿量。
\item 相空间里一团初态可以变形而保持体积。正则变换还要求保持各变量间的基本括号，体积正确仍不足以通过这项检查。
\item 守恒要把量本身显含时间的变化与沿轨道的变化一起算，即 $\partial_tf+\{f,H\}=0$。
\item 生成函数只有在所选独立变量可用的区域内才能使用。若转折点处不能反解，就换一组变量接着描述同一条轨道。
\item 用主函数重建运动需要足够多的独立参数。作用角坐标适合规则的周期轨道，接近分离曲线或静止点时应检查是否仍可用。
\item 用辛算法时，分别检查步长、能量和相位。能量曲线即使很平稳，许多周期后也可能走到错误的位置。
\end{itemize}

\section*{本章小结}
\addcontentsline{toc}{section}{本章小结}
以正则动量代替速度后，动力学在相空间由反对称结构生成。作用量边界项说明可用哪些端点与生成函数；泊松括号联系时间演化、对称性和守恒量；HJ完全积分提供回到轨道的常量反解，作用角把正则周期族组织成匀速相位运动。正则性、紧致性、分支和数值误差都需在使用这些工具时检查。

第\ref{ch:relativity}章保留作用量与正则方法，但重新检验牛顿绝对时空。哈密顿形式本身不规定洛伦兹或伽利略对称性，必须由具体物理作用量和时空结构决定。非相对论 $H=p^2/(2m)$ 的适用范围也将在相对论粒子的能量动量关系中重新界定。

\section*{分层习题}
\addcontentsline{toc}{section}{分层习题}
先分清位置、速度和正则动量，再写选定变量下的方程。换变量时，可以用基本泊松括号或代回旧方程复核。

\begin{bookexercises}
\item \exlevel{基础} \textbf{速度一次项。} 对 $L=m\dot q^2/2+a(q)\dot q-V(q)$ 求 $p,H$，用正则方程恢复牛顿方程。证明一维光滑 $a(q)\dot q$ 在局部只改变固定构型端点作用量的边界项。
\item \exlevel{基础} \textbf{局部与全局。} 对正文四次速度模型列出 $p=0$ 的三个速度及相应 $H$。非零海森矩阵说明什么？对 $L=(\dot x-y)^2/2$，为什么普通变换不能反解两速度？
\item \exlevel{基础} \textbf{边界项。} 给相空间作用量加入 $-[pq]_{t_1}^{t_2}$，完整写出变分，说明固定动量端点如何得到正则方程；任意固定全部$q,p$两端为什么一般过约束？
\item \exlevel{综合} \textbf{移动系。} 对正文两种生成函数分别求新坐标、动量、哈密顿量和轨道，比较原系机械能与移动系机械能，说明为什么不能交叉混用。
\item \exlevel{综合} \textbf{罗斯约化。} 对平面中心势在固定 $\ell$ 下导出约化径向方程和角度重构；若相应角坐标不循环，为什么不能仍将其动量固定为常数？
\item \exlevel{综合} \textbf{密度与刘维尔。} 自由粒子的初密度为 $C\exp[-q^2/(2\sigma_q^2)-p^2/(2\sigma_p^2)]$。求 $\rho(q,p,t)$、位置方差及相协方差行列式，解释扩散形状不等于相体积增加。
\item \exlevel{综合} \textbf{含时守恒与生成元。} 对 $G=tp-mx$ 验证守恒并写出所生的 $\delta x,\delta p$。再验证角动量括号 $\{J_x,J_y\}=J_z$，说明泊松定理为何不保证新量独立。
\item \exlevel{综合} \textbf{磁括号。} 在 $\bA=(0,Bx,0)$ 下算 $\{\pi_x,\pi_y\}$，比较 $p_y$ 与 $mv_y$ 的演化；规范改变后哪个动量保持物理数值不变？
\item \exlevel{综合} \textbf{体积与辛。} 核验正文四维反例的行列式和两对基本括号。改成 $Q_a=a_aq_a,P_a=p_a/a_a$ 后是否正则？
\item \exlevel{综合} \textbf{含时缩放。} 对谐振子取 $F_2(q,P,t)=e^{\lambda t}qP$，常数$\lambda$单位为时间的负一次方。求 $Q,P,K$，再由原正则方程独立验证新方程中的含时项。
\item \exlevel{综合} \textbf{生成函数图。} 恒等变换为何有常规$F_2$却不能用常规独立$(q,Q)$作同一$F_1$？谐振子$F_1(q,\theta)$在何处失效，物理轨道是否在那里终止？
\item \exlevel{综合} \textbf{完全积分。} 对恒力例逐项验证HJ方程与 $S_{x\alpha}=1$，再由 $\beta=S_\alpha$ 重构 $x,p$。只有一个不含可调参数的HJ特解是否足够？
\item \exlevel{综合} \textbf{HJ分支。} 对谐振子算 $W_E,W_{qE}$，说明正根分支的相位范围，并解释转折后的负根为何必需。
\item \exlevel{综合} \textbf{作用与相位。} 用相面积和直接积分两种方式算谐振子 $J$，验证 $\{\theta,J\}=1$。若 $J_{\rm alt}=2\pi J$，其共轭角周期及演化速率是什么？
\item \exlevel{拓展} \textbf{分离曲线。} 算单摆摆动侧、正转动侧的分离极限作用，说明因子2来自何种回路差别；证明摆动周期发散而作用有限不矛盾。
\item \exlevel{拓展} \textbf{开普勒可积性。} 从 $J_r=\kappa\sqrt{m/(-2E)}-\ell$ 反解 $H$，求两频率。计算二维 $H(J_r,J_\varphi)$ 的海森矩阵秩，能否据其退化就说开普勒不可积？能否直接套正文列出的非退化KAM版本？
\item \exlevel{拓展} \textbf{辛算法与稳定。} 验证辛欧拉法、Verlet的矩阵辛条件及各自修正二次型。辛欧拉法在$h=2$时对$(q,p)=(1,1)$与$(1,0)$的行为是否相同？
\item \exlevel{拓展} \textbf{计算检验。} 运行伴随脚本，分别检查 $\dd J/\dd E=T/(2\pi)$、HJ残差、积分器有限区间收敛和长期能量/相位。为什么低能量误差或单位行列式不能单独证明算法轨道准确？
\end{bookexercises}

\section*{习题解答}
先用同一题号找到解答。核对时，从第一条方程开始比较自己的做法，再检查最后的符号、单位和适用条件。
\begin{bookexercises}
\item $p=m\dot q+a,H=(p-a)^2/(2m)+V$，故 $\dot q=(p-a)/m,\dot p=a'\dot q-V'$，相减得 $m\ddot q=-V'$。$a(q)\dot q=\dd[\int^qa(s)\dd s]/\dd t$，局部固定$q$端点不改内部方程。
\item $u=0,\pm1$，对应$H=V,V+E_*/4,V+E_*/4$。每点局部可反解，但三个分支不能拼成任意全局单值关系。奇异例的$p_y=0$与$\dot y$无关，需约束理论。

令 $u=\dot q/v_*$。由 $p=(E_*/v_*)(u^3-u)$，$p=0$ 时 $u(u-1)(u+1)=0$。代入 $H=E_*(3u^4/4-u^2/2)+V$ 得三支能量。速度海森矩阵在这一维为 $(E_*/v_*^2)(3u^2-1)$，在这三点均非零，所以可各自局部反解。第二例的海森矩阵为 $\operatorname{diag}(1,0)$，$p_x$ 只确定 $\dot x$。
\item 边界项为$-[q\delta p]$，内部仍为$(\dot q-H_p)\delta p-(\dot p+H_q)\delta q$。$2s$个初始数据已定一阶系统；任意额外两端条件通常不相容，不能由边界项消失推断存在唯一解。

原作用量边界项为 $[p\,\delta q]$。减去的端点乘积贡献 $-[p\,\delta q+q\,\delta p]$，相加正好只剩 $-[q\,\delta p]$。固定两端动量后，独立的内部变分分别要求 $\dot q=H_p$、$\dot p=-H_q$。
\item 第一种$P=p,K=P^2/(2m)-VP$；第二种$\widehat P=p-mV,\widehat K=\widehat P^2/(2m)$，两者$Q=x-Vt$相同。原系能量为$p^2/(2m)$，移动系机械能为$(p-mV)^2/(2m)$。

第一种 $F_2=(x-Vt)P$ 给 $\dot P=0$、$\dot Q=P/m-V$，所以 $Q=Q_0+(P/m-V)t$。第二种 $\widehat F_2=(x-Vt)(\widehat P+mV)+mV^2t/2$ 给 $\dot{\widehat P}=0$、$\dot Q=\widehat P/m$，所以 $Q=Q_0+\widehat P t/m$。当 $\widehat P=P-mV$ 时两轨迹重合，换回 $x=Q+Vt$ 得同一原系轨道。
\item $\mathcal R=m\dot r^2/2-V-\ell^2/(2mr^2)$，$m\ddot r=-V'+\ell^2/(mr^3)$，$\dot\varphi=-\mathcal R_\ell$。非循环时$\dot\ell=L_\varphi$一般非零，须联立其演化。

从 $L=m(\dot r^2+r^2\dot\varphi^2)/2-V$ 得 $\ell=mr^2\dot\varphi$。先从 $L$ 减去 $\ell\dot\varphi$ 再代入 $\dot\varphi=\ell/(mr^2)$，角动能与被减项合为 $-\ell^2/(2mr^2)$，由此得到正确的离心项。给定径向轨道后用 $\varphi(t)=\varphi(0)+\int_0^t\ell/[mr(s)^2]\dd s$ 重构。
\item $\rho=C\exp[-(q-pt/m)^2/(2\sigma_q^2)-p^2/(2\sigma_p^2)]$，位置方差为$\sigma_q^2+t^2\sigma_p^2/m^2$，协方差为$t\sigma_p^2/m$，行列式仍$\sigma_q^2\sigma_p^2$。精细密度沿流不变，位置边缘却展宽。

用逆流 $q_0=q-pt/m$、$p_0=p$ 代入初密度。初始两变量独立，故 $\operatorname{Var}(q_0+tp_0/m)$ 无交叉项，$\operatorname{Cov}(q,p)=t\operatorname{Var}(p_0)/m$。相协方差行列式中 $t^2\sigma_p^4/m^2$ 的两项相消。
\item $G_t=p,\{G,H\}=-p$，$\delta x=\epsilon t,\delta p=\epsilon m$。角动量括号循环封闭在已有三个分量中；括号也可能为零，不保证函数独立。

写 $J_x=yp_z-zp_y$、$J_y=zp_x-xp_z$。只有 $(z,p_z)$ 一对同时贡献两项，故 $\{J_x,J_y\}=(-p_y)(-x)-y p_x=xp_y-yp_x=J_z$。
\item $\{\pi_x,\pi_y\}=eB$；$p_y$守恒而$\dot\pi_y=-eBv_x$。机械动量$\pi$规范不变，正则$p$随$e\nabla\chi$改变。

此规范下 $\pi_x=p_x$、$\pi_y=p_y-eBx$，由 $\{p_x,x\}=-1$ 立得磁括号。$y$ 循环给 $\dot p_y=0$，所以 $\dot\pi_y=-eB\dot x$；规范变化使 $p$ 与 $eA$ 增加相同的梯度项。
\item 原反例行列式1，括号分别2、1/2，不正则。互逆缩放则全部基本括号保持，是正则变换。

按 $(q_1,p_1,q_2,p_2)$ 排列，导数矩阵对角为 $(2,1,1,1/2)$。行列式为1，但两个基本括号分别乘以2和 $1/2$。互逆缩放时 $\{a_i q_i,p_j/a_j\}=(a_i/a_j)\delta_{ij}=\delta_{ij}$（$a_i\ne0$），其余基本括号仍为零。
\item $Q=e^{\lambda t}q,p=e^{\lambda t}P$，
\[K=\frac{e^{2\lambda t}P^2}{2m}+\frac{m\omega^2e^{-2\lambda t}Q^2}{2}+\lambda QP.\]
从而$\dot Q=e^{2\lambda t}P/m+\lambda Q,\dot P=-m\omega^2e^{-2\lambda t}Q-\lambda P$。末两项来自$F_{2,t}$。

直接微分 $Q=e^{\lambda t}q$、$P=e^{-\lambda t}p$，得 $\dot Q=\lambda Q+e^{\lambda t}p/m$、$\dot P=-\lambda P-e^{-\lambda t}m\omega^2q$。再用反解代入，和 $K$ 的正则方程逐项一致。
\item 恒等变换的$q,Q=q$不能独立变化，投影图退化；$F_2=qP$没有此问题。谐振子的$\sin\theta=0$处$F_1$失效，但$J>0$的原轨道可光滑穿过$q=0$。
\item 以正文的 $S$ 为候选完全积分。需一个带足够独立参数且混合海森矩阵可逆的完全积分，不能仅凭一个特解覆盖全部初值。

将 $p=S_x=\alpha+Ft$ 代入 $H=p^2/(2m)-Fx$，展开为 $\alpha^2/(2m)+\alpha Ft/m+F^2t^2/(2m)-Fx$，恰与 $S_t$ 相消。对 $S_x$ 求 $\alpha$ 偏导即为1。由 $\beta=S_\alpha=x-\alpha t/m-Ft^2/(2m)$ 解出 $x=\beta+\alpha t/m+Ft^2/(2m)$，可见 $\alpha=p(0),\ \beta=x(0)$，两个初值均可独立指定。
\item 正支$W_E=\arcsin(q/A)/\omega$、$W_{qE}=m/p$，相位限于$(-\pi/2,\pi/2)$的选定区间。转折后$p$符号改变，换支并匹配累计相位，不能重置成反正弦主值。

写 $W(q,E)=\int_0^q\sqrt{2mE-m^2\omega^2s^2}\dd s$。在开区间 $|q|<A$ 内对 $E$ 求导，被积函数变为 $m/\sqrt{2mE-m^2\omega^2s^2}$，积分即得上式，再对 $q$ 求导给 $m/p$。
\item $J=E/\omega$，直接对$\theta=\operatorname{atan2}(m\omega q,p)$求偏导给括号1。$J_{\rm alt}$的共轭角为$\theta/(2\pi)$，周期1，速率$\omega/(2\pi)$。

椭圆半轴为 $A=\sqrt{2E/(m\omega^2)}$、$p_0=m\omega A$，故相面积 $\pi A p_0=2\pi E/\omega$。直接沿运动取 $q=A\sin\theta,p=p_0\cos\theta$，有 $\oint p\,\dd q=A p_0\int_0^{2\pi}\cos^2\theta\,\dd\theta=\pi A p_0$。设 $D=p^2+(m\omega q)^2$，则 $\theta_q=m\omega p/D,\ \theta_p=-m\omega q/D$。由 $J_p=p/(m\omega),\ J_q=m\omega q$，得 $\theta_qJ_p-\theta_pJ_q=1$。
\item 分别为$8\sqrt{IK_0}/\pi$和$4\sqrt{IK_0}/\pi$。摆动闭合曲线包含正负两动量支，正转动一周只用正支；有限$J$可以具有发散的能量导数，正如$\dd J/\dd E=T/(2\pi)$。

在势垒能量 $E=2K_0$、$-\pi\le q\le\pi$ 时，$p_\pm=\pm2\sqrt{IK_0}\cos(q/2)$。正转动侧取正支积分，$\int_{-\pi}^{\pi}p_+\dd q=8\sqrt{IK_0}$；摆动闭环的回程积分与去程同号，面积加倍，再除 $2\pi$ 得所列作用。周期积分的时间元为 $I\dd q/|p|$，分离极限下近 $q=\pi$ 如 $\dd q/|\pi-q|$，故周期对数发散，作用积分本身仍收敛。
\item $H=-m\kappa^2/[2(J_r+J_\varphi)^2]$，两频率相同。海森矩阵的每个元素都等于同一个非零二阶导，秩1、行列式0；这是频率退化而非不可积，正文所列非退化KAM版本不能直接用。

令 $S=J_r+J_\varphi$，由 $S=\kappa\sqrt{m/(-2E)}$ 解得上述 $H$。两频率均为 $m\kappa^2/S^3$，海森矩阵为 $-3m\kappa^2/S^4$ 乘以全1的 $2\times2$ 矩阵，因此两行相同且非零。
\item 直接计算 $\det A_S=1-h^2+h^2=1$，$\det A_V=(1-h^2/2)^2+h^2(1-h^2/4)=1$。在二维，任意矩阵满足 $A^{\mathsf T}\Omega A=(\det A)\Omega$，故两步均辛。对辛欧拉，用 $q'=q+hp'$、$p=p'+hq$，有
\[
q'^2+p'^2-hq'p'=q^2+p'^2+hqp'=q^2+p^2-hqp.
\]
对 Verlet 记 $d=1-h^2/4$、$a=1-h^2/2$，则 $q'=aq+hp$、$p'=-hdq+ap$。展开 $dq'^2+p'^2$ 时交叉项相消，其余为 $(a^2+h^2d)(dq^2+p^2)=dq^2+p^2$。在 $h=2$，$A_S=-I+N$、$N=\bigl(\begin{smallmatrix}-2&2\\-2&2\end{smallmatrix}\bigr)$，$N^2=0$，所以 $A_S^n=(-1)^n(I-nN)$。$(1,1)^{\mathsf T}$ 被 $N$ 消去而保持有界，$(1,0)^{\mathsf T}$ 则出现线性增长项。
\item 前两项核验解析模型/变量关系，有限时间步长比较核验数值阶；长期相位另记录。保持体积不是保持所有辛结构，更不控制离散频率偏差。

例如单摆 $E=0.8$ 时脚本给 $T=7.11008$、$\dd J/\dd E=1.13160$，与 $T/(2\pi)$ 一致。谐振子末态误差在 $h=0.05$ 到 $0.025$ 时，辛欧拉从 $0.01423$ 降到 $0.006950$，Verlet 从 $9.0388\times10^{-4}$ 降到 $2.2595\times10^{-4}$。约百周期的 $h=0.1$ 测试中，Verlet 最大相对原能量误差约 $0.00250$，累计相位差仍约 $0.262\,\mathrm{rad}$。检验 HJ 式时将 $S_t$ 和 $H(q,S_q)$ 分别计算后相加，确认抵消。
\end{bookexercises}

\section*{进一步阅读与来源}
链接复核：2026-11-08。模型按本书正则括号、生成函数和$J=(2\pi)^{-1}\oint p\dd q$约定重新推导，数值均为教学设定。
\begin{bookreading}
\item David Tong，\href{https://davidtong.org/pdfs/teaching/classical-dynamics/clas4.pdf}{The Hamiltonian Formalism}，可对照正则方程、括号、作用角和HJ方法，比较时核对作用归一化与生成元符号。
\item Iain Stewart，MIT 8.09，\href{https://ocw.mit.edu/courses/8-09-classical-mechanics-iii-fall-2014/f00f7f68ac7ba346a0868efb7430582c_MIT8_09F14_Chapter_4.pdf}{Chapter 4: Canonical Transformations, Hamilton--Jacobi Equations, and Action--Angle Variables}，用于进一步研究正则结构与完全积分。
\item MIT 8.09，\href{https://ocw.mit.edu/courses/8-09-classical-mechanics-iii-fall-2014/a55cb49d8706f671ccfdac7c6937cddc_MIT8_09F14_pset5.pdf}{Problem Set 5}，训练含时正则变换、括号及HJ重构。本章题目按当前模型另行编写。
\item E. Hairer、C. Lubich、G. Wanner，\href{https://www.cambridge.org/core/journals/acta-numerica/article/geometric-numerical-integration-illustrated-by-the-stormerverlet-method/E55395D5DD7A4E0526D10EA74DA6C46B}{Geometric numerical integration illustrated by the Störmer--Verlet method}，Acta Numerica原始研究入口；继续学习数值辛性、时间对称与后向误差分析。
\end{bookreading}

\chapter{狭义相对论}\label{ch:relativity}
\begin{guidebox}
设想两站各有一盏灯。要判断灯是不是同时闪，得先约定两站怎样对表。狭义相对论从这种可操作的测量出发，重新问时间、距离和速度怎样在不同惯性系之间换算。

做题先写出正在比较的两个事件，例如杆头和杆尾各经过探测器一次。再标出每个事件的位置和时刻由哪一套钟和标尺测量。掌握这一步后，再用间隔与四动量处理旅行、衰变和碰撞。
\end{guidebox}

\section{事件、同步与相对性原理}
想比较两地的闪光是否同时，先得让两地的钟有共同的对时办法。本节用光信号约定怎样给每一次闪光标上时间。

\subsection{先规定怎样测量}
一盏灯何时闪，与我们何时收到它的光，是两次不同的事件。要给第一次事件定时间，可以读灯旁边已经同步的钟。

一次闪光、两粒子相遇或一只钟的一次滴答，都发生在确定地点和时刻，称为\textbf{事件}（event）。惯性系 $S$ 用 $(t,x,y,z)$ 标记事件：空间坐标由静止标尺网络确定，$t$ 由事件发生地点的静止钟读取。这里的时间不是远处观察者看到闪光的时刻；后者还包含信号传播所需的时间。

怎样使空间分离的钟同步？设钟 $A,B$ 相对静止，$A$ 在 $t_1$ 发光，光在 $B$ 反射后于 $t_2$ 返回 $A$。采用\textbf{爱因斯坦同步约定}，把反射事件的 $B$ 钟读数设为 $(t_1+t_2)/2$，两钟距离为 $c(t_2-t_1)/2$。在惯性系中，这与各方向相同的真空光速相容。同步约定针对同一惯性系的静止钟网络；不能不加说明地沿用到转动参考系。

\begin{figure}[htbp]\centering
\begin{tikzpicture}[>=Stealth,x=1.2cm,y=0.9cm]
\draw[->] (-.5,0)--(4.7,0) node[right]{$x$};
\draw[->] (0,-.2)--(0,4.8) node[above]{$ct$};
\draw[thick] (3,0)--(3,4.5) node[above]{钟 $B$};
\node[below left] at (0,0){钟 $A$};
\draw[blue,very thick,->] (0,.6)--(3,2.1);
\draw[blue,very thick,->] (3,2.1)--(0,3.6);
\draw[dashed] (0,2.1)--(3,2.1);
\fill (0,.6) circle(2pt) node[left]{$ct_1$};
\fill (0,3.6) circle(2pt) node[left]{$ct_2$};
\fill (3,2.1) circle(2pt) node[right]{反射};
\node[left] at (0,2.1){$c(t_1+t_2)/2$};
\node[blue,above,sloped] at (1.6,1.4){发出};
\node[blue,above,sloped] at (1.6,2.9){返回};
\end{tikzpicture}
\caption{静止钟的雷达同步。此图为光路示意，横纵绘图比例不同；不以图上斜率读取光速。}\label{fig:sr-sync}
\end{figure}

\subsection{为什么要改变时间变换}
如果照常认为不同惯性系共用同一时间，追着光跑的人会算出较小的光速。要让每个惯性系都测得同一真空光速，就必须重新考虑时间和位置怎样换算。

第\ref{ch:particle}章的伽利略变换（式 \eqref{eq:galilean}）$x'=x-vt,t'=t$ 导出 $u'_x=u_x-v$。若 $u_x=c$，另一惯性系就得到 $c-v$，与真空光速不变冲突。需要改变的是时空变换，而不是在牛顿方程中临时更换一个质量函数。
\begin{keybox}
\textbf{相对性原理：}物理定律在所有惯性系中具有相同形式。\quad
\textbf{光速不变原理：}真空光速为 $c$，与光源运动状态无关，在所有惯性系中相同。

本章讨论平直时空。$c$ 是时空结构的不变速度，也是局域因果信号的极限速度；介质中的光速不用于定义洛伦兹变换。
\end{keybox}
狭义相对论可以描述加速粒子，只需在惯性坐标中写出它的世界线。排除引力时空曲率，不等于排除加速度；涉及引力钟差时则需要额外理论。

\section{洛伦兹变换：从原理到公式}
变换公式要告诉我们，同一次事件在两套钟和标尺网络中各得到什么读数。先约定两个原点何时重合，以及它们相对运动的方向。

\subsection{标准构型与逐步推导}
把两个惯性系的轴对齐，只让它们沿一条轴相对滑动。先用移动原点的轨迹确定位置变换的形式，再用光信号确定时间变换。

取坐标轴平行的惯性系 $S,S'$，$S'$ 沿 $S$ 的 $+x$ 方向以恒速 $v$ 运动；两原点在 $t=t'=0$ 重合。时空均匀性要求坐标差的映射不依赖基点，故变换为线性；原点世界线 $x=vt$ 又要求
\[
x'=A(v)(x-vt),\qquad t'=B(v)t+D(v)x.
\]
原点沿正、负 $x$ 发出的光分别满足 $x=\pm ct,x'=\pm ct'$，所以
\[
A(c-v)=cB+c^2D,\qquad A(c+v)=cB-c^2D.
\]
相加、相减得 $B=A,D=-Av/c^2$。反转 $x$ 轴使 $v\to-v$，空间各向同性给出 $A(-v)=A(v)$。沿 $S$ 的原点世界线 $x=0$，有 $x'/t'=-v$。所以由相对性原理，逆变换应使用同一形式而取速度 $-v$，即 $x=A(-v)(x'+vt')$。将 $x'$ 与 $t'$ 代入，并用 $A(-v)=A(v)$，得到
\[
x=A^2(1-v^2/c^2)x.
\]
选择从 $v=0$ 恒等变换连续延伸的正根，定义 $\beta=v/c$ 与 $\gamma=(1-\beta^2)^{-1/2}$，得到
\begin{equation}\label{eq:sr-lorentz}
\boxed{\begin{aligned}
x'&=\gamma(x-vt),& ct'&=\gamma(ct-\beta x),\\
y'&=y,&z'&=z.
\end{aligned}}
\end{equation}
其中 $|v|<c$。这种只改变惯性系相对速度的变换称为\textbf{洛伦兹推进}（boost）。横向变换由绕运动轴的对称性可写为 $y'=a(v)y,z'=a(v)z$；互易与各向同性给出 $a^2=1$，连续性选 $a=1$。这里采用相同的标尺单位，并排除了额外的空间转动。

逆变换只需交换撇号并令 $v\to-v$：
\[
x=\gamma(x'+vt'),\qquad t=\gamma(t'+vx'/c^2).
\]
不能只改空间公式的符号，而让时间保持不变。对事件坐标差也使用同样的线性变换。

\begin{example}[先变换，再反变换]
用 $\ell$ 为长度单位，设事件在 $S$ 中为 $(ct,x)=(5\ell,3\ell)$，取 $v=0.6c,\gamma=5/4$。则 $ct'=4\ell,x'=0$；这正是运动原点上的事件。逆变换给出 $ct=5\ell,x=3\ell$。两系均有 $c^2t^2-x^2=16\ell^2$，已显露出比坐标更稳定的几何量。

直接代入为 $ct'=(5/4)(5-0.6\times3)\ell$、$x'=(5/4)(3-0.6\times5)\ell$。在逆变换中 $x'=0$，于是 $ct=\gamma ct'$、$x=\gamma\beta ct'$。
\end{example}

\subsection{低速极限需要什么条件}
日常车速与光速相比很小，许多相对论修正可以忽略。但若比较相隔很远的两次事件，位置对时间变换的贡献也应单独估计。

用 $\gamma-1=\beta^2/2+O(\beta^4)$，写出相对于伽利略变换的差：
\[
x'-(x-vt)=(\gamma-1)(x-vt),\qquad
t'-t=-\frac{vx}{c^2}+(\gamma-1)(t-vx/c^2).
\]
第一项 $-vx/c^2$ 是同时性修正。后一项来自 $\gamma$，若精度达到 $\beta^2$ 量级，也须保留。若研究过程的典型长度、时间为 $L,T$，除了 $|v|/c\ll1$，还须在要求的精度内有 $|v|L/c^2\ll T$，才能把时间也近似看作绝对时间。对长距离精密授时，即使速度不大，同时性修正仍可能重要。

\section{同时性、钟慢和尺缩：每次都选择事件}
钟慢、尺缩和同时性用的是不同的事件对。先写清两次事件发生在哪里、由哪套钟读数，再决定用哪条关系。

\subsection{同时不等于同时看见}
近处和远处的灯即使同时发光，光也可能先后到达同一观察者。比较发光时刻时，应查看各盏灯旁已经同步的钟。

对两个事件作差：
\[
\Delta t'=\gamma(\Delta t-v\Delta x/c^2).
\]
$S$ 中同时而异地的事件满足 $\Delta t=0,\Delta x\ne0$，通常有 $\Delta t'\ne0$。这是同步钟读数的差异，已经与光到达眼睛的先后区分开。
\begin{example}[两灯闪光]
实验室两灯位于 $x=0,L$，在 $t=0$ 闪光。对向右运动的 $S'$，右灯事件减左灯事件的时间差为 $-\gamma vL/c^2$，右灯先闪。若改问运动观察者何时收到两束光，则还须求其世界线与光线的交点，不能直接用这个时间差回答。
\end{example}

\subsection{同一只钟：时间膨胀}
让同一只钟在两次滴答时留下记录。钟自己的参考系里两次滴答发生在同一位置，这就是推导时间膨胀需要的事件条件。

静止于 $S'$ 原点的一只钟，两次滴答的事件差满足 $\Delta x'=0$。这只钟自身记录的时间称为\textbf{固有时}（proper time）$\Delta\tau=\Delta t'$。由逆变换
\begin{equation}\label{eq:sr-time}
\boxed{\Delta t=\gamma\Delta\tau.}
\end{equation}
实验室的 $\Delta t$ 由经过的不同位置处的同步钟比较，$\Delta\tau$ 则由同一只移动钟记录。两种测量操作并不对称。

光钟给出直观推导。两镜在钟的静止系中沿横向相距 $h$，一次往返周期为 $\Delta\tau=2h/c$。实验室看到钟横向运动，半周期光路为斜线，故
\[
(c\Delta t/2)^2=h^2+(v\Delta t/2)^2
\quad\Longrightarrow\quad \Delta t=\gamma\Delta\tau.
\]
这里比较下镜发射与下一次返回下镜的两个事件。
\begin{figure}[htbp]\centering
\begin{tikzpicture}[>=Stealth,scale=.95]
\begin{scope}
\draw[thick] (-.4,0)--(.4,0);\draw[thick] (-.4,2.4)--(.4,2.4);
\draw[blue,->,very thick] (-.1,.1)--(-.1,2.3);
\draw[blue,->,very thick] (.1,2.3)--(.1,.1);
\draw[<->] (-.8,0)--(-.8,2.4) node[midway,left]{$h$};
\node at (0,-.6){钟静止系};
\end{scope}
\begin{scope}[xshift=4cm]
\draw[dashed] (0,0)--(0,2.4);\draw[dashed] (1.8,0)--(1.8,2.4);
\draw[thick] (-.3,0)--(.3,0);\draw[thick] (1.5,2.4)--(2.1,2.4);
\draw[thick] (3.3,0)--(3.9,0);
\draw[blue,->,very thick] (0,0)--(1.8,2.4);
\draw[blue,->,very thick] (1.8,2.4)--(3.6,0);
\draw[<->] (0,-.35)--(1.8,-.35) node[midway,below]{$v\Delta t/2$};
\node[above left,blue] at (.9,1.3){$c\Delta t/2$};
\node at (1.8,-1.05){实验室中的空间光路};
\end{scope}
\end{tikzpicture}
\caption{横向光钟。右图按 $v=0.6c$ 绘制，半周期光路、横移、镜距构成 $5:3:4$ 直角三角形；它是空间路径图，不是时空图。}\label{fig:sr-lightclock}
\end{figure}

\begin{example}[寿命与飞行距离]
一种模型粒子的固有平均寿命为 $\tau_0$，实验室速率为 $0.8c$，忽略能损。实验室平均寿命为 $5\tau_0/3$，平均飞行距离为 $4c\tau_0/3$。若固有时衰变律为 $N=N_0e^{-\tau/\tau_0}$，实验室飞行距离 $d$ 后的存活比例为 $e^{-d/(\gamma u\tau_0)}$。平均距离不是所有粒子的共同飞行终点；此例未采用实际实验数据。
\end{example}

\subsection{同一时刻的两个端点：长度收缩}
测运动的杆有多长，应在测量者的同一时刻读出两端位置。换到杆的静止系，这两次读数对应的事件一般不再同时。

杆静止于 $S'$，端点为 $x'_L=0,x'_R=L_0$，$L_0$ 为\textbf{固有长度}。实验室必须在同一 $t$ 记录两端位置。令两测量事件差 $\Delta t=0$，便有 $L_0=\gamma L$，所以
\[
\boxed{L=L_0/\gamma.}
\]
表中选 $t=0$ 的测量事件，清楚显示它们在杆系并不同时：
\begin{center}\begin{tabular}{c c c c c}\toprule
事件&$t$&$x$&$t'$&$x'$\\\midrule
$A$：记录左端&$0$&$0$&$0$&$0$\\
$B$：记录右端&$0$&$L_0/\gamma$&$-vL_0/c^2$&$L_0$\\\bottomrule
\end{tabular}\end{center}
\begin{figure}[htbp]\centering
\begin{tikzpicture}[>=Stealth,scale=1.05]
\draw[->] (-.5,0)--(5.1,0) node[right]{$x$};
\draw[->] (0,-.4)--(0,4.1) node[above]{$ct$};
\draw[thick,blue] (0,0)--(2,4) node[above]{$x'_L=0$};
\draw[thick,blue] (2,0)--(4,4) node[above]{$x'_R=L_0$};
\draw[red,thick] (0,0)--(2.667,1.333) node[right]{$t'=0$};
\fill (0,0) circle(2pt) node[below left]{$A$};
\fill (2,0) circle(2pt) node[below]{$B$};
\fill (2.667,1.333) circle(2pt) node[above left]{$C$};
\draw[<->] (0,-.25)--(2,-.25) node[midway,below]{$L=L_0/\gamma$};
\node[right] at (3.4,2.2){杆端世界线};
\end{tikzpicture}
\caption{运动杆的两种事件选择。$AB$ 是实验室同时测量，$AC$ 是杆系同时测量；不能把 $AC$ 的水平投影当作实验室杆长。取 $v=0.5c$，横纵轴等尺度。}\label{fig:sr-rod}
\end{figure}

\begin{example}[杆通过一个探测器]
固有长 $L_0$ 的杆以 $0.6c$ 通过实验室固定探测器。实验室杆长为 $0.8L_0$，两端经过探测器的时间差为 $\Delta t=L/u=4L_0/(3c)$。这两个事件在实验室同地，故探测器记录的固有时就是 $\Delta t$。在杆系它们的时间差为 $\gamma\Delta t=5L_0/(3c)=L_0/u$。辨认同地条件后，钟慢关系没有矛盾。
\end{example}
\begin{physicsbox}
“各看对方钟慢”比较的是不同同时切片选出的事件对。若两只钟最后真正重逢，则比较沿两条完整世界线累积的固有时。长度收缩也由同时切片定义；一张照片所记录的是到达相机同一时刻的光，仍需处理光行时，不能直接读作收缩长度。
\end{physicsbox}

\section{闵可夫斯基时空与固有时}
同一次旅行的时间差与空间位移会随参考系改变，但可以组合成大家都算出同一值的间隔。它帮助判断两事件能否用光或物体的运动联系起来。

\subsection{不变间隔与因果分类}
先计算由时间差和空间位移组成的量，再根据它的正、零、负分类。类时、类光和类空分别对应不同的因果联系可能性。

代入洛伦兹变换并展开，有
\[
\begin{aligned}
c^2\Delta t'^2-\Delta x'^2
&=\gamma^2\big[(c\Delta t-\beta\Delta x)^2-(\Delta x-\beta c\Delta t)^2\big]\\
&=c^2\Delta t^2-\Delta x^2.
\end{aligned}
\]
因此四维间隔
\begin{equation}\label{eq:sr-interval}
\boxed{\Delta s^2=c^2\Delta t^2-|\Delta\br|^2}
\end{equation}
不依赖惯性系。本章固定采用 $(+,-,-,-)$ 号差。粒子的事件轨迹称\textbf{世界线}（worldline）。沿 $x$ 的时空图取 $x,ct$ 相同尺度时，真空光线为 $45^\circ$，有质量粒子的未来世界线位于未来光锥内。

\begin{figure}[htbp]\centering
\begin{tikzpicture}[>=Stealth,scale=.85]
\fill[blue!8] (0,0)--(-3,3)--(3,3)--cycle;
\fill[blue!8] (0,0)--(-3,-3)--(3,-3)--cycle;
\draw[->] (-3.7,0)--(3.8,0) node[right]{$x$};
\draw[->] (0,-3.4)--(0,3.7) node[above]{$ct$};
\draw[orange,thick] (-3.2,-3.2)--(3.2,3.2);
\draw[orange,thick] (-3.2,3.2)--(3.2,-3.2);
\draw[blue,very thick,->] (0,0)--(1.1,2.8);
\fill (1.1,2.8) circle(2pt) node[above]{$T$：类时};
\fill (2.8,.9) circle(2pt) node[right]{$Q$：类空};
\node at (-.5,2.4){未来};\node at (.6,-2.5){过去};
\node[below left] at (0,0){$O$};
\node[orange,rotate=-45] at (-2.5,2.8){类光};
\end{tikzpicture}
\caption{光锥按相同横纵尺度绘制。$O$ 与 $T$ 可由亚光速信号连接；$O$ 与 $Q$ 不能建立这样的因果联系。图仅显示一个空间方向。}\label{fig:sr-cone}
\end{figure}

沿分离方向取 $\Delta x>0$，三种情形可用变换直接判断：
\begin{center}\begin{tabular}{p{3.1cm}p{4.1cm}p{6cm}}\toprule
间隔&可选择的参考系&物理含义\\\midrule
类时：$\Delta s^2>0$&$v=\Delta x/\Delta t$ 使 $\Delta x'=0$&存在两事件同地的惯性系；未来先后对所有惯性系相同。\\
类空：$\Delta s^2<0$&$v=c^2\Delta t/\Delta x$ 使 $\Delta t'=0$&存在同时系；先后可反转，但不存在亚光速因果连接。\\
类光：$\Delta s^2=0$&使其同地需 $|v|=c$，不是惯性系&没有光子的静止系；沿光线固有时间隔为零。\\\bottomrule
\end{tabular}\end{center}
因果次序不变还需限制于保持时间方向的洛伦兹变换，本章使用的连续标准推进满足此条件。

\begin{example}[事件先后何时能反转？]
设 $\Delta x=4\ell,c\Delta t=2\ell$，则 $\Delta s^2=-12\ell^2$，在 $v=0.5c$ 的系中同时，$v>0.5c$ 时 $\Delta t'<0$。若改为 $\Delta x=2\ell,c\Delta t=4\ell$，则类时；$\Delta t'=\gamma(4-2\beta)\ell/c>0$ 对所有 $|\beta|<1$ 成立。坐标先后能变的恰是没有亚光速因果关联的事件对。
\end{example}

\subsection{倾斜坐标轴也需要校准}
时空图上的斜线可以画出运动观察者的坐标轴。为了用它量长度或时间，还要确定斜轴上一格对应多少单位。

在 $S$ 的时空图上，$ct'$ 轴满足 $x'=0$，即 $x=\beta ct$；$x'$ 轴满足 $t'=0$，即 $ct=\beta x$。二者向光线方向倾斜。图上并非欧氏直角坐标，不能作垂足读出撇号坐标。

若原系长度单位为 $\ell$，则 $ct'=\ell,x'=0$ 的单位点为 $(ct,x)=(\gamma\ell,\gamma\beta\ell)$；$ct'=0,x'=\ell$ 的单位点为 $(\gamma\beta\ell,\gamma\ell)$。前者位于 $c^2t^2-x^2=\ell^2$，后者位于 $x^2-c^2t^2=\ell^2$。不变双曲线校准了新坐标轴。
\begin{figure}[htbp]\centering
\begin{tikzpicture}[>=Stealth,scale=2.25]
\draw[->] (-.15,0)--(2.2,0) node[right]{$x/\ell$};
\draw[->] (0,-.15)--(0,2.2) node[above]{$ct/\ell$};
\draw[dashed,gray] (0,0)--(2,2);
\draw[blue,thick,->] (0,0)--(1,2) node[above]{$ct'$};
\draw[red,thick,->] (0,0)--(2,1) node[right]{$x'$};
\draw[blue!60,domain=0:1.5,samples=60] plot (\x,{sqrt(1+\x*\x)});
\draw[red!60,domain=0:1.5,samples=60] plot ({sqrt(1+\x*\x)},\x);
\fill[blue] (.57735,1.1547) circle(.025) node[left]{$ct'=\ell$};
\fill[red] (1.1547,.57735) circle(.025) node[below right]{$x'=\ell$};
\fill (0,1) circle(.018) node[left]{$1$};
\fill (1,0) circle(.018) node[below]{$1$};
\end{tikzpicture}
\caption{取 $\beta=1/2$ 的校准坐标图。新轴单位点的欧氏距离并非原轴单位长度；虚线为光线，双曲线分别校准时间与空间单位。}\label{fig:sr-axes}
\end{figure}

\subsection{变速世界线与双生子}
变速旅行可分成许多小段，每一段用当时的速率计算旅行者钟走了多少。把这些时间加起来，就是出发到重逢的固有时间。

对有质量粒子，把相邻事件的间隔除以 $c^2$，定义
\begin{equation}\label{eq:sr-proper}
\dd\tau=\dd t\sqrt{1-u^2(t)/c^2},\qquad
\boxed{\tau=\int_{t_A}^{t_B}\sqrt{1-u^2(t)/c^2}\,\dd t.}
\end{equation}
这使用粒子相对所选惯性系的瞬时速率 $u(t)$。理想钟测量此固有时；积分对任意亚光速光滑世界线成立，分段光滑运动可分段积分。

\begin{example}[双生子的年龄和同时性切换]
地球钟留在 $x=0$，旅行者以 $u=\beta c$ 出发，在地球时刻 $T/2$、位置 $uT/2$ 瞬时掉头，再于 $T$ 回到原点。忽略掉头所占时间是极限模型。地球累积 $T$，旅行者累积 $T/\gamma$；取 $T=20$ 年、$\beta=0.8$，分别为 $20$ 年与 $12$ 年。

为何旅行者不能在整个旅程中使用一个惯性系，反说地球年轻？掉头事件处，出发段共动系的同时线给地球时刻
\[
t_-=\frac{T}{2}-\frac{u(uT/2)}{c^2}=\frac{T}{2}(1-\beta^2),
\]
返回段共动系的同时线则给
\[
t_+=\frac{T}{2}+\frac{u(uT/2)}{c^2}=\frac{T}{2}(1+\beta^2).
\]
二者相差 $\beta^2T$，是更换惯性系时对远处同时事件的重新归属，不是地球钟突然跳动。$\beta=0.8$ 时为 $3.6$ 年与 $16.4$ 年。实际看到的光信号连续传播，也不会出现这种读数跳跃。
\end{example}

\begin{figure}[htbp]\centering
\begin{tikzpicture}[>=Stealth,scale=1.2]
\draw[->] (-.4,0)--(2.8,0) node[right]{$x/(cT)$};
\draw[->] (0,-.2)--(0,4.5) node[above]{$ct/(cT)$};
\draw[blue,very thick] (0,0)--(0,4);
\draw[red,very thick,->] (0,0)--(1.6,2);
\draw[red,very thick,->] (1.6,2)--(0,4);
\draw[dashed] (0,.72)--(2.3,2.56);
\draw[dashed] (0,3.28)--(2.3,1.44);
\fill (1.6,2) circle(2pt) node[right]{掉头 $(0.4,0.5)$};
\fill (0,.72) circle(2pt) node[left]{$0.18$};
\fill (0,3.28) circle(2pt) node[left]{$0.82$};
\node[left] at (0,4){$1$};\node[below left] at (0,0){$0$};
\node[red,right] at (1,3.3){返回段};
\end{tikzpicture}
\caption{双生子世界线与掉头处两套同时线，$\beta=0.8$。坐标以 $cT$ 归一化，图上每单位对应四个绘图单位；两轴等尺度。地球钟两交点为 $0.18T,0.82T$。}\label{fig:sr-twins}
\end{figure}

\begin{mathbox}[title={数学注记：直线世界线使固有时最大}]
在同一惯性系中，固定两端事件、总坐标时间 $T$ 和位移 $\Delta\br$。函数 $f(\bu)=\sqrt{1-|\bu|^2/c^2}$ 在 $|\bu|<c$ 内严格凹。由凹函数的积分不等式，
\[
\tau=\int_0^T f(\bu)\,\dd t
\le T f\!\left(\frac1T\int_0^T\bu\,\dd t\right)
=T\sqrt{1-\frac{|\Delta\br|^2}{c^2T^2}}.
\]
等号仅在速度恒定时成立。因此平直时空中固定类时端点之间，惯性直线具有最大固有时。双生子年龄差取决于路径；加速度使路径改变，却不是在钟旁额外扣除时间的独立项。
\end{mathbox}

\begin{modernbox}[title={进一步探索：快度把同向速度合成变成加法}]
令\textbf{快度} $\chi=\operatorname{artanh}\beta$，则 $\gamma=\cosh\chi,\gamma\beta=\sinh\chi$，推进写为
\[
\begin{pmatrix}ct'\\x'\end{pmatrix}
=\begin{pmatrix}\cosh\chi&-\sinh\chi\\-\sinh\chi&\cosh\chi\end{pmatrix}
\begin{pmatrix}ct\\x\end{pmatrix}.
\]
双曲函数加法公式表明，同一直线上的连续推进快度相加，速度则为 $c\tanh(\chi_1+\chi_2)$。它揭示洛伦兹推进是保持双曲型间隔的变换。只需矩阵乘法和双曲函数即可验证；非共线推进还会涉及空间转动，不能直接推广成三维速度的逐分量相加。
\end{modernbox}

\section{速度合成、光行差与多普勒效应}
改变参考系后，速度、光的传播方向和频率都可能变。下面分别从事件坐标的变化和光波相位的变化求出它们。

\subsection{对坐标变换取微分}
速度是位移除以时间。洛伦兹变换会同时改变这个分式的分子和分母，因此连横向速度分量也需要重新计算。

设粒子在 $S$ 中速度为 $\bu$，用 $\dd t'=\gamma(1-vu_x/c^2)\dd t$ 作分母，得到
\begin{equation}\label{eq:sr-velocity}
\boxed{u'_x=\frac{u_x-v}{1-vu_x/c^2},\quad
u'_y=\frac{u_y}{\gamma(1-vu_x/c^2)},\quad
u'_z=\frac{u_z}{\gamma(1-vu_x/c^2)}.}
\end{equation}
横向速度也会改变，因为时间间隔改变。反解时 $v\to-v$。沿粒子世界线，间隔可写为 $\dd s^2=c^2\dd t^2(1-u^2/c^2)$，在 $S'$ 中则为 $c^2\dd t'^2(1-u'^2/c^2)$。两式相等，再代入 $\dd t'/\dd t$，得到
\[
1-\frac{u'^2}{c^2}=\frac{1-u^2/c^2}{\gamma^2(1-vu_x/c^2)^2}.
\]
因此亚光速仍为亚光速，光速仍为光速。若有人使用 $u_x+v$ 作为高速合成，便破坏这个性质。
\begin{example}[飞船发射探测器]
飞船相对实验室以 $0.6c$ 向右运动，探测器在飞船系中也以 $0.6c$ 向右运动。使用逆速度变换，实验室速度为
\[
u_x=\frac{0.6c+0.6c}{1+0.36}=\frac{15}{17}c,
\]
小于 $c$。对应快度为两倍 $\operatorname{artanh}(0.6)$，结果相同。
\end{example}

\subsection{光的方向怎样改变}
一束光在某个惯性系中垂直于运动方向，在另一个系中可能变斜。把它的速度分量重新算一遍，合起来的大小仍为光速。

光在 $S$ 中沿与 $+x$ 夹角 $\theta$ 的方向传播，取 $u_x=c\cos\theta,u_y=c\sin\theta$，则
\[
\cos\theta'=\frac{\cos\theta-\beta}{1-\beta\cos\theta},\qquad
\sin\theta'=\frac{\sin\theta}{\gamma(1-\beta\cos\theta)}.
\]
这是\textbf{光行差}（aberration），两个角各自属于不同参考系。对一般象限应同时使用正弦、余弦判别方向。
\begin{example}[横向光并非在另一系仍横向]
在 $S$ 中沿 $+y$ 传播的光有 $\theta=\pi/2$。对 $v=0.6c$ 的 $S'$，$\cos\theta'=-0.6,\sin\theta'=0.8$。合成后仍有 $u'^2=c^2$，但光线向 $-x'$ 偏斜；只变纵向速度而保留横向速度会给出错误光速。
\end{example}

\subsection{由波相位推导频率变换}
同一波峰到达某个事件时，两系应认出同一个波峰。利用相位一致，就能把频率和传播方向放在同一套变换中求解。

平面光波相位为 $\Phi=\omega t-\boldsymbol{k}\cdot\br$。同一事件的相位不依赖坐标标记，把逆洛伦兹变换代入得
\[
\omega'=\gamma(\omega-vk_x),\quad
k'_x=\gamma(k_x-v\omega/c^2),\quad k'_y=k_y,\quad k'_z=k_z.
\]
真空光波 $|\boldsymbol{k}|=\omega/c$，若 $\theta$ 是\emph{光传播方向在 $S$ 中}与 $+x$ 的夹角，便有
\begin{equation}\label{eq:sr-doppler}
\boxed{\omega'=\gamma\omega(1-\beta\cos\theta).}
\end{equation}
这里 $S'$ 中静止的观察者测得 $\omega'$。若沿 $+x$ 追赶光，即 $\theta=0,v>0$，则
\[
\omega'=\omega\sqrt{\frac{1-\beta}{1+\beta}}.
\]
迎光运动则频率升高。波矢与角度也要一起变换。
\begin{example}[两种“横向”条件]
若 $\theta=\pi/2$，则 $\omega'=\gamma\omega$。若要求接收系中光线横向，即 $\theta'=\pi/2$，光行差给出 $\cos\theta=\beta$，从而 $\omega'=\omega/\gamma$。两结果并不冲突：横向条件属于不同参考系，所选光线不同。使用“横向多普勒”一词时必须注明这个条件。
\end{example}

\section{四矢量、四动量与不变质量}
能量和动量在换参考系时会相互混合。把它们合起来处理，能找出所有观察者都同意的质量。

\subsection{不定内积与四速度}
在普通空间里，向量的各分量平方相加。时空中时间和空间项的符号不同，这使光的方向能够有零间隔，也使四速度有固定的长度平方。

定义位置四矢量 $x^\mu=(ct,x,y,z)$，指标 $\mu=0,1,2,3$。对任意四矢量 $A,B$，约定
\[
A\cdot B=\eta_{\mu\nu}A^\mu B^\nu=A^0B^0-\boldsymbol{A}\cdot\boldsymbol{B},\qquad
\eta=\operatorname{diag}(1,-1,-1,-1).
\]
重复指标求和，空间三维点积仍为欧氏点积。矩阵条件 $\Lambda^{\mathsf T}\eta\Lambda=\eta$ 保证四维内积不变。降指标 $A_\mu=\eta_{\mu\nu}A^\nu$，所以 $A_\mu=(A^0,-\boldsymbol{A})$；指标位置在此有实际符号意义。

有质量粒子的四速度为
\[
U^\mu=\frac{\dd x^\mu}{\dd\tau}=\gamma_u(c,\bu),\qquad U\cdot U=c^2,
\quad \gamma_u=(1-u^2/c^2)^{-1/2}.
\]
不变量 $\tau$ 作分母保证 $U$ 按四矢量变换。光线没有以自身固有时定义的四速度，不能代入 $u=c$ 强行计算。

\subsection{能量与动量属于同一个四矢量}
把四速度乘以粒子的质量，时间分量与能量有关，空间分量就是相对论动量。把它们组成不变间隔式，就得到能量和动量的关系。

对自由有质量粒子定义不变质量 $m$ 和四动量
\begin{equation}\label{eq:sr-p4}
P^\mu=mU^\mu=(E/c,\bp),\quad
\boxed{E=\gamma_umc^2,\quad \bp=\gamma_um\bu.}
\end{equation}
$m$ 与惯性系无关。本书不把 $\gamma_um$ 称作随速度变化的质量。四动量的时间分量为 $P^0=E/c$，其物理归一化由静止能 $mc^2$ 确定；后面将从作用量重新得到这些量。
\begin{equation}\label{eq:sr-mass-shell}
P\cdot P=m^2c^2\quad\Longrightarrow\quad
\boxed{E^2=p^2c^2+m^2c^4,\quad\bu=\frac{c^2\bp}{E}.}
\end{equation}
固定 $m$ 时，式中的能量--动量约束曲面称为\textbf{质量壳}。对正能量自由粒子取正根。动能 $K=E-mc^2$，低速展开为
\[
K=\tfrac12mu^2+\tfrac38mu^4/c^2+O(mu^6/c^4),\qquad
\bp=m\bu\,[1+u^2/(2c^2)+O(u^4/c^4)].
\]
动量小参数形式则为 $E=mc^2+p^2/(2m)-p^4/(8m^3c^2)+\cdots$，不要把两个不同自变量的展开混用。

\begin{figure}[htbp]\centering
\begin{tikzpicture}
\begin{axis}[width=.81\textwidth,height=5.3cm,xlabel={$p/(mc)$},ylabel={$E/(mc^2)$},xmin=0,xmax=3,xtick={0,0.5,1,1.5,2,2.5,3},ymin=0,ymax=4,legend pos=north west,grid=major]
\addplot[blue,thick,domain=0:3,samples=100]{sqrt(1+x*x)};\addlegendentry{有质量粒子}
\addplot[orange,thick,dashed,domain=0:3]{x};\addlegendentry{光子 $E=pc$}
\addplot[red,dotted,thick,domain=0:2.3]{1+x*x/2};\addlegendentry{低动量近似}
\end{axis}\end{tikzpicture}
\caption{正能量质量壳。图用一个固定的非零参考质量 $m$ 归一化，光子曲线也使用该参考单位；光子自身的质量为零。低动量近似只在 $p\ll mc$ 时可靠。}\label{fig:sr-mass-shell}
\end{figure}

\begin{example}[动能和动量不能使用同一个修正因子]
粒子速率为 $0.8c$，则 $E=5mc^2/3,K=2mc^2/3,p=4mc/3$。若把牛顿动能 $mu^2/2$ 简单乘 $\gamma$，得到 $8mc^2/15$，并非正确动能。正确方式是先用总能量，再减静止能。

由 $\gamma_u=(1-0.8^2)^{-1/2}=5/3$，分别代入 $E=\gamma_umc^2$、$K=(\gamma_u-1)mc^2$、$p=\gamma_umu$。将 $E$ 和 $p$ 放回 $E^2-p^2c^2$，得到 $(25-16)m^2c^4/9=m^2c^4$。
\end{example}

\subsection{能动量变换与观察者的能量}
迎着粒子运动的观察者与追着粒子运动的观察者，测得的能量可能不同。选定观察者的速度后，就能从粒子的四动量读出这个能量。

四动量变换与位置相同：
\[
E'=\gamma(E-vp_x),\quad p'_x=\gamma(p_x-vE/c^2),\quad p'_y=p_y,\quad p'_z=p_z.
\]
对指定观察者，设其四速度为 $W=\gamma_w(c,\boldsymbol{w})$，它测得的粒子能量可写为
\[
E_{\rm obs}=P\cdot W=\gamma_w(E-\bp\cdot\boldsymbol{w}).
\]
在观察者静止系中 $W=(c,0,0,0)$，这个式子就是该系能量，所以结论成立。$P\cdot W$ 对同时变换两个四矢量不变；更换观察者则改变 $W$，能量仍会改变。

\begin{example}[追赶一个光子]
自由光子取 $m=0$，故 $E=pc$，四动量为 $(E/c,(E/c)\boldsymbol{n})$。沿光方向以 $0.6c$ 运动的观察者测得 $E'=\gamma E(1-0.6)=E/2$。加上量子关系 $E=\hbar\omega$，即得到同方向多普勒结果。这个量子关系是额外物理输入，并非单由洛伦兹变换推出。
\end{example}

\subsection{系统的不变质量与质心系}
研究一群粒子时，先把它们的能量和动量分别相加，再由总量求不变质量。若存在总动量为零的参考系，系统的不变质量就是该系总能量除以光速平方。

孤立系统的总四动量守恒。对可视为互不作用的碰撞入射态或出射态，将全部粒子与辐射的四动量相加：$P_{\rm tot}=\sum_aP_a$。相互作用期间须按完整模型计入场或相互作用的能动量，不能只加粒子的机械四动量。与第\ref{ch:particle}章的动量收支相同，只取外场中的测试粒子通常得不到孤立系统。定义系统不变质量 $M$：
\[
\boxed{M^2c^4=E_{\rm tot}^2-c^2|\bP_{\rm tot}|^2.}
\]
若 $M>0$，取 $\bV=c^2\bP_{\rm tot}/E_{\rm tot}$ 可到总三动量为零的\textbf{质心系}（也称动量中心系）；此系总能量为 $Mc^2$。$M$ 包含内部相对运动能、辐射能和相互作用能，通常不等于各静质量的和。相互作用体系的总能动量还必须计入场。

\begin{example}[两个无质量粒子组成有质量系统]
两光子能量为 $E_1,E_2$，传播夹角为 $\theta$。将总能动量平方，得到
\[
M^2c^4=(E_1+E_2)^2-|E_1\boldsymbol{n}_1+E_2\boldsymbol{n}_2|^2
=2E_1E_2(1-\cos\theta).
\]
等能反向光子有 $M=2E_1/c^2$，尽管每个光子质量都为零。同向光子系统 $M=0$，不存在总动量为零的惯性系；不能对所有粒子系统无条件指定质心静止系。
\end{example}

\begin{figure}[htbp]\centering
\begin{tikzpicture}[>=Stealth,scale=1]
\fill (0,0) circle(2pt);
\draw[blue,very thick,->] (0,0)--(2.8,0) node[right]{$E_1\boldsymbol{n}_1/c$};
\draw[red,very thick,->] (0,0)--(-1.4,2.425) node[above]{$E_2\boldsymbol{n}_2/c$};
\draw[->] (.65,0) arc[start angle=0,end angle=120,radius=.65];
\node at (.5,.7){$\theta$};
\draw[dashed,gray] (2.8,0)--(1.4,2.425)--(-1.4,2.425);
\draw[purple,thick,->] (0,0)--(1.4,2.425) node[right]{$\bP_{\rm tot}$};
\end{tikzpicture}
\caption{双光子三动量的矢量合成，图取等能光子与 $\theta=120^\circ$。系统不变质量由能量和动量共同决定，不能只将两个光子的静质量相加。}\label{fig:sr-photons}
\end{figure}

\section{相对论动力学与分析力学的衔接}
在高速情况下，仍能用力求动量的变化，也能用功率求能量的变化。需要换掉的是低速下速度与动量之间的简单比例。

\subsection{三力、功率和四力}
先用实验室的时间定义三力和功率，再用粒子自己的钟定义四力。写方程时要分清导数对哪一种时间求。

在所选惯性系中定义三力 $\bF=\dd\bp/\dd t$。对不变质量恒定的粒子，$\dd\gamma_u/\dd t=\gamma_u^3\bu\cdot\ba/c^2$，因而
\begin{equation}\label{eq:sr-force}
\bF=m\gamma_u\ba+\frac{m\gamma_u^3}{c^2}(\bu\cdot\ba)\bu.
\end{equation}
分解加速度为沿速度和垂直速度的分量，有
\[
F_\parallel=m\gamma_u^3a_\parallel,\qquad
\bF_\perp=m\gamma_u\ba_\perp.
\]
同样大小的力在不同方向产生不同的加速度，所以不能保留 $\bF=m_{\rm eff}\ba$ 并用一个标量替代所有动力学。由能量动量关系微分，
\[
2E\dot E=2c^2\bp\cdot\dot{\bp}
\quad\Longrightarrow\quad\boxed{\dot E=\bF\cdot\bu.}
\]
在 $m$ 不变时总能量和动能导数相同，功率关系保持熟悉的形式。

四力定义为
\[
f^\mu=\frac{\dd P^\mu}{\dd\tau}
=\gamma_u\left(\frac{\bF\cdot\bu}{c},\bF\right).
\]
由于 $P=mU$ 与 $U\cdot U=c^2$，有 $f\cdot U=0$。若粒子因吸热、辐射等改变自身不变质量，$\dd P/\dd\tau=m\dd U/\dd\tau+\frac{\dd m}{\dd\tau}U$，正交结论需要相应修正；此处动力学公式均以恒定 $m$ 为条件。

\begin{example}[为什么磁场能弯曲轨迹而不做功]
带电粒子处于均匀静磁场，忽略辐射与其他作用。洛伦兹力 $e\bu\times\bB$ 与速度垂直，故 $\dot E=0$，$u,\gamma_u$ 均不变。在垂直磁场的平面内运动时，$p=|e|B\rho$，角频率为 $|e|B/(\gamma_um)$。磁场改变动量方向，不能给粒子补充动能；加速器增加能量还需要电场。
\end{example}

\subsection{恒力运动与固有加速度}
恒力能让动量随实验室时间匀速增加，但速度越接近光速，继续增加就越慢。先求动量，再换回速度和位置。

粒子从实验室原点静止出发，受沿 $+x$ 的恒定三力 $F>0$。先积分动量方程而非猜测加速度：$p=Ft$。由质量壳求速度并积分：
\begin{equation}\label{eq:sr-constant-force}
\begin{aligned}
u(t)&=\frac{Ft/m}{\sqrt{1+(Ft/mc)^2}},\\
x(t)&=\frac{mc^2}{F}\left[\sqrt{1+(Ft/mc)^2}-1\right],\\
E(t)&=mc^2\sqrt{1+(Ft/mc)^2}.
\end{aligned}
\end{equation}
短时间 $Ft\ll mc$ 时恢复 $u=Ft/m,x=Ft^2/(2m)$；任意有限时间内 $u<c$。坐标加速度为 $\dd u/\dd t=(F/m)[1+(Ft/mc)^2]^{-3/2}$，逐渐减小。

定义\textbf{固有加速度} $\alpha=\sqrt{-A\cdot A}$，其中 $A=\dd U/\dd\tau$。一维运动令 $u=c\tanh\chi$，则 $U=c(\cosh\chi,\sinh\chi)$，故 $\alpha=c|\dd\chi/\dd\tau|$。从静止开始、沿正向且恒定 $\alpha$ 的运动满足
\[
\chi=\alpha\tau/c,\quad
t=\frac{c}{\alpha}\sinh\frac{\alpha\tau}{c},\quad
x=\frac{c^2}{\alpha}\left(\cosh\frac{\alpha\tau}{c}-1\right).
\]
消去 $\tau$ 得 $(x+c^2/\alpha)^2-c^2t^2=(c^2/\alpha)^2$。与式 \eqref{eq:sr-constant-force} 比较，\emph{本例一维、恒质量情形}有 $\alpha=F/m$；恒固有加速度并不意味着恒坐标加速度。

\begin{figure}[htbp]\centering
\begin{tikzpicture}
\begin{axis}[width=.79\textwidth,height=5.2cm,xlabel={$z=Ft/(mc)$},ylabel={归一化物理量},xmin=0,xmax=4,ymin=0,ymax=4,legend pos=north west,grid=major]
\addplot[blue,thick,domain=0:4,samples=100]{x/sqrt(1+x*x)};\addlegendentry{$u/c$}
\addplot[red,thick,domain=0:4,samples=100]{sqrt(1+x*x)-1};\addlegendentry{$Fx/(mc^2)$}
\addplot[gray,dashed,domain=0:4]{1};\addlegendentry{光速上界}
\end{axis}\end{tikzpicture}
\caption{恒力使动量持续线性增加，但速度趋近 $c$。横轴 $z$ 无量纲，纵轴分别表示速率和位移的无量纲值，不表示两者具有相同量纲。}\label{fig:sr-constant-force}
\end{figure}

\begin{example}[恒力做功的闭合核查]
令 $Ft/(mc)=3/4$，则 $\gamma_u=5/4,u=3c/5,x=mc^2/(4F)$，于是 $Fx=mc^2/4=E-mc^2$。同一时刻固有时满足 $F\tau/(mc)=\operatorname{arsinh}(3/4)=\ln2$。由双曲形式计算 $t,x$ 又得到同一结果，验证了坐标时间与固有时两种描述。
\end{example}

\subsection{自由粒子的作用量与哈密顿量}
前几章从作用量求运动方程。在相对论里，先用各惯性系都同意的固有时间写自由粒子的作用量，再用同一套变分和变量变换求动力学。

前面使用了能动量的四矢量结构；现在回到第2、3、6章的方法。自由粒子作用量应是洛伦兹标量，取世界线固有长度的一次积分：
\[
S=-mc^2\int\dd\tau=\int L\,\dd t,\qquad
L=-mc^2\sqrt{1-u^2/c^2}.
\]
系数和符号由低速展开 $L=-mc^2+mu^2/2+\cdots$ 与牛顿动能匹配确定。对坐标速度求偏导，得到正则动量
\[
p_i=\frac{\partial L}{\partial u_i}=m\gamma_uu_i.
\]
自由粒子的正则动量与机械动量相同。勒让德变换给出
\[
H=\bp\cdot\bu-L
=m\gamma_uu^2+\frac{mc^2}{\gamma_u}
=\gamma_umc^2=\sqrt{m^2c^4+c^2p^2}.
\]
于是哈密顿方程直接导出 $\dot{\br}=c^2\bp/H=\bu,\dot{\bp}=0$，与自由直线世界线一致。式 \eqref{eq:hamilton-momentum} 定义的速度海森矩阵在这里为
\[
\frac{\partial^2L}{\partial u_i\partial u_j}
=m\gamma_u\delta_{ij}+\frac{m\gamma_u^3}{c^2}u_iu_j.
\]
在 $m>0,u<c$ 内正定，纵横本征值分别为 $m\gamma_u^3,m\gamma_u$，所以此区域的勒让德反解良好。它同时解释了式 \eqref{eq:sr-force} 的方向差别。

\begin{physicsbox}[title={物理图像与理解：作用量驻值与最大固有时}]
自由惯性世界线在固定类时端点之间使固有时最大，而 $S=-mc^2\tau$ 有负系数，因此该作用量取最小值，当然也满足驻值条件。低速时常数 $-mc^2$ 不影响欧拉--拉格朗日方程，可以从计算运动的拉格朗日量中略去；若把自由粒子哈密顿量解释为总能量 $E=cP^0$，就要恢复静止能的归一化。不能一边任意平移能量零点，一边仍把平移后的量直接代入质量壳。
\end{physicsbox}

\begin{modernbox}[title={进一步探索：带电粒子的哈密顿量}]
沿用第6章电荷符号 $e$ 和给定外场势 $(\phi,\bA)$，最小耦合为
\[
L=-mc^2\sqrt{1-u^2/c^2}+e\bA\cdot\bu-e\phi.
\]
正则动量（即第\ref{ch:hamilton}章的 $\bp$） $\bp_{\rm can}=\gamma_um\bu+e\bA$ 与机械动量 $\boldsymbol{\pi}=\gamma_um\bu$ 不同。勒让德变换给出
\[
H=\sqrt{m^2c^4+c^2|\bp_{\rm can}-e\bA|^2}+e\phi.
\]
由欧拉--拉格朗日方程可得
\[
\dot\pi_i=e(-\partial_i\phi-\partial_t A_i)
+e\sum_j u_j(\partial_i A_j-\partial_j A_i),
\]
即 $\dot{\boldsymbol{\pi}}=e(\bE+\bu\times\bB)$，机械能变化为 $\dd(\gamma_umc^2)/\dd t=e\bE\cdot\bu$。这里只处理给定外场且忽略辐射反作用；外场可对粒子交换能量，粒子单独未必是孤立系统。所需数学是偏导与向量恒等式，可结合第\ref{ch:hamilton}章的规范变换进一步学习。
在 $|\boldsymbol\pi|\ll mc$ 时，根号展开为 $mc^2+|\boldsymbol\pi|^2/(2m)+O(|\boldsymbol\pi|^4/(m^3c^2))$。减去静止能后，恢复式 \eqref{eq:em-hamiltonian}。
\end{modernbox}
自由光子的 $H=c|\bp|$ 可给出速率 $c$，但不能直接在 $-mc^2\int\dd\tau$ 中令 $m=0$ 来推导光子运动：该作用量此时退化。无质量粒子的变分描述需要另引约束形式，留待后续课程。

\section{碰撞、衰变与反应阈能}\label{sec:sr-collision}
碰撞和衰变题通常可先用总能量和总动量求一个不变量，再到总动量为零的系里做较简单的计算。

第\ref{sec:newton-collision}节用质心系分离整体运动；高速时仍可用这一思路，但守恒量要换成总四动量。
\subsection{先求不变量，再选最简单的系}
若题目给的是实验室里的入射量，先在实验室写出总四动量。它的平方在换系时保持不变，因而可直接用于求复合质量或反应阈值。

碰撞前后把所有粒子与逸出辐射计入，使用
\[
\sum_{\rm in}P_a^\mu=\sum_{\rm out}P_b^\mu.
\]
实际求解可按四步进行：列出初态四动量；求总不变质量；在质心系或某粒子静止系求解；需要时再变回实验室。若把不变质量守恒误写成各粒子静质量总和守恒，就会排除内部能和新粒子的产生。

\begin{example}[完全非弹性碰撞：动能去了哪里？]
两个静质量均为 $m$ 的粒子在质心系以速率 $u$ 相向运动，碰后黏成一个静止复合体，不释放辐射。初态总动量为零，总能量 $2\gamma_umc^2$，所以复合体质量
\[
M=2\gamma_um>2m.
\]
相对运动动能成为内部能，增加复合体的不变质量；不是能量消失。若有光子逸出，则应将光子的能动量列入守恒，剩余复合体质量会不同。

具体地，$\Delta E_{\rm int}=(M-2m)c^2=2(\gamma_u-1)mc^2$，正好等于两粒子碰前在质心系的总动能。
\end{example}

\begin{example}[固定靶黏合与复合体速度]
一粒子质量 $m$、总能量 $E$，与实验室静止的同质量粒子黏合，无辐射。初态总能量为 $E+mc^2$，入射动量满足 $p^2c^2=E^2-m^2c^4$。因此
\[
M^2=2m^2+\frac{2mE}{c^2},\qquad V=\frac{pc^2}{E+mc^2}.
\]
取 $E=5mc^2/3$，得 $p=4mc/3,M=4m/\sqrt3,V=c/2$。复合体的机械动量为 $\gamma_VM V=4mc/3$，总能量为 $\gamma_VM c^2=8mc^2/3$，可复核两条守恒。
\end{example}

\subsection{两体衰变的通用解}
一个静止粒子变成两个粒子时，末态动量必须等大反向。若两个粒子质量不同，用各自的能量关系就能确定怎样分配总能量。

静止母粒子质量为 $M$，衰变为自由粒子 $1,2$，质量 $m_1,m_2$。母粒子静止系用星号标记，三动量反向且等大，$\bp_1^*=-\bp_2^*$。由 $P_2=P_0-P_1$ 两边各自平方，右侧为 $P_0^2+P_1^2-2P_0\cdot P_1$。在母粒子静止系，$P_0=(Mc,\boldsymbol0)$，所以 $P_0\cdot P_1=ME_1^*$。将三个质量壳代入得
\[
m_2^2c^2=M^2c^2+m_1^2c^2-2ME_1^*,
\]
所以
\begin{equation}\label{eq:sr-decay}
E_1^*=\frac{M^2+m_1^2-m_2^2}{2M}c^2,\quad
E_2^*=\frac{M^2+m_2^2-m_1^2}{2M}c^2.
\end{equation}
再将 $p_*^2c^2=(E_1^*)^2-m_1^2c^4$ 因式分解，得到
\[
p_* =\frac{c}{2M}\sqrt{[M^2-(m_1+m_2)^2][M^2-(m_1-m_2)^2]}.
\]
阈值直接来自 $Mc^2=E_1^*+E_2^*$ 与 $E_i^*\ge m_i c^2$，故需 $M\ge m_1+m_2$。根号中的乘积即使在 $M<|m_1-m_2|$ 时也可为正，但那不满足此正能量条件。若有无质量末态，等号处的零能量仅表示阈值极限，不能认为一个非零能量光子在静止。

\begin{example}[运动母粒子衰变为两光子]
母粒子静止时两光子能量均为 $Mc^2/2$。若母粒子在实验室以 $V=\beta c$ 沿 $+x$ 运动，第一束光在母系的方向余弦为 $\cos\theta^*$，用逆变换得到
\[
E_1=\frac{\gamma Mc^2}{2}(1+\beta\cos\theta^*),\quad
E_2=\frac{\gamma Mc^2}{2}(1-\beta\cos\theta^*).
\]
两能量之和为母粒子实验室总能量。若沿运动轴发射且 $\beta=0.6$，两光子分别为 $Mc^2$ 与 $Mc^2/4$；能量不等却仍满足总四动量守恒。仅由运动学不能推断发射角的概率分布。
\end{example}

\subsection{两体到两体散射的通用步骤}\label{sec:sr-two-to-two}
两辆小车碰撞时，跟着质心一起走，就能把整个系统的平动扣掉。对高速散射也可以这样做。在质心系，两粒子从相反方向入射，散射后又沿相反方向飞出。能量守恒决定末态动量的大小。它们朝哪个方向飞，则用散射角来描述。

把过程记为 $1+2\to3+4$，$P_i=(E_i/c,\bp_i)$，$m_i$ 为各粒子质量。设实验室系中粒子2静止，入射粒子1沿 $+x$ 方向。定义能量平方不变量 $\mathcal S$，并用 $W=\sqrt{\mathcal S}$ 表示质心系总能量。由 $(P_1+P_2)^2$ 得
\begin{align}
\mathcal S&=c^2(P_1+P_2)^2
=(E_1+m_2c^2)^2-p_1^2c^2
=m_1^2c^4+m_2^2c^4+2m_2c^2E_1,\nonumber\\
\beta_{\rm cm}&=\frac{p_1 c}{E_1+m_2c^2},
\qquad
\gamma_{\rm cm}=\frac{E_1+m_2c^2}{W}.
\label{eq:sr-scatter-fixed}
\end{align}
这里 $E_1$ 是入射粒子的\emph{总能量}，光子取 $m_1=0$、$p_1c=E_1$。$\beta_{\rm cm}$ 是使总纵向动量变成零的参考系速度除以 $c$。

求末态时，把总四动量看成一个不变质量为 $W/c^2$ 的两体系统。用式 \eqref{eq:sr-decay} 的同一计算，有
\begin{align}
E_3^*&=\frac{\mathcal S+m_3^2c^4-m_4^2c^4}{2W},
& E_4^*&=W-E_3^*,\nonumber\\
p_f^*c&=\frac{\sqrt{[\mathcal S-(m_3+m_4)^2c^4]
                              [\mathcal S-(m_3-m_4)^2c^4]}}{2W}.
\label{eq:sr-scatter-cm}
\end{align}
公式用于正能量允许区 $W\ge(m_3+m_4)c^2$。用同一公式对初态质量 $m_1,m_2$ 计算，可得入射质心动量大小 $p_i^*$。若是弹性散射，即 $m_3=m_1,\ m_4=m_2$ 且粒子种类不变，则 $p_f^*=p_i^*=p_*$，散射只改变质心系动量的方向。

以入射方向为基准，令粒子3在质心系的角为 $\theta^*$，选散射平面为 $xy$ 平面，使粒子3的横向动量非负。记 $\sigma_3=+1,\ \sigma_4=-1$，则两个末态粒子变回实验室系后满足
\begin{align}
 E_j&=\gamma_{\rm cm}\bigl(E_j^*
        +\sigma_j\beta_{\rm cm}cp_f^*\cos\theta^*\bigr),\nonumber\\
 p_{jx}&=\gamma_{\rm cm}\bigl(\sigma_jp_f^*\cos\theta^*
                          +\beta_{\rm cm}E_j^*/c\bigr),&
 p_{jy}&=\sigma_jp_f^*\sin\theta^*,\qquad j=3,4.
\label{eq:sr-scatter-boost}
\end{align}
横向分量仍相消，纵向分量各自多了一项。因此在实验室系，两个粒子一般不再背对背。实验室极角可用 $\theta_j=\operatorname{atan2}(|p_{jy}|,p_{jx})$ 求出，范围为 $[0,\pi]$。这里 $\operatorname{atan2}(y,x)$ 按两个分量的符号选择象限。动量为零时没有出射方向。

\begin{figure}[htbp]
\centering
\begin{tikzpicture}[>=Stealth,scale=1.1]
\node at (0,1.58) {质心系};
\node at (5.7,1.58) {实验室系};
\draw[gray,->] (-1.9,0)--(2.05,0) node[right]{$+x$};
\draw[gray,->] (4.15,0)--(8.65,0) node[right]{$+x$};
\fill (0,0) circle (1.5pt);
\fill (5.7,0) circle (1.5pt);
\draw[very thick,blue!75!black,->] (0,0)--(.72,1.247) node[above]{$\bp_3^*$};
\draw[very thick,red!70!black,->] (0,0)--(-.72,-1.247) node[below]{$\bp_4^*$};
\draw (.46,0) arc (0:60:.46);
\node at (.73,.35) {$\theta^*$};
\draw[very thick,blue!75!black,->] (5.7,0)--(8.345,1.247) node[above]{$\bp_3$};
\draw[very thick,red!70!black,->] (5.7,0)--(6.582,-1.247) node[below]{$\bp_4$};
\draw[gray,->] (2.6,.55)--(4,.55) node[midway,above]{$+\beta_{\rm cm}c$};
\end{tikzpicture}
\caption{同一次等质量弹性散射在两个参考系中的末态三动量。图中取 $K=m_Nc^2$、$\theta^*=60^\circ$。转回实验室系时，横向分量不变，两个纵向分量都朝前。}
\label{fig:sr-two-to-two}
\end{figure}

\paragraph{用不变量记录动量转移}
有时实验给的是四动量差的平方。按本书保留 $c$ 的单位，定义
\[
\mathcal T=c^2(P_1-P_3)^2,\qquad
\mathcal U=c^2(P_1-P_4)^2,\qquad
\mathcal S+\mathcal T+\mathcal U=(m_1^2+m_2^2+m_3^2+m_4^2)c^4 .
\]
最后一式由展开平方并用 $P_3+P_4=P_1+P_2$ 得到。在自然单位 $c=1$ 下，这三个量通常写成曼德尔斯塔姆变量 $s,t,u$。弹性散射时，粒子1与3在质心系的能量相同，故
\begin{equation}\label{eq:sr-elastic-transfer}
\mathcal T=-2(p_*c)^2(1-\cos\theta^*),
\qquad -4(p_*c)^2\le\mathcal T\le0 .
\end{equation}
因而 $\mathcal T=0$ 是正向散射端点，$-\mathcal T$ 随偏转角增大。这里用 $\mathcal T$ 表示能量平方，用 $K$ 表示动能。

\subsection{光子--核子弹性散射与康普顿运动学}
光子偏转后，靶会得到反冲动量。用能量和动量守恒求出射光子的能量，若还要换成波长，再用光子的量子关系。

一个能量为 $\varepsilon$ 的光子与初始静止、质量为 $m$ 的自由粒子散射，出射光子能量 $\varepsilon'$、实验室散射角为 $\theta$。光子四动量记为 $Q,Q'$，靶初态为 $P=(mc,\boldsymbol{0})$；末态靶满足 $(P+Q-Q')^2=m^2c^2$。因 $Q^2=Q'^2=0$，展开得
\[
2P\cdot(Q-Q')=2Q\cdot Q'
\quad\Longrightarrow\quad
m(\varepsilon-\varepsilon')=\frac{\varepsilon\varepsilon'}{c^2}(1-\cos\theta).
\]
因此
\[
\boxed{\varepsilon'=\frac{\varepsilon}{1+[\varepsilon/(mc^2)](1-\cos\theta)}}.
\]
再使用量子关系 $\varepsilon=hc/\lambda$，才得到 $\lambda'-\lambda=h(1-\cos\theta)/(mc)$。相对论给出能动量守恒，光子关系则来自量子物理；两者的作用要分清。这里不讨论散射概率，也不把束缚粒子无条件视作静止自由靶。
\begin{example}[反向散射与低能极限]
若 $\varepsilon=mc^2$ 且 $\theta=\pi$，则 $\varepsilon'=mc^2/3$，反冲靶的动能为 $2mc^2/3$，动量大小为 $(\varepsilon+\varepsilon')/c=4mc/3$，正好满足靶的质量壳。若 $\varepsilon\ll mc^2$，相对能量转移很小，与经典低能散射的运动学极限相容。

出射与入射反向，故两光子动量差沿入射方向，大小为 $(\varepsilon+\varepsilon')/c$。靶的末能量是 $5mc^2/3$，其质量壳为 $[(5/3)^2-(4/3)^2]m^2c^4=m^2c^4$。
\end{example}

\paragraph{自由核子靶及其反冲}
以自由质子或中子 $N$ 作固定靶，选定弹性道 $\gamma+N\to\gamma+N$。上面的康普顿运动学可直接使用，只需取 $m=m_N$。若核子束缚在原子核内，或末态另外产生粒子，须按实际初态与末态重列四动量。

为看清两个参考系怎样对应，记 $a=m_Nc^2$。由式 \eqref{eq:sr-scatter-fixed} 和 \eqref{eq:sr-scatter-cm}，
\[
W=\sqrt{a^2+2a\varepsilon},\qquad
\beta_{\rm cm}=\frac{\varepsilon}{a+\varepsilon},\qquad
\varepsilon^*=p_*c=\frac{a\varepsilon}{W}.
\]
实验室中的光子角记为 $\theta_\gamma$。对光子的纵向动量除以能量，有
\[
\cos\theta_\gamma=\frac{\cos\theta^*+\beta_{\rm cm}}
                       {1+\beta_{\rm cm}\cos\theta^*},\qquad
\varepsilon'=\gamma_{\rm cm}\varepsilon^*(1+\beta_{\rm cm}\cos\theta^*)
=\frac{\varepsilon}{1+(\varepsilon/a)(1-\cos\theta_\gamma)}.
\]
最后一式可通过消去 $\cos\theta^*$ 得到，与靶质量壳的结果一致。靶的反冲动能 $K_N=\varepsilon-\varepsilon'$，反冲动量是入射光子动量减去出射光子动量，因此
\[
p_{Nx}c=\varepsilon-\varepsilon'\cos\theta_\gamma,\qquad
p_{Ny}c=-\varepsilon'\sin\theta_\gamma .
\]
另一个实用关系来自 $\mathcal T=c^2(P_N'-P_N)^2$。由于初态核子静止，展开得
\begin{equation}\label{eq:sr-photon-recoil}
-\mathcal T=2a K_N=2\varepsilon\varepsilon'(1-\cos\theta_\gamma).
\end{equation}
所以测得反冲核子的动能也能确定动量转移。

\begin{example}[光子散射到 $60^\circ$ 时的核子反冲]
取入射能量 $\varepsilon=a/3$，测得光子在实验室偏转 $60^\circ$。出射光子能量为 $(a/3)/(1+1/6)=2a/7$，核子动能为 $a/3-2a/7=a/21$。其纵、横动量分别为 $4a/(21c)$、$-\sqrt3a/(7c)$，所以 $|\bp_N|c=\sqrt{43}\,a/21$，反冲角为 $\arctan(3\sqrt3/4)\simeq52.4^\circ$，位于入射轴的另一侧。质量壳给 $(22a/21)^2-(\sqrt{43}a/21)^2=a^2$。再用 $\beta_{\rm cm}=1/4$ 反变换光子的方向，得 $\cos\theta^*=2/7$，可见 $60^\circ$ 只属于实验室系。
\end{example}

\subsection{核子--核子弹性散射}\label{sec:sr-nn-scatter}
让一个核子打到静止核子上。两粒子在质心系中动量大小相同，散射后只改变方向。实验室还看到两者共同向前的运动，因此同一个质心散射角会对应不同的两支实验室角。

以两个质子或两个中子的等质量情形为例，仍记 $a=m_Nc^2$，入射核子动能为 $K$，故 $E_1=a+K$。代入式 \eqref{eq:sr-scatter-fixed} 得
\begin{align*}
\mathcal S&=4a^2+2aK,& p_1c&=\sqrt{K(K+2a)},\\
E_3^*&=E_4^*=\frac{W}{2},&
p_*c&=\sqrt{aK/2},\\
\beta_{\rm cm}&=\sqrt{\frac{K}{K+2a}},&
\gamma_{\rm cm}&=\sqrt{1+\frac{K}{2a}} .
\end{align*}
这里 $p_*^2c^2=W^2/4-a^2$，并且 $\beta_{\rm cm}E_3^*/c=p_*$。因此式 \eqref{eq:sr-scatter-boost} 的两个纵向分量化为 $\gamma_{\rm cm}p_*(1\pm\cos\theta^*)$，均非负。两粒子动能和实验室极角为
\begin{align}
K_3&=\frac K2(1+\cos\theta^*),&
K_4&=\frac K2(1-\cos\theta^*),\nonumber\\
\tan\theta_3&=\frac{1}{\gamma_{\rm cm}}\tan\frac{\theta^*}{2},&
\tan\theta_4&=\frac{1}{\gamma_{\rm cm}}\cot\frac{\theta^*}{2}.
\label{eq:sr-nn-lab}
\end{align}
能量式可用 $\gamma_{\rm cm}W/2=(2a+K)/2$ 和 $\gamma_{\rm cm}\beta_{\rm cm}cp_*=K/2$ 复核。极角式适用于 $K>0$、$0<\theta^*<\pi$。在两个端点，有一个末态核子静止，不能给它定义散射角。对有两个出射方向的事件，横向动量反向，实验室两轨迹的夹角是 $\theta_3+\theta_4$。由于 $\tan\theta_3\tan\theta_4=1/\gamma_{\rm cm}^2<1$，夹角小于 $90^\circ$，在低速极限趋于熟悉的直角关系。

\begin{example}[质心系直角散射怎样分配动能]
取 $K=a$ 且 $\theta^*=90^\circ$。先得 $W=\sqrt6\,a$、$p_*c=a/\sqrt2$。回到实验室，两个核子各得动能 $a/2$，各有 $p_xc=\sqrt3\,a/2$，横向分量 $p_yc=\pm a/\sqrt2$。两极角均为 $\arctan\sqrt{2/3}\simeq39.2^\circ$，轨迹夹角约 $78.5^\circ$。两粒子纵向动量之和为 $\sqrt3\,a/c$，等于入射动量。每个粒子的 $E^2-p^2c^2=(3a/2)^2-5a^2/4=a^2$，也符合各自的质量壳。
\end{example}

若是质子--中子散射且需保留质量差，在通式中分别代入 $m_p,m_n$，不能沿用等质量时的半角和能量分配公式。对于 $pp$ 或 $nn$，末态3、4只是给两条出射轨迹贴的标签。交换两个标签相当于用 $\pi-\theta^*$ 描述另一支轨迹，$\mathcal T$ 和 $\mathcal U$ 也随之互换。上述计算给出允许的能量和角度。核子内部结构、自旋及相互作用决定各方向出现的概率，需要用散射动力学处理。

\subsection{固定靶阈能与对撞机}
同样的实验室输入能量，固定靶和对撞机可用于产生新粒子的能量不同。先算总动量，再比较质心系中的总能量。

设质量 $m_a$ 的入射粒子总能量为 $E_a$，靶粒子质量 $m_b$ 且在实验室静止。沿用本节的能量平方不变量 $\mathcal S=E_{\rm tot}^2-c^2P_{\rm tot}^2$。利用入射粒子质量壳，
\[
\mathcal S=(E_a+m_bc^2)^2-p_a^2c^2
=m_a^2c^4+m_b^2c^4+2m_bc^2E_a.
\]
若末态全为自由有质量粒子，质量和为 $M_f$，其质心系总能量在相对速度均趋零时最小，为 $M_fc^2$。因此吸能反应的运动学阈值为
\begin{equation}\label{eq:sr-threshold}
E_{a,\rm th}=\frac{M_f^2-m_a^2-m_b^2}{2m_b}c^2,\qquad
K_{a,\rm th}=E_{a,\rm th}-m_ac^2.
\end{equation}
若代数阈值为负且末态包含至少两个自由有质量粒子，运动学上可取最低入射动能 $0$，剩余能量进入末态相对运动。若末态只有一个粒子，则必须恰有 $\mathcal S=M_f^2c^4$，不能用相对运动吸收剩余能量。含光子末态需单独讨论零能极限；“所有末态粒子在质心系静止”不能照搬给非零能量光子。满足运动学阈值也不保证实际反应会发生，还需相互作用与其他守恒量允许。

\begin{example}[达到同一产生阈值，固定靶和对撞各需多少动能？]
取 $m_a=m_b=m$，假想允许产生总静质量 $M_f=4m$ 的全有质量末态。固定靶需要 $E_a=7mc^2$，即入射动能 $6mc^2$；等能反向对撞只需每束总能量 $2mc^2$，总输入动能 $2mc^2$。固定靶的大部分额外能量被整个末态系统的平动占用，对撞可让总三动量为零，更多能量用于内部产生。反过来，若给定同样的总动能 $2mc^2$，固定靶的质心能量为 $\sqrt{8}\,mc^2$，等能对撞为 $4mc^2$。此处只比较碰撞运动学。

固定靶的阈条件为 $2m^2c^4+2mc^2E_a=(4mc^2)^2$。若入射动能只有 $2mc^2$，则 $E_a=3mc^2$，左端为 $8m^2c^4$。对撞时总动量相消，质心能量就是两束总能量之和。
\end{example}

\begin{figure}[htbp]\centering
\begin{tikzpicture}[>=Stealth,scale=.95]
\node[anchor=west] at (-.3,2.3){固定靶：总动量不为零};
\draw[blue,very thick,->] (0,1.5)--(2.2,1.5) node[midway,above]{$E_a,p_a$};
\fill (2.6,1.5) circle(4pt) node[right]{靶 $m$ 静止};
\draw[->,thick] (4.7,1.5)--(6.4,1.5) node[right]{末态整体运动};
\node[anchor=west] at (-.3,.4){等能对撞：总动量为零};
\draw[blue,very thick,->] (0,-.4)--(2.3,-.4) node[midway,above]{$E,p$};
\draw[red,very thick,->] (5.2,-.4)--(2.9,-.4) node[midway,above]{$E,-p$};
\node at (2.6,-1){质心系即实验室系};
\end{tikzpicture}
\caption{固定靶与等能反向对撞的初态动量。箭头只表示入射方向，不按两装置所需动量比例绘制；定量比较由 $\mathcal S$ 完成。}\label{fig:sr-collider}
\end{figure}

\begin{modernbox}[title={进一步探索：光子产生粒子对为什么需要第二个物体？}]
单个自由光子在真空中不能变成两个有质量粒子：其初态总四动量平方为零，而两未来指向的有质量粒子总四动量为类时，无法相等。两光子可提供非零系统不变质量。产生两个质量均为 $m$ 的粒子，运动学条件为
\[
2E_1E_2(1-\cos\theta)\ge4m^2c^4.
\]
若光子打到质量为 $M_T$ 的静止靶上，靶参与反冲、末态靶质量保持 $M_T$，则总末态静质量至少为 $M_T+2m$，由同一不变量方法得到光子阈能
\[
\varepsilon_{\rm th}=2mc^2\left(1+\frac{m}{M_T}\right).
\]
有限靶反冲使阈值高于 $2mc^2$；反应概率仍需量子理论确定。
\end{modernbox}

\section{计算与模拟：沿世界线积累固有时}
知道每个时刻的速率，就可以把旅行切成许多短时间片，逐片算旅行者的钟走了多少。加密时间片，检查总固有时间是否收敛。

对平滑往返速度 $u(t)=u_0\sin(2\pi t/T)$，令 $s=t/T$，用复合辛普森求积计算
\[
\frac{\tau}{T}=\int_0^1\sqrt{1-(u_0/c)^2\sin^2(2\pi s)}\,\dd s.
\]
用偶数个子区间，减半步长比较结果，并检查位移积分为零。伴随脚本 \texttt{review/verify\_ch07.py} 给出核查值；使用 $c=1$ 时速度仍须小于1。

\section*{关键概念与常见误区}

写式子之前，先在草图上标明是谁相对谁运动。$v$ 表示两参考系的相对速度，$u$ 表示选定系中粒子的速度。做碰撞题时，再用 $m$ 和 $M$ 区分单粒子质量与系统质量，注明用的是总能量还是动能。若题里有角度，也应标出由哪一系测量。算完可用量纲、守恒和低速极限检验结果。

\section*{本章小结}
\addcontentsline{toc}{section}{本章小结}
\begin{keybox}
\textbf{事件与测量：}先选事件，再写变换。钟慢要求同一只钟的两个事件；测长要求所选系的同时端点事件；收到信号的时刻还要计算光行时。

\textbf{时空与不变量：}$\Delta s^2=c^2\Delta t^2-\Delta r^2$，固有时沿世界线积分；惯性系对坐标值可有不同判断，对因果类型和同一条世界线的固有时一致。

\textbf{动力学与守恒：}$P=(E/c,\bp)$，$E^2=p^2c^2+m^2c^4$。总四动量守恒，系统不变质量包括内部能；作用量和勒让德变换延续第2、3、6章的方法。
\textbf{两体散射。}先由初态求 $W=\sqrt{\mathcal S}$，再在质心系求末态动量大小。选定质心散射角后，变回实验室系即可预测两条轨迹的能量和方向。
\end{keybox}

\Needspace{7\baselineskip}
\section*{分层习题}
\addcontentsline{toc}{section}{分层习题}
先列事件及所属参考系。碰撞或衰变题可先把能量和三动量分别相加，再选择能简化计算的参考系。

以下均忽略引力；有质量粒子的 $m$ 恒定，除非题目另有说明。
\begin{bookexercises}
\item \exlevel{基础} \textbf{事件与接收。}两盏实验室静止灯在 $x=\pm L$、$t=0$ 闪光。观察者从原点以 $v>0$ 匀速出发。求其接收左右光的实验室时刻及自身钟读数；说明这与两发光事件的坐标时间差有何不同。
\item \exlevel{基础} \textbf{同步钟。}静止钟 $A$ 发光时为 $2t_0$，收回反射光时为 $8t_0$。求远处反射事件的同步时间及两钟距离。若改由运动观察者描述，是否仍可把这两个 $A$ 钟读数直接当作其坐标时间？
\item \exlevel{基础} \textbf{洛伦兹变换。}取 $v=0.6c$，将 $(ct,x)=(5\ell,\ell)$ 变换到 $S'$，核查间隔并作逆变换。
\item \exlevel{基础} \textbf{同时系与同地系。}分别对事件差 $(c\Delta t,\Delta x)=(3\ell,5\ell)$ 和 $(5\ell,3\ell)$，判断类型并找出同时系或同地系，求不变量大小。
\item \exlevel{基础} \textbf{长度与通过时间。}静止长 $L_0$ 的杆以 $0.8c$ 经过实验室固定探测器，求实验室杆长和探测器记录的通过时间；再到杆系复核。
\item \exlevel{基础} \textbf{寿命。}粒子按固有时作指数衰变，平均寿命为 $\tau_0$，实验室速率 $0.6c$。求经过距离 $d=c\tau_0$ 后的存活比例，并解释平均寿命与确定寿命的区别。
\item \exlevel{基础} \textbf{双生子。}地球总历时 $10$ 年，旅行者以 $0.6c$ 对称往返。求旅行者固有时及瞬时掉头前后同时线在地球钟上的两个交点。
\item \exlevel{基础} \textbf{速度与光行差。}飞船沿实验室 $+x$ 以 $0.6c$ 运动，两系坐标轴平行。在飞船系沿 $+y'$ 以 $0.8c$ 发射探测器，求实验室速度分量和速率。若改成光束，求其实验室方向。
\item \exlevel{基础} \textbf{多普勒角度。}在 $S$ 中光线与 $+x$ 成 $60^\circ$，$S'$ 沿 $+x$ 以 $0.6c$ 运动，两系坐标轴平行。求 $\omega'/\omega,\cos\theta'$，并用反变换复核。
\item \exlevel{基础} \textbf{能动量。}粒子动能为 $mc^2$，求 $E,p,u$。沿其运动方向以 $u$ 运动的观察者测得什么能量？
\item \exlevel{基础} \textbf{双光子。}两光子等能 $\varepsilon$，夹角 $120^\circ$。求总能量、总动量大小、系统不变质量及质心系速度；若夹角趋零会怎样？
\item \exlevel{基础} \textbf{恒力。}粒子从 $x=0$ 处静止释放后受沿 $+x$ 的恒力 $F>0$，在 $Ft/(mc)=1$ 时求速率、位移、动能、固有时与坐标加速度。
\item \exlevel{综合} \textbf{作用量。}从自由粒子 $L$ 推导速度海森矩阵的纵横本征值，并从 $H$ 求 $\partial H/\partial p_i$，核查两种速度表达式互逆。解释为什么直接令 $m=0$ 的固有时作用量失效。
\item \exlevel{综合} \textbf{黏合。}质量 $m$、总能量 $3mc^2$ 的粒子打到静止同质量靶并黏合，无辐射。求复合体质量与实验室速度，在复合体静止系检验总能量。
\item \exlevel{综合} \textbf{不等质量衰变。}母粒子质量 $5m$，衰变为质量 $m,2m$ 的两个粒子。求母系两能量和动量大小，并说明两个能量一般不相等的原因。
\item \exlevel{综合} \textbf{康普顿散射。}入射光子能量 $mc^2/2$、散射角 $90^\circ$。求出射光子能量、靶反冲动能与动量大小，核查质量壳。
\item \exlevel{综合} \textbf{阈能。}质量均为 $m$ 的入射粒子和静止靶，产生允许的全有质量末态，末态静质量和为 $3m$。求固定靶入射动能阈值；比较等能反向对撞所需的总动能。
\item \exlevel{拓展} \textbf{粒子对阈值。}两反向等能光子产生质量均为 $m$ 的一对粒子，各光子至少需多少能量？若光子改打质量 $M_T=10m$ 的静止靶且靶保持原质量，求光子阈能。
\item \exlevel{拓展} \textbf{平滑掉头。}构造世界线 $u(t)=u_0\sin(2\pi t/T)$，$0<u_0<c$，$0\le t\le T$。证明回到出发位置；用数值积分算 $u_0=0.8c$ 时 $\tau/T$，比较同端点惯性钟及同样峰值速度的瞬时掉头模型。缩小积分步长检查收敛。
\item \exlevel{拓展} \textbf{加速器。}均匀磁场中粒子垂直入射，忽略辐射。由轨迹半径 $\rho$ 推导 $E,\gamma_u$ 和回旋角频率；固定 $B$ 时增大能量，半径和角频率怎样变化？指出电场补充能量与磁场转向的区别。
\item \exlevel{综合} \textbf{从反冲动能还原光子方向。} 能量 $\varepsilon=m_Nc^2/2$ 的光子打到自由静止核子，弹性散射后测得核子反冲动能为 $m_Nc^2/8$。求出射光子能量、实验室与质心系的散射角，以及 $-\mathcal T$。对同一入射能量，核子反冲动能最大是多少？
\item \exlevel{综合} \textbf{固定靶核子--核子散射。} 等质量核子的入射动能为 $K=2m_Nc^2$，出射粒子3的质心散射角为 $60^\circ$。求两个末态核子的实验室动能和极角，计算 $\mathcal T,\mathcal U$。交换末态标签会怎样？
\end{bookexercises}

\section*{习题解答}
先用同一题号找到解答。核对时，从第一条方程开始比较自己的做法，再检查最后的符号、单位和适用条件。
\begin{bookexercises}
\item 右灯光线 $x=L-ct$ 与 $x=vt$ 相交，$t_R=L/(c+v)$；左灯光线 $x=-L+ct$ 给 $t_L=L/(c-v)$。各接收时刻的固有时为 $t_{R,L}/\gamma$。发光事件则有 $t'_R=-\gamma vL/c^2,t'_L=+\gamma vL/c^2$，它们是另一组事件。

上述固有时以观察者经过原点的出发事件为零点。由 $t'=\gamma(t-vx/c^2)$，在观察者的接收世界线上代入 $x=vt$ 得 $t'=t/\gamma$；在灯的发光事件上则应代入 $t=0,x=\pm L$。
\item 反射时间 $5t_0$，距离 $3ct_0$。$A$ 钟只在其静止系中直接给该位置坐标时间；运动系需变换事件，不能照搬读数。

静止系中去程与回程等长，所用时间均为 $(8-2)t_0/2=3t_0$。故反射时刻是两读数的中点 $(2+8)t_0/2$，距离是单程时间乘以 $c$。
\item $ct'=11\ell/2,x'=-5\ell/2$，间隔为 $24\ell^2$；逆变换恢复原坐标。

这里 $\gamma=5/4$。直接检查 $(11\ell/2)^2-(-5\ell/2)^2=24\ell^2$。逆变换给 $ct=(5/4)(11/2-3/2)\ell=5\ell$、$x=(5/4)(-5/2+33/10)\ell=\ell$。
\item 第一组类空，$v=3c/5$ 使同时，此系空间间隔 $4\ell$；第二组类时，$v=3c/5$ 使同地，固有时为 $4\ell/c$。

第一组令 $c\Delta t'=\gamma(c\Delta t-\beta\Delta x)=0$，故 $\beta=3/5$。第二组令 $\Delta x'=\gamma(\Delta x-\beta c\Delta t)=0$，也得 $\beta=3/5$。两组 $\Delta s^2$ 分别为 $-16\ell^2$ 和 $16\ell^2$，所以分别取纯空间和纯时间间隔时得到上述数值。
\item $L=3L_0/5$，$\Delta t=3L_0/(4c)$。杆系探测器以 $0.8c$ 通过距离 $L_0$，$\Delta t'=5L_0/(4c)=\gamma\Delta t$。

由 $\gamma=5/3$ 得 $L=L_0/\gamma$。固定探测器处两端经过同一位置，故实验室时间差为 $L/(0.8c)$，也是探测器的固有时间。杆系中的两次通过发生在不同位置，故其坐标时间差为 $\gamma\Delta t$。
\item $\gamma u\tau_0=3c\tau_0/4$，故比例 $e^{-4/3}\simeq0.2636$。平均寿命是指数分布的均值，单个粒子衰变时刻不确定。

实验室飞行时间为 $d/u=5\tau_0/3$，粒子钟只经过 $(d/u)/\gamma=4\tau_0/3$。对固有时指数衰变模型，存活概率为 $\exp[-(d/u)/(\gamma\tau_0)]$。
\item 旅行者 $8$ 年；地球交点为 $5(1-0.36)=3.2$ 年与 $5(1+0.36)=6.8$ 年。

地球系中掉头事件为 $(t_A,x_A)=(5\,\mathrm{yr},3\,\text{光年})$。旅行者的两段固有时各为 $5/\gamma=4$ 年。设地球钟交点为 $(t_E,0)$，与掉头事件同时要求 $(t_E-t_A)-v_{\rm leg}(0-x_A)/c^2=0$，分别用 $v_{\rm leg}=+0.6c$ 和 $-0.6c$，得到 $t_E=t_A\mp1.8\,\mathrm{yr}$。
\item 探测器 $u_x=0.6c,u_y=0.64c$，速率 $\sqrt{0.7696}c\simeq0.8773c$。光束则 $u_x=0.6c,u_y=0.8c$，方向余弦为 $0.6$，仍为光速。

逆速度变换给 $u_x=(u_x'+v)/(1+vu_x'/c^2)$、$u_y=u_y'/[\gamma(1+vu_x'/c^2)]$。这里 $u_x'=0$，所以分母分别为1和 $5/4$。光束情形用 $u_y'=c$，其在 $xy$ 平面内与 $+x$ 成 $\arccos(0.6)\simeq53.13^\circ$。
\item 频率比 $7/8$，$\cos\theta'=-1/7$。反变换因子 $\gamma(1+\beta\cos\theta')=8/7$，乘积为 $1$。

四维波矢变换给 $\omega'/\omega=\gamma(1-\beta\cos\theta)=(5/4)(1-0.6/2)$，而 $\cos\theta'=(\cos\theta-\beta)/(1-\beta\cos\theta)=(1/2-0.6)/(1-0.6/2)$。反向角变换给 $(\cos\theta'+\beta)/(1+\beta\cos\theta')=1/2$。
\item $E=2mc^2,p=\sqrt3mc,u=\sqrt3c/2$。共动观察者测 $mc^2$，可用 $P\cdot U$ 直接核查。

用 $p^2c^2=E^2-m^2c^4$ 得 $p$，再用 $u=pc^2/E$ 得速率。共动观察者的洛伦兹因子为2，故 $E'=2(E-up)=mc^2$。
\item 总能量 $2\varepsilon$，总动量 $\varepsilon/c$，$M=\sqrt3\varepsilon/c^2,V=c/2$。夹角趋零时 $M\to0,V\to c$，不存在极限的惯性质心系。

两动量沿角平分线叠加，大小为 $2(\varepsilon/c)\cos60^\circ=\varepsilon/c$。由 $M^2c^4=(2\varepsilon)^2-|\bp_{\rm tot}|^2c^2$ 得不变质量。沿角平分线作速度 $V=c^2|\bp_{\rm tot}|/E_{\rm tot}$ 的洛伦兹变换可使总动量为零。
\item $u=c/\sqrt2,x=(\sqrt2-1)mc^2/F,K=(\sqrt2-1)mc^2$，$\tau=(mc/F)\ln(1+\sqrt2)$，$a=F/(2\sqrt2m)$。

令 $z=Ft/(mc)$，有 $p=Ft$、$u/c=z/\sqrt{1+z^2}$。由 $x(0)=0$ 积分得
\[
x=\frac{mc^2}{F}\bigl(\sqrt{1+z^2}-1\bigr),
\]
动能也等于 $Fx$。因 $\gamma_u=\sqrt{1+z^2}$，积分 $\dd\tau=\dd t/\gamma_u$ 得 $\tau=(mc/F)\operatorname{arsinh}z$。对 $u$ 求导得 $a=(F/m)(1+z^2)^{-3/2}$，最后令 $z=1$。
\item 海森矩阵为 $m\gamma_u I+m\gamma_u^3\bu\bu^{\mathsf T}/c^2$，作用在纵向给 $m\gamma_u^3$，横向给 $m\gamma_u$。$\partial H/\partial p_i=c^2p_i/H=u_i$。$m=0$ 的原作用量恒零，不能产生光子运动方程。

从 $L=-mc^2\sqrt{1-|\bu|^2/c^2}$ 先求 $p_i=m\gamma_u u_i$。再用 $\partial\gamma_u/\partial u_j=\gamma_u^3u_j/c^2$ 即得海森矩阵。纵向本征值化简依赖 $1+\gamma_u^2|\bu|^2/c^2=\gamma_u^2$。反过来，$H=\sqrt{m^2c^4+c^2|\bp|^2}$ 给 $\bu=c^2\bp/H$，因此 $1-|\bu|^2/c^2=m^2c^4/H^2$，恢复 $\bp=m\gamma_u\bu$。
\item $p=\sqrt8mc,M=\sqrt8m,V=c/\sqrt2$。质心能量为 $Mc^2$，$\gamma_V=\sqrt2$ 给回实验室总能量 $4mc^2$。

碰前总能量 $4mc^2$、总动量 $\sqrt8mc$，故 $Mc^2=\sqrt{(4mc^2)^2-(\sqrt8mc^2)^2}$，而 $V=c^2p_{\rm tot}/E_{\rm tot}$。推动到复合体静止系，$p^*=\gamma_V(p_{\rm tot}-VE_{\rm tot}/c^2)=0$，$E^*=\gamma_V(E_{\rm tot}-Vp_{\rm tot})=\sqrt8mc^2$。
\item $E_1^*=11mc^2/5,E_2^*=14mc^2/5,p_*=4\sqrt6mc/5$。相反等大的是三动量；质量不同使能量不同，且两能量和为 $5mc^2$。

母粒子静止时两动量为 $\pm\bp_*$，故质量壳相减给 $(E_2^*)^2-(E_1^*)^2=3m^2c^4$，同时 $E_2^*+E_1^*=5mc^2$。解得 $E_2^*-E_1^*=3mc^2/5$，再由 $p_*c=\sqrt{(E_1^*)^2-m^2c^4}$ 得动量大小。
\item $\varepsilon'=mc^2/3$，反冲 $K=mc^2/6$；$p=\sqrt{13}mc/6$，于是 $(mc^2+K)^2-p^2c^2=m^2c^4$。

取入射沿 $+x$，出射沿 $+y$，则靶动量分量为 $(\varepsilon/c,-\varepsilon'/c)$，末能量为 $mc^2+\varepsilon-\varepsilon'$。代入质量壳并约去公共项，得 $mc^2(\varepsilon-\varepsilon')=\varepsilon\varepsilon'$，即 $\varepsilon'=\varepsilon/(1+\varepsilon/(mc^2))$。数值代入后，能量平方与动量项之差为 $(49-13)m^2c^4/36$。
\item 固定靶 $E_a=7mc^2/2,K_a=5mc^2/2$；对撞每束 $E=3mc^2/2$，总动能 $mc^2$。两方案总初态静质量均为 $2m$。

固定靶初态的不变量为 $E_{\rm cm}^2=2m^2c^4+2mc^2E_a$。令它达到 $(3mc^2)^2$ 求 $E_a$，再减去 $mc^2$ 得阈动能。对撞总动量为零，阈值是 $2E=3mc^2$，总动能由 $2E-2mc^2$ 给出。
\item 两等能反向光子的总动量为零，阈条件 $2\varepsilon=2mc^2$ 给出每个光子至少需 $mc^2$。固定靶初态的不变量为
\[
E_{\rm cm}^2=(M_Tc^2+\varepsilon)^2-\varepsilon^2
=M_T^2c^4+2M_Tc^2\varepsilon .
\]
阈值处末态各粒子在质心系没有相对运动，故令 $E_{\rm cm}=(M_T+2m)c^2$，解得 $\varepsilon_{\rm th}=2mc^2+2m^2c^2/M_T=2.2mc^2$。其中 $2mc^2$ 用于新增静质量，额外的 $0.2mc^2$ 是末态整体在实验室系运动所需的动能，包含原靶和新生粒子的动能。
\item $\int_0^T u\,\dd t=0$。$\tau/T=\int_0^1\sqrt{1-0.64\sin^2(2\pi s)}\,\dd s\simeq0.81255$；介于地球钟 $1$ 和恒速绝对值 $0.8c$ 的往返模型 $0.6$ 之间。峰值相同并不意味着整段速率相同。

以 $f(s)=\sqrt{1-0.64\sin^2(2\pi s)}$，等分 $[0,1]$，复合 Simpson 公式用权重 $1,4,2,\ldots,4,1$，加权和乘 $1/(3N)$。所得数值为
\[
\begin{array}{c|ccc}
N & 16 & 32 & 64\\ \hline
\tau/T & 0.812628877930 & 0.812549714011 & 0.812549610283
\end{array}
\]
相邻差值由约 $7.92\times10^{-5}$ 缩小到 $1.04\times10^{-7}$。
\item $p=|e|B\rho$，$E=\sqrt{m^2c^4+e^2B^2\rho^2c^2}$，$\gamma_u=E/(mc^2)$，$\omega_c=|e|B/(\gamma_um)$。固定 $B$ 下能量增加使半径增大、角频率降低；磁场本身不做功。

磁力垂直速度，故 $E$、$p$ 的大小恒定。匀速转向使 $|\dot{\bp}|=pu/\rho$，与 $|e|uB$ 相等给磁刚度关系。再用 $\omega_c=u/\rho$、$p=\gamma_um u$ 得角频率。电场若沿运动方向有分量，功率为 $e\bE\cdot\bu$，可改变能量。
\item 记 $a=m_Nc^2$。能量守恒给 $\varepsilon'=a/2-a/8=3a/8$。由康普顿关系，
\[
1-\cos\theta_\gamma=\frac{aK_N}{\varepsilon\varepsilon'}=\frac23 ,
\qquad
-\mathcal T=2aK_N=\frac{a^2}{4}.
\]
所以 $\theta_\gamma=\arccos(1/3)\simeq70.5^\circ$。由 $\beta_{\rm cm}=\varepsilon/(a+\varepsilon)=1/3$，逆光行差给 $\cos\theta^*=(\cos\theta_\gamma-\beta_{\rm cm})/(1-\beta_{\rm cm}\cos\theta_\gamma)=0$，即 $\theta^*=90^\circ$。光子在实验室反向散射时 $1-\cos\theta_\gamma$ 最大，$\varepsilon'$ 最小，为 $(a/2)/(1+1)=a/4$，故 $K_{N,\max}=a/4$。
\item 记 $a=m_Nc^2$。有 $\gamma_{\rm cm}=\sqrt2$、$p_*c=a$。将 $\cos\theta^*=1/2$ 代入式 \eqref{eq:sr-nn-lab}，得 $K_3=3a/2,\ K_4=a/2$，两者之和为 $2a$。两极角满足
\[
\tan\theta_3=\frac1{\sqrt6},\qquad
\tan\theta_4=\sqrt{\frac32},
\]
分别约为 $22.2^\circ$ 与 $50.8^\circ$，位于束轴两侧。由式 \eqref{eq:sr-elastic-transfer} 得 $\mathcal T=-a^2$，对粒子4用 $\pi-\theta^*$ 得 $\mathcal U=-3a^2$。亦可由 $\mathcal S=8a^2$ 核对 $\mathcal S+\mathcal T+\mathcal U=4a^2$。交换标签会交换两份动能、两个角度和 $\mathcal T,\mathcal U$，对应同一对末态轨迹。
\end{bookexercises}

\section*{进一步阅读与来源}
本章采用事件表、世界线和不变量组织推导。以下为可核实的官方或作者原始资料，正文例题和插图按本书符号独立构造（链接复核：2026-11-08）：
\begin{bookreading}
\item \href{https://ocw.mit.edu/courses/8-033-introduction-to-relativity-and-spacetime-physics-fall-2024/mit8_033_f24_lec04.pdf}{MIT 8.033，2024，Lecture 4}：时空图与同时性。可比较不同事件选择如何改变时间和长度测量。
\item \href{https://ocw.mit.edu/courses/8-033-introduction-to-relativity-and-spacetime-physics-fall-2024/mit8_033_f24_lec08.pdf}{MIT 8.033，2024，Lecture 8}：四动量应用与观察者能量。该资料采用 $(-,+,+,+)$，与本书号差相反；转换时需同时改变内积与质量壳的符号。
\item \href{https://davidtong.org/pdfs/teaching/dynamics-and-relativity/dynrel7.pdf}{David Tong：Special Relativity}：本科相对论讲义。可选读快度与碰撞；后部旋量等内容超出本书要求。
\item \href{https://pdg.lbl.gov/2024/reviews/rpp2024-rev-kinematics.pdf}{Particle Data Group，Kinematics，第49章}。可对照固定靶的质心系变换和两体散射不变量。该资料使用自然单位 $c=1$。
\end{bookreading}
牛顿力学并未被丢弃，而是在 $u/c\ll1$、测量精度允许时成为相对论的近似；变分法、拉格朗日和哈密顿方法则通过改变作用量继续适用。全书的共同线索是：先明确物理对象和可测量量，再选择合适变量，用对称性与不变量建立可检验的运动规律。


\end{document}
